\xchapter{Conspiracy of Catiline.}

\xsection{The Argument.}

The Introduction, I.-IV. The character of Catiline, V. Virtues of the
ancient Romans, VI.-IX. Degeneracy of their posterity, X.-XIII.
Catiline's associates and supporters, and the arts by which he
collected them, XIV. His crimes and wretchedness, XV. His tuition of
his accomplices, and resolution to subvert the government, XVI. His
convocation of the conspirators, and their names, XVII. His concern in
a former conspiracy, XVIII., XIX. Speech to the conspirators, XX. His
promises to them, XXI. His supposed ceremony to unite them, XXII. His
designs discovered by Fulvia, XXIII. His alarm on the election of
Cicero to the consulship, and his design in engaging women in his
cause, XXIV. His accomplice, Sempronia, characterized, XXV. His
ambition of the consulship, his plot to assassinate Cicero, and his
disappointment in both, XXVI. His mission of Manlius into Etruria, and
his second convention of the conspirators, XXVII. His second attempt
to kill Cicero; his directions to Manlius well observed, XXVIII. His
machinations induce the Senate to confer extraordinary power on the
consuls, XXIX. His proceedings are opposed by various precautions,
XXX. His effrontery in the Senate, XXXI. He sets out for Etruria,
XXXII. His accomplice, Manlius, sends a deputation to Marcius, XXXIII.
His representations to various respectable characters, XXXIV. His
letter to Catulus, XXXV. His arrival at Manlius's camp; he is declared
an enemy by the Senate; his adherents continue faithful and resolute,
XXXVI. The discontent and disaffection of the populace in Rome,
XXXVII. The old contentions between the patricians and plebeians,
XXXVIII. The effect which a victory of Catiline would have produced,
XXXIX. The Allobroges are solicited to engage in the conspiracy, XL.
They discover it to Cicero, XLI. The incaution of Catiline's
accomplices in Gaul and Italy, XLII. The plans of his adherents at
Rome, XLIII. The Allobroges succeed in obtaining proofs of the
conspirators' guilt, XLIV. The Allobroges and Volturcius are arrested
by the contrivance of Cicero, XLV. The principal conspirators at Rome
are brought before the Senate, XLVI. The evidence against them, and
their consignment to custody, XLVII. The alteration in the minds of
the populace, and the suspicions entertained against Crassus, XLVIII.
The attempts of Catulus and Piso to criminate Caesar, XLIX. The plans
of Lentulus and Cethegus for their rescue, and the deliberations of
the Senate, L. The speech of Caesar on the mode of punishing the
conspirators, LI. The speech of Cato on the same subject, LII. The
condemnation of the prisoners; the causes of Roman greatness, LIII.
Parallel between Caesar and Cato, LIV. The execution of the criminals,
LV. Catiline's warlike preparations in Etruria, LVI. He is compelled
by Metullus and Antonius to hazard an action, LVII. His exhortation to
his men, LVIII. His arrangements, and those of his opponents, for the
battle, LIX. His bravery, defeat, and death, LX., LXI.

* * * * *

I. It becomes all men, who desire to excel other animals,\footnote{I. Desire to excel other animals—\emph{Sese student praestare
caeteris animalibus.} The pronoun, which is usually omitted, is, says
Cortius, not without its force; for it is equivalent to \emph{ut ipsi}:
student \emph{ut ipsi praestent}. In support of his opinion he quotes, with
other passages, Plaut. Asinar. i. 3, 31: Vult placere sese amicae,
i.e. vult \emph{ut ipse amicae placeat}; and Coelius Antipater apud Festum
in ``Topper,'' Ita uti sese quisque vobis studeat aemulari, i.e.
\emph{studeat ut ipse aemuletur.} This explanation is approved by Bernouf.
Cortius might have added Cat. 7: \emph{sese} quisque hostem \emph{ferire
—properabat.} ``Student,'' Cortius interprets by ``cupiunt.''} to strive,
to the utmost of their power,\footnote{To the utmost of their power—\emph{Summâ ope}, with their utmost
ability. ``A Sallustian mode of expression. Cicero would have said
\emph{summâ operâ, summo studio, summâ contentione.} Ennius has '\emph{Summa
nituntur opum vi}.''' Colerus.} not to pass through life in obscurity,
\footnote{In obscurity—\emph{Silentio.} So as to have nothing said of them,
either during their lives or at their death. So in c. 2: \emph{Eorum ego
vitam mortemque juxta aestumo, quoniam de utrâque siletur}. When Ovid
says, \emph{Bene qui latuit, bene vixit,} and Horace, \emph{Nec vixit malè, qui
vivens moriensque fefellit,} they merely signify that he has some
comfort in life, who, in ignoble obscurity, escapes trouble and
censure. But men thus undistinguished are, in the estimation of
Sallust, little superior to the brute creation. ``Optimus quisque,''
says Muretus, quoting Cicero, ``honoris et gloriae studio maximè
ducitur;'' the ablest men are most actuated by the desire of honor and
glory, and are more solicitous about the character which they will
bear among posterity. With reason, therefore, does Pallas, in the
Odyssey, address the following exhortation to Telemachus:} like the beasts of the field,\footnote{Like the beasts of the field—\emph{Veluti pecora.} Many translators
have rendered \emph{pecora} ``brutes'' or ``beasts;'' \emph{pecus}, however, does
not mean brutes in general, but answers to our English word \emph{cattle}.} which nature has formed groveling\footnote{Groveling—\emph{Prona.} I have adopted \emph{groveling} from Mair's
old translation. \emph{Pronus}, stooping \emph{to the earth}, is applied to
\emph{cattle}, in opposition to \emph{erectus}, which is applied to \emph{man}; as
in the following lines of Ovid, Met. i.:}
and subservient to appetite.

All our power is situate in the mind and in the body.\footnote{All our power is situate in the mind and in the body—\emph{Sed
omnis nostra vis in animo et corpore sita}. All our power is placed,
or consists, in our mind and our body. The particle \emph{sed,} which is
merely a connective, answering to the Greek \emph{dé}, and which would be
useless in an English translation, I have omitted.} Of the mind
we rather employ the government;\footnote{Of the mind we—employ the government—\emph{Animi imperio—utimur}.
``What the Deity is in the universe, the mind is in man; what matter
is to the universe, the body is to us; let the worse, therefore,
serve the better.''—Sen. Epist. lxv. \emph{Dux et imperator vitae mortalium
animus est,} the mind is the guide and ruler of the life of mortals.
—Jug. c. 1. ``An animal consists of mind and body, of which the one
is formed by nature to rule, and the other to obey.''—Aristot. Polit.
i. 5. Muretus and Graswinckel will supply abundance of similar passages.} of the body the service.\footnote{Of the mind we rather employ the government; of the body, the
service—\emph{Animi imperio, corporis servitio, magis utimur}. The word
\emph{magis} is not to be regarded as useless. ``It signifies,'' says Cortius,
``that the mind rules, and the body obeys, \emph{in general}, and \emph{with
greater reason}.'' At certain times the body may \emph{seem to have the
mastery}, as when we are under the irresistible influence of hunger
or thirst.} The
one is common to us with the gods; the other with the brutes. It
appears to me, therefore, more reasonable\footnote{It appears to me, therefore, more reasonable, etc.—\emph{Quo mihi
rectius videtur}, etc. I have rendered \emph{quo} by \emph{therefore}. ``\emph{Quo},''
observes Cortius, ``is \emph{propter quod}, with the proper force of the
ablative case. So Jug. c. 84: \emph{Quo} mihi acrius adnitendum est, etc;
c. 2, \emph{Quo} magis pravitas eorum admiranda est. Some expositors would
force us to believe that these ablatives are inseparably connected
with the comparative degree, as in \emph{quo minus, eo major}, and similar
expressions; whereas common sense shows that they can not be so
connected.'' Kritzius is one of those who interprets in the way to
which Cortius alludes, as if the drift of the passage were, \emph{Quanto
magis animus corpori praestat, tanto rectius ingenii opibus gloriam
quaerere}. But most of the commentators and translators rightly follow
Cortius. ``\emph{Quo},'' says Pappaur, ``is for \emph{quocirca}.''}to pursue glory by means
of the intellect than of bodily strength, and, since the life which we
enjoy is short, to make the remembrance of us as lasting as possible.
For the glory of wealth and beauty is fleeting and perishable; that of
intellectual power is illustrious and immortal.\footnote{\emph{That of} intellectual power is illustrious and immortal—\emph{Virtus
clara aeternaque habetur}. The only one of our English translators who
has given the right sense of \emph{virtus} In this passage, is Sir Henry
Steuart, who was guided to it by the Abbé Thyvon and M. Beauzée.
``It appears somewhat singular,'' says Sir Henry, ``that none of the
numerous translators of Sallust, whether among ourselves or among
foreign nations—the Abbé Thyvon and M. Beauzée excepted—have thought
of giving to the word \emph{virtus}, in this place, what so obviously is the
meaning intended by the historian; namely, 'genius, ability,
distinguished talents.''' Indeed, the whole tenor of the passage, as well
as the scope of the context, leaves no room to doubt the fact. The main
objects of comparison, throughout the three first sections of this
Proemium, or introductory discourse, are not vice and virtue, but body
and mind; a listless indolence, and a vigorous, honorable activity.
On this account it is pretty evident, that by \emph{virtus} Sallust could
never mean the [Greek \emph{aretae}], 'virtue or moral worth,' but that he
had in his eye the well-known interpretation of Varro, who considers it
\emph{ut viri vis} (De Ling. Lat. iv.), as denoting the useful energy which
ennobles a man, and should chiefly distinguish him among his
fellow-creatures. In order to be convinced of the justice of this
rendering, we need only turn to another passage of our author, in the
second section of the Proemium to the Jugurthine War, where the same
train of thought is again pursued, although he gives it somewhat a
different turn in the piece last mentioned. The object, notwithstanding,
of both these dissertations is to illustrate, in a striking manner, the
pre-eminence of the mind over extrinsic advantage, or bodily endowments,
and to show that it is by genius alone that we may aspire to a reputation
which shall never die. ``\emph{Igitur praeclara facies, magnae divitiae,
adhuc vis corporis, et alia hujusmondi omnia, brevi dilabuntur: at
ingenii egregia facinora, sicut anima, immortalia sunt}''.}

Yet it was long a subject of dispute among mankind, whether military
efforts were more advanced by strength of body, or by force of
intellect. For, in affairs of war, it is necessary to plan before
beginning to act,\footnote{It is necessary to plan before beginning to act—\emph{Priusquam
incipias, consulto—opus est}. Most translators have rendered
\emph{consulto} ``deliberation,'' or something equivalent; but it is
\emph{planning} or \emph{contrivance} that is signified. Demosthenes, in his
Oration \emph{de Pace}, reproaches the Athenians with acting without any
settled plan: [Greek: \emph{Oi men gar alloi puntes anthropoi pro ton
pragmatonheiothasi chraesthai to Bouleuesthai, umeis oude meta ta
pragmata}.]} and, after planning, to act with promptitude
and vigor.\footnote{To act with promptitude and vigor—\emph{Maturè facto opus est}.
``Maturè facto'' seems to include the notions both of promptitude and
vigor, of force as well as speed; for what would be the use of acting
expeditiously, unless expedition be attended with power and effect?} Thus, each\footnote{Each—\emph{Utrumque}. The corporeal and mental faculties.} being insufficient of itself, the one
requires the assistance of the other.\footnote{The one requires the assistance of the other—\emph{Alterum
alterius auxilio eget}. ``\emph{Eget},'' says Cortius, ``is the reading of all
the MSS.'' \emph{Veget}, which Havercamp and some others have adopted, was
the conjecture of Palmerius, on account of \emph{indigens} occurring in the
same sentence. But \emph{eget} agrees far better with \emph{consulto et—maturè
facto opus est}, in the preceding sentence.}

II. In early times, accordingly, kings (for that was the first title
of sovereignty in the world) applied themselves in different ways;\footnote{II. Applied themselves in different ways—\emph{Diversi}. ``Modo
et instituto diverso, diversa sequentes.'' \emph{Cortius}.}
some exercised the mind, others the body. At that period, however,\footnote{At that period, however—\emph{Et jam tum}. ``Tunc temporis
\emph{praecisè}, at that time \emph{precisely}, which is the force of the
particle \emph{jam}. as donatus shows. I have therefore written \emph{et jam}
separately. Virg. Aen. vii. 737. Late \emph{jam tum} ditione premebat
Sarrastes populos.'' \emph{Cortius}.}
the life of man was passed without covetousness;\footnote{Without covetousness—Sine cupiditate\_. ``As in the famous
golden age. See Tacit. Ann. iii. 28.'' \emph{Cortius}. See also Ovid. Met.
i. 80, \emph{seq}. But ``such times were never,'' as Cowper says.} every one was
satisfied with his own. But after Cyrus in Asia\footnote{But after Cyrus in Asia, etc.—\emph{Postea verò quàm in, Asiâ
Cyrus}, etc. Sallust writes as if he had supposed that kings were more
moderate before the time of Cyrus. But this can hardly have been the
case. ``The Romans,'' says De Brosses, whose words I abridge, ``though
not learned in antiquity, could not have been ignorant that there were
great conquerors before Cyrus; as Ninus and Sesostris. But as their
reigns belonged rather to the fabulous ages, Sallust, in entering upon
a serious history, wished to confine himself to what was certain, and
went no further back than the records of Herodotus and Thucydides.''
Ninus, says Justin. i. 1, was the first to change, through inordinate
ambition, the \emph{veterem et quasi avitum gentibus morem}, that is, to
break through the settled restraints of law and order. Gerlach agrees
in opinion with De Brosses.} and the
Lacedaemonians and Athenians in Greece, began to subjugate cities and
nations, to deem the lust of dominion a reason for war, and to imagine
the greatest glory to be in the most extensive empire, it was then at
length discovered, by proof and experience,\footnote{Proof and experience—\emph{Periculo atque negotiis}. Gronovius
rightly interprets \emph{periculo} ``experiundo, experimentis,'' by
experiment or trial. Cortius takes \emph{periculo atque negotiis} for
\emph{periculosis negotiis}, by hendyadys; but to this figure, as Kritzius
remarks, we ought but sparingly to have recourse. It is better, he
adds, to take the words in their ordinary signification, understanding
by \emph{negotia} ``res graviores.'' Bernouf judiciously explains \emph{negotiis}
by ``ipsa negotiorum tractatione,'' \emph{i. e.} by the management of affairs,
or by experience in affairs. Dureau Delamalle, the French translator,
has ``l'expérience et la pratique.'' Mair has ``trial and experience.''
which, I believe, faithfully expresses Sallust's meaning. Rose gives
only ``experience'' for both words.} that mental power has
the greatest effect in military operations. And, indeed,\footnote{And, indeed, if the intellectual ability, etc.—\emph{Quod
si—animi virtus}, etc. ``Quod si'' can not here be rendered \emph{but if;}
it is rather equivalent to \emph{quapropter si}, and might be expressed by
\emph{wherefore if, if therefore, if then, so that if}.} if the
intellectual ability\footnote{Intellectual ability—\emph{Animi virtus}. See the remarks on
\emph{virtus}, above noted.} of kings and magistrates\footnote{Magistrates—\emph{Imperatorum}. ``Understand all who govern
states, whether in war or in peace.'' \emph{Bernouf}. Sallust calls the
consuls \emph{imperatores}, c. 6.} were exerted to
the same degree in peace as in war, human affairs would be more
orderly and settled, and you would not see governments shifted from
hand to hand,\footnote{Governments shifted from hand to hand—\emph{aliud aliò ferri}.
Evidently alluding to changes in government.} and things universally changed and confused. For
dominion is easily secured by those qualities by which it was at first
obtained. But when sloth has introduced itself in the place of industry,
and covetousness and pride in that of moderation and equity, the fortune
of a state is altered together with its morals; and thus authority is
always transferred from the less to the more deserving.\footnote{Less to the more deserving—\emph{Ad optimum quemque à minus
bono}. ``From the less good to the best.''}

Even in agriculture,\footnote{Even in agriculture, etc.—\emph{Quae homines arant, navigant,
aedificant, virtuti omnia parent}. Literally, \emph{what men plow, sail},
etc. Sallust's meaning is, that agriculture, navigation, and
architecture, though they may seem to be effected by mere bodily
exertion, are as much the result of mental power as the highest of
human pursuits.} in navigation, and in architecture, whatever
man performs owns the dominion of intellect. Yet many human beings,
resigned to sensuality and indolence, un-instructed and unimproved,
have passed through life like travellers in a strange country\footnote{Like travelers in a strange country—\emph{Sicuti peregrinantes}.
``Vivere nesciunt; igitur in vita quasi hospites sunt:'' they know not
how to use life, and are therefore, as it were, strangers in it.
\emph{Dietsch}. ``\emph{Peregrinantes}, qui, qua transeunt, nullum sui vestigium
relinquunt;'' they are as travelers who do nothing to leave any trace
of their course. Pappaur.}; to
whom, certainly, contrary to the intention of nature, the body was a
gratification, and the mind a burden. Of these I hold the life and
death in equal estimation\footnote{Of these I hold the life and death in equal estimation—\emph{Eorum
ego vitam mortemque juxta aestimo}. I count them of the same value dead
as alive, for they are honored in the one state as much as in the other.
``Those who are devoted to the gratification of their appetites,'' as
Sallust says, ``let us regard as inferior animals, not as men; and some,
indeed, not as living, but as dead animals.'' Seneca, Ep. lx.}; for silence is maintained concerning
both. But he only, indeed, seems to me to live, and to enjoy life,
who, intent upon some employment, seeks reputation from some ennobling
enterprise, or honorable pursuit.

But in the great abundance of occupations, nature points out different
paths to different individuals. III. To act well for the Commonwealth
is noble, and even to speak well for it is not without merit\footnote{III. Not without merit—\emph{Haud absurdum}. I have borrowed
this expression from Rose, to whom Muretus furnished ``sua laude non
caret.'' ``The word \emph{absurdus} is often used by the Latins as an epithet
for sounds disagreeable to the ear; but at length it came to be
applied to any action unbecoming a rational being.'' \emph{Kunhardt}.}. Both
in peace and in war it is possible to obtain celebrity; many who have
acted, and many who have recorded the actions of others, receive their
tribute of praise. And to me, assuredly, though by no means equal
glory attends the narrator and the performer of illustrious deeds, it
yet seems in the highest degree difficult to write the history of
great transactions; first, because deeds must be adequately
represented\footnote{Deeds must be adequately represented, etc.—\emph{Facta dictis
sunt exaequanda}. Most translators have regarded these words as
signifying \emph{that the subject must be equaled by the style}. But it is
not of mere style that Sallust is speaking. ``He means that the matter
must be so represented by the words, that honorable actions may not be
too much praised, and that dishonorable actions may not be too much
blamed; and that the reader may at once understand what was done and
how it was done.'' \emph{Kunhardt}.} by words; and next, because most readers consider that
whatever errors you mention with censure, are mentioned through
malevolence and envy; while, when you speak of the great virtue and
glory of eminent men, every one hears with acquiescence\footnote{Every one hears with acquiescence, etc.—\emph{Quae sibi—aequo
animo accipit}, etc. This is taken from Thucydides, ii. 35. ``For
praises spoken of others are only endured so far as each one thinks
that he is himself also capable of doing any of the things he hears;
but that which exceeds their own capacity, men at once envy and
disbelieve.'' Dale's Translation: Bohn's Classical Library.} only that
which he himself thinks easy to be performed; all beyond his own
conception he regards as fictitious and incredible\footnote{Regards as fictitious and incredible—\emph{Veluti ficta, pro
falsis ducit. Ducit pro falsis}, he considers as false or incredible,
\emph{veluti ficta}, as if invented.}.

I myself, however, when a young man\footnote{When a young man—\emph{Adolescentulus}. ``It is generally admitted
that all were called \emph{adolescentes} by the Romans, who were between
the fifteenth or seventeenth year of their age and the fortieth.
The diminutive is used in the same sense, but with a view to contrast
more strongly the ardor and spirit of youth with the moderation,
prudence, and experience of age. So Caesar is called \emph{adolescentulus},
in c. 49, at a time when he was in his thirty-third year.'' \emph{Dietsch}.
And Cicero, referring to the time of his consulship, says, \emph{Defendi
rempublicam adolescens}, Philipp. ii. 46.}, was at first led by
inclination, like most others, to engage in political affairs\footnote{To engage in political affairs—\emph{Ad rempublicam}. ``In the phrase
of Cornelius Nepos, \emph{honoribus operam dedi}, I sought to obtain some
share in the management of the Republic. All public matters were
comprehended under the term \emph{Respublica}.'' \emph{Cortius}.}; but
in that pursuit many circumstances were unfavorable to me; for,
instead of modesty, temperance, and integrity\footnote{Integrity—\emph{Virtute}. Cortius rightly explains this word as
meaning\_justice, equity\_, and all other virtues necessary in those who
manage the affairs of a state. Observe that it is here opposed to
\emph{avaritia}, not, as some critics would have it, to \emph{largitio}.}, there prevailed
shamelessness, corruption, and rapacity. And although my mind,
inexperienced in dishonest practices, detested these vices, yet, in
the midst of so great corruption, my tender age was insnared and
infected\footnote{Was ensnared and infected—\emph{Corrupta, tenebatur}. As
\emph{obsessus tenetur}, Jug., c. 24.} by ambition; and, though I shrunk from the vicious
principles of those around me, yet the same eagerness for honors, the
same obloquy and jealousy\footnote{The same eagerness for honors, the same obloquy and
jealousy, etc.—\emph{Honoris cupido eadem quae caeteros, fama atque
invidia vexabat}. I follow the interpretation of Cortius: ``Me vexabat
honoris cupido, et vexabat \emph{propterea} etiam eadem, quae caeteros,
fama atqua invidia.'' He adds, from a gloss in the Guelferbytan MS.,
that it is a \emph{zeugma}. ``\emph{Fama atque invidia},'' says Gronovius, ``is
[Greek: \emph{en dia duoin}], for \emph{invidiosa et maligna fama}.'' Bernouf,
with Zanchius and others, read \emph{fama atque invidia} in the ablative
case; and the Bipont edition has \emph{eadem qua—fama, etc.}; but the
method of Cortius is, to me, by far the most straightforward and
satisfactory. Sallust, observes De Brosses, in his note on this
passage, wrote the account of Catiline's conspiracy shortly after his
expulsion from the Senate, and wishes to make it appear that he
suffered from calumny on the occasion; though he took no trouble, in
the subsequent part of his life, to put such calumny to silence.}, which disquieted others, disquieted
myself.

IV. When, therefore, my mind had rest from its numerous troubles and
trials, and I had determined to pass the remainder of my days
unconnected with public life, it was not my intention to waste my
valuable leisure in indolence and inactivity, or, engaging in servile
occupations, to spend my time in agriculture or hunting\footnote{IV. Servile occupations—agriculture or hunting—\emph{Agrum
colendo, aut venando, servilibus officiis intentum}. By calling
agriculture and hunting \emph{servilia officia}, Sallust intends, as is
remarked by Graswinckelius, little more than was expressed in the
saying of Julian the emperor, \emph{Turpe est sapienti, cum habeat animum,
captare laudes ex corpore}. ``Ita ergo,'' adds the commentator,
``agricultura et venatio servilio officia sunt, quum in solo consistant
corporis usu, animum, vero nec meliorem nec prudentiorem reddant. Quia
labor in se certe est illiberalis, ei praesertim cui facultas sit ad
meliora.'' Symmachus (1 v. Ep. 66) and some others, whose remarks the
reader may see in Havercamp, think that Sallust might have spoken of
hunting and agriculture with more respect, and accuse him of not
remembering, with sufficient veneration, the kings and princes that
have amused themselves in hunting, and such illustrious plowmen as
Curius and Cincinnatus. Sallust, however, is sufficiently defended
from censure by the Abbé Thyvon, in a dissertation much longer than
the subject deserves, and much longer than most readers are willing to
peruse.}; but,
returning to those studies\footnote{Returning to those studies, etc.—\emph{A quo incepto studio me
ambitio mala detinuerat, eodem regressus}. ``The study, namely, of
writing history, to which he signifies that he was attached in c. 3.''
\emph{Cortius}.} from which, at their commencement, a
corrupt ambition had allured me, I determined to write, in detached
portions\footnote{In detached portions—\emph{Carptim}. ``Plin. Ep. viii., 47:
Respondebis non posse perinde \emph{carptim}, ut \emph{contexta} placere: et vi.
22: Egit \emph{carptim} et [Greek: \emph{kata kephulaia}],'' \emph{Dietsch}.}, the transactions of the Roman people, as any occurrence
should seem worthy of mention; an undertaking to which I was the
rather inclined, as my mind was uninfluenced by hope, fear, or
political partisanship. I shall accordingly give a brief account, with
as much truth as I can, of the Conspiracy of Catiline; for I think it
an enterprise eminently deserving of record, from the unusual nature
both of its guilt and of its perils. But before I enter upon my
narrative, I must give a short description of the character of the
man.

V. Lucius Catiline was a man of noble birth\footnote{V. Of noble birth—\emph{Nobili genere natus}. His three names
were Lucius \emph{Sergius} Catilina, he being of the family of the Sergii,
for whose antiquity Virgil is responsible, Aen. v. 121: \emph{Sergestusque,
domus tenet a quo Sergia nomen}. And Juvenal says, Sat. viii. 321:
\emph{Quid, Catilino, tuis natalibus atque Cethegi Inveniet quisquam
sublimius?} His great grandfather, L. Sergius Silus, had eminently
distinguished himself by his services in the second Punic war. See
Plin. Hist. Nat. vii. 29. ``Catiline was born A.U.C. 647, A.C. 107.''
\emph{Dietsch}. Ammianus Marcellinus (lib. xxv.) says that he was the last
of the Sergii.}, and of eminent mental
and personal endowments; but of a vicious and depraved disposition.
His delight, from his youth, had been civil commotions, bloodshed,
robbery, and sedition\footnote{\emph{Sedition—Discordia civilis}.}; and in such scenes he had spent his early
years.\footnote{And in such scenes he had spent his early years—\emph{Ibique
juventutem suam exercuit}. ``It is to be observed that the Roman
writers often used an adverb, where we, of modern times, should
express ourselves more specifically by using a noun.'' \emph{Dietsch} on c.
3, \emph{ibique multa mihi advorsa fuere}. \emph{Juventus} properly signified
the time between thirty and forty-five years of age; \emph{adolescentia}
that between fifteen and thirty. But this distinction was not always
accurately observed. Catiline had taken an active part in supporting
Sylla, and in carrying into execution his cruel proscriptions and
mandates. ``Quis erat hujus (Syllae) imperii minister? Quis nisi
Catilina jam in omne facinus manus exercens?'' Sen. de Ira, iii. 18.} His constitution could endure hunger, want of sleep, and
cold, to a degree surpassing belief. His mind was daring, subtle, and
versatile, capable of pretending or dissembling whatever he wished.\footnote{Capable of pretending or dissembling whatever he wished
—\emph{Cujuslibet, rei simulator ac dissimulator}. ``Dissimulation is
the negative, when a man lets fall signs and arguments, that he is not
that he is; simulation is the affirmative, when a man industriously
and expressly feigns and pretends to be that he is not.'' Bacon,
Essay vi.}
He was covetous of other men's property, and prodigal of his own. He
had abundance of eloquence,\footnote{Abundance of eloquence—\emph{Satis eloquentiae}. Cortius reads
\emph{loquentiae} ``\emph{Loquentia} is a certain facility of speech not
necessarily attended with sound sense; called by the Greeks [Greek:
\emph{lalia}].'' \emph{Bernouf}. ``Julius Candidus used excellently to observe
that \emph{eloquentia} was one thing, and \emph{loquentia} another; for
eloquence is given to few, but what Candidus called \emph{loquentia}, or
fluency of speech, is the talent of many, and especially of the most
impudent.'' Plin. Ep. v. 20. But \emph{eloquentiae} is the reading of most
of the MSS., and \emph{loquentiae}, if Aulus Gellius (i. 15) was rightly
informed, was a correction of Valerius Probus, the grammarian, who
said that Sallust \emph{must} have written so, as \emph{eloquentiae} could not
agree with \emph{sapientiae parum}. This opinion of Probus, the grammarian,
who said that Sallust \emph{must} have written so, as \emph{eloquentiae} could
not agree with \emph{sapientiae parum}. This opinion of Probus, however,
may be questioned. May not Sallust have written \emph{eloquentiae}, with
the intention of signifying that Catiline had abundance of eloquence
to work on the minds of others, though he wanted prudence to regulate
his own conduct? Have there not been other men of whom the same may be
said, as Mirabeau, for example? The speeches that Sallust puts into
Catiline's mouth (c. 20, 58) are surely to be characterized rather as
\emph{eloquentia}, than \emph{loquentia}. On the whole, and especially from the
concurrence of MSS., I prefer to read \emph{eloquentiae}, with the more
recent editors, Gerlach, Kritz and Dietsch.} though but little wisdom. His
insatiable ambition was always pursuing objects extravagant, romantic,
and unattainable.

Since the time of Sylla's dictatorship,\footnote{Since the time of Sylla's dictatorship—\emph{Post dominationem
Lucii Syllae}. ``The meaning is not the same as if it were \emph{finitâ
dominatione} but is the same as \emph{ab eo tempore quo dominari caeperat}.
In French, therefore, \emph{post} should be rendered by \emph{depuis}, not, as
it is commonly translated, \emph{après}.'' \emph{Bernouf}. As \emph{dictator} was the
title that Sylla assumed, I have translated \emph{dominatio}, ``dictatorship''.
Rose, Gordon, and others, render it ``usurpation''.} a strong desire of seizing
the government possessed him, nor did he at all care, provided that he
secured power\footnote{Power—\emph{Regnum}. Chief authority, rule, dominion.} for himself, by what means he might arrive at it.
His violent spirit was daily more and more hurried on by the
diminution of his patrimony, and by his consciousness of guilt; both
which evils he had increased by those practices which I have mentioned
above. The corrupt morals of the state, too, which extravagance and
selfishness, pernicious and contending vices, rendered thoroughly
depraved,\footnote{Rendered thoroughly depraved—\emph{Vexabant}. ``Corrumpere et
pessundare studebant.'' \emph{Bernouf}. \emph{Quos vexabant}, be it observed,
refers to \emph{mores}, as Gerlach and Kritz interpret, not to \emph{cives}
understood in \emph{civitatis}, which is the evidently erroneous method of
Cortius.} furnished him with additional incentives to action.

Since the occasion has thus brought public morals under my notice, the
subject itself seems to call upon me to look back, and briefly to
describe the conduct of our ancestors\footnote{Conduct of our ancestors—\emph{Instituta majorum}. The principles
adopted by our ancestors, with regard both to their own conduct, and
to the management of the state. That this is the meaning, is evident
from the following account.} in peace and war; how they
managed the state, and how powerful they left it; and how, by gradual
alteration, it became, from being the most virtuous, the most vicious
and depraved.

VI. Of the city of Rome, as I understand,\footnote{VI. As I understand—\emph{Sicuti ego accepi}. ``By these words he
plainly shows that nothing certain was known about the origin of Rome.
The reader may consult Livy, lib. i.; Justin, lib. xliii.; and Dionys.
Halicar., lib.i.; all of whom attribute its rise to the Trojans.''
\emph{Bernouf}.} the founders and
earliest inhabitants were the Trojans, who, under the conduct of
Aeneas, were wandering about as exiles from their country, without any
settled abode; and with these were joined the Aborigines,\footnote{Aborigines—\emph{Aborigines}. The original inhabitants of Italy;
the same as \emph{indigenae}, or the [Greek: \emph{Autochthones}].} a savage
race of men, without laws or government, free, and owning no control.
How easily these two tribes, though of different origin, dissimilar
language, and opposite habits of life, formed a union when they met
within the same walls, is almost incredible.\footnote{: Almost incredible—\emph{Incredibile memoratu}. ``Non credi potest,
si memoratur; superat omnem fidem.'' \emph{Pappaur}. Yet that which
actually happened, can not be absolutely incredible; and I have,
therefore, inserted \emph{almost}.} But when their state,
from an accession of population and territory, and an improved
condition of morals, showed itself tolerably flourishing and powerful,
envy, as is generally the case in human affairs, was the consequence
of its prosperity. The neighboring kings and people, accordingly,
began to assail them in war, while a few only of their friends came to
their support; for the rest, struck with alarm, shrunk from sharing
their dangers. But the Romans, active at home and in the field,
prepared with alacrity for their defense.\footnote{Prepared with alacrity for there defense—\emph{Festinare, parare}.
``Made haste, prepared.'' ``\emph{Intenti ut festinanter pararent} ea, quae
defensioni aut bello usui essent.'' \emph{Pappaur}.} They encouraged one
another, and hurried to meet the enemy. They protected, with their
arms, their liberty, their country, and their homes. And when they had
at length repelled danger by valor, they lent assistance to their
allies and supporters, and procured friendships rather by
bestowing\footnote{Procured friendships rather by bestowing, etc;—\emph{Magisque
dandis, quam accipiundis beneficiis amicitias parabant}. Thucyd. ii.,
40: [Greek: \emph{Ou paschontes eu, alla drontes, ktometha tous philous}]} favors than by receiving them.

They had a government regulated by laws. The denomination of their
government was monarchy. Chosen men, whose bodies might be enfeebled
by years, but whose minds were vigorous in understanding, formed the
council of the state; and these, whether from their age, or from the
similarity of their duty, were called FATHERS.\footnote{FATHERS—PATRES. ``(Romulus) appointed that the direction of
the state should be in the hands of the old men, who, from their
authority, were called \emph{Fathers}; from their age, \emph{Senatus}.'' Florus,
i. 1. \emph{Senatus} from \emph{senex}. ``\emph{Patres} ab honore—appellati.''
\emph{Livy}.} But afterward, when
the monarchical power, which had been originally established for the
protection of liberty, and for the promotion of the public interest,
had degenerated into tyranny and oppression, they changed their plan,
and appointed two magistrates,\footnote{Two magistrates—\emph{Binos imperatores}. The two consuls. They
were more properly called \emph{imperatores} at first, when the law, which
settled their power, said ``\emph{Regio imperio} duo sunto'' (Cic. de Legg.
iii. 4), than afterward, when the people and tribunes had made
encroachments on their authority.} with power only annual; for they
conceived that, by this method, the human mind would be least likely
to grow overbearing for want of control.

VII. At this period every citizen began to seek distinction, and to
display his talents with greater freedom; for, with princes, the
meritorious are greater objects of suspicion than the undeserving, and
to them the worth of others is a source of alarm. But when liberty was
secured, it is almost incredible\footnote{VII. Almost incredible—\emph{Incredibile memoratu}. See above, c. 6.} how much the state strengthened
itself in a short space of time, so strong a passion for distinction
had pervaded it. Now, for the first time, the youth, as soon as they
were able to bear the toil of war,\footnote{Able to bear the toils of war—\emph{Laboris ac belli patiens}.
As by \emph{laboris} the labor of war is evidently intended, I have thought
it better to render the words in this manner. The reading is Cortius'.
Havercamp and others have ``simul \emph{ac belli} patiens erat, in castris
\emph{per laborem usu} militiam discebat;'' but \emph{per laborem usu} is
assuredly not the hand of Sallust.} acquired military skill by
actual service in the camp, and took pleasure rather in splendid arms
and military steeds than in the society of mistresses and convivial
indulgence. To such men no toil was unusual, no place was difficult or
inaccessible, no armed enemy was formidable; their valor had overcome
every thing. But among themselves the grand rivalry was for glory;
each sought to be first to wound an enemy, to scale a wall, and to be
noticed while performing such an exploit. Distinction such as this
they regarded as wealth, honor, and true nobility.\footnote{Honor and true nobility—\emph{Bonam famam magnamque nobilitatem}.} They were
covetous of praise, but liberal of money; they desired competent
riches but boundless glory. I could mention, but that the account
would draw me too far from my subject, places in which the Roman
people, with a small body of men, routed vast armies of the enemy; and
cities, which, though fortified by nature, they carried by assault.

VIII. But, assuredly, Fortune rules in all things. She makes every
thing famous or obscure rather from caprice than in conformity with
truth. The exploits of the Athenians, as far as I can judge, were very
great and glorious,\footnote{VIII. Very great and glorious—\emph{Satis amplae magnificaeque}.
In speaking of this amplification of the Athenian exploits, he
alludes, as Colerus observes, to the histories of Thucydides,
Xenophen, and perhaps Herodotus; not, as Wasse seems to imagine,
to the representations of the poets.} something inferior to what fame has represented
them. But because writers of great talent flourished there, the actions
of the Athenians are celebrated over the world as the most splendid
achievements. Thus, the merit of those who have acted is estimated at
the highest point to which illustrious intellects could exalt it in
their writings.

But among the Romans there was never any such abundance of writers;\footnote{There was never any such abundance of writers—\emph{Nunquam ea
copia fuit}. I follow Kuhnhardt, who thinks \emph{copia} equivalent to
\emph{multitudo}. Others render it \emph{advantage}, or something similar;
which seems less applicable to the passage. Compare c.28:
\emph{Latrones}—\emph{quorum}—magna copia \emph{erat}.}
for, with them, the most able men were the most actively employed. No
one exercised the mind independently of the body: every man of ability
chose to act rather than narrate,\footnote{Chose to act rather than narrate—``For,'' as Cicero says,
``neither among those who are engaged in establishing a state, nor
among those carrying on wars, nor among those who are curbed and
restrained under the rule of kings, is the desire of distinction in
eloquence wont to arise.'' \emph{Graswinckelius}.} and was more desirous that his
own merits should be celebrated by others, than that he himself should
record theirs.

IX. Good morals, accordingly, were cultivated in the city and in the
camp. There was the greatest possible concord, and the least possible
avarice. Justice and probity prevailed among the citizens, not more
from the influence of the laws than from natural inclination. They
displayed animosity, enmity, and resentment only against the enemy.
Citizens contended with citizens in nothing but honor. They were
magnificent in their religious services, frugal in their families,
and steady in their friendships.

By these two virtues, intrepidity in war, and equity in peace, they
maintained themselves and their state. Of their exercise of which
virtues, I consider these as the greatest proofs; that, in war,
punishment was oftener inflicted on those who attacked an enemy
contrary to orders, and who, when commanded to retreat, retired too
slowly from the contest, than on those who had dared to desert their
standards, or, when pressed by the enemy,\footnote{IX. Pressed by the enemy—\emph{Pulsi}. In the words \emph{pulsi loco
cedere ausi erant}, \emph{loco} is to be joined, as Dietsch observes, with
cedere\_, not, as Kritzius puts it, with \emph{pulsi}. ``To retreat,'' adds
Dietsch, ``is disgraceful only to those \emph{qui ab hostibus se pelli
patiantur}, who suffer themselves to be \emph{repulsed by the enemy}.''} to abandon their posts;
and that, in peace, they governed more by conferring benefits than by
exciting terror, and, when they received an injury, chose rather to
pardon than to revenge it.

X. But when, by perseverance and integrity, the republic had increased
its power; when mighty princes had been vanquished in war;\footnote{X. When mighty princes had been vanquished in war—Perses,
Antiochus, Mithridates, Tigranes, and others.} when
barbarous tribes and populous states had been reduced to subjection;
when Carthage, the rival of Rome's dominion, had been utterly
destroyed, and sea and land lay every where open to her sway, Fortune
then began to exercise her tyranny, and to introduce universal
innovation. To those who had easily endured toils, dangers, and
doubtful and difficult circumstances, ease and wealth, the objects of
desire to others, became a burden and a trouble. At first the love of
money, and then that of power, began to prevail, and these became, as
it were, the sources of every evil. For avarice subverted honesty,
integrity, and other honorable principles, and, in their stead,
inculcated pride, inhumanity, contempt of religion, and general
venality. Ambition prompted many to become deceitful; to keep one
thing concealed in the breast, and another ready on the tongue;\footnote{To keep one thing concealed in the breast, and another ready
on the tongue—\_Aliud clausum in pectore, aliud in lingua promptum,} to
estimate friendships and enmities, not by their worth, but according
to interest; and to carry rather a specious countenance than an honest
heart. These vices at first advanced but slowly, and were sometimes
restrained by correction; but afterward, when their infection had
spread like a pestilence, the state was entirely changed, and the
government, from being the most equitable and praiseworthy, became
rapacious and insupportable.

XI. At first, however, it was ambition, rather than avarice,\footnote{XI. At first, however, it was ambition, rather than avarice,
etc.—\emph{Sed primo magis ambitio quam avaritia animos hominum
exercebat}. Sallust has been accused of having made, in this passage,
an assertion at variance with what he had said before (c.10), \emph{Igitur
primo pecuniae, deinde imperii cupido, crevit}, and it will be hard to
prove that the accusation is not just. Sir H. Steuart, indeed,
endeavors to reconcile the passages by giving them the following
``meaning'', which, he says, ``seems perfectly evident'': ``Although
avarice was the first to make its appearance at Rome, yet, after both
had had existence, it was ambition that, of the two vices, laid the
stronger hold on the minds of men, and more speedily grew to an
inordinate height''. To me, however, it ``seems perfectly evident'' that
the Latin can be made to yield no such ``meaning''. ``How these passages
agree,'' says Rupertus, ``I do not understand: unless we suppose that
Sallust, by the word \emph{primo}, does not always signify order''.} that
influenced the minds of men; a vice which approaches nearer to virtue
than the other. For of glory, honor, and power, the worthy is as
desirous as the worthless; but the one pursues them by just methods;
the other, being destitute of honorable qualities, works with fraud
and deceit. But avarice has merely money for its object, which no wise
man has ever immoderately desired. It is a vice which, as if imbued
with deadly poison, enervates whatever is manly in body or mind.\footnote{Enervates whatever is manly in body or mind—\emph{Corpus
virilemque animum effaeminat}. That avarice weakens the mind, is
generally admitted. But how does it weaken the body? The most
satisfactory answer to this question is, in the opinion of Aulus
Gellius (iii. 1), that those who are intent on getting riches devote
themselves to sedentary pursuits, as those of usurers and
money-changers, neglecting all such exercises and employments as
strengthen the body. There is, however, another explanation by
Valerius Probus, given in the same chapter of Aulus Gellius, which
perhaps is the true one; namely, that Sallust, by \emph{body and mind},
intended merely to signify \emph{the whole man}.}
It is always unbounded and insatiable, and is abated neither by
abundance nor by want.

But after Lucius Sylla, having recovered the government\footnote{Having recovered the government—\emph{Receptâ republicâ}. Having
wrested it from the hands of Marius and his party.} by force
of arms, proceeded, after a fair commencement, to a pernicious
termination, all became robbers and plunderers;\footnote{All became robbers and plunderers—\emph{Rapere omnes, trahere}.
He means that there was a general indulgence in plunder among Sylla's
party, and among all who, in whatever character, could profit by
supporting it. Thus he says immediately afterward, ``neque modum neque
modestiam \emph{victores} habere.''} some set their
affections on houses, others on lands; his victorious troops knew
neither restraint nor moderation, but inflicted on the citizens
disgraceful and inhuman outrages. Their rapacity was increased by the
circumstance that Sylla, in order to secure the attachment of the
forces which he had commanded in Asia,\footnote{which he had commanded in Asia—\emph{Quem in Asiâ dustaverat}. I
have here deserted Cortius, who gives \emph{in Asiam}, ``into Asia,'' but this,
as Bernouf justly observes, is incompatible with the frequentative verb
\emph{ductaverat}.} had treated them, contrary
to the practice of our ancestors, with extraordinary indulgence, and
exemption from discipline; and pleasant and luxurious quarters had
easily, during seasons of idleness, enervated the minds of the
soldiery. Then the armies of the Roman people first became habituated
to licentiousness and intemperance, and began to admire statues,
pictures, and sculptured vases; to seize such objects alike in public
edifices and private dwellings;\footnote{in public edifices and private dwellings—\emph{Privatim ac
publice}. I have translated this according to the notion of Burnouf.
Others, as Dietsch and Pappaur, consider \emph{privatim} as signifying
\emph{each on his own account}, and \emph{publice}, \emph{in the name of the
Republic}.} to spoil temples; and to cast off
respect for every thing, sacred and profane. Such troops, accordingly,
when once they obtained the mastery, left nothing to be vanquished.
Success unsettles the principles even of the wise, and scarcely would
those of debauched habits use victory with moderation.

XII. When wealth was once considered an honor, and glory, authority,
and power attended on it, virtue lost her influence, poverty was
thought a disgrace, and a life of innocence was regarded as a life of
ill-nature.\footnote{XII. A life of innocence was regarded as a life of ill-nature
—\emph{Innocentia pro malivolentiâ duci caepit}. ``Whoever continued honest
and upright, was considered by the unprincipled around him as their
enemy; for a good man among the bad can never be regarded as of their
party.'' \emph{Bernouf}.} From the influence of riches, accordingly, luxury,
avarice, and pride prevailed among the youth; they grew at once
rapacious and prodigal; they undervalued what was their own, and
coveted what was another's; they set at naught modesty and continence;
they lost all distinction between sacred and profane, and threw off
all consideration and self-restraint.

It furnishes much matter for reflection,\footnote{It furnishes much matter for reflection—\emph{Operae pretium est}.} after viewing our modern
mansions and villas extended to the size of cities, to contemplate the
temples which our ancestors, a most devout race of men, erected to the
gods. But our forefathers adorned the fanes of the deities with devotion,
and their homes with their own glory, and took nothing from those whom
they conquered but the power of doing harm; their descendants, on the
contrary, the basest of mankind,\footnote{Basest of mankind—\emph{Ignavissumi mortales}. It is opposed to
\emph{fortissumi viri}, which follows, ``Qui nec fortiter nec bene quidquam
fecere.'' \emph{Cortius}.} have even wrested from their allies,
with the most flagrant injustice, whatever their brave and victorious
ancestors had left to their vanquished enemies; as if the only use of
power were to inflict injury.

XIII. For why should I mention those displays of extravagance, which
can be believed by none but those who have seen them; as that mountains
have been leveled, and seas covered with edifices,\footnote{XIII. Seas covered with edifices—\emph{Maria constructa esse}.} by many private
citizens; men whom I consider to have made a sport of their wealth,\footnote{To have made a sport of their wealth—\emph{Quibus mihi videntur
ledibrio fuisse divitiae}. ``They spent their riches on objects which,
in the judgment of men of sense, are ridiculous and contemptible.''
\emph{Cortius}.}
since they were impatient to squander disreputably what they might have
enjoyed with honor.

But the love of irregular gratification, open debauchery, and all
kinds of luxury,\footnote{Luxury—\emph{Cultûs}. ``Deliciarum in victu\_, luxuries of the table;
for we must be careful not to suppose that apparel is meant.''
\emph{Cortius}.} had spread abroad with no less force. Men forgot
their sex; women threw off all the restraints of modesty. To gratify
appetite, they sought for every kind of production by land and by sea;
they slept before there was any inclination for sleep; they no longer
waited to feel hunger, thirst, cold,\footnote{Cold—\emph{Frigus}. It is mentioned by Cortius that this word is
wanting in one MS.; and the English reader may possibly wish that it
were away altogether. Cortius refers it to cool places built of stone,
sometimes underground, to which the luxurious retired in the hot
weather; and he cites Pliny, Ep., v. 6, who speaks of \emph{crytoporticus},
a gallery from which the sun was excluded, almost as if it were
underground, and which, even in summer was cold nearly to freezing.
He also refers to Ambros., Epist. xii., and Casaubon. Ad Spartian.
Adrian., c. x., p. 87.} or fatigue, but anticipated
them all by luxurious indulgence. Such propensities drove the youth,
when their patrimonies were exhausted, to criminal practices; for
their minds, impregnated with evil habits, could not easily abstain
from gratifying their passions, and were thus the more inordinately
devoted in every way to rapacity and extravagance.

XIV. In so populous and so corrupt a city, Catiline, as it was very
easy to do, kept about him, like a body-guard, crowds of the
unprincipled and desperate. For all those shameless, libertine, and
profligate characters, who had dissipated their patrimonies by
gaming,\footnote{XIV. Gaming—\emph{Manu}. Gerlach, Dietsch, Kritzius, and all the
recent editors, agree to interpret \emph{manu} by \emph{gaming}.} luxury, and sensuality; all who had contracted heavy
debts, to purchase immunity for their crimes or offenses; all
assassins\footnote{Assassins—\emph{Parricidae}. ``Not only he who had killed his father
was called a \emph{parricide}, but he who had killed any man; as is
evident from a law of Numa Pompilius: If any one unlawfully and
knowingly bring a free man to death, let him be \emph{a parricide}.''
\emph{Festus} sub voce \emph{Parrici}.} or sacrilegious persons from every quarter, convicted or
dreading conviction for their evil deeds; all, besides, whom their
tongue or their hand maintained by perjury or civil bloodshed; all, in
fine, whom wickedness, poverty, or a guilty conscience disquieted,
were the associates and intimate friends of Catiline. And if any one,
as yet of unblemished character, fell into his society, he was
presently rendered, by daily intercourse and temptation, similar and
equal to the rest. But it was the young whose acquaintance he chiefly
courted; as their minds, ductile and unsettled from their age, were
easily insnared by his stratagems. For as the passions of each,
according to his years, appeared excited, he furnished mistresses to
some, bought horses and dogs for others, and spared, in a word,
neither his purse nor his character, if he could but make them his
devoted and trustworthy supporters. There were some, I know, who
thought that the youth, who frequented the house of Catiline, were
guilty of crimes against nature; but this report arose rather from
other causes than from any evidence of the fact\footnote{Than from any evidence of the fact—\emph{Quam quod cuiquam id
compertum foret}.}.

XV. Catiline, in his youth, had been guilty of many criminal
connections, with a virgin of noble birth\footnote{XV. With a virgin of noble birth—\emph{Cum virgine nobili}. Who
this was is not known. The name may have been suppressed from respect
to her family. If what is found in a fragment of Cicero be true,
Catiline had an illicit connection with some female, and afterward
married the daughter who was the fruit of the connection: \emph{Ex eodem
stupro et uxorem et filiam invenisti}; Orat. in Tog. Cand. (Oration
xvi., Ernesti's edit.) On which words Asconius Pedianus makes this
comment: ``Dicitur Catilinam adulterium commisisse cum ea quae ci
postea socrus fuit, et ex eo stupro duxisse uxorem, cum filia ejus
esset. Haec Lucceius quoque Catilinae objecit in orationibus, quas in
eum scripsit. Nomina harum mulierum nondum inveni.'' Plutarch, too
(Life of Cicero, c. 10), says that Catiline was accused of having
corrupted his own daughter.}, with a priestess of
Vesta\footnote{With a priestess of Vesta—\emph{Cum sacerdote Vestae}. This
priestess of Vesta was Fabia Terentia, sister to Terentia, Cicero's
wife, whom Sallust, after she was divorced by Cicero, married. Clodius
accused her, but she was acquitted, either because she was thought
innocent, or because the interest of Catulus and others, who exerted
themselves in her favor, procured her acquittal. See Orosius, vi. 3;
the Oration of Cicero, quoted in the preceding note; and Asconius's
commentary on it.}, and of many other offenses of this nature, in defiance
alike of law and religion. At last, when he was smitten with a passion
for Aurelia Orestilla\footnote{Aurelia Orestilla—See c. 35. She was the sister or daughter,
as De Brosses thinks, of Cneius Aurelius Orestis, who had been praetor,
A.U.C. 677.}, in whom no good man, at any time of her
life, commended any thing but her beauty, it is confidently believed
that because she hesitated to marry him, from the dread of having a
grown-up step-son\footnote{A grown-up step-son—\emph{Privignum adulta aetate}. A son of
Catiline's by a former marriage.}, he cleared the house for their nuptials by
putting his son to death. And this crime appears to me to have been
the chief cause of hurrying forward the conspiracy. For his guilty
mind, at peace with neither gods nor men, found no comfort either
waking or sleeping; so effectually did conscience desolate his
tortured spirit.\footnote{Desolate his tortured spirit—\emph{Mentem exciteam vastabat}.
``Conscience desolates the mind, when it deprives it of its proper
power and tranquillity, and introduces into it perpetual disquietude.''
\emph{Cortius}. Many editions have \emph{vexabat}.} His complexion, in consequence, was pale, his
eyes haggard, his walk sometimes quick and sometimes slow, and
distraction was plainly apparent in every feature and look.

XVI. The young men, whom, as I said before, he had enticed to join
him, he initiated, by various methods, in evil practices. From among
them he furnished false witnesses,\footnote{XVI. He furnished false witnesses, etc. \emph{Testis signatoresque
falsos commodare}. ``If any one wanted any such character, Catiline was
ready to supply him from among his troop.''\emph{Bernouf}.} and forgers of signatures; and
he taught them all to regard, with equal unconcern, honor, property,
and danger. At length, when he had stripped them of all character and
shame, he led them to other and greater enormities. If a motive for
crime did not readily occur, he incited them, nevertheless, to
circumvent and murder inoffensive persons\footnote{Inoffensive persons, etc.—\emph{Insontes, sicuti sontes.} Most
translators have rendered these words ``innocent'' and ``guilty,'' terms
which suggest nothing satisfactory to the English reader. The
\emph{insontes} are those who had given Catiline no cause of offens; the
\emph{sontes} those who had in some way incurred his displeasure, or become
objects of his rapacity.} just as if they had
injured him; for, lest their hand or heart should grow torpid for want
of employment, he chose to be gratuitously wicked and cruel.

Depending on such accomplices and adherents, and knowing that the load
of debt was every where great, and that the veterans of Sylla,\footnote{Veterans of Sylla, etc.—Elsewhere called the colonists of
Sylla; men to whom Sylla had given large tracts of land as rewards for
their services, but who, having lived extravagantly, had fallen into
such debt and distress, that, as Cicero said, nothing could relieve
them but the resurrection of Sylla from the dead. Cic. ii. Orat. in
Cat.}
having spent their money too liberally, and remembering their spoils
and former victory, were longing for a civil war, Catiline formed the
design of overthrowing the government. There was no army in Italy;
Pompey was fighting in a distant part of the world;\footnote{Pompey was fighting in a distant part of the world—\emph{In extremis
terris}. Pompey was then conducting the war against Mithridates and
Tigranes, in Pontus and Armenia.} he himself had
great hopes of obtaining the consulship; the senate was wholly off its
guard;\footnote{The senate was wholly off its guard—\emph{Senatus nihil sane intentus}.
The senate was \emph{regardless}, and unsuspicious of any danger.} every thing was quiet and tranquil; and all those
circumstances were exceedingly favorable for Catiline.

XVII. Accordingly, about the beginning of June, in the consulship of
Lucius Caesar\footnote{XVII. Lucius Caesar—He was a relation of Julius Caesar; and his
sister was the wife of M. Antonius, the orator, and mother of Mark
Antony, the triumvir.} and Caius Figulus, he at first addressed each of his
accomplices separately, encouraged some, and sounded others, and
informed them, of his own resources, of the unprepared condition of
the state, and of the great prizes to be expected from the conspiracy.
When he had ascertained, to his satisfaction, all that he required, he
summoned all whose necessities were the most urgent, and whose spirits
were the most daring, to a general conference.

At that meeting there were present, of senatorial rank, Publius
Lentulus Sura,\footnote{Publius Lentulus Sura—He was of the same family with Sylla,
that of the Cornelii. He had filled the office of consul, but his
conduct had been afterward so profligate, that the censors expelled
him from the senate. To enable him to resume his seat, he had
obtained, as a qualification, the office of praetor, which he held at
the time of the conspiracy. He was called Sura, because, when he had
squandered the public money in his quaestorship, and was called to
account by Sylla for his dishonesty, he declined to make any defense,
but said, ``I present you the calf of my leg (\emph{sura});'' alluding to a
custom among boys playing at ball, of inflicting a certain number of
strokes on the leg of an unsuccessful player. Plutarch, Life of
Cicero, c.17.} Publius Autronius,\footnote{Publius Autronius—He had been a companion of Cicero in his
boyhood, and his colleague in the quaestorship. He was banished in the
year after the conspiracy, together with Cassius, Laeca, Vargunteius,
Servius Sylla, and Caius Cornelius, under the Plautian law. \emph{De
Brosses}.} Lucius Cassius Longinus,\footnote{Lucius Cassius Longinus.—He had been a competitor with Cicero
for the consulship. Ascon. Ped., in Cic. Orat. in Tog. Cand. His
corpulence was such that Cassius's fat (\emph{Cassii adeps}) became
proverbial. Cic. Orat. in Catil., iii. 7.}
Caius Cethegus,\footnote{Caius Cethegus—He also was one of the Cornelian family. In the
civil wars, says De Brosses, he had first taken the side of Marius,
and afterward that of Sylla. Both Cicero (Orat. in Catil., ii.7) and
Sallust describe him as fiery and rash.} Publius and Servius Sylla\footnote{Publius and Servius Sylla—These were nephews of Sylla the
dictator. Publius, though present on this occasion, seems not to have
joined in the plot, since, when he was afterward accused of having
been a conspirator, he was defended by Cicero and acquitted. See Cic.
Orat. pro P. Sylla. He was afterward with Caesar in the battle of
Pharsalia. Caes. de B.C., iii. 89.} the sons of
Servius Sylla, Lucius Vargunteius,\footnote{Lucius Vargunteius—``Of him or his family little is known.
He had been, before this period, accused of bribery, and defended by
Hortensius. Cic. pro P. Sylla, c. 2.'' \emph{Bernouf}.} Quintus Annius,\footnote{Quintus Annius—He is thought by De Brosses to have been the same
Annius that cut off the head of M. Antonius the orator, and carried it
to Marius. Plutarch, Vit. Marii, c. 44.} Marcus
Porcius Laeca,\footnote{Marcus Porcius Laeca—He was one of the same \emph{gens} with the
Catones, but of a different family.} Lucius Bestia,\footnote{Lucius Bestia—Of the Calpurnian \emph{gens}. He escaped death
on the discovery of the conspiracy, and was afterward aedile, and
candidate for the praetorship, but was driven into exile for bribery.
Being recalled by Caesar, he became candidate for the consulship, but
was unsuccessful. \emph{De Brosses}.} Quintus Curius;\footnote{Quintus Curius—He was a descendant of M. Curius Dentatus, the
opponent of Pyrrhus. He was so notorious as a gamester and a profligate,
that he was removed from the senate, A.U.C. 683. See c. 23. As he had
been the first to give information of the conspiracy to Cicero, public
honors were decreed him, but he was deprived of them by the influence of
Caesar, whom he had named as one of the conspirators. Sueton. Caes. 17;
Appian. De Bell. Civ., lib. ii.} and, of
the equestrian order, Marcus Fulvius Nobilior,\footnote{M. Fulvius Nobilior—``He was not put to death, but exiled,
A.U.C. 699. Cic. ad Att. iv., 16.'' \emph{Bernouf}.} Lucius
Statilius,\footnote{Lucius Statilius—of him nothing more is known than is told by
Sallust.} Publius Gabinius Capito,\footnote{Publius Gabinius Capito—Cicero, instead of Capito, calls him
Cimber. Orat. in Cat., iii. 3. The family was originally from Gabii.} Caius Cornelius;\footnote{Caius Cornelius—There were two branches of the \emph{gens Cornelia},
one patrician, the other plebeian, from which sprung this conspirator.}
with many from the colonies and municipal towns,\footnote{Municipal towns—\emph{Municipiis}. ``The \emph{municipia} were towns
of which the inhabitants were admitted to the rights of Roman citizens,
but which were allowed to govern themselves by their own laws, and to
choose their own magistrates. See Aul. Gell, xvi. 13; Beaufort, Rep.
Rom., vol. v.'' \emph{Bernouf}.} persons of
consequence in their own localities. There were many others, too,
among the nobility, concerned in the plot, but less openly: men whom
the hope of power, rather than poverty or any other exigence, prompted
to join in the affair. But most of the young men, and especially the
sons of the nobility, favored the schemes of Catiline; they who had
abundant means of living at ease, either splendidly or voluptuously,
preferred uncertainties to certainties, war to peace. There were some,
also, at that time, who believed that Marcus Licinius Crassus\footnote{Marcus Licinius Crassus—The same who, with Pompey and Caesar,
formed the first triumvirate, and who was afterward killed in his
expedition against the Parthians. He had, before the time of the
conspiracy, held the offices of praetor and consul.} was
not unacquainted with the conspiracy; because Cneius Pompey, whom he
hated, was at the head of a large army, and he was willing that the
power of any one whomsoever should raise itself against Pompey's
influence; trusting, at the same time, that if the plot should
succeed, he would easily place himself at the head of the
conspirators.

XVIII. But previously\footnote{XVIII. But previously, etc.—Sallust here makes a digression,
to give an account of a conspiracy that was formed three years before
that of Catiline.} to this period, a small number of persons,
among whom was Catiline, had formed a design against the state: of
which affair I shall here give as accurate account as I am able. Under
the consulship of Lucius Tullus and Marcus Lepidus, Publius Autronius
and Publius Sylla,\footnote{Publius Autronius and Publius Sylla—The same who are mentioned
in the preceding chapter. They were consuls elect, and some editions
have the words \emph{designati consules}, immediately following their names.} having been tried for bribery under the laws
against it,\footnote{Having been tried for bribery under the laws against it
—\emph{Legibus ambitus interrogati}. \emph{Bribery at their election}, is the
meaning of the word \emph{ambitus}, for \emph{ambire}, as Cortius observes, is
\emph{circumeundo favorem et suffragia quaerere}. De Brosses translates the
passage thus: ``Autrone et Sylla, convaincus d'avoir obtenu le consulat
par corruption des suffrages, avaient été punis selon la rigueur de la
loi''. There were several very severe Roman laws against bribery.
Autronius and Sylla were both excluded from the consulship.} had paid the penalty of the offense. Shortly after
Catiline, being brought to trial for extortion,\footnote{For extortion—\emph{Pecuniarum repetundarum}. Catiline had been
praetor in Africa, and, at the expiration of his office, was accused
of extortion by Publius Clodius, on the part of the Africans. He
escaped by bribing the prosecutor and judges.} had been
prevented from standing for the consulship, because he had been unable
to declare himself a candidate within the legitimate number of
days.\footnote{To declare himself a candidate within the legitimate number
of days—\emph{Prohibitus erat consulatum petere, quod intra legitimos
dies profiteri} (se candidatum, says Cortius, citing Suet. Aug. 4)
\emph{nequiverit}. A person could not be a candidate for the consulship,
unless he could declare himself free from accusation within a certain
number of days before the time of holding the \emph{comitia centuriata}.
That number of days was \emph{trinundinum spatium}, that is, the time
occupied by three market-days, \emph{tres nundinae}, with seven days
intervening between the first and second, and between the second and
third; or \emph{seventeen days}. The \emph{nundinae} (from \emph{novem} and \emph{dies})
were held, as it is commonly expressed, every ninth day; whence
Cortius and others considered \emph{trinundinum spatium} to be twenty-seven,
or even thirty days; but this way of reckoning was not that of the
Romans, who made the last day of \emph{the first ennead} to be also the first
day \emph{of the second}. Concerning the \emph{nundinae} see Macrob., Sat. i. 16.
``Muller and Longius most erroneously supposed the \emph{trinundinum} to be
about thirty days; for that it embraced only seventeen days has been
fully shown by Ernesti. Clav. Cic., sub voce; by Scheller in Lex. Ampl.,
p. 11, 669; by Nitschius Antiquitt. Romm. i. p. 623: and by Drachenborch
(cited by Gerlach) ad Liv. iii. 35.'' \emph{Kritzius}.} There was at that time, too, a young patrician of the most
daring spirit, needy and discontented, named Cneius Piso,\footnote{Cneius Piso—Of the Calpurnian gens. Suetonius (Vit. Caes., c. 9)
mentions three authors who related that Crassus and Caesar were both
concerned in this plot; and that, if it had succeeded, Crassus was to
have assumed the dictatorship, and made Caesar his master of the horse.
The conspiracy, as these writers state, failed through the remorse or
irresolution of Crassus.} whom
poverty and vicious principles instigated to disturb the government.
Catiline and Autronius,\footnote{Catiline and Autronius—After these two names, in Havercamp's
and many other editions, follow the words \emph{circiter nonas Decembres},
\emph{i.e.}, about the fifth of December.} having concerted measures with this Piso,
prepared to assassinate the consuls, Lucius Cotta and Lucius Torquatus,
in the Capitol, on the first of January,\footnote{On the first of January—\emph{Kalendis Januariis}. On this day the
consuls were accustomed to enter on their office. The consuls whom
they were going to kill, Cotta and Torquatus, were those who had been
chosen in the place of Antronius and Sylla.} when they, having seized
on the fasces, were to send Piso with an army to take possession of the
two Spains.\footnote{The two Spains—Hither and Thither Spain. \emph{Hispania Citerior}
and \emph{Ulterior}, as they were called by the Romans.} But their design being discovered, they postponed the
assassination to the fifth of February; when they meditated the
destruction, not of the consuls only, but of most of the senate. And had
not Catiline, who was in front of the senate-house, been too hasty to
give the signal to his associates, there would that day have been
perpetrated the most atrocious outrage since the city of Rome was
founded. But as the armed conspirators had not yet assembled in
sufficient numbers, the want of force frustrated the design.

XIX. Some time afterward, Piso was sent as quaestor, with Praetorian
authority, into Hither Spain; Crassus promoting the appointment,
because he knew him to be a bitter enemy to Cneius Pompey. Nor were
the senate, indeed, unwilling\footnote{XIX. Nor were the senate, indeed, unwilling, etc.—See Dio Cass.
xxxvi. 27.} to grant him the province; for they
wished so infamous a character to be removed from the seat of
government; and many worthy men, at the same time, thought that there
was some security in him against the power of Pompey, which was then
becoming formidable. But this Piso, on his march toward his province,
was murdered by some Spanish cavalry whom he had in his army. These
barbarians, as some say, had been unable to endure his unjust,
haughty, and cruel orders; but others assert that this body of
cavalry, being old and trusty adherents of Pompey, attacked Piso at
his instigation; since the Spaniards, they observed, had never before
committed such an outrage, but had patiently submitted to many severe
commands. This question we shall leave undecided. Of the first
conspiracy enough has been said.

XX. When Catiline saw those, whom I have just above mentioned,\footnote{XX. Just above mentioned—In c. 17.}
assembled, though he had often discussed many points with them singly,
yet thinking it would be to his purpose to address and exhort them in
a body, retired with them into a private apartment of his house,
where, when all witnesses were withdrawn, he harangued them to the
following effect:

``If your courage and fidelity had not been sufficiently proved by me,
this favorable opportunity\footnote{Favorable opportunity—\emph{Opportuna res}. See the latter part
of c. 16.} would have occurred to no purpose;
mighty hopes, absolute power, would in vain be within our grasp; nor
should I, depending on irresolution or ficklemindedness, pursue
contingencies instead of certainties. But as I have, on many remarkable
occasions, experienced your bravery and attachment to me, I have
ventured to engage in a most important and glorious enterprise. I am
aware, too, that whatever advantages or evils affect you, the same
affect me; and to have the same desires and the same aversions, is
assuredly a firm bond of friendship.

``What I have been meditating you have already heard separately. But my
ardor for action is daily more and more excited, when I consider what
our future condition of life must be, unless we ourselves assert our
claims to liberty.\footnote{Assert our claims to liberty—\emph{Nosmet ipsi vindicamus in
libertatem}.Unless we vindicate ourselves into liberty. See below,
``En illa, illa, quam saepe optastis, libertas,'' etc.} For since the government has fallen under the
power and jurisdiction of a few, kings and princes\footnote{Kings and princes—\emph{Reges tetrarchae}. \emph{Tetrarchs} were
properly those who had the government of the fourth part of the
country; but at length, the signification of the word being extended,
it was applied to any governors of any country who were possessed of
supreme authority, and yet were not acknowledged as kings by the
Romans. See Hirt. Bell. Alex. c. 67: ``Deiotarus, at that time
\emph{tetrarch} of almost all Gallograecia, a supremacy which the other
\emph{tetrarchs} would not allow to be granted him either by the laws or by
custom, but indisputably acknowledged as king of Armenia Minor by the
senate,'' etc. \emph{Dietsch.} ``Hesychius has, [Greek: \emph{Tetrarchas,
basileis}]. See Isidor., ix. 8; Alex. ab. Alex., ii. 17.'' \emph{Colerus}.
``Cicero, Phil. II., speaks of Reges Tetrarchas Dynastasque. And Lucan
has (vii. 46) Tetrarchae regesque tenent, magnique tyranni.'' \emph{Wasse.}
Horace also says,} have constantly
been their tributaries; nations and states have paid them taxes; but all
the rest of us, however brave and worthy, whether noble or plebeian,
have been regarded as a mere mob, without interest or authority, and
subject to those, to whom, if the state were in a sound condition, we
should be a terror. Hence, all influence, power, honor, and wealth, are
in their hands, or where they dispose of them: to us they have left only
insults,\footnote{Insults—\emph{Repulsas}. Repulses in standing for office.} dangers, persecutions, and poverty. To such indignities,
O bravest of men, how long will you submit? Is it not better to die in
a glorious attempt, than, after having been the sport of other men's
insolence, to resign a wretched and degraded existence with ignominy?

``But success (I call gods and men to witness!) is in our own hands.
Our years are fresh, our spirit is unbroken; among our oppressors, on
the contrary, through age and wealth a general debility has been
produced. We have therefore only to make a beginning; the course of
events\footnote{The course of events, etc.—\emph{Caetera res expediet}.—``Of. Cic.
Ep. Div. xiii. 26: \emph{explicare et expedire negotia}.'' Gerlach.} will accomplish the rest.

``Who in the world, indeed, that has the feelings of a man, can endure
that they should have a superfluity of riches, to squander in building
over seas\footnote{Building over seas—See c. 13.} and leveling mountains, and that means should be wanting
to us even for the necessaries of life; that they should join together
two houses or more, and that we should not have a hearth to call our
own? They, though they purchase pictures, statues, and embossed plate;
\footnote{Embossed plate—\emph{Toreumata}. The same as \emph{vasa coelata},
sculptured vases, c. 11. Vessels ornamented in bas-relief; from
[Greek: \emph{toreuein}], \emph{sculpere}; see Bentley ad Hor. A. P., 441.
``Perbona toreumata, in his pecula duo,'' etc. Cic. in Verr. iv. 18.} though they pull down now buildings and erect others, and lavish
and abuse their wealth in every possible method; yet can not, with the
utmost efforts of caprice, exhaust it. But for us there is poverty at
home, debts abroad; our present circumstances are bad, our prospects
much worse; and what, in a word, have we left, but a miserable existence?

``Will you not, then, awake to action? Behold that liberty, that
liberty for which you have so often wished, with wealth, honor, and
glory, are set before your eyes. All these prizes fortune offers to
the victorious. Let the enterprise itself, then, let the opportunity,
let your poverty, your dangers, and the glorious spoils of war,
animate you far more than my words. Use me either as your leader or
your fellow-soldier; neither my heart nor my hand shall be wanting to
you. These objects I hope to effect, in concert with you, in the
character of consul; unless, indeed, my expectation deceives me, and
you prefer to be slaves rather than masters.''

XXI. When these men, surrounded with numberless evils, but without any
resources or hopes of good, had heard this address, though they
thought it much for their advantage to disturb the public tranquillity,
yet most of them called on Catiline to state on what terms they were to
engage in the contest; what benefits they were to expect from taking up
arms; and what support and encouragement they had, and in what quarters.
\footnote{XXI. What support or encouragement they had, and in what
quarters.—\emph{Quid ubique opis aut spei haberent; i.e.} quid opis aut
So c. 27, \emph{init.} Quem ubique opportunum credebat, \emph{i.e.}, says
Cortius, ``quem, et ubi \emph{illum}, opportunum credebat''.} Catiline then promised them the abolition of their debts;\footnote{Abolition of their debts—\emph{Tabulas novas.} Debts were
registered on tablets; and, when the debts were paid, the score was
effaced, and the tablets were ready to be used \emph{as new.} See Ernesti's
Clav. in Cio.\emph{sub voce}.} a
proscription of the wealthy citizens;\footnote{Proscription of the wealthy citizens—\emph{Proscriptionem
locupletium.} The practice of proscription was commenced by Sylla, who
posted up, in public places of the city, the names of those whom he
doomed to death, offering rewards to such as should bring him their
heads. Their money and estates he divided among his adherents, and
Catiline excited his adherents with hopes of similar plunder.} offices, sacerdotal dignities,
plunder, and all other gratifications which war, and the license of
conquerors, can afford. He added that Piso was in Hither Spain, and
Publius Sittius Nucerinus with an army in Mauritania, both of whom were
privy to his plans; that Caius Antonius, whom he hoped to have for a
colleague, was canvassing for the consulship, a man with whom he was
intimate, and who was involved in all manner of embarrassments; and that,
in conjunction with him, he himself, when consul, would commence
operations. He, moreover, assailed all the respectable citizens with
reproaches, commended each of his associates by name, reminded one of
his poverty, another of his ruling passion,\footnote{Another of his ruling passion—\emph{Admonebat—alium cupiditatis
suae}. Rose renders this passage, ``Some he put in mind of their
poverty, others of their amours.'' De Brosses renders it, ``Il
remontre à l'un sa pauvreté, à l'autre son ambition.'' \emph{Ruling
passion}, however, seems to be the proper sense of \emph{cupiditatis};
as it is said, in c. 14, ``As the passions of each, according to his
years, appeared excited, he furnished mistresses to some, bought
horses and dogs for others'', etc.} several others of their
danger or disgrace, and many of the spoils which they had obtained by
the victory of Sylla. When he saw their spirits sufficiently elevated,
he charged them to attend to his interest at the election of consuls,
and dismissed the assembly.

XXII. There were some, at the time, who said that Catiline, having
ended his speech, and wishing to bind his accomplices in guilt by an
oath, handed round among them, in goblets, the blood of a human body
mixed with wine; and that when all, after an imprecation, had tasted
of it, as is usual in sacred rites, he disclosed his design; and they
asserted\footnote{XXII. They asserted—\emph{Dictitare}. In referring this word to
the circulators of the report, I follow Cortius, Gerlach, Kritzius,
and Bernouf. Wasse, with less discrimination, refers it to Catiline.
This story of the drinking of human blood is copied by Florus, iv 1,
and by Plutarch in his Life of Cicero. Dio Cassius (lib. xxxvii.) says
that the conspirators were reported to have killed a child on the
occasion.} that he did this, in order that they might be the more
closely attached to one another, by being mutually conscious of such
an atrocity. But some thought that this report, and many others, were
invented by persons who supposed that the odium against Cicero, which
afterward arose, might be lessened by imputing an enormity of guilt to
the conspirators who had suffered death. The evidence which I have
obtained, in support of this charge, is not at all in proportion to
its magnitude.

XXIII. Among those present at this meeting was Quintus Curius,\footnote{XXIII. Quintus Curius—the same that is mentioned in c. 17.} a
man of no mean family, but immersed in vices and crimes, and whom the
censors had ignominiously expelled from the senate. In this person
there was not less levity than impudence; he could neither keep secret
what he heard, not conceal his own crimes; he was altogether heedless
what he said or what he did. He had long had a criminal intercourse
with Fulvia, a woman of high birth; but growing less acceptable to her,
because, in his reduced circumstances, he had less means of being
liberal, he began, on a sudden, to boast, and to promise her seas and
mountains;\footnote{To promise her seas and mountains—\emph{Maria montesque polliceri}.
A proverbial expression. Ter. Phorm., i. 2, 18: \emph{Modò non montes auri
pollicens}. Perc., iii. 65: \emph{Et quid opus Cratero magnos promittere
emontes.}} threatening her, at times, with the sword, if she were
not submissive to his will; and acting, in his general conduct, with
greater arrogance than ever.\footnote{With greater arrogance than ever—\emph{Ferocius quam solitus erat.}} Fulvia, having learned the cause of
his extravagant behavior, did not keep such danger to the state a
secret; but, without naming her informant, communicated to several
persons what she had heard and under what circumstances, concerning
Catiline's conspiracy. This intelligence it was that incited the
feelings of the citizens to give the consulship to Marcus Tullius
Cicero.\footnote{To Marcus Tullius Cicero—Cicero was now in his forty-third
year, and had filled the office of quaestor, aedile, and praetor.} For before this period, most of the nobility were moved
with jealousy, and thought the consulship in some degree sullied, if a
man of no family,\footnote{A man of no family—\emph{Novus homo.} A term applied to such as
could not boast of any ancestor that had held any curule magistracy,
that is, had been consul, praetor, censor, or chief aedile.} however meritorious, obtained it. But when
danger showed itself, envy and pride were laid aside. XXIV.
Accordingly, when the comitia were held, Marcus Tullius and Caius
Antonius were declared consuls; an event which gave the first shock to
the conspirators. The ardor of Catiline, however, was not at all
diminished; he formed every day new schemes; he deposited arms, in
convenient places, throughout Italy; he sent sums of money borrowed on
his own credit, or that of his friends, to a certain Manlius,\footnote{XXIV. Manlius—He had been an officer in the army of Sylla,
and, having been distinguished for his services, had been placed at
the head of a colony of veterans settled about Faesulae: but he had
squandered his property in extravagance. See Plutarch, Vit. Cic., Dio
Cassius, and Appian.} at
Faesulae,\footnote{Faesulae—A town of Etruria, at the foot of the Appennines,} who was subsequently the first to engage in hostilities.
At this period, too, he is said to have attached to his cause great
numbers of men of all classes, and some women, who had, in their earlier
days, supported an expensive life by the price of their beauty, but who,
when age had lessened their gains but not their extravagance, had
contracted heavy debts. By the influence of these females, Catiline
hoped to gain over the slaves in Rome, to get the city set on fire, and
either to secure the support of their husbands or take away their lives.

XXV. In the number of those ladies was Sempronia,\footnote{XXV. Sempronia—Of the same \emph{gens} as the two Gracchi. She
was the wife of Decimus Brutus.} a woman who had
committed many crimes with the spirit of a man. In birth and beauty,
in her husband and her children, she was extremely fortunate; she was
skilled in Greek and Roman literature; she could sing, play, and
dance,\footnote{Sing, play, and dance—\emph{Psallere, saltare.} As \emph{psallo}
signifies both to play on a musical instrument, and to sing to it
while playing, I have thought it necessary to give both senses in the
translation.} with greater elegance than became a woman of virtue, and
possessed many other accomplishments that tend to excite the passions.
But nothing was ever less valued by her than honor or chastity.
Whether she was more prodigal of her money or her reputation, it would
have been difficult to decide. Her desires were so ardent that she
oftener made advances to the other sex than waited for solicitation.
She had frequently, before this period, forfeited her word, forsworn
debts, been privy to murder, and hurried into the utmost excesses by
her extravagance and poverty. But her abilities were by no means
despicable;\footnote{By no means despicable—\emph{Haud absurdum.} Compare, \emph{Bene dicere
haud absurdum est,} c. 8.} she could compose verses, jest, and join in
conversation either modest, tender, or licentious. In a word, she was
distinguished\footnote{She was distinguished, etc.—\emph{Multae facetiae, multusque lepos
inerat.} Both \emph{facetiae} and \emph{lepos} mean ``agreeableness, humor,
pleasantry,'' but \emph{lepos} here seems to refer to diction, as in Cic.
Orat. i. 7: \emph{Magnus in jocando lepos.}} by much refinement of wit, and much grace of
expression.

XXVI. Catiline, having made these arrangements, still canvassed for
the consulship for the following year; hoping that, if he should be
elected, he would easily manage Antonius according to his pleasure.
Nor did he, in the mean time remain inactive, but devised schemes, in
every possible way, against Cicero, who, however, did not want skill
or policy to guard, against them. For, at the very beginning of his
consulship, he had, by making many promises through Fulvia, prevailed
on Quintus Curius, whom I have already mentioned, to give him secret
information of Catiline's proceedings. He had also persuaded his
colleague, Antonius, by an arrangement respecting their provinces,\footnote{XXVI. By an arrangement respecting their provinces—\emph{Pactione
provinciae}. This passage has been absurdly misrepresented by most
translators, except De Brosses. Even Rose, who was a scholar, translated
\emph{pactione provinciae}, ``by promising a province to his colleague.''
Plutarch, in his Life of Cicero, says that the two provinces, which
Cicero and his colleague Antonius shared between them, were Gaul and
Macedonia, and that Cicero, in order to retain Antonius in the interest
of the senate, exchanged with him Macedonia, which had fallen to himself,
for the inferior province of Gaul. See Jug., c. 27.}
to entertain no sentiment of disaffection toward the state; and he kept
around him, though without ostentation, a guard of his friends and
dependents.

When the day of the comitia came, and neither Catiline's efforts for
the consulship, nor the plots which he had laid for the consuls in the
Campus Martius,\footnote{Plots which he had laid for the consuls in the Campus Martius
—\emph{Insidiae quas consuli in campo fecerat}. I have here departed from
the text of Cortius, who reads \emph{consulibus}, thinking that Catiline, in
his rage, might have extended his plots even to the consuls-elect. But
\emph{consuli}, there is little doubt, is the right reading, as it is favored
by what is said at the beginning of the chapter, \emph{insidias parabat
Ciceroni}, by what follows in the next chapter, \emph{consuli insidias tendere},
and by the words, \emph{sperans, si designatus foret, facile se ex voluntate
Antonio usurum}; for if Catiline trusted that he should be able to use
his pleasure with Antonius, he could hardly think it necessary to form
plots against his life. I have De Brosses on my side, who translates the
phrase, \emph{les pieges où il comptait faire périr le consul}. The words \emph{in
campo}, which look extremely like an intruded gloss, I wonder that
Cortius should have retained. ``\emph{Consuli},'' says Gerlach, ``appears the
more eligible, not only on account of \emph{consuli insidias tendere}, c. 27,
but because nothing but the death of Cicero was necessary to make
everything favorable for Catiline.'' Kritzius, Bernouf, Dietsch, Pappaur,
Allen, and all the modern editors, read \emph{Consuli}. See also the end of
c. 27: \emph{Si prius Ciceronem oppressisset}.] [note 144: Had ended in
confusion and disgrace—\emph{Aspera faedaque evenerant}. I have borrowed
from Murphy.} were attended with success, he determined to
proceed to war, and resort to the utmost extremities, since what he
had attempted secretly had ended in confusion and disgrace.

XXVII. He accordingly dispatched Caius Manlius to Faesulae, and the
adjacent parts of Etruria; one Septimius, of Carinum,\footnote{XXVII. Of Camerinum—Camertem. ``That is, a native of Camerinum,
a town on the confines of Umbria and Picenum. Hence the noun \emph{Camers},
as Cic. Pro. Syll., c. 19, \emph{in agro Camerti}.'' Cortius.} into the
Picenian territory; Caius Julius into Apulia; and others to various
places, wherever he thought each would be most serviceable.\footnote{Wherever he thought each would be most serviceable—\_Ubi
quemque opportunum credebat. ``Proprie reddas: quam, \emph{et ubi} illum,
\emph{opportunum credebat},'' Cortius. See c. 23.} He
himself, in the mean time, was making many simultaneous efforts at
Rome; he laid plots for the consul; he arranged schemes for burning
the city; he occupied suitable posts with armed men; he went constantly
armed himself, and ordered his followers to do the same; he exhorted
them to be always on their guard and prepared for action; he was active
and vigilant by day and by night, and was exhausted neither by
sleeplessness nor by toil. At last, however, when none of his
numerous projects succeeded,\footnote{When none of his numerous projects succeeded—\emph{Ubi multa
agilanti nihil procedit}.} he again, with the aid of Marcus
Porcius Laeca, convoked the leaders of the conspiracy in the dead of
night, when, after many complaints of their apathy, he informed them
that he had sent forward Manlius to that body of men whom he had
prepared to take up arms; and others of the confederates into other
eligible places, to make a commencement of hostilities; and that he
himself was eager to set out to the army, if he could but first cut
off Cicero, who was the chief obstruction to his measures.

XXVIII. While, therefore, the rest were in alarm and hesitation, Caius
Cornelius, a Roman knight, who offered his services, and Lucius
Vargunteius, a senator, in company with him, agreed to go with an
armed force, on that very night, and with but little delay,\footnote{XXVIII. On that very night, and with but little delay—\emph{Ea
nocte, paulo post}. They resolved on going soon after the meeting
broke up, so that they might reach Cicero's house early in the
morning, which was the usual time for waiting on great men. \emph{Ingentem
foribus domus alla superbis} Mane \emph{salutantûm totis vomit aedibus
undam}. Virg. Georg., ii. 461.} to
the house of Cicero, under pretense of paying their respects to him,
and to kill him unawares, and unprepared for defense, in his own
residence. But Curius, when he heard of the imminent danger that
threatened the consul, immediately gave him notice, by the agency of
Fulvia, of the treachery which was contemplated. The assassins, in
consequence, were refused admission, and found that they had
undertaken such an attempt only to be disappointed.

In the mean time, Manlius was in Etruria, stirring up the populace,
who, both from poverty, and from resentment for their injuries (for,
under the tyranny of Sylla, they had lost their lands and other
property) were eager for a revolution. He also attached to himself all
sorts of marauders, who were numerous in those parts, and some of
Sylla's colonists, whose dissipation and extravagance had exhausted
their enormous plunder.

XXIX. When these proceedings were reported to Cicero, he, being
alarmed at the twofold danger, since he could no longer secure the
city against treachery by his private efforts, nor could gain
satisfactory intelligence of the magnitude or intentions of the army
of Manlius, laid the matter, which was already a subject of discussion
among the people, before the senate. The senate, accordingly, as is
usual in any perilous emergency, decreed that THE CONSULS SHOULD MAKE
IT THEIR CARE THAT THE COMMONWEALTH SHOULD RECEIVE NO INJURY. This is
the greatest power which, according to the practice at Rome, is
granted\footnote{XXIX. This is the greatest power which—is granted, etc.
—\emph{Ea potestas per senatum, more Romano, magistratui maxima
permittitur}. Cortius, \emph{mirâ judicii peversitate}, as Kritzius
observes, makes \emph{ea} the ablative case, understanding ``decretione,''
``formula,'' or some such word; but, happily, no one has followed him.} by the senate to the magistrate, and which authorizes him
to raise troops; to make war; to assume unlimited control over the
allies and the citizens; to take the chief command and jurisdiction at
home and in the field; rights which, without an order of the people,
the consul is not permitted to exercise.

XXX. A few days afterward, Lucius Saenius, a senator, read to the
senate a letter, which, he said, he had received from Faesulae, and in
which, it was stated that Caius Manlius, with a large force, had taken
the field by the 27th of October.\footnote{XXX. By the 27th of October—\emph{Ante diem VI. Kalendas
Novembres}. He means that they were in arms on or before that day.} Others at the same time, as is
not uncommon in such a crisis, spread reports of omens and prodigies;
others of meetings being held, of arms being transported, and of
insurrections of the slaves at Capua and in Apulia. In consequence of
these rumors, Quintus Marcius Rex\footnote{Quintus Marcius Rex—He had been proconsul in Cilicia, and
was expecting a triumph for his successes.} was dispatched, by a decree of
the senate, to Faesulae, and Quintus Metellus Creticus\footnote{Quintus Metellus Creticus—He had obtained the surname of
Creticus from having reduced the island of Crete.} into
Apulia and the parts adjacent; both which officers, with the title of
commanders,\footnote{Both which officers, with the title of commanders, etc.
—\emph{hi utrique ad urbem imperatores erant; impediti ne triumpharent
calumniâ paucorum quibus omnia honesta atque inhonesta vendere mos
erat}. ``Imperator'' was a title given by the army, and confirmed by the
senate, to a victorious general, who had slain a certain number of the
enemy. What the number was is not known. The general bore this title
as an addition to his name, until he obtained (if it were granted him)
a triumph, for which he was obliged to wait \emph{ad urbem}, near the city,
since he was not allowed to enter the gates as long as he held any
military command. These \emph{imperatores} had been debarred from their
expected honor by a party who would sell \emph{any thing honorable}, as a
triumph, or \emph{any thing dishonorable}, as a license to violate the laws.} were waiting near the city, having been prevented
from entering in triumph, by the malice of a cabal, whose custom it
was to ask a price for every thing, whether honorable or infamous. The
praetors, too, Quintus Pompeius Rufus, and Quintus Metellus Celer, were
sent off, the one to Capua, the other to Picenum, and power was given
them to levy a force proportioned to the exigency and the danger. The
senate also decreed, that if any one should give information of the
conspiracy which had been formed against the state, his reward should
be, if a slave, his freedom and a hundred sestertia; if a freeman, a
complete pardon and two hundred sestertia\footnote{A hundred sestertia—two hundred sestertia—A hundred sestertia
were about 807£. 5s. 10d. of our money.}. They further appointed
that the schools of gladiators\footnote{Schools of gladiators—\emph{Gladiatoriae familiae}. Any number of
gladiators under one teacher, or trainer (\emph{lanista}), was called
\emph{familia}. They were to be distributed in different parts, and to be
strictly watched, that they might not run off to join Catiline. See
Graswinckelius, Rupertus, and Gerlach.} should be distributed in Capua and
other municipal towns, according to the capacity of each; and that, at
Rome, watches should be posted throughout the city, of which the
inferior magistrates\footnote{The inferior magistrates—The aediles, tribunes, quaestors,
and all others below the consuls, censors, and praetors. Aul. Cell.,
xiii. 15.} should have the charge.

XXXI. By such proceedings as these the citizens were struck with
alarm, and the appearance of the city was changed. In place of that
extreme gayety and dissipation,\footnote{XXXI. Dissipation—Lascivia. ``Devotion to public amusements
and gayety. The word is used in the same sense as in Lucretius, v.} to which long tranquillity\footnote{Long tranquillity—\emph{Diuturna quies}. ``Since the victory of
Sylla to the time of which Sallust is speaking, that is, for about
twenty years, there had been a complete cessation from civil discord
and disturbance'' \emph{Bernouf}.}
had given rise, a sudden gloom spread over all classes; they became
anxious and agitated; they felt secure neither in any place, nor with
any person; they were not at war, yet enjoyed no peace; each measured
the public danger by his own fear. The women, also, to whom, from the
extent of the empire, the dread of war was new, gave way to lamentation,
raised supplicating hands to heaven, mourned over their infants, made
constant inquiries, trembled at every thing, and, forgetting their pride
and their pleasures, felt nothing but alarm for themselves and their
country.

Yet the unrelenting spirit of Catiline persisted in the same purposes,
notwithstanding the precautions that were adopted against him, and
though he himself was accused by Lucius Paullus under the Plautian
law.\footnote{The Plautian law—\emph{Lege Plautia}. ``This law was that of M.
Plautius Silanus, a tribune of the people, which was directed against
such as excited a sedition in the state, or formed plots against the
life of any individual.'' \emph{Cyprianus Popma}. See Dr. Smith's Dict. of
Gr. and Rom. Antiquities, sub Vis.} At last, with a view to dissemble, and under pretense of
clearing his character, as if he had been provoked by some attack, he
went into the senate-house. It was then that Marcus Tullius, the
consul, whether alarmed at his presence, or fired with indignation
against him, delivered that splendid speech, so beneficial to the
republic, which he afterward wrote and published.\footnote{Which he afterward wrote and published—\emph{Quam postea scriptam
edidit}. This was the first of Cicero's four Orations against
Catiline. The epithet applied to it by Sallust, which I have rendered
``splendid,'' is \emph{luculentam}; that is, says Gerlach, ``luminibus
verborum et sententiarum ornatam,'' distinguished by much brilliancy of
words and thoughts. And so say Kritzius, Bernouf, and Dietsch. Cortius,
who is followed by Dahl, Langius, and Muller, makes the word equivalent
merely to \emph{lucid}, in the supposition that Sallust intended to bestow
on the speech, as on other performances of Cicero, only very cool praise.
\emph{Luculentus}, however, seems certainly to mean something more than
\emph{lucidus}.}

When Cicero sat down, Catiline, being prepared to pretend ignorance of
the whole matter, entreated, with downcast looks and suppliant voice,
that ``the Conscript Fathers would not too hastily believe any thing
against him;'' saying ``that he was sprung from such a family, and had
so ordered his life from his youth, as to have every happiness in
prospect; and that they were not to suppose that he, a patrician,
whose services to the Roman people, as well as those of his ancestors,
had been so numerous, should want to ruin the state, when Marcus
Tullius, a mere adopted citizen of Rome,\footnote{A mere adopted citizen of Rome—\emph{Inquilinus civis urbis Romae}.
``Inquilinus'' means properly a lodger, or tenant in the house of another.
Cicero was born at Arpinum, and is therefore called by Catiline a
citizen of Rome merely by adoption or by sufferance. Appian, in
repeating this account (Bell. Civ., ii. 104), says, [Greek:
\emph{Ingkouilinon, phi raemati kalousi tous enoikountas en allotriais
oikiais}.]} was eager to preserve
it.'' When he was proceeding to add other invectives, they all raised
an outcry against him, and called him an enemy and a traitor.\footnote{Traitor—\emph{Parricidam}. See c. 14. ``An oppressor or betrayer
of his country is justly called a parricide; for our country is the
common parent of all. Cic. ad Attic.'' \emph{Wasse}.}
Being thus exasperated, ``Since I am encompassed by enemies,'' he
exclaimed,\footnote{Since I am encompassed, by enemies, he exclaimed, etc.—``It
was not on this day, nor indeed to Cicero, that this answer was made
by Catiline. It was a reply to Cato, uttered a few days before the
comitia for electing consuls, which were held on the 22d day of
October. See Cic. pro Muraeno, c. 25. Cicero's speech was delivered on
the 8th of November. Sallust is, therefore, in error on this point, as
well as Florus and Valerius Maximus, who have followed him.''
\emph{Bernouf}. From other accounts we may infer that no reply was made to
Cicero by Catiline on this occasion. Plutarch, in his Life of Cicero,
says that Catiline, before Cicero rose, seemed desirous to address the
senate in defense of his proceedings, but that the senators refused to
listen to him. Of any answer to Cicero's speech, on the part of
Catiline, he makes no mention. Cicero himself, in his second Oration
against Catiline, says that Catiline \emph{could not endure his voice},
but, when he was ordered to go into exile, ``paruit, quievit,'' \emph{obeyed
and submitted in silence}. And in his Oration, c. 37, he says, ``That
most audacious of men, Catiline, when he was accused by me in the
senate, was dumb.''} ``and driven to desperation, I will extinguish the
flame kindled around me in a general ruin.''

XXXII He then hurried from the senate to his own house; and then,
after much reflection with himself, thinking that, as his plots
against the consul had been unsuccessful, and as he knew the city to
be secured from fire by the watch, his best course would be to augment
his army, and make provision for the war before the legions could be
raised, he set out in the dead of night, and with a few attendants, to
the camp of Manlius. But he left in charge to Lentulus and Cethegus,
and others of whose prompt determination he was assured, to strengthen
the interests of their party in every possible way, to forward the
plots against the consul, and to make arrangements for a massacre, for
firing the city, and for other destructive operations of war;
promising that he himself would shortly advance on the city with a
large army.

During the course of these proceedings at Rome, Caius Manlius
dispatched some of his followers as deputies to Quintus Marcius Rex,
with directions to address him\footnote{XXXII. With directions to address him, etc.—\emph{Cum mandatis
hujuscemodi}. The communication, as Cortius observes, was not an
epistle, but a verbal message.} to the following effect:

XXXIII. ``We call gods and men to witness, general, that we have taken
up arms neither to injure our country, nor to occasion peril to any
one, but to defend our own persons from harm; who, wretched and in
want, have been deprived most of us, of our homes, and all of us of
our character and property, by the oppression and cruelty of usurers;
nor has any one of us been allowed, according to the usage of our
ancestors, to have the benefit of the law,\footnote{XXXIII. To have the benefit of the law—\emph{Lege uti}. The law
here meant was the Papirian law, by which it was provided, contrary to
the old law of the Twelve Tables, that no one should be confined in
prison for debt, and that the property of the debtor only, not his
person, should be liable for what he owed. Livy (viii. 28) relates the
occurrence which gave rise to this law, and says that it ruptured one
of the strongest bonds of credit.} or, when our property
was lost to keep our persons free. Such has been the inhumanity of the
usurers and of the praetor.\footnote{The praetor—The \emph{praetor urbanus}, or city praetor, who
decided all causes between citizens, and passed sentence on debtors.}

Often have your forefathers, taking compassion on the commonalty at
Rome, relieved their distress by decrees;\footnote{Relieved their distress by decrees—\emph{Decretis suis inopiae
opitulati sunt}. In allusion to the laws passed at various times for
diminishing the rate of interest.} and very lately, within
our own memory, silver, by reason of the pressure of debt, and with
the consent of all respectable citizens, was paid with brass.\footnote{Silver—was paid with brass—\emph{Agentum aere solutum est}.
Thus a \emph{sestertius}, which was of silver, and was worth four \emph{asses},
was paid with one \emph{as}, which was of brass; or \emph{the fourth part only
of the debt was paid}. See Plin. H. N. xxxiii. 3; and Velleius
Paterculus, ii. 23; who says, \emph{quadrantem solvi}, that \emph{a quarter} of
their debts were paid by the debtors, by a law of Valerius Flaccus,
when he became consul on the death of Marius.}

Often too, we must own, have the commonalty themselves, driven by
desire of power, or by the arrogance of their rulers, seceded\footnote{Often—have the commonalty—seceded, etc.—``This happened
three times: 1. To the Mons Sacer, on account of debt; Liv. ii. 32. 2.
To the Aventine, and thence to the Mons Sacer, through the tyranny of
Appius Claudius, the decemvir; Liv. iii. 50. 3. To the Janiculum, on
account of debt; Liv. Epist. xi.'' \emph{Bernouf}.}
under arms from the patricians. But at power or wealth, for the sake
of which wars, and all kinds of strife, arise among mankind, we do not
aim; we desire only our liberty, which no honorable man relinquishes
but with life. We therefore conjure you and the senate to befriend
your unhappy fellow-citizens; to restore us the protection of the law,
which the injustice of the praetor has taken from us; and not to lay
on us the necessity of considering how we may perish, so as best to
avenge our blood.''

XXXIV. To this address Quintus Marcius replied, that, ``if they wished
to make any petition to the senate, they must lay down their arms, and
proceed as suppliants to Rome;'' adding, that ``such had always been the
kindness\footnote{XXXIV. That such had always been the kindness, etc.—\emph{Ea,
mansuetudine atque misericordia senatum populumque Romanum, semper
fuisse.} ``That the senate, etc., had always been of such kindness.'' I
have deserted the Latin for the English idiom.} and humanity of the Roman senate and people, that none
had ever asked help of them in vain.''

Catiline, on his march, sent letters to most men of consular dignity,
and to all the most respectable citizens, stating that ``as he was
beset by false accusations, and unable to resist the combination of
his enemies, he was submitting to the will of fortune, and going into
exile at Marseilles; not that he was guilty of the great wickedness
laid to his charge, but that the state might be undisturbed, and that
no insurrection might arise from his defense of himself.''

Quintus Catulus, however, read in the senate a letter of a very
different character, which, he said, was delivered to him in the
name of Catiline, and of which the following is a copy.

\footnote{XXXV. The commencement of this letter is different in different
editions. In Havercamp it stands thus: \emph{Egregiatua fides, re cognita,
grata mihi, magnis in meis periculis, fiduciam commendationi meae
tribuit.} Cortius corrected it as follows: \emph{Egregia tua fides, re
cognita, gratam in magnis periculis fiduciam commendationi meae
tribuit.} Cortius's reading has been adopted by Kritzius, Bernouf, and
most other editors. Gerlach and Dietsch have recalled the old text.
That Cortius's is the better; few will deny; for it can hardly be
supposed that Sallust used \emph{mihi, meis}, and \emph{meae} in such close
succession. Some, however, as Rupertus and Gerlach, defend Havercamp's
text, by asserting, from the phrase \emph{earum exemplam infra scriptum,}
that this is a true copy of the letter, and that the style is,
therefore, not Sallust's, but Catiline's. But such an opinion is
sufficiently refuted by Cortius, whose remarks I will transcribe:
``Rupertus,'' says he, ``quod in promptu erat, Catilinae culpam tribuit,
qui non eo, quo Crispus, stilo scripserit. Sed cur oratio ejus tam
apta et composita suprà, c. 20 refertur? At, inquis, hic ipsum
litterarum exemplum exhibetur. At vide mihi exemplum litterarum
Lentuli, c. 44; et lege Ciceronem, qui idem exhibet, et senties sensum
magis quam verba referri. Quare inanis haec quidem excusatio.'' Yet it
is not to be denied that \emph{grata mihi} is the reading of all the
manuscripts.}XXXV. ``Lucius Catiline to Quintus Catulus, wishing health. Your
eminent integrity, known to me by experience,\footnote{Known—by experience.—\emph{Re cognita.} ``Cognita'' be it observed,
\emph{tironum gratia,} is the nominative case. ``Catiline had experienced
the friendship of Catulus in his affair with Fabia Terentia; for it
was by his means that he escaped when he was brought to trial, as is
related by Orosius.'' \emph{Bernouf.}} gives a pleasing
confidence, in the midst of great perils, to my present recommendation.
\footnote{Recommendation—\emph{Commendationi.} His recommendation of his
affairs, and of Orestilla, to the care of Catulus.} I have determined, therefore, to make no formal defense\footnote{Formal defense—\emph{Defensionem.} Opposed to \emph{satisfactionem},
which follows, and which means a private apology or explanation.
``\emph{Defensio}, a defense, was properly a statement or speech to be made
against an adversary, or before judges; \emph{satisfactio} was rather an
excuse or apology made to a friend, or any other person, in a private
communication.'' \emph{Cortius.}} with
regard to my new course of conduct; yet I was resolved, though conscious
of no guilt,\footnote{Though conscious of no guilt—\emph{Ex nullâ conscientiâ de culpâ}.
This phrase is explained by Cortius as equivalent to ``Propter
conscientam denullâ culpâ,'' or ``inasmuch as I am conscious of no
fault.'' ``\emph{De culpâ}, he adds, is the same as \emph{culpae}; so in the ii.
Epist. to Caesar, c. 1: Neque \emph{de futuro} quisquam satix callidus;
and c. 9: \emph{de illis} potissimum jactura fit.''} to offer you some explanation,\footnote{To make no formal defense—to offer you some explanation
—\emph{Defensionem—parare; satisfactionem—proponere}. ``Parare,'' says
Cortius, ``is applied to a defense which might require some study and
premeditation; \emph{proponere} to such a statement as it was easy to make
at once''.} which, on my word
of honor,\footnote{On my word of honor—\emph{Me dius fidius}, sc. juvet. So may the
god of faith help me, as I speak truth. But who is the god of faith?
\emph{Dius}, say some, is the same as \emph{Deus} (Plautus has \emph{Deus} fidius,
Asin i. 1, 18); and the god here meant is probably Jupiter (\emph{sub dio}
being equivalent to \emph{sub Jove}); so that \emph{Dius fidius} (\emph{fidius} being
an adjective from \emph{fides}) will be the [Greek: \emph{Zeus pistios}] of the
Greeks. ``\emph{Me dius fidius}'' will therefore be, ``May Jupiter help me!''
This is the mode of explication adopted by Gerlach, Bernouf, and
Dietsch. Others, with Festus (sub voce \emph{Medius fidius}) make \emph{fidius}
equivalent to \emph{filius}, because the ancients, according to Festus,
often used D for L, and \emph{dius fidius} will then be the same as [Greek:
\emph{Dios}] or Jovis filius, or Hercules, and \emph{medius fidius} will be the
same as \emph{mehercules} or \emph{mehercule}. Varro de L. L. (v. 10, ed.
Sprengel) mentions a certain Aelius who was of this opinion. Against
this derivation there is the quantity of \emph{fidius}, of which the first
syllable is short: \emph{Quaerebam Nonas Sanco fidone referrem}, Ov. Fast.
vi. 213. But if we consider \emph{dius} the same as \emph{deus}, we may as well
consider \emph{dius fidius} to be the god Hercules as the god Jupiter, and
may thus make \emph{medius fidius} identical with \emph{mehercules}, as it
probably is. ``Tertullian, de Idol. 20, says that \emph{medius fidius} is a
form of swearing by Hercules.'' Schiller's Lex. sub \emph{Fidius}. This
point will be made tolerably clear if we consider (with Varro, v. 10,
and Ovid, \emph{loc. cit.}) Dius Fidius to be the same with the Sabine
Sancus, or Semo Sancus, and Semo Sancus to be the same with Hercules.} you may receive as true.\footnote{You may receive as true—\emph{Veram licet cognoscas}. Some
editions, before that of Cortius, have \emph{quae—licet vera mecum
recognoscas}; which was adopted from a quotation of Servius ad Aen.
iv. 204. But twenty of the best MSS., according to Certius, have
\emph{veram licet cognoscas}.} Provoked by injuries and
indignities, since, being robbed of the fruit of my labor and exertion,
\footnote{Robbed of the fruit of my labor and exertion—\emph{Fructu laboris
industriaeque meae privatus}. ``The honors which he sought he
elegantly calls the \emph{fruit} of his labor, because the one is obtained
by the other.'' \emph{Cortius}.} I did not obtain the post of honor due to me,\footnote{Post of honor due to me—\emph{Statum dignitatis}. The consulship.} I have
undertaken, according to my custom, the public cause of the distressed.
Not but that I could have paid, out of my own property, the debts
contracted on my own security;\footnote{On my own security—\emph{Meis nominibus}. ``He uses the plural,''
says Herzogius, ``because he had not borrowed once only, or from one
person, but oftentimes, and from many.'' No other critic attempts to
explain this point. For \emph{alienis nominibus}, which follows, being in
the plural, there is very good reason. My translation is in conformity
with Bernouf's comment.} while the generosity of Orestilla,
out of her own fortune and her daughter's, would discharge those
incurred on the security of others. But because I saw unworthy men
ennobled with honors, and myself proscribed\footnote{Proscribed—\emph{Alienatum}. ``Repulsed from all hope of the
consulship.'' \emph{Bernouf}.} on groundless suspicion,
I have for this very reason, adopted a course,\footnote{Adopted a course—\emph{Spes—secutus sum}. ``\emph{Spem sequi} is a
phrase often used when the direction of the mind to any thing, action,
or course of conduct, and the subsequent election and adoption of what
appears advantageous, is signified.'' \emph{Cortius}.} amply justifiable
in my present circumstances, for preserving what honor is left to me.
When I was proceeding to write more, intelligence was brought that
violence is preparing against me. I now commend and intrust Orestilla
to your protection;\footnote{Protection—\emph{Fidei}.} intreating you, by your love for your own
children, to defend her from injury.\footnote{Intreating you, by your love for your own children, to defend
her from injury—\emph{Eam ab injuria defendas, per liberos tuos rogatus}.
``Defend her from injury, being intreated [to do so] by [or for the
sake of] your own children.''} Farewell.''

XXXVI. Catiline himself, having stayed a few days with Caius Flaminius
Flamma in the neighborhood of Arretium,\footnote{XXXVI. In the neighborhood of Arretium—\emph{In agro Arretino}.
Havercamp, and many of the old editions, have \emph{Reatino}; ``but,'' says
Cortius, ``if Catiline went the direct road to Faesulae, as is rendered
extremely probable by his pretense that he was going to Marseilles,
and by the assertion of Cicero, made the day after his departure, that
he was on his way to join Manlius, we must certainly read \emph{Arretino}.''
Arretium (now \emph{Arezzo}) lay in his road to Faesulae; Reate was many
miles out of it.} while he was supplying
the adjacent parts, already excited to insurrection, with arms,
marched with his fasces, and other ensigns of authority, to join
Manlius in his camp.

When this was known at Rome, the senate declared Catiline and Manlius
enemies to the state, and fixed a day as to the rest of their force,
before which they might lay down their arms with impunity, except such
as had been convicted of capital offenses. They also decreed that the
consuls should hold a levy; that Antonius, with an army, should hasten
in pursuit of Catiline; and that Cicero should protect the city.

At this period the empire of Rome appears to me to have been in an
extremely deplorable condition;\footnote{In an extremely deplorable condition—\emph{Multo maxime miserabile}.
\emph{Multe} is added to superlatives, like \emph{longe}. So c. 52, \emph{multo
pulcherrimam} eam nos haberemus. Cortius gives several other instances.} for though every nation, from the
rising to the setting of the sun, lay in subjection to her arms, and
though peace and prosperity, which mankind think the greatest
blessings, were hers in abundance, there yet were found, among her
citizens, men who were bent with obstinate determination, to plunge
themselves and their country into ruin; for, notwithstanding the two
decrees of the senate,\footnote{Notwithstanding the two decrees of the senate—\emph{Duobus senati
decretis.} I have translated it ``\emph{the} two decrees,'' with Rose.
One of the two was that respecting the rewards mentioned in c. 30; the
other was that spoken of in c. 36., allowing the followers of Catiline
to lay down their arms before a certain day.} not one individual, out of so vast a
number, was induced by the offer of reward to give information of the
conspiracy; nor was there a single deserter from the camp of Catiline.
So strong a spirit of disaffection had, like a pestilence, pervaded
the minds of most of the citizens.

XXXVII. Nor was this disaffected spirit confined to those who were
actually concerned in the conspiracy; for the whole of the common
people, from a desire of change, favored the projects of Catiline.
This they seemed to do in accordance with their general character;
for, in every state, they that are poor envy those of a better class,
and endeavor to exalt the factious;\footnote{XXXVII. Endeavor to exalt the factious—\emph{Malos extollunt}.
They strive to elevate into office those who resemble themselves.} they dislike the established
condition of things, and long for something new; they are discontented
with their own circumstances, and desire a general alteration; they
can support themselves amid tumult and sedition, without anxiety,
since poverty does not easily suffer loss.\footnote{Poverty does not easily suffer loss—\emph{Egestas facile habetur
sine damna} He that has nothing, has nothing to lose. Petron.
Sat., c. 119: \emph{Inops audacia tuta est}.}

As for the populace of the city, they had become disaffected\footnote{Had become disaffected—Praeceps abierat. Had grown demoralized,
sunk in corruption, and ready to join in any plots against the state.
So Sallust says of Sempronia, \emph{praeceps abierat}, c. 25.} from
various causes. In the first place,\footnote{In the first place—Primum omnium. ``These words refer, not to
\emph{item} and \_postremo in the same sentence, but to \emph{deinde} at the
commencement of the next.'' \emph{Bernouf}.} such as every where took the
lead in crime and profligacy, with others who had squandered their
fortunes in dissipation, and, in a word, all whom vice and villainy
had driven from their homes, had flocked to Rome as a general
receptacle of impurity. In the next place, many, who thought of the
success of Sylla, when they had seen some raised from common soldiers
into senators, and others so enriched as to live in regal luxury and
pomp, hoped, each for himself, similar results from victory, if they
should once take up arms. In addition to this, the youth, who, in the
country, had earned a scanty livelihood by manual labor, tempted by
public and private largesses, had preferred idleness in the city to
unwelcome toil in the field. To these, and all others of similar
character, public disorders would furnish subsistence. It is not at
all surprising, therefore, that men in distress, of dissolute
principles and extravagant expectations, should have consulted the
interest of the state no further than as it was subservient to their
own. Besides, those whose parents, by the victory of Sylla, had been
proscribed, whose property had been confiscated, and whose civil
rights had been curtailed,\footnote{Civil rights had been curtailed—\emph{Jus libertatis imminutum
erat}. ``Sylla, by one of his laws, had rendered the children of
proscribed persons incapable of holding any public office; a law
unjust, indeed, but which, having been established and acted upon for
more than twenty years, could not be rescinded without inconvenience
to the government. Cicero, accordingly, opposed the attempts which
were made, in his consulship, to remove this restriction, as he
himself states in his Oration against Piso, c. 2.'' \emph{Bernouf}. See
Vell. Patere., ii., 28; Plutarch, Vit. Syll.; Quintil., xi. 1, where a
fragment of Cicero's speech, \emph{De Proscriptorum Liberis}, is preserved.
This law of Sylla was at length abrogated by Julius Caesar, Suet. J.
Caes. 41; Plutarch Vit. Caes.; Dio Cass., xli. 18.} looked forward to the event of a war
with precisely the same feelings.

All those, too, who were of any party opposed to that of the senate,
were desirous rather that the state should be embroiled, than that
they themselves should be out of power. This was an evil, which, after
many years, had returned upon the community to the extent to which it
now prevailed.\footnote{This was an evil—to the extent to which it now prevailed—\emph{Id
adeò malum multos post annos in civitatem reverterat}. ``\emph{Adeò}, says
Cortius, ``\emph{in particula elegantissima}'' Allen makes it equivalent to
\emph{eò usque}.}

XXXVIII. For after the powers of the tribunes, in the consulate of
Cneius Pompey and Marcus Crassus, had been fully restored,\footnote{XXXVIII. The powers of the tribunes—had been fully restored
—\emph{Tribunicia potestas restituta}. Before the time of Sylla,
the power of the tribunes had grown immoderate, but Sylla diminished
and almost annihilated it, by taking from them the privileges of
holding any other magistracy after the tribunate, of publicly
addressing the people, of proposing laws, and of listening to appeals.
But in the consulship of Cotta, A.U.C. 679, the first of these
privileges had been restored; and in that of Pompey and Crassus,
A.U.C. 683, the tribunes were reinstated in all their former powers.}
certain young men, of an ardent age and temper, having obtained that
high office,\footnote{Having obtained that high office—\emph{Summam potestatem nacti}.
Cortius thinks these words spurious.} began to stir up the populace by inveighing against
the senate, and proceeded, in course of time, by means of largesses
and promises, to inflame them more and more; by which methods they
became popular and powerful. On the other hand, the most of the
nobility opposed their proceedings to the utmost; under pretense,
indeed, of supporting the senate, but in reality for their own
aggrandizement. For, to state the truth in few words, whatever
parties, during that period, disturbed the republic under plausible
pretexts, some, as if to defend the rights of the people, others, to
make the authority of the senate as great as possible, all, though
affecting concern for the public good, contended every one for his own
interest. In such contests there was neither moderation nor limit;
each party made a merciless use of its successes.

XXXIX. After Pompey, however, was sent to the maritime and Mithridatic
wars, the power of the people was diminished, and the influence of the
few increased. These few kept all public offices, the administration
of the provinces, and every thing else, in their own hands; they
themselves lived free from harm,\footnote{XXXIX. Free from harm—\emph{Innoxii}. In a passive sense.} in flourishing circumstances,
and without apprehension; overawing others, at the same time, with
threats of impeachment,\footnote{Overawing others—with threats of impeachment—\emph{Caeteros
judiciis terrere}. ``Accusationibus et judiciorum periculis.''
\emph{Bernouf}.} so that when in office, they might be
less inclined to inflame the people. But as soon as a prospect of
change, in this dubious state of affairs, had presented itself, the
old spirit of contention awakened their passions; and had Catiline, in
his first battle, come off victorious, or left the struggle undecided,
great distress and calamity must certainly have fallen upon the state,
nor would those, who might at last have gained the ascendency, have
been allowed to enjoy it long, for some superior power would have
wrested dominion and liberty from them when weary and exhausted.

There were some, however, unconnected with the conspiracy, who set out
to join Catiline at an early period of his proceedings. Among these
was Aulus Fulvius, the son of a senator, whom, being arrested on his
journey, his father ordered to be put to death.\footnote{His father ordered to be put to death—\emph{Parens necari jussit}.
``His father put him to death, not by order of the consuls, but by his
own private authority; nor was he the only one who, at the same
period, exercised similar power.'' Dion. Cass., lib. xxxvii. The
father observed on the occasion, that, ``he had begotten him, not for
Catiline against his country, but for his country against Catiline''.
Val. Max., v.8. The Roman laws allowed fathers absolute control over
the lives of their children.} In Rome, at the
same time, Lentulus, in pursuance of Catiline's directions, was
endeavoring to gain over, by his own agency or that of others, all
whom he thought adapted, either by principles or circumstances, to
promote an insurrection; and not citizens only, but every description
of men who could be of any service in war.

XL. He accordingly commissioned one Publius Umbrenus to apply to
certain deputies of the Allobroges,\footnote{XL. Certain deputies of the Allobroges—\emph{Legatos Allobrogum}.
Plutarch, in his Life of Cicero, says that there were then at Rome
\emph{two deputies} from this Gallic nation, sent to complain of oppression
on the part of the Roman governors.} and to lead them, if he
could, to a participation in the war; supposing that as they were
nationally and individually involved in debt, and as the Gauls were
naturally warlike, they might easily be drawn into such an enterprise.
Umbrenus, as he had traded in Gaul, was known to most of the chief men
there, and personally acquainted with them; and consequently, without
loss of time, as soon as he noticed the deputies in the Forum, he
asked them, after making a few inquiries about the state of their
country, and affecting to commiserate its fallen condition, ``what
termination they expected to such calamities?'' When he found that they
complained of the rapacity of the magistrates, inveighed against the
senate for not affording them relief, and looked to death as the only
remedy for their sufferings, ``Yet I,'' said he, ``if you will but act as
men, will show you a method by which you may escape these pressing
difficulties.'' When he had said this, the Allobroges, animated with
the highest hopes, besought Umbrenus to take compassion on them;
saying that there was nothing so disagreeable or difficult, which they
would not most gladly perform, if it would but free their country from
debt. He then conducted them to the house of Decimus Brutus, which was
close to the Forum, and, on account of Sempronia, not unsuitable to
his purpose, as Brutus was then absent from Rome.\footnote{As Brutus was then absent from Borne—\emph{Nam tum Brutus ab
Româ, aberat}. From this remark, say Zanchius and Omnibonus, it is
evident that Brutus was not privy to the conspiracy. ``What sort of
woman \emph{Sempronia} was, has been told in c. 25. Some have thought that
she was the wife of Decimus Brutus; but since Sallust speaks of her as
being in the decay of her beauty at the time of the conspiracy, and
since Brutus, as may be seen in Caesar (B. G. vii., sub fin.), was
then very young, it is probable that she had only an illicit
connection with him, but had gained such an ascendency over his
affections, by her arts of seduction, as to induce him to make her his
mistress, and to allow her to reside in his house.'' \emph{Beauzée}. I have,
however, followed those who think that Brutus was the husband of
Sempronia. Sallust (c. 24), speaking of the woman, of whom Sempronia
was one, says that Catiline \emph{credebat posse—viros earum vel adjungere
sibi, vel interficere}. The truth, on such a point, is of little
importance.} In order, too,
to give greater weight to his representations, he sent for Gabinius,
and, in his presence, explained the objects of the conspiracy, and
mentioned the names of the confederates, as well as those of many
other persons, of every sort, who were guiltless of it, for the
purpose of inspiring the embassadors with greater confidence. At
length, when they had promised their assistance, he let them depart.

XLI. Yet the Allobroges were long in suspense what course they should
adopt. On the one hand, there was debt, an inclination for war, and
great advantages to be expected from victory;\footnote{XLI. To be expected from victory—\emph{In spe victoriae}.} on the other,
superior resources, safe plans, and certain rewards\footnote{Certain rewards—\emph{Certa praemia}. ``Offered by the senate to
those who should give information of the conspiracy. See c. 30.''
\emph{Kuhnhardt}.} instead of
uncertain expectations. As they were balancing these considerations,
the good fortune of the state at length prevailed. They accordingly
disclosed the whole affair, just as they had learned it, to Quintus
Fabius Sanga,\footnote{Quintus Fabius Sanga—``A descendent of that Fabius who, for
having subdued the Allobroges, was surnamed Allobrogicus.'' \emph{Bernouf}.
Whole states often chose patrons as well as individuals.} to whose patronage their state was very greatly
indebted. Cicero, being apprized of the matter by Sanga, directed the
deputies to pretend a strong desire for the success of the plot, to
seek interviews with the rest of the conspirators, to make them fair
promises, and to endeavor to lay them open to conviction as much as
possible.

XLII. Much about the same time there were commotions\footnote{XLII. There were commotions—\emph{Motus erat}. ``\emph{Motus} is also
used by Cicero and Livy in the singular number for \emph{seditiones} and
\emph{tumultus}. No change is therefore to be made in the text.'' \emph{Gerlach}.
``Motus bellicos intelligit, \emph{tumultus}; ut Flor., iii. 13.'' \emph{Cortius}.} in Hither
and Further Gaul, in the Picenian and Bruttian territories, and in
Apulia. For those, whom Catiline had previously sent to those parts,
had begun, without consideration, and seemingly with madness, to
attempt every thing at once; and, by nocturnal meetings, by removing
armor and weapons from place to place, and by hurrying and confusing
every thing, had created more alarm than danger. Of these, Quintus
Metellus Celer, the praetor, having brought several to trial,\footnote{Having brought several to trial—\emph{Complures—caussâ cognitâ.
``Caussum cognoscere} is the legal phrase for examining as to the
authors and causes of any crime.'' \emph{Dietsch}.}
under the decree of the senate, had thrown them into prison, as had
also Caius Muraena in Further Gaul,\footnote{Caius Muraena in Further Gaul—\emph{In Ulteriore Galliâ C. Muraena}.
All the editions, previous to that of Cortius, have \emph{in citeriore
Galliâ}. ``But C. Muraena,'' says the critic, ``commanded in Gallia
Transalpina, or Ulterior Gaul, as appears from Cic. pro Muraena,
c. 41. To attribute such an error to a lapse or memory in Sallust,
would be absurd. I have, therefore, confidently altered \emph{citeriore}
into \emph{ulteriore}.'' The praise of having first discovered the error,
however, is due, not to Cortius, but to Felicius Durantinus, a friend
of Rivius, in whose note on the passage his discovery is recorded.} who governed that province in
quality of legate.

XLIII. But at Rome, in the mean time, Lentulus, with the other leaders
of the conspiracy, having secured what they thought a large force, had
arranged, that as soon as Catiline should reach the neighborhood of
Faesulae, Lucius Bestia, a tribune of the people, having called an
assembly, should complain of the proceedings of Cicero, and lay the
odium of this most oppressive war on the excellent consul;\footnote{XLIII. The excellent consul—\emph{Optimo consuli}. With the
exception of the slight commendation bestowed on his speech,
\emph{luculentam} atque \emph{utilem reipublicae}, c. 31, this is the only
epithet of praise that Sallust bestows on the consul throughout his
narrative. That it could be regarded only as frigid eulogy, is
apparent from a passage in one of Cicero's letters to Atticus (xii.
21), in which he speaks of the same epithet having been applied to him
by Brutus: ``Brutus thinks that he pays me a great compliment when he
calls me an excellent consul (optimum consulem); but what enemy could
speak more coldly of me?''} and
that the rest of the conspirators, taking this as a signal, should, on
the following night, proceed to execute their respective parts.

These parts are said to have been thus distributed. Statilius and
Gabinius, with a large force, were to set on fire twelve places of the
city, convenient for their purpose,\footnote{Twelve places of the city, convenient for their purpose—\emph{Duodecim—opportuna loca}. Plutarch, in his Life of Cicero, says a
hundred places. Few narratives lose by repetition.} at the same time; in order
that, during the consequent tumult,\footnote{In order that, during the consequent tumult—\emph{Quò tumultu}.
``It is best,'' says Dietsch, ``to take \emph{quo} as the \emph{particula finalis}
(to the end that), and \emph{tumultu} as the ablative of the instrument''.} an easier access might be
obtained to the consul, and to the others whose destruction was
intended; Cethegus was to beset the gate of Cicero, and attack him
personally with violence; others were to single out other victims;
while the sons of certain families, mostly of the nobility, were to
kill their fathers; and, when all were in consternation at the
massacre and conflagration, they were to sally forth to join Catiline.

While they were thus forming and settling their plans, Cethegus was
incessantly complaining of the want of spirit in his associates;
observing, that they wasted excellent opportunities through hesitation
and delay;\footnote{Delay—\emph{Dies prolatando}. By putting off from day to day.} that, in such an enterprise, there was need, not of
deliberation, but of action; and that he himself, if a few would
support him, would storm the senate-house while the others remained
inactive. Being naturally bold, sanguine, and prompt to act, he
thought that success depended on rapidity of execution.

XLIV. The Allobroges, according to the directions of Cicero, procured
interviews, by means of Gabinius, with the other conspirators; and
from Lentulus, Cethegus, Statilius, and Cassius, they demanded an
oath, which they might carry under seal to their countrymen, who
otherwise would hardly join in so important an affair. To this the
others consented without suspicion; but Cassius promised them soon to
visit their country,\footnote{XLIV. Soon to visit their country—\emph{Semet eò brevi venturum}.
``It is plain that the adverb relates to what precedes (\emph{ad cives});
and that Cassius expresses an intention to set out for Gaul.'' \emph{Dietsch}.} and, indeed, left the city a little before
the deputies.

In order that the Allobroges, before they reached home, might confirm
their agreement with Catiline, by giving and receiving pledges of
faith, Lentulus sent with them one Titus Volturcius, a native of
Crotona, he himself giving Volturcius a letter for Catiline, of which
the following is a copy:

``Who I am, you will learn from the person whom I have sent to you.
Reflect seriously in how desperate a situation you are placed, and
remember that you are a man.\footnote{Remember that you are a man—\emph{Memineris te virum}. Remember
that you are a man, and ought to act as one. Cicero, in repeating this
letter from memory (Orat. in Cat., iii. 5), gives the phrase, \emph{Cura ut
vir sis}.} Consider what your views demand, and
seek aid from all, even the lowest.'' In addition, he gave him this
verbal message: ``Since he was declared an enemy by the senate, for
what reason should he reject the assistance of slaves? That, in the
city, every thing which he had directed was arranged; and that he
should not delay to make nearer approaches to it.''

XLV. Matters having proceeded thus far, and a night being appointed
for the departure of the deputies, Cicero, being by them made
acquainted with every thing, directed the praetors,\footnote{XLV. The praetors—\emph{Praetoribus} urbanis, the praetors of the city.} Lucius
Valerius Flaccus, and Caius Pomtinus, to arrest the retinue of the
Allobroges, by laying in wait for them on the Milvian Bridge;\footnote{The Milvian Bridge—\emph{Ponte Mulvio}. Now \emph{Ponte Molle}.} he
gave them a full explanation of the object with which they were
sent,\footnote{Of the object with which they were sent—\emph{Rem—cujus gratiâ
mittebantur}.} and left them to manage the rest as occasion might require.
Being military men, they placed a force, as had been directed, without
disturbance, and secretly invested the bridge; when the deputies, with
Volturcius, came to the place, and a shout was raised from each side
of the bridge,\footnote{From each side of the bridge—\emph{Utrinque}. ``Utrinque,'' observes
Cortius, ``glossae MSS. exponunt \_ex utrâque parte pontis,'' and there
is little doubt that the exposition is correct. No translator, however,
before myself, has availed himself of it.} the Gauls, at once comprehending the matter,
surrendered themselves immediately to the praetors. Volturcius, at
first, encouraging his companions, defended himself against numbers
with his sword; but afterward, being unsupported by the Allobroges, he
began earnestly to beg Pomtinus, to whom he was known, to save his
life, and at last, terrified and despairing of safety, he surrendered
himself to the praetors as unconditionally as to foreign enemies.

XLVI. The affair being thus concluded, a full account of it was
immediately transmitted to the consul by messengers. Great anxiety,
and great joy, affected him at the same moment. He rejoiced that, by
the discovery of the conspiracy, the state was freed from danger; but
he was doubtful how he ought to act, when citizens of such eminence
were detected in treason so atrocious. He saw that their punishment
would be a weight upon himself, and their escape the destruction of
the Commonwealth. Having, however, formed his resolution, he ordered
Lentulus, Cethegus, Statilius, Gabinius, and one Quintus Coeparius of
Terracina, who was preparing to go to Apulia to raise the slaves, to
be summoned before him. The others came without delay; but Coeparius,
having left his house a little before, and heard of the discovery of
the conspiracy, had fled from the city. The consul himself conducted
Lentulus, as he was praetor, holding him by the hand, and ordered the
others to be brought into the Temple of Concord, under a guard. Here
he assembled the senate, and in a very full attendance of that body,
introduced Volturcius with the deputies. Hither also he ordered
Valerius Flaccus, the praetor, to bring the box with the letters\footnote{XLVI. The box with the letters—\emph{Scrinium cum literis. Litterae}
may be rendered either \emph{letter} or \emph{letters}. There is no mention made
previously of more letters than that of Lentulus to Catiline, c. 44.
But as it is not likely that the deputies carried a box to convey only
one letter, I have followed other translators by putting the word in
the plural. The oath of the conspirators, too, which was a written
document, was probably in the box.}
which he had taken from the deputies.

XLVII. Volturcius, being questioned concerning his journey, concerning
his letter,\footnote{XLVII. His letter—\emph{Litteris.} His own letter to Catiline, c. 44.
So \emph{praeter litteras} a little below.} and lastly, what object he had had in view,\footnote{What object he had had in view, etc.—\emph{Quid, aut quâ de causâ,
consilli habuisset}. What design he had entertained, and from what
motive \emph{he had entertained it}.} and
from what motives he had acted, at first began to prevaricate,\footnote{To prevaricate.—\emph{Fingere alia.} ``To pretend other things
than what had reference to the conspiracy.'' \emph{Bernouf.}}
and to pretend ignorance of the conspiracy; but at length, when he was
told to speak on the security of the public faith,\footnote{On the security of the public faith—\emph{Fide publicá.}
``Cicero pledged to him the public faith, with the consent of the
senate; or engaged, in the name of the republic, that his life
should be spared, if he would but speak the truth.'' \emph{Bernouf.}} he disclosed
every circumstance as it had really occurred, stating that he had been
admitted as an associate, a few days before, by Gabinius and Coeparius;
that he knew no more than the deputies, only that he used to hear from
Gabinius, that Publius Autronius, Servius Sylla, Lucius Vargunteius,
and many others, were engaged in the conspiracy. The Gauls made a
similar confession, and charged Lentulus, who began to affect ignorance,
not only with the letter to Catiline, but with remarks which he was in
the habit of making, ``that the sovereignty of Rome, by the Sibylline
books, was predestined to three Cornelii; that Cinna and Sylla had ruled
already;\footnote{That Cinna and Sylla had ruled already—\emph{Cinnam atque Syllam
antea.} ``Had ruled,'' or something similar, must be supplied. Cinna
had been the means of recalling Marius from Africa, in conjunction
with whom he domineered over the city, and made it a scene of
bloodshed and desolation.} and that he himself was the third, whose fate it would be
to govern the city; and that this, too, was the twentieth year since the
Capitol was burned; a year which the augurs, from certain omens, had
often said would be stained with the blood of civil war.''

The letter then being read, the senate, when all had previously
acknowledged their seals,\footnote{Their seals—\emph{Signa sua}. ``Leurs cachets, leurs sceaux.''
Bernouf. The Romans tied their letters round with a string, the knot
of which they covered with wax, and impressed with a seal. To open the
letter it was necessary to cut the string: ``\emph{nos linum incidimus}.''
Cic. Or. in Cat. iii. 5. See also C. Nep. Panc. 4, and Adam's \emph{Roman
Antiquities}. The seal of Lentulus had on it a likeness of one of his
ancestors; see Cicero, \emph{loc. cit.}} decreed that Lentulus, being deprived
of his office, should, as well as the rest, be placed in private
custody.\footnote{In private custody—\emph{In liberis custodiis.} Literally, in
``free custody,'' but ``private custody'' conveys a better notion of the
arrangement to the mind of the English reader. It was called \emph{free}
because the persons in custody were not confined in prison. Plutarch
calls it [Greek: \emph{adeomon phylakin}] as also Dion., cap. lviii. 3. See
Tacit. Ann. vi. 8. It was adopted in the case of persons of rank and
consideration.} Lentulus, accordingly, was given in charge to Publius
Lentulus Spinther, who was then aedile; Cethegus, to Quintus
Cornificius; Statilius, to Caius Caesar; Gabinius, to Marcus Crassus;
and Coeparius, who had just before been arrested in his flight, to
Cneius Terentius, a senator.

XLVIII. The common people, meanwhile, who had at first, from a desire
of change in the government, been too much inclined to war, having, on
the discovery of the plot, altered their sentiments, began to execrate
the projects of Catiline, to extol Cicero to the skies; and, as if
rescued from slavery, to give proofs of joy and exultation. Other
effects of war they expected as a gain rather than a loss; but the
burning of the city they thought inhuman, outrageous, and fatal,
especially to themselves, whose whole property consisted in their
daily necessaries and the clothes which they wore.

On the following day, a certain Lucius Tarquinius was brought before
the senate, who was said to have been arrested as he was setting out
to join Catiline. This person, having offered to give information of
the conspiracy, if the public faith were pledged to him,\footnote{XLVIII. If the public faith were pledged to him—\emph{Si fides
publica data, esset}. See c. 47.} and
being directed by the consul to state what he knew, gave the senate
nearly the same account as Volturcius had given, concerning the
intended conflagration, the massacre of respectable citizens, and the
approach of the enemy, adding that ``he was sent by Marcus Crassus to
assure Catiline that the apprehension of Lentulus, Cethegus, and
others of the conspirators, ought not to alarm him, but that he should
hasten, with so much the more expedition to the city, in order to
revive the courage of the rest, and to facilitate the escape of those
in custody''.\footnote{And to facilitate the escape of those in custody—\emph{Et illi
facilius è periculo eriperentur}.} When Tarquinius named Crassus, a man of noble birth,
of very great wealth, and of vast influence, some, thinking the
statement incredible, others, though they supposed it true, yet,
judging that at such a crisis a man of such power\footnote{A man of such power—\emph{Tanta vis hominis}. So great power of
the man.} was rather to
be soothed than irritated (most of them, too, from personal reasons,
being; under obligation to Crassus), exclaimed that he was ``a false
witness,'' and demanded that the matter should be put to the vote.
Cicero, accordingly, taking their opinions, a full senate decreed
``that the testimony of Tarquinius appeared false; that he himself
should be kept in prison; and that no further liberty of speaking\footnote{Liberty of speaking—\emph{Potestatem}. ``Potestatem loquendi.''
\emph{Cyprianus Popma}. As it did not appear that he spoke the truth, the
pledge which the senate had given him, \emph{on condition that he spoke the
truth}, went for nothing; he was not allowed to continue his evidence,
and was sent to prison.}
should be granted him, unless he should name the person at whose
instigation he had fabricated so shameful a calumny.''

There were some, at that time, who thought that this affair was
contrived by Publius Autronius, in order that the interest of Crassus,
if he were accused, might, from participation in the danger, more
readily screen the rest. Others said that Tarquinius was suborned by
Cicero, that Crassus might not disturb the state, by taking upon him,
as was his custom,\footnote{As was his custom—\emph{More suo}. Plutarch, in his Life of Crassus,
relates that frequently when Pompey, Caesar, and Cicero, had refused
to undertake the defense of certain persons, as being unworthy of
their support, Crassus would plead in their behalf; and that he thus
gained great popularity among the common people.} the defense of the criminals. That this attack
on his character was made by Cicero, I afterward heard Crassus himself
assert.

XLIX. Yet, at the same time, neither by interest, nor by solicitation,
nor by bribes, could Quintus Catulus, and Caius Piso, prevail upon
Cicero to have Caius Caesar falsely accused, either by means of the
Allobroges, or any other evidence. Both of these men were at bitter
enmity with Caesar; Piso, as having been attacked by him, when he was
on\footnote{XLIX. Piso, as having been attacked by him, when he was on,
etc.—\emph{Piso oppugnatus in judicio repetundarum propter cujusdam
Transpadani supplicium injustum}. Such is the reading and punctuation
of Cortius. Some editions insert \emph{pecuniarum} before \emph{repetundarum},
and some a comma after it. I have interpreted the passage in
conformity with the explanation of Kritzius, which seems to me the
most judicious that has been offered. \emph{Oppugnatus}, says he, is
equivalent to \emph{gravitur vexatus}, or violently assailed; and Piso was
thus assailed by Caesar on account of his unjust execution of the
Gaul; the words \emph{in judicio repetundarum} merely mark the time when
Caesar's attack was made. While he was on his trial for one thing, he
was attacked by Caesar for another. Gerlach, observing that the words
\emph{in judicio} are wanting in one MS., would emit them, and make
\emph{oppugnatus} govern \emph{pecuniarum repetundarum}, as if it were
\emph{accusatus}; a change which would certainly not improve the passage.
The Galli Transpadani seem to have been much attached to Caesar; see
Cic. Ep. ad Att., v. 2; ad Fam. xvi. 12.} his trial for extortion, on a charge of having illegally put
to death a Transpadane Gaul; Catulus, as having hated him ever since
he stood for the pontificate, because, at an advanced age, and after
filling the highest offices, he had been defeated by Caesar, who was
then comparatively a youth.\footnote{Comparatively a youth—\emph{Adolescentalo}. Caesar was then in
the thirty-third, or, as some say, the thirty-seventh year of his age.
See the note on this word, c. 3.} The opportunity, too, seemed
favorable for such an accusation; for Caesar, by extraordinary
generosity in private, and by magnificent exhibitions in public,\footnote{By magnificent exhibitions in public—\emph{Publicè maximis muneribus}.
Shows of gladiators.}
had fallen greatly into debt. But when they failed to persuade the
consul to such injustice, they themselves, by going from one person to
another, and spreading fictions of their own, which they pretended to
have heard from Volturcius or the Allobroges, excited such violent
odium against him, that certain Roman knights, who were stationed as
an armed guard round the Temple of Concord, being prompted, either by
the greatness of the danger, or by the impulse of a high spirit, to
testify more openly their zeal for the republic, threatened Caesar
with their swords as he went out of the senate-house.

L. While these occurrences were passing in the senate, and while
rewards were being voted, an approbation of their evidence, to the
Allobrogian deputies and to Titus Volturcius, the freedmen, and some
of the other dependents of Lentulus, were urging the artisans and
slaves, in various directions throughout the city,\footnote{L. In various directions throughout the city—\emph{Variis itineribus
—in vicis}. Going hither and thither through the streets.} to attempt his
rescue; some, too, applied to the ringleaders of the mob, who were
always ready to disturb the state for pay. Cethegus, at the same time,
was soliciting, through his agents, his slaves\footnote{Slaves—\emph{Familiam}. ``Servos suos, qui proprie \emph{familia},''
Cortius. \emph{Familia} is a number of \emph{famuli}.} and freedmen, men
trained to deeds of audacity, to collect themselves into an armed
body, and force a way into his place of confinement.

The consul, when he heard that these things were in agitation, having
distributed armed bodies of men, as the circumstances and occasion
demanded, called a meeting of the senate, and desired to know ``what
they wished to be done concerning those who had been committed to
custody.'' A full senate, however, had but a short time before\footnote{A full senate, however, had but a short time before,
etc.—The senate had already decreed that they were enemies to their
country; Cicero now calls a meeting to ascertain what sentence should
be passed on them.}
declared them traitors to their country. On this occasion, Decimus
Junius Silanus, who, as consul elect, was first asked his opinion,
moved\footnote{On this occasion—moved—\emph{Tunc—decreverat}. The \emph{tunc}
(or, as most editors have it, \emph{tum}) must be referred to the second
meeting or the senate, for it does not appear that any proposal
concerning the punishment of the prisoners was made at the first
meeting. There would be no doubt on this point, were it not for the
pluperfect tense, \emph{decreverat}. I have translated it as the perfect.
We must suppose that Sallust had his thoughts on Caesar's speech,
which was to follow, and signifies that all this business \emph{had been
done} before Caesar addressed the house. Kritzius thinks that the
pluperfect was referred by Sallust, not to Caesar's speech, but to the
decree of the senate which was finally made; but this is surely a less
satisfactory method of settling the matter. Sallust often uses the
pluperfect, where his reader would expect the perfect; see, for
instance, \emph{concusserat}, at the beginning of c. 24.} that capital punishment should be inflicted, not only on
those who were in confinement, but also on Lucius Cassius, Publius
Furius, Publius Umbrenus, and Quintus Annius, if they should be
apprehended; but afterward, being influenced by the speech of Caius
Caesar, he said that he would go over to the opinion of Tiberius
Nero,\footnote{That he would go over to the opinion of Tiberius Nero—\emph{Pedibus
in sententian Tib. Neronis—iturum}. Any question submitted to the
senate was decided by the majority of votes, which was ascertained
either by \emph{numeratio}, a counting of the votes, or by \emph{discessio},
when those who were of one opinion, at the direction of the presiding
magistrate, passed over to one side of the house, and those who were
of the contrary opinion, to the other. See Aul. Gell. xiv. 7; Suet.
Tib. 31; Adam's Rom. Ant.; Dr. Smith's Dictionary, Art. \emph{Senatus}.} who had proposed that the guards should be increased, and
that the senate should deliberate further on the matter. Caesar, when
it came to his turn, being asked his opinion by the consul, spoke to
the following effect:

LI. ``It becomes all men,\footnote{LI. It becomes all men, etc.—The beginning of this speech,
attributed to Caesar, is imitated from Demosthenes, [Greek: \emph{Peri ton
hen Chersonaeso pragmaton: Edei men, o andres Athaenaioi, tous
legontas apantas en umin maete pros echthran poieisthai logon maedena,
maete pros charin}]. ``It should be incumbent on all who speak before
you, O Athenians, to advance no sentiment with any view either to
enmity or to favor.''} Conscript Fathers, who deliberate on
dubious matters, to be influenced neither by hatred, affection, anger,
nor pity. The mind, when such feelings obstruct its view, can not
easily see what is right; nor has any human being consulted, at the
same moment, his passion and his interest. When the mind is freely
exerted, its reasoning is sound; but passion, if it gain possession of
it, becomes its tyrant, and reason is powerless.

I could easily mention, Conscript Fathers, numerous examples of kings
and nations, who, swayed by resentment or compassion, have adopted
injudicious courses of conduct; but I had rather speak of these
instances in which our ancestors, in opposition, to the impulse of
passion, acted with wisdom and sound policy.

In the Macedonian war, which we carried on against king Perses, the
great and powerful state of Rhodes, which had risen by the aid of the
Roman people, was faithless and hostile to us; yet, when the war was
ended, and the conduct of the Rhodians was taken into consideration,
our forefathers left them unmolested lest any should say that war was
made upon them for the sake of seizing their wealth, rather than of
punishing their faithlessness. Throughout the Punic war, too, though
the Carthaginians, both during peace and in suspension of arms, were
guilty of many acts of injustice, yet our ancestors never took
occasion to retaliate, but considered rather what was worthy of
themselves, than what might be justly inflicted on their enemies.

Similar caution, Conscript Fathers, is to be observed by yourselves,
that the guilt of Lentulus, and the other conspirators, may not have
greater weight with you than your own dignity, and that you may not
regard your indignation more than your character. If, indeed, a
punishment adequate to their crimes be discovered, I consent to
extraordinary measures;\footnote{I consent to extraordinary measures—\emph{Novum consilium adprobo}.
``That is, I consent that you depart from the usage of your ancestors,
by which Roman citizens were protected from death.'' \emph{Bernouf}.} but if the enormity of their crime
exceeds whatever can be devised,\footnote{Whatever can be devised—\emph{Omnium ingenia}.} I think that we should inflict
only such penalties as the laws have provided.

Most of those, who have given their opinions before me, have
deplored, in studied and impressive language,\footnote{Studied and impressive language—\emph{Compositè atque magnificè.
Compositè}, in language nicely put together; elegantly. \emph{Magnificè},
in striking or imposing terms. \emph{Compositè} is applied to the speech
of Caesar, by Cato, in the following chapter.} the sad fate that
threatens the republic; they have recounted the barbarities of war,
and the afflictions that would fall on the vanquished; they have told
us that maidens would be dishonored, and youths abused; that children
would be torn from the embraces of their parents; that matrons would
be subjected to the pleasure of the conquerors; that temples and
dwelling-houses would be plundered; that massacres and fires would
follow; and that every place would be filled with arms, corpses,
blood, and lamentation. But to what end, in the name of the eternal
gods! was such eloquence directed? Was it intended to render you
indignant at the conspiracy? A speech, no doubt, will inflame him whom
so frightful and monstrous a reality has not provoked! Far from it:
for to no man does evil, directed against himself, appear a light
matter; many, on the contrary, have felt it more seriously than was
right.

But to different persons, Conscript Fathers, different degrees of
license are allowed. If those who pass a life sunk in obscurity,
commit any error, through excessive anger, few become aware of it, for
their fame is as limited as their fortune; but of those who live
invested with extensive power, and in an exalted station, the whole
world knows the proceedings. Thus in the highest position there is the
least liberty of action; and it becomes us to indulge neither
partiality nor aversion, but least of all animosity; for what in
others is called resentment, is in the powerful termed violence and
cruelty.

I am indeed of opinion, Conscript Fathers, that the utmost degree of
torture is inadequate to punish their crime; but the generality of
mankind dwell on that which happens last, and, in the case of
malefactors, forget their guilt, and talk only of their punishment,
should that punishment have been inordinately severe. I feel assured,
too, that Decimus Silanus, a man of spirit and resolution, made the
suggestions which he offered, from zeal for the state, and that he had
no view, in so important a matter, to favor or to enmity; such I know
to be his character, and such his discretion.\footnote{Such I know to be his character, such his discretion—\emph{Eos
mores, eam modestiam viri cognovi}. I have translated \emph{modestiam,
discretion}, which seems to be the proper meaning of the word. Beauzée
renders it \emph{prudence}, and adds a note upon it, which may be worth
transcription. ``I translate \emph{modestia},'' says he, ``by \emph{prudence}, and
think myself authorized to do so. \emph{Sic definitur a Stoicis}, says
Cicero (De Off. i. 40), \emph{ut modestia sit sicentia earum rerum, quae
agentur, aut dicentur, loco suo collocandarum}; and shortly afterward,
\emph{Sic fit ut modestia scientia sit opportunitatis idoneorum ad agendum
temporum}. And what is understood in French by prudence? It is,
according to the Dictionary of the Academy, 'a virtue by which we
discern and practice what is proper in the conduct of life.' This is
almost a translation of the words of Cicero''.} Yet his proposal
appears to me, I will not say cruel (for what can be cruel that is
directed against such characters?), but foreign to our policy. For
assuredly, Silanus, either your fears, or their treason, must have
induced you, a consul elect, to propose this new kind of punishment.
Of fear it is unnecessary to speak, when by the prompt activity of
that distinguished man our consul, such numerous forces are under
arms; and as to the punishment, we may say, what is indeed the truth,
that in trouble and distress, death is a relief from suffering, and
not a torment;\footnote{That—death is a relief from suffering, not a torment, etc.
—This Epicurean doctrine prevailed very much at Rome in Caesar's, and
afterward. We may very well suppose Caesar to have been a sincere
convert to it. Cato alludes to this passage in the speech which
follows; as also Cicero, in his fourth Oration against Catiline, c. 4.
See, for opinions on this point, the first book of Cicero's Tusculan
Questions.} that it puts an end to all human woes; and that,
beyond it, there is no place either for sorrow or joy.

But why, in the name of the immortal gods, did you not add to your
proposal, Silanus, that, before they were put to death, they should be
punished with the scourge? Was it because the Porcian law\footnote{The Porcian Law—\emph{Lex Porcia}. A law proposed by P. Porcius
Loeca, one of the tribunes, A.U.O. 454, which enacted that no one
should bind, scourge or kill a Roman citizen. See Liv., x. 9; Cic.
pro. Rabir., 3, 4: Verr., v 63; de Rep., ii, 31.} forbids
it? But other laws\footnote{Other laws—\emph{Aliae leges}. So Caesar says below, ``Tum lex
Porcia aliaeque paratae, quibus legibus auxilium damnatis permissum;''
what other laws these were is uncertain. One of them, however, was the
Sempronian law, proposed by Caius Gracchus, which ordained that
sentence should not be passed on the life of a Roman citizen without
the order of the people. See Cic. pro Rabir. 4. So ``O lex Porcia
legesque Semproniae!'' Cic. in. Verr., v. 63.} forbid condemned citizens to be deprived of
life, and allow them to go into exile. Or was it because scourging is
a severer penalty than death? Yet what can be too severe, or too
harsh, toward men convicted of such an offense? But if scourging be a
milder punishment than death, how is it consistent to observe the law
as to the smaller point, when you disregard it as to the greater?

But who it may be asked, will blame any severity that shall be
decreed against these parricides\footnote{Parricides—See c. 14, 32.} of their country? I answer that
time, the course of events,\footnote{The course of events—\emph{Dies}. ``Id est, temporis momentum
(\emph{der veränderte Zeitpunkt}).'' \emph{Dietsch}. Things change, and that
which is approved at one period, is blamed at another. \emph{Tempus} and
\emph{dies} are sometimes joined (Liv., xxii. 39, ii. 45), as if not only
time in general, but particular periods, as \emph{from day to day}, were
intended.} and fortune, whose caprice governs
nations, may blame it. Whatever shall fall on the traitors, will fall
on them justly; but it is for you, Conscript Fathers, to consider well
what you resolve to inflict on others. All precedents productive of
evil effects,\footnote{All precedents productive of evil effects—\emph{Omnia mala exempla}.
Examples of severe punishments are meant.} have had their origin from what was good; but when
a government passes into the hands of the ignorant or unprincipled,
any new example of severity,\footnote{Any new example of severity, etc.—\emph{Novum illud exemplum ab
dignis et idoneis ad indignos et non idoneos transferetur}. Gerlach,
Kritzius, Dietsch, and Bernouf, agree to giving to this passage the
sense which is given in the translation. \emph{Digni} and \emph{idonei} are
here used in a bad sense, for \emph{digni et idonei qui poena afficiantur},
deserving and fit objects for punishment.} inflicted on deserving and suitable
objects, is extended to those that are improper and undeserving of it.
The Lacedaemonians, when they had conquered the Athenians,\footnote{When they had conquered the Athenians—At the conclusion of
the Peloponnesian war.}
appointed thirty men to govern their state. These thirty began their
administration by putting to death, even without a trial, all who were
notoriously wicked, or publicly detestable; acts at which the people
rejoiced, and extolled their justice. But afterward, when their
lawless power gradually increased, they proceeded, at their pleasure,
to kill the good and the bad indiscriminately, and to strike terror
into all; and thus the state, overpowered and enslaved, paid a heavy
penalty for its imprudent exultation.

Within our own memory, too, when the victorious Sylla ordered
Damasippus,\footnote{Damasippus—``He, in the consulship of Caius Marius, the younger,
and Cneius Carbo, was city praetor, and put to death some of the most
eminent senators, a short time before the victory of Sylla. See Vell.
Paterc. ii. 26.'' \emph{Bernouf}.} and others of similar character, who had risen by
distressing their country, to be put to death, who did not commend the
proceeding? All exclaimed that wicked and factious men, who had
troubled the state with their seditious practices, had justly
forfeited their lives. Yet this proceeding was the commencement of
great bloodshed. For whenever anyone coveted the mansion or villa, or
even the plate or apparel of another, he exerted his influence to have
him numbered among the proscribed. Thus they, to whom the death of
Damasippus had been a subject of joy, were soon after dragged to death
themselves; nor was there any cessation of slaughter, until Sylla had
glutted all his partisans with riches.

Such excesses, indeed, I do not fear from Marcus Tullius, or in these
times. But in a large state there arise many men of various
dispositions. At some other period, and under another consul, who,
like the present, may have an army at his command, some false
accusation may be credited as true; and when, with our example for a
precedent, the consul shall have drawn the sword on the authority of
the senate, who shall stay its progress, or moderate its fury?

Our ancestors, Conscript Fathers, were never deficient in conduct or
courage; nor did pride prevent them from imitating the customs of
other nations, if they appeared deserving of regard. Their armor, and
weapons of war, they borrowed from the Samnites; their ensigns of
authority,\footnote{Ensigns of authority—\emph{Insignia magistratum}. ``The fasces and
axes of the twelve lictors, the robe adorned with purple, the curule
chair, and the ivory scepter. For the Etrurians, as Dionysius
Halicarnassensis relates, having been subdued, in a nine years' war,
by Tarquinius Priscus, and having obtained peace on condition of
submitting to him as their sovereign, presented him with the
\emph{insignia} of their own monarchs. See Strabo, lib. V.; Florus, i. 5,''
\emph{Kuhnhardt}.} for the most part, from the Etrurians; and, in short,
whatever appeared eligible to them, whether among allies or among
enemies, they adopted at home with the greatest readiness, being more
inclined to emulate merit than to be jealous of it. But at the same
time, adopting a practice from Greece, they punished their citizens
with the scourge, and inflicted capital punishment on such as were
condemned. When the republic, however, became powerful, and faction
grew strong from the vast number of citizens, men began to involve the
innocent in condemnation, and other like abuses were practiced; and it
was then that the Porcian and other laws were provided, by which
condemned citizens were allowed to go into exile. This lenity of our
ancestors, Conscript Fathers, I regard as a very strong reason why we
should not adopt any new measures of severity. For assuredly there was
greater merit and wisdom in those, who raised so mighty an empire from
humble means, than in us, who can scarcely preserve what they so
honorably acquired. Am I of opinion, then, you will ask, that the
conspirators should be set free, and that the army of Catiline should
thus be increased? Far from it; my recommendation is, that their
property be confiscated, and that they themselves be kept in custody
in such of the municipal towns as are best able to bear the
expense;\footnote{Best able to bear the expense—\emph{Maxime opibus valent}. Are
possessed of most resources.} that no one hereafter bring their case before the
senate, or speak on it to the people; and that the senate now give
their opinion, that he who shall act contrary to this, will act
against the republic and the general safety.''

LII. When Caesar had ended his speech, the rest briefly expressed
their assent,\footnote{LII. The rest briefly expressed their assent, etc.—\emph{Caeteri
verbo, alius alii, varie assentiebantur. Verbo assentiebantur}
signifies that they expressed their assent merely by a word or two,
as \emph{assentior Silano, assentior Tiberio Neroni, aut Caesari}, the
three who had already spoken. \emph{Varie}, ``in support of their different
proposals.''} some to one speaker, and some to another, in
support of their different proposals; but Marcius Porcius Cato, being
asked his opinion, made a speech to the following purport:

``My feelings, Conscript Fathers, are extremely different,\footnote{My feelings, Conscript Fathers, are extremely different,
etc.—\emph{Longe mihi alia mens est, P. C.}, etc. The commencement of
Cato's speech is evidently copied from the beginning of the third
Olynthiac of Demosthenes: [Greek: \emph{Ouchi tauta paristatai moi
ginoskein, o andres Athaenaioi, otan te eis ta pragmata apoblepso kai
otan pros tous logous ous akouo tous men gar logous peri tou
timoraesasthai Philippon oro gignomenous, ta de pragmata eis touto
proaekonta oste opos mae peisometha autoi proteron kakos skepsasthai
deon}.] ``I am by no means affected in the same manner. Athenians, when
I review the state of our affairs, and when I attend to those speakers
who have now declared their sentiments. They insist that we should
punish Philip; but our affairs, situated as they now appear, warn us
to guard against the dangers with which we ourselves are threatened.''
\emph{Leland}.} when I
contemplate our circumstances and dangers, and when I revolve in my
mind the sentiments of some who have spoken before me. Those speakers,
as it seems to me, have considered only how to punish the traitors who
have raised war against their country, their parents, their altars,
and their homes;\footnote{Their altars and their homes—\emph{Aris atque focis suis.}
``When \emph{arae} and \emph{foci} are joined, beware of supposing that they are
to be distinguished as referring the one (\_arae) to the public
temples, and the other (\emph{foci}) to private dwellings. Both are to be
understood of private houses, in which the \emph{ara} belonged to the \emph{Dii
Penates}, and was placed in the \emph{impluvium} in the inner part of the
house; the \emph{focus} was dedicated to the \emph{lares}, and was in the hall.''
Ernesti, Clav. Cic., sub. v. \emph{Ara}. Of the commentators on Sallust,
Kritzius is, I believe, the only one who has concurred in this notion
of Ernesti; Langins and Dietsch (with Cortius) adhere to the common
opinion that \emph{arae} are the public altars. Dietsch refers, for a
complete refutation of Ernesti, to G. A. B. Hertzberg \emph{de Diis
Romanorum Penatibus}, Halae, 1840, p. 64; a book which I have not
seen. Certainly, in the observation of Cicero ad Att., vii. 11, ``Non
est respublica in parietibus, sed in aris et focis,'' \emph{arae} must be
considered (as Schiller observes) to denote the public altars and
national religion. See Schiller's Lex. v. \emph{Ara}.} but the state of affairs warns us rather to
secure ourselves against them, than to take counsel as to what
sentence we should pass upon them. Other crimes you may punish after
they have been committed; but as to this, unless you prevent its
commission, you will, when it has once taken effect, in vain appeal to
justice.\footnote{In vain appeal to justice—\emph{Frusta judicia implores. Judicia},
trials, to procure the inflictions of legal penalties.} When the city is taken, no power is left to the
vanquished. But, in the name of the immortal gods, I call upon you,
who have always valued your mansions and villas, your statues and
pictures, at a higher price than the welfare of your country; if you
wish to preserve those possessions, of whatever kind they are, to
which you are attached; if you wish to secure quiet for the enjoyment
of your pleasures, arouse yourselves, and act in defense of your
country. We are not now debating on the revenues, or on injuries done
to our allies, but our liberty and our life is at stake.

Often, Conscript Fathers, have I spoken at great length in this
assembly; often have I complained of the luxury and avarice of our
citizens, and, by that very means, have incurred the displeasure of
many. I, who never excused to myself, or to my own conscience, the
commission of any fault, could not easily pardon the misconduct,\footnote{Could not easily pardon the misconduct, etc.—\emph{Haud facile
alterius lubidini malefacta condonabam}. ``Could not easily forgive the
licentiousness of another its evil deeds.''}
or indulge the licentiousness, of others. But though you little
regarded my remonstrances, yet the republic remained secure; its own
strength\footnote{Yet the republic remained secure; its own strength, etc.
—\emph{Tamen respublica firma, opulentia neglegentiam tolerabat}. This is
Cortius's reading; some editors, as Havercamp, Kritzius, and Dietsch,
insert \emph{erat} after \emph{firma}. Whether \emph{opulentia} is the nominative or
ablative, is disputed. ``\emph{Opulentia},'' says Allen, ``casum sextum
intellige, et repete \emph{respublica} (ad \emph{tolerabat}).'' ``\emph{Opulentia},''
says Kritzius, ``melius nominativo capiendum videtur; nam quae
sequuntur verba novam enunciationem efficiunt.'' I have preferred to
take it as a nominative.} was proof against your remissness. The question, however,
at present under discussion, is not whether we live in a good or a bad
state of morals; nor how great, or how splendid, the empire of the
Roman people is; but whether these things around us, of whatever value
they are, are to continue our own, or to fall, with ourselves, into the
hands of the enemy.

In such a case, does any one talk to me of gentleness and compassion?
For some time past, it is true, we have lost the real name of things;
\footnote{We have lost the real names of things, etc.—Imitated from
Thucydides, iii. 32: [Greek: \emph{Kai taen eiothuian axiosin ton onomaton
es ta erga antaellaxan tae dikaiosei. Tolma men gar alogistos, andria
philetairos enomisthae, mellasis te promaethaes, deilia euprepaes to
de sophron. Tou anandrou proschaema, kai to pros apan syneton, epi pan
argon}.] ``The ordinary meaning of words was changed by them as they
thought proper. For reckless daring was regarded as courage that was
true to its friends; prudent delay, as specious cowardice; moderation,
as a cloak for unmanliness; being intelligent in every thing, as being
useful for nothing.'' \emph{Dale's} translation; Bohn's Classical Library.} for to lavish the property of others is called generosity, and
audacity in wickedness is called heroism; and hence the state is reduced
to the brink of ruin. But let those, who thus misname things, be liberal,
since such is the practice, out of the property of our allies; let them
be merciful to the robbers of the treasury; but let them not lavish our
blood, and, while they spare a few criminals, bring destruction on all
the guiltless.

Caius Caesar, a short time ago, spoke in fair and elegant language,
\footnote{Elegant language—\emph{Compositè}. See above, c. 51.} before this assembly, on the subject of life and death; considering
as false, I suppose, what is told of the dead; that the bad, going a
different way from the good, inhabit places gloomy, desolate, dreary, and
full of horror. He accordingly proposed \emph{that the property of the
conspirators should be confiscated, and themselves kept in custody in
the municipal towns}; fearing, it seems, that, if they remain at Rome,
they may be rescued either by their accomplices in the conspiracy, or by
a hired mob; as if, forsooth, the mischievous and profligate were to be
found only in the city, and not through the whole of Italy, or as if
desperate attempts would not be more likely to succeed where there is
less power to resist them. His proposal, therefore, if he fears any
danger from them, is absurd; but if, amid such universal terror, he
alone is free from alarm, it the more concerns me to fear for you and
myself.

Be assured, then, that when you decide on the fate of Lentulus and
the other prisoners, you at the same time determine that of the army
of Catiline, and of all the conspirators. The more spirit you display
in your decision, the more will their confidence be diminished; but if
they shall perceive you in the smallest degree irresolute, they will
advance upon you with fury.

Do not suppose that our ancestors, from so small a commencement,
raised the republic to greatness merely by force of arms. If such had
been the case, we should enjoy it in a most excellent condition;\footnote{In a most excellent condition—\emph{Multo pulcherrumam.} See c. 36.}
for of allies and citizens,\footnote{For of allies and citizens, etc.—Imitated from Demosthenes,
Philipp. III.4.} as well as arms and horses, we have a
much greater abundance than they had. But there were other things
which made them great, but which among us have no existence; such as
industry at home, equitable government abroad, and minds impartial in
council, uninfluenced by any immoral or improper feeling. Instead of
such virtues, we have luxury and avarice; public distress, and private
superfluity; we extol wealth, and yield to indolence; no distinction
is made between good men and bad; and ambition usurps the honors due
to virtue. Nor is this wonderful; since you study each his individual
interest, and since at home you are slaves to pleasure, and here to
money or favor; and hence it happens that an attack is made on the
defenseless state.

But on these subjects I shall say no more. Certain citizens, of the
highest rank, have conspired to ruin their country; they are engaging
the Gauls, the bitterest foes of the Roman name, to join in a war
against us; the leader of the enemy is ready to make a descent upon
us; and do you hesitate; even in such circumstances, how to treat
armed incendiaries arrested within your walls? I advise you to have
mercy upon them;\footnote{I advise you to have mercy upon them—\emph{Misereamini censeo,
i.e.}, censeo \emph{ut} misereanum, spoken ironically. Most translators
have taken the words in the sense of ``You would take pity on them, I
suppose,'' or something similar.} they are young men who have been led astray by
ambition; send them away, even with arms in their hands. But such
mercy, and such clemency, if they turn those arms against you, will
end in misery to yourselves. The case is, assuredly, dangerous, but
you do not fear it; yes, you fear it greatly, but you hesitate how to
act, through weakness and want of spirit, waiting one for another, and
trusting to the immortal gods, who have so often preserved your
country in the greatest dangers. But the protection of the gods is not
obtained by vows and effeminate supplications; it is by vigilance,
activity, and prudent measures, that general welfare is secured. When
you are once resigned to sloth and indolence, it is in vain that you
implore the gods; for they are then indignant and threaten vengeance.

In the days of our forefathers, Titus Manlius Torquatus, during a war
with the Gauls, ordered his own son to be put to death, because he had
fought with an enemy contrary to orders. That noble youth suffered for
excess of bravery; and do you hesitate what sentence to pass on the
most inhuman of traitors? Perhaps their former life is at variance
with their present crime. Spare, then, the dignity of Lentulus, if he
has ever spared his own honor or character, or had any regard for gods
or for men. Pardon the youth of Cethegus, unless this be the second
time that he has made war upon his country.\footnote{Unless this be the second time that he has made war upon
his country—``Cethegus first made war on his country in conjunction
with Marius.'' \emph{Bernouf}. Whether Sallust alludes to this, or intimates
(as Gerlach thinks) that he was engaged in the first conspiracy, is
doubtful.} As to Gabinius,
Slatilius, Coeparius, why should I make any remark upon them? Had they
ever possessed the smallest share of discretion, they would never have
engaged in such a plot against their country.

In conclusion, Conscript Fathers, if there were time to amend an
error, I might easily suffer you, since you disregard words, to be
corrected by experience of consequences. But we are beset by dangers on
all sides; Catiline, with his army, is ready to devour us;\footnote{Is ready to devour us—\emph{Faucibus urget}. Cortius, Kritzius,
Gerlach, Burnouf, Allen, and Dietsch, are unanimous in interpreting
this as a metaphorical expression, alluding to a wild beast with open
jaws ready to spring upon its prey. They support this interpretation
by Val. Max., v. 3: ``Faucibus apprehensam rempublicam;'' Cic. pro.
Cluent., 31: ``Quum faucibus premetur;'' and Plaut. Casin. v. 3,4,
``Manifesto faucibus teneor.'' Some, editors have read \emph{in faucibus},
and understood the words as referring to the jaws or narrow passes of
Etruria, where Catiline was with his army.} while
there are other enemies within the walls, and in the heart of the
city; nor can any measures be taken, or any plans arranged, without
their knowledge. The more necessary is it, therefore, to act with
promptitude. What I advise, then, is this: that since the state, by a
treasonable combination of abandoned citizens, has been brought into
the greatest peril; and since the conspirators have been convicted on
the evidence of Titus Volturcius, and the deputies of the Allobroges,
and on their own confession, of having concerted massacres,
conflagrations, and other horrible and cruel outrages, against their
fellow-citizens and their country, punishment be inflicted, according
to the usage of our ancestors, on the prisoners who have confessed
their guilt, as on men convicted of capital crimes.''

LIII. When Cato had resumed his seat, all the senators of consular
dignity, and a great part of the rest,\footnote{LIII. All the senators of consular dignity, and a great
part of the rest—\emph{Consulares omnes, itemque senatus magna pars}. ``As
the consulars were senators, the reader would perhaps expect Sallust
to have said \emph{reliqui senatús} but \emph{itemque} is equivalent to \emph{et
praeter eos}.'' \emph{Dietsch}.} applauded his opinion, and
extolled his firmness of mind to the skies. With mutual reproaches,
they accused one another of timidity, while Cato was regarded as the
greatest and noblest of men; and a decree of the senate was made as he
had advised.

After reading and hearing of the many glorious achievements which the
Roman people had performed at home and in the field, by sea as well as
by land, I happened to be led to consider what had been the great
foundation of such illustrious deeds. I knew that the Romans had
frequently, with small bodies of men, encountered vast armies of the
enemy; I was aware that they had carried on wars\footnote{That they had carried on wars—\emph{Bella gesta}. That wars had
been carried on \emph{by them}.} with limited
forces against powerful sovereigns; that they had often sustained,
too, the violence of adverse fortune; yet that, while the Greeks
excelled them in eloquence, the Gauls surpassed them in military
glory. After much reflection, I felt convinced that the eminent virtue
of a few citizens had been the cause of all these successes; and hence
it had happened that poverty had triumphed over riches, and a few over
a multitude. And even in later times, when the state had become
corrupted by luxury and indolence, the republic still supported
itself, by its own strength, under the misconduct of its generals and
magistrates; when, as if the parent stock were exhausted,\footnote{As if the parent stock were exhausted—\emph{Sicuti effoeta
parentum}. This is the reading of Cortius, which he endeavors to
explain thus: ``Ac sicuti \emph{effoeta parens}, inter parentes, \emph{sese
habere solet}, ut nullos amplius liberas proferat, sic Roma sese
habuit, ubi multis tempestatibus nemo virtute magnus fuit.'' ``\emph{Est},''
he adds, ``or \emph{solet esse}, or \emph{sese habere solet}, may very well be
understood from the \emph{fuit} which follows.'' But all this only serves to
show what a critic may find to say in defense of a reading to which he
is determined to adhere. All the MSS., indeed, have \emph{parentum}, except
one, which has \emph{parente}. Dietsch thinks that some word has been lost
between \emph{effoeta} and \emph{parentum}, and proposes to read \_sicuti effoetá
aetate parentum, with the sense, \emph{as if the age of the parents were
too much exhausted to produce strong children}. Kritzius, from a
suggestion of Cortius (or rather of his predecessor, Rupertus), reads
\emph{effoetae parentum} (the effoetae agreeing with Romae which follows),
considering the sense to be the same as as \emph{effoetae parentis}—as
\emph{divina dearum} for \emph{divina dea}, etc. Gerlach retains the rending of
Cortius, and adopts his explanation (4to. ed., 1827), but says that
the \emph{explicatio} may seem \emph{durior}, and that it is doubtful whether we
ought not to have recourse to the \emph{effoeta parente} of the old critics.
Assuredly if we retain \emph{parentum}, \emph{effoetae} is the only reading that
we can well put with it. We may compare with it \emph{loca nuda gignentium},
(Jug. c. 79), i.e. ``places bare of objects producing any thing.''
Gronovius know not what to do with the passage, called it \emph{locus
intellectus nemini}, and at last decided on understanding \emph{virtute}
with \emph{effoetae parentum}, which, \emph{pace tarti viri}, and although Allen
has followed him, is little better than folly. The concurrence of the
majority of manuscripts in giving \emph{parentum} makes the scholar
unwilling to set it aside. However, as no one has explained it
satisfactorily even to himself, I have thought it better, with Dietsch,
to regard it a \emph{scriptura non ferenda}, and to acquiesce, with
Glareanus, Rivius, Burnouf, and the Bipont edition, in the reading
\emph{effoetâ parente}.} there
was certainly not produced at Rome, for many years, a single citizen
of eminent ability. Within my recollection, however, there arose two
men of remarkable powers, though of very different character, Marcus
Cato and Caius Caesar, whom, since the subject has brought them before
me, it is not my intention to pass in silence, but to describe, to the
best of my ability, the disposition and manners of each.

LIV. Their birth, age, and eloquence, were nearly on an equality;
their greatness of mind similar, as was also their reputation, though
attained by different means.\footnote{LIV. Though attained by different means—\emph{Sed alia alii}.
``Alii alia \emph{gloria},'' for \emph{altera alteri}. So Livy, i. 21: \emph{Duo
reges}, alius alia via.} Caesar grew eminent by generosity
and munificence; Cato by the integrity of his life. Caesar was
esteemed for his humanity and benevolence; austereness had given
dignity to Cato. Caesar acquired renown by giving, relieving, and
pardoning; Cato by bestowing nothing. In Caesar, there was a refuge
for the unfortunate; in Cato, destruction for the bad. In Caesar, his
easiness of temper was admired; in Cato, his firmness. Caesar, in
fine, had applied himself to a life of energy and activity; intent
upon the interest of his friends, he was neglectful of his own; he
refused nothing to others that was worthy of acceptance, while for
himself he desired great power, the command of an army, and a new war
in which his talents might be displayed. But Cato's ambition was that
of temperance, discretion, and, above all, of austerity; he did not
contend in splendor with the rich, or in faction with the seditious,
but with the brave in fortitude, with the modest in simplicity,\footnote{Simplicity—\emph{Pudore}. The word here seems to mean the absence
of display and ostentation.}
with the temperate\footnote{With the temperate—\emph{Cum innocente}. ``That is \emph{cum integro
et abstinente}. For \emph{innocentia} is used for \emph{abstinentia}, and
opposed to \emph{avaritia}. See Cic. pro Lego Manil., c. 13.'' \emph{Burnouf}.} in abstinence; he was more desirous to be,
than to appear, virtuous; and thus, the less he courted popularity,
the more it pursued him.

LV. When the senate, as I have stated, had gone over to the opinion of
Cato, the counsel, thinking it best not to wait till night, which was
coming on, lest any new attempts should be made during the interval,
ordered the triumvirs\footnote{LV. The triumvirs—\emph{Triumviros}. The \emph{triumviri capitales},
who had the charge of the prison and of the punishment of the
condemned. They performed their office by deputy, Val. Max., v. 4. 7.} to make such preparations as the execution
of the conspirators required. He himself, having posted the necessary
guards, conducted Lentulus to the prison; and the same office was
performed for the rest by the praetors. There is a place in the
prison, which is called the Tullian dungeon,\footnote{The Tullian dungeon—\emph{Tullianum}. ``Tullianum'' is an adjective,
with which \emph{robur} must be understood, as it was originally
constructed, wholly or partially, with oak. See Festus, sub voce
\emph{Robum} or \emph{Robur}: his words are \emph{arcis robustis includebatur}, of
which the sense is not very clear. The prison at Rome was built by
Ancus Marcius, and enlarged by Servius Tullius, from whom this part of
it had its name; Varro de L. L., iv. 33. It is now transformed into a
subterranean chapel, beneath a small church erected over it, called
\emph{San Pietro in Carcere}. De Brosses and Eustace both visited it; See
Eustace's Classical Tour, vol. i. p. 260, in the \emph{Family Library}. See
also Wasse's note on this passage.} and which, after a
slight ascent to the left, is sunk about twelve feet under ground.
Walls secure it on every side, and over it is a vaulted roof connected
with stone arches;\footnote{A vaulted roof connected with stone arches—\emph{Camera lapideis
fornicibus vincta}. ``That \emph{camera} was a roof curved in the form of
a \emph{testudo}, is generally admitted; see Vitruv. vii. 3; Varr.,
R. R. iii. 7, init.'' \emph{Dietsch}. The roof is now arched in the usual
way.} but its appearance is disgusting and horrible,
by reason of the filth, darkness, and stench. When Lentulus had been
let down into this place, certain men, to whom orders had been
given,\footnote{Certain men, to whom orders had been given—\emph{Quibus praeceptum
erat}. The editions of Havercamp, Gerlach, Kritzius, and Dietsch,
have \emph{vindices rerum capitalium, quibus}, etc. Cortius ejected the
first three words from his text, as an intruded gloss. If the words
be genuine, we must consider these \emph{vindices} to have been the
deputies, or lictors, of the ``triumvirs'' mentioned above.} strangled him with a cord. Thus this patrician, who was of
the illustrious family of the Cornelii, and who filled the office of
consul at Rome, met with an end suited to his character and conduct.
On Cethegus, Statilius, Gabinius, and Coeparius, punishment was
inflicted in a similar manner.

LVI. During these proceedings at Rome, Catiline, out of the entire
force which he himself had brought with him, and that which Manlius
had previously collected, formed two legions, filling up the cohorts
as far as his number would allow;\footnote{LVI. As far as his numbers would allow—\emph{Pro numero militum}.
He formed his men into two bodies, which he called legions, and
divided each legion, as was usual, into ten cohorts, putting into
each cohort as many men as he could. The cohort of a full legion
consisted of three maniples, or six hundred men; the legion would then
be six thousand men. But the legions were seldom so large as this;
they varied at different periods, from six thousand to three thousand;
in the time of Polybius they were usually four thousand two hundred.
See Adam's Rom. Ant., and Lipsius de Mil. Rom Dial. iv.} and afterward, as any
volunteers, or recruits from his confederates,\footnote{From his confederates—\emph{Ex sociis}. ``Understand, not only
the leaders in the conspiracy, but those who, in c. 35, are said to
have set out to join Catiline, though not at that time exactly
implicated in the plot.'' \emph{Kritzius}. It is necessary to notice this,
because Cortius erroneously supposes ``sociis'' to mean the \emph{allies of
Rome}. Dahl, Longius, Müller, Burnouf, Gerlach, and Dietsch, all
interpret in the same manner as Kritzius.} arrived in his
camp, he distributed them equally throughout the cohorts, and thus
filled up his legions, in a short time, with their regular number of
men, though at first he had not more than two thousand. But, of his
whole army, only about a fourth part had the proper weapons of
soldiers; the rest, as chance had equipped them, carried darts,
spears, or sharpened stakes.

As Antonius approached with his army, Catiline directed his march over
the hills, encamping, at one time, in the direction of Rome, at
another in that of Gaul. He gave the enemy no opportunity of fighting,
yet hoped himself shortly to find one,\footnote{Hoped himself shortly to find one—\emph{Sperabat propediem sese
habiturum}. Other editions, as those of Havercamp, Gerlach, Kritzius,
Dietsch, and Burnouf, have the words \emph{magnas copias} before \emph{sese}.
Cortius struck them out, observing that \emph{copiae} occurred too often in
this chapter, and that in one MS. they were wanting. One manuscript,
however, was insufficient authority for discarding them; and the
phrase suits much better with what follows, \emph{si Romae socii incepta
patravissent}, if they are retained.} if his accomplices at Rome
should succeed in their objects. Slaves, meanwhile, of whom vast numbers
\footnote{Slaves—of whom vast numbers, etc.—\emph{Servitia—cujus magnae
copiae}. ``\emph{Cujus},'' says Priscian (xvii. 20, vol. ii., p. 81, cd. Krehl),
``is referred \emph{ad rem}, that is \emph{cujus rei servitiorum}.'' \emph{Servorum} or
\emph{hominum genus}, is, perhaps, rather what Sallust had in his mind, as
the subject of his relation. Gerlach adduces as an expression most
nearly approaching to Sallust's, Thucyd., iii. 92; [Greek: \emph{Kai dorieis,
hae maetropolis ton Lakedaimonion}].} had at first flocked to him, he continued to reject, not only as
depending on the strength of the conspiracy, but as thinking it impolitic
\footnote{Impolitic—\emph{Alienum suis rationibus}. Foreign to his views;
inconsistent with his policy.} to appear to share the cause of citizens with runagates.

LVII. When it was reported in his camp, however, that the conspiracy
had been discovered at Rome, and that Lentulus, Cethegus, and the rest
whom I have named, had been put to death, most of those whom the hope
of plunder, or the love of change, had led to join in the war, fell
away. The remainder Catiline conducted, over rugged mountains, and by
forced marches, into the neighborhood of Pistoria, with a view to
escape covertly, by cross roads, into Gaul.

But Quintus Metellus Celer, with a force of three legions, had at that
time, his station in Picenum, who suspected that Catiline, from the
difficulties of his position, would adopt precisely the course which
we have just described. When, therefore, he had learned his route from
some deserters, he immediately broke up his camp, and took his post at
the very foot of the hills, at the point where Catiline's descent
would be, in his hurried march into Gaul\footnote{LVII. In his hurried march into Gaul—\emph{In Galliam properanti}.
These words Cortius inclosed in brackets, pronouncing them as a useless
gloss. But all editors have retained them as genuine, except the Bipont
and Burnouf, who wholly omitted them.}. Nor was Antonius far
distant, as he was pursuing, though with, a large army, yet through
plainer ground, and with fewer hinderances, the enemy in retreat.\footnote{As he was pursuing, though with a large army, yet through
plainer ground, and with fewer hinderances; the enemy in
retreat—\emph{Utpote qui magna exercitu, locis aequioribus, expeditus, in
fuga sequeretur}. It would be tedious to notice all that has been
written upon this passage of Sallust. All the editions, before that of
Cortius, had \emph{expeditos, in fugam}, some joining \emph{expeditos} with
\emph{locis aequioribus}, and some with \emph{in fugam}. \emph{Expeditos in fugam}
was first condemned by Wasse, no negligent observer of phrases, who
said that no expression parallel to it could be found in any Latin
writer. Cortius, seeing that the \emph{expedition}, of which Sallust is
speaking, is on the part of Antonius, not of Catiline, altered
\emph{expeditos}, though found in all the manuscripts, into \emph{expeditus};
and \emph{in fugam}, at the same time, into \emph{in fuga}; and in both these
emendations he has been cordially followed by the subsequent editors,
Gerlach, Kritzius, and Dietsch. I have translated \emph{magno exercitu},
``\emph{though} with a large army,'' although, according to Dietsch and some
others, we need not consider a large army as a cause of slowness, but
may rather regard it as a cause of speed; since the more numerous were
Metellus's forces, the less he would care how many he might leave
behind through fatigue, or to guard the baggage; so that he might be
the more \emph{expeditus}, unincumbered. With \emph{sequeretur} we must
understand \emph{hostes}. The Bipont, Burnouf's, which often follows it,
and Havercamp's, are now the only editions of any note that retain
\emph{expeditos in fugam}.}

Catiline, when he saw that he was surrounded by mountains and by
hostile forces, that his schemes in the city had been unsuccessful,
and that there was no hope either of escape or of succor, thinking it
best, in such circumstances, to try the fortune of a battle, resolved
upon engaging, as speedily as possible, with Antonius. Having,
therefore, assembled his troops, he addressed them in the following
manner:

LVIII. ``I am well aware, soldiers, that words can not inspire courage;
and that a spiritless army can not be rendered active,\footnote{LVIII. That a spiritless army can not be rendered active,
etc.—\emph{Neque ex ignavo strenuum, neque fortem ex timido exercitum
oratione imperatoris fieri}. I have departed a little from the literal
reading, for the sake of ease.} or a timid
army valiant, by the speech of its commander. Whatever courage is in
the heart of a man, whether from nature or from habit, so much will be
shown by him in the field; and on him whom neither glory nor danger
can move, exhortation is bestowed in vain; for the terror in his
breast stops his ears.

I have called you together, however, to give you a few instructions,
and to explain to you, at the same time, my reasons for the course
which I have adopted. You all know, soldiers, how severe a penalty the
inactivity and cowardice of Lentulus has brought upon himself and us;
and how, while waiting for reinforcements from the city, I was unable
to march into Gaul.

In what situation our affairs now are, you all understand as well as
myself. Two armies of the enemy, one on the side of Rome, and the
other on that of Gaul, oppose our progress; while the want of corn,
and of other necessaries, prevents us from remaining, however strongly
we may desire to remain, in our present position. Whithersoever we
would go, we must open a passage with our swords. I conjure you,
therefore, to maintain a brave and resolute spirit; and to remember,
when you advance to battle, that on your own right hands depend\footnote{That on your own right hands depend, etc.—\emph{In dextris
portare}. ``That you carry in your right hands.''}
riches, honor, and glory, with the enjoyment of your liberty and of
your country. If we conquer, all will be safe; we shall have
provisions in abundance; and the colonies and corporate towns will
open their gates to us. But if we lose the victory through want of
courage, these same places\footnote{Those same places—\emph{Eadem illa}. ``Coloniae atque municipia
portas claudent.'' \emph{Burnouf}.} will turn against us; for neither
place nor friend will protect him whom his arms have not protected.
Besides, soldiers, the same exigency does not press upon our
adversaries, as presses upon us; we fight for our country, for our
liberty, for our life; they contend for what but little concerns
them,\footnote{They contend for what but little concerns them—\emph{Illis
supervacaneum est pugnare}. It is but of little concern to the great
body of them personally: they may fight, but others will have the
advantages of their efforts.} the power of a small party. Attack them, therefore, with so
much the greater confidence, and call to mind your achievements of
old.

We might,\footnote{We might, etc.—\emph{Licuit nobis}. The editions vary between
\emph{nobis} and \emph{robis}; but most, with Cortius, have \emph{nobis}.} with the utmost ignominy, have passed the rest of our
days in exile. Some of you, after losing your property, might have
waited at Rome for assistance from others. But because such a life, to
men of spirit, was disgusting and unendurable, you resolved upon your
present course. If you wish to quit it, you must exert all your
resolution, for none but conquerors have exchanged war for peace. To
hope for safety in flight, when you have turned away from the enemy
the arms by which the body is defended, is indeed madness. In battle,
those who are most afraid are always in most danger; but courage is
equivalent to a rampart. When I contemplate you, soldiers, and when I
consider your past exploits, a strong hope of victory animates me.
Your spirit, your age, your valor, give me confidence; to say nothing
of necessity, which makes even cowards brave. To prevent the numbers
of the enemy from surrounding us, our confined situation is
sufficient. But should Fortune be unjust to your valor, take care not
to lose your lives unavenged; take care not to be taken and butchered
like cattle, rather than fighting like men, to leave to your enemies a
bloody and mournful victory.''

LIX. When he had thus spoken, he ordered, after a short delay, the
signal for battle to be sounded, and led down his troops, in regular
order, to the level ground. Having then sent away the horses of all
the cavalry, in order to increase the men's courage by making their
danger equal, he himself, on foot, drew up his troops suitably to
their numbers and the nature of the ground. As a plain stretched
between the mountains on the left, with a rugged rock on the right, he
placed eight cohorts in front, and stationed the rest of his force, in
close order, in the rear.\footnote{LIX. In the rear—\emph{In subsidio}. Most translators have
rendered this, ``as a body of reserve;'' but such can not well be the
signification. It seems only to mean the part behind the front:
Catiline places the eight cohorts \emph{in front}, and the rest of his
force \emph{in subsidio}, to support the front. \emph{Subsidia}, according to
Varro (de L. L., iv. 16) and Festus (v. \emph{Subsidium}), was a term
applied to the Triarii, because they \emph{subsidebant}, or sunk down on
one knee, until it was their turn to act. See Sheller's Lex. v.
\emph{Subsidium}. ``Novissimi ordines ita dicuntur.'' \emph{Gerlach}. \emph{In
subsidiis}, which occurs a few lines below, seems to signify \emph{in lines
in the rear}; as in Jug. 49, \emph{triplicibus subsidiis aciem intruxit,
i.e. with three lines behind the front}. ``Subsidium ea pars aciei
vocabatur quae reliquis submitti posset; Caes. B. G., ii. 25.''
\emph{Dietsch}.} From among these he removed all the
ablest centurions,\footnote{All the ablest centurions—\emph{Centuriones omnes lectos}.
``\emph{Lectos} you may consider to be the same as \emph{eximios, praestantes},
centurionum praestantissimum quemque.'' \emph{Kritzius}. Cortius and others
take it for a participle, \emph{chosen}.} the veterans,\footnote{Into the foremost ranks—\emph{In primam aciem}. Whether Sallust
means that he ranged them with the eight cohorts, or only in the first
line of the \emph{subsidia}, is not clear.} and the stoutest of the
common soldiers that were regularly armed, into the foremost
ranks.\footnote{Into the foremost ranks—\emph{In primam aciem}. Whether Sallust
means that he ranged them with the eight cohorts, or only in the first
line of the \emph{subsidia}, is not clear.} He ordered Caius Manlius to take the command on the right,
and a certain officer of Faesulae\footnote{A certain officer of Faesulae—\emph{Faesulanum quemdam}. ``He is
thought to have been that P. Furius, whom Cicero (Cat., iii. 6, 14)
mentions as having been one of the colonists that Sylla settled at
Faesulae, and who was to have been executed, if he had been
apprehended, for having been concerned in corrupting the Allobrogian
deputies.'' \emph{Dietsch}. Plutarch calls this officer Furius.} on the left; while he himself,
with his freedmen\footnote{His freedmen—\emph{Libertis}. ``His own freedmen, whom he probably
had about him as a body-guard, deeming them the most attached of his
adherents. Among them was, possibly, that Sergius, whom we find
from Cic. pro Domo, 5, 6, to have been Catiline's armor bearer.''
\emph{Dietsch}.} and the colonists,\footnote{The colonists—\emph{Colonis}. ``Veterans of Sylla, who had been
settled by him as colonists in Etruria, and who had now been induced
to join Catiline.'' \emph{Gerlach}. See c. 28.} took his station by the
eagle,\footnote{By the eagle—\emph{Propter aquilam}. See Cic. in Cat., i. 9.} which Caius Marius was said to have had in his army in the
Cimbrian war.

On the other side, Caius Antonius, who, being lame,\footnote{Being lame—\emph{Pedibus aeger}. It has been common among
translators to render \emph{pedibus aeger} afflicted with the gout, though
a Roman might surely be lame without having the gout. As the lameness
of Antonius, however, according to Dion Cassius (xxxvii. 39), was only
pretended, it may be thought more probable that he counterfeited the
gout than any other malady. It was with this belief, I suppose, that
the writer of a gloss on one of the manuscripts consulted by Cortius,
interpreted the words, \emph{ultroneam passus est podogram}, ``he was
affected with a voluntary gout.'' Dion Cassius says that he preferred
engaging with Antonius, who had the larger army, rather than with
Metellus, who had the smaller, because he hoped that Antonius would
designedly act in such a way as to lose the victory.} was unable to
be present in the engagement, gave the command of the army to Marcus
Petreius, his lieutenant-general. Petreius, ranged the cohorts of
veterans, which he had raised to meet the present insurrection,\footnote{To meet the present insurrection—\emph{Tumulti causa}. Any sudden
war or insurrection in Italy or Gaul was called \emph{tumultus}. See
Cic. Philipp. v. 12.}
in front, and behind them the rest of his force in lines. Then, riding
round among his troops, and addressing his men by name, he encouraged
them, and bade them remember that they were to fight against unarmed
marauders, in defense of their country, their children, their temples,
and their homes.\footnote{Their temples and their homes—\emph{Aris atque focis suis}. See
c. 52.} Being a military man, and having served with
great reputation, for more than thirty years, as tribune, praefect,
lieutenant, or praetor, he knew most of the soldiers and their
honorable actions, and, by calling these to their remembrance, roused
the spirits of the men.

LX. When he had made a complete survey, he gave the signal with the
trumpet, and ordered the cohorts to advance slowly. The army of the
enemy followed his example; and when they approached so near that the
action could be commenced by the light-armed troops, both sides, with
a loud shout, rushed together in a furious charge.\footnote{LX. In a furious charge—\emph{Infestis siqnis}.} They threw
aside their missiles, and fought only with their swords. The veterans,
calling to mind their deeds of old, engaged fiercely in the closest
combat. The enemy made an obstinate resistance; and both sides
contended with the utmost fury. Catiline, during this time, was
exerting himself with his light troops in the front, sustaining such
as were pressed, substituting fresh men for the wounded, attending to
every exigency, charging in person, wounding many an enemy, and
performing at once the duties of a valiant soldier and a skillful
general.

When Petreius, contrary to his expectation, found Catiline attacking
him with such impetuosity, he led his praetorian cohort against the
centre of the enemy, among whom, being thus thrown into confusion, and
offering but partial resistance,\footnote{Offering but partial resistance—\emph{Alios alibi resistentes}.
Not making a stand in a body, but only some in one place, and some in
another.} he made great slaughter, and
ordered, at the same time, an assault on both flanks. Manlius and the
Faesulan, sword in hand, were among the first\footnote{Among the first, etc.—\emph{In primis pugnantes cadunt}. Cortius
very properly refers \emph{in primis} to \emph{cadunt}.} that fell; and
Catiline, when he saw his army routed, and himself left with but few
supporters, remembering his birth and former dignity, rushed into the
thickest of the enemy, where he was slain, fighting to the last.

LXI. When the battle was over, it was plainly seen what boldness, and
what energy of spirit, had prevailed throughout the army of Catiline;
for, almost every where, every soldier, after yielding up his breath,
covered with his corpse the spot which he had occupied when alive. A
few, indeed, whom the praetorian cohort had dispersed, had fallen
somewhat differently, but all with wounds in front. Catiline himself
was found, far in advance of his men, among the dead bodies of the
enemy; he was not quite breathless, and still expressed in his
countenance the fierceness of spirit which he had shown during his
life. Of his whole army, neither in the battle, nor in flight, was any
free-born citizen made prisoner, for they had spared their own lives
no more than those of the enemy.

Nor did the army of the Roman people obtain a joyful or bloodless
victory; for all their bravest men were either killed in the battle,
or left the field severely wounded.

Of many who went from the camp to view the ground, or plunder the
slain, some, in turning over the bodies of the enemy, discovered a
friend, others an acquaintance, others a relative; some, too,
recognized their enemies. Thus, gladness and sorrow, grief and joy,
were variously felt throughout the whole army.

\xchapter{Chronology of the Conspiracy of Catiline.}

A.U.C 685.—COSS. L. CAECILIUS METELLUS, Q. MARCIIS REX.—Catiline is

Praetor.

686.—C. CALPURNIUS PISO, M. ACILIUS GLABRIO.—Catiline Governor of
Africa.

687.—L. VOLCATIUS TULLUS, M. AEMILIUS LEPIDUS.—Deputies from Africa
accuse Catiline of extortion, through the agency of Clodius. He is
obliged to desist from standing for the consulship, and forms the
project of the first conspiracy. See Sall. Cat., c. 18.

688.—L. MANLIUS TORQUATUS, L. AURELIUS COTTA.—\emph{Jan}. 1: Catiline's
project of the first conspiracy becomes known, and he defers the
execution of it to the 5th of February, when he makes an unsuccessful
attempt to execute it. \emph{July} 17: He is acquitted of extortion, and
begins to canvass for the consulship for the year 690.

689.—L. JULIUS CAESAR, C. MARCIUS FIGULUS THERMUS.—\emph{June} 1:
Catiline convokes the chiefs of the second conspiracy. He is
disappointed in his views on the consulship.

690—M. TULLIUS CICERO, C. ANTONIUS HYBRIDA.—\emph{Oct}. 19: Cicero lays
the affair of the conspiracy before the senate, who decree plenary
powers to the consuls for defending the state. \emph{Oct}. 21: Silanus and
Muraena are elected consuls for the next year, Catiline, who was a
candidate, being rejected. \emph{Oct}. 22: Catiline is accused under the
Plautian Law \emph{de vi}. Sall. Cat., c. 31. \emph{Oct}. 24: Manlius takes up
arms in Etruria. \emph{Nov}. 6: Catiline assembles the chief conspirators,
by the agency of Porcius Laeca Sall. Cat., c. 27. \emph{Nov}. 7: Vargunteius
and Cornelius undertake to assassinate Cicero. Sall. Cat., c. 28.
\emph{Nov}. 8: Catiline appears in the senate; Cicero delivers his first
Oration against him; he threatens to extinguish the flame raised
around him in a general destruction, and quits Rome. Sall. Cat., c.
31. \emph{Nov}. 9: Cicero delivers his second Oration against Catiline,
before an assembly of the people, convoked by order of the senate.
\emph{Nov}. 20, \emph{or thereabouts}: Catiline and Manlius are declared public
enemies. Soon after this the conspirators attempt to secure the
support of the Allobrogian deputies. \emph{Dec}. 3: About two o'clock in
the morning the Allobroges are apprehended. Toward evening Cicero
delivers his third Oration against Catiline, before the people. \emph{Dec}.
5. Cicero's fourth Oration against Catiline, before the senate. Soon
after, the conspirators are condemned to death, and great honors are
decreed by the senate to Cicero. 691.—D. JUNIUS SILANUS, L. LICINIUS
MURAENA—\emph{Jan}. 5: Battle of Pistoria, and death of Catiline.

* * * * *

The narrative of Sallust terminates with the account of the battle of
Pistoria. There are a few other particulars connected with the history
of the conspiracy, which, for the sake of the English reader, it may
not be improper to add.

When the victory was gained, Antonius caused Catiline's head to be cut
off, and sent it to Rome by the messengers who carried the news.
Antonius himself was honored, by a public decree, with the title of
\emph{Imperator}, although he had done little to merit the distinction, and
although the number of slain, which was three thousand, was less than
that for which the title was generally given. See Dio Cass. xxxvii.,
40, 41.

The remains of Catiline's army, after the death of their leader,
continued to make efforts to raise another insurrection. In August,
eight months after the battle, a party, under the command of Lucius
Sergius, perhaps a relative or freedman of Catiline, still offered
resistance to the forces of the government in Etruria. \emph{Reliquiae
conjuratorum, cum L. Sergio, tumultuantur in Hetruria}. Fragm. Act.
Diurn. The responsibility of watching these marauders was left to the
proconsul Metellus Celer. After some petty encounters, in which the
insurgents were generally worsted, Sergius, having collected his force
at the foot of the Alps, attempted to penetrate into the country of
the Allobroges, expecting to find them ready to take up arms; but
Metellus, learning his intention, pre-occupied the passes, and then
surrounded and destroyed him and his followers.

At Rome, in the mean time, great honors were paid to Cicero. A
thanksgiving of thirty days was decreed in his name, an honor which
had previously been granted to none but military men, and which was
granted to him, to use his own words, because \emph{he had delivered the
city from fire, the citizens from slaughter, and Italy from war}. ``If
my thanksgiving,'' he also observes, ``be compared with those of others,
there will be found this difference, that theirs were granted them for
having managed the interests of the republic successfully, but that
mine was decreed to me for having preserved the republic from ruin.''
See Cic. Orat. iii., in Cat., c. 6. Pro Sylla, c. 30. In Pison. c. 3.
Philipp. xiv., 8. Quintus Catulus, then \emph{princeps senatus}, and Marcus

Roma parentem,

Roma patrem patriae Ciceronem libera dixit.

Juv. Sat., viii. 244.

Of the inferior conspirators, who did not follow Sergius, and who were
apprehended at Rome, or in other parts of Italy, after the death of
the leaders in the plot, some were put to death, chiefly on the
testimony of Lucius Vettius, one of their number, who turned informer
against the rest. But many whom he accused were acquitted; others,
supposed to be guilty, were allowed to escape.

\xchapter{The Jugurthine War.}

\xsection{The Argument.}

The Introduction, I.-IV. The author's declaration of his design, and
prefatory account of Jugurtha's family, V. Jugurtha's character, VI.
His talents excite apprehensions in his uncle Micipsa, VII. He is sent
to Numantia. His merits, his favor with Scipio, and his popularity in
the army, VIII. He receives commendation and advice from Scipio and is
adopted by Micipsa, who resolves that Jugurtha, Adherbal, and
Hiempsal, shall, at his death, divide his kingdom equally between
them, IX. He is addressed by Micipsa on his death-bed, X. His
proceedings, and those of Adherbal and Hiempsal, after the death of
Micipsa, XI. He murders Hiempsal, XII. He defeats Adherbal, and drives
him for refuge to Rome. He dreads the vengeance of the senate, and
sends embassadors to Rome, who are confronted with those of Adherbal
in the senate-house, XIII. The speech of Adherbal, XIV. The reply of
Jugurtha's embassadors, and the opinions of the senators, XV. The
prevalence of Jugurtha's money, and the partition of the kingdom
between him and Adherbal, XVI. A description of Africa, XVII. An
account of its inhabitants, and of its principal divisions at the
commencement of the Jugurthine war, XVIII., XIX. Jugurtha invades
Adherbal's part of the kingdom, XX. He defeats Adherbal, and besieges
him in Cirta, XXI. He frustrates the intentions of the Roman deputies,
XXII. Adherbal's distresses, XXIII. His letter to the senate, XXIV.
Jugurtha disappoints a second Roman deputation, XXV. He takes Cirta,
and puts Adherdal to death, XXVI. The senate determine to make war
upon him, and commit the management of it to Calpurnius, XXVII. He
sends an ineffectual embassy to the senate. His dominions are
vigorously invaded by Calpurnius, XXVIII. He bribes Calpurnius, and
makes a treaty with him, XXIX. His proceedings are discussed at Rome,
XXX. The speech of Memmius concerning them, XXXI. The consequences of
it, XXXII. The arrival of Jugurtha at Rome, and his appearance before
the people, XXXIII., XXXIV. He procures the assassination of Massiva,
and is ordered to quit Italy, XXXV. Albinus, the successor of
Calpurnius, renews the war. He returns to Rome, and leaves his brother
Aulus to command in his absence, XXXVI. Aulus miscarries in the siege
of Suthul, and concludes a dishonorable treaty with Jugurtha, XXXVII,
XXXVIII. His treaty is annulled by the senate. His brother, Albinus,
resumes the command, XXXIX. The people decree an inquiry into the
conduct of those who had treated with Jugurtha, XL. Consideration on
the popular and senatorial factions, XLI., XLII. Metellus assumes the
conduct of the war, XLIII. He finds the army in Numidia without
discipline, XLIV. He restores subordination, XLV. He rejects
Jugurtha's offers of submission, bribes his deputies, and marches into
the country, XLVI. He places a garrison in Vacca, and seduces other
deputies of Jugurtha, XLVII. He engages with Jugurtha, and defeats
him. His lieutenant, Rutilius, puts to flight Bomilcar, the general of
Jugurtha, XLVIII.-LIII. He is threatened with new opposition. He lays
waste the country. His stragglers are cut off by Jugurtha, LIV. His
merits are celebrated at Rome. His caution. His progress retarded, LV.
He commences the siege of Zama, which is reinforced by Jugurtha. His
lieutenant, Marius, repulses Jugurtha at Sicca, LVI. He is joined by
Marius, and prosecutes the siege. His camp is surprised, LVII., LVIII.
His struggles with Jugurtha, and his operations before the town, LIX.,
LX. He raises the siege, and goes into winter quarters. He attaches
Bomilcar to his interest, LXI. He makes a treaty with Jugurtha, who
breaks it, LXII. The ambition of Marius. His character. His desire of
the consulship, LXIII. His animosity toward Metellus. His intrigues to
supplant him, LXIV, LXV. The Vaccians surprise the Roman garrison, and
kill all the Romans but Turpilius, the governor, LXVI., LXVII.
Metellus recovers Vacca, and puts Turpilius to death, LXVIII., LXIX.
The conspiracy of Bomilcar, and Nabdalsa against Jugurtha, and the
discovery of it. Jugurtha's disquietude, LXX.-LXXII. Metellus makes
preparations for a second campaign. Marius returns to Rome, and is
chosen consul, and appointed to command the army in Numidia, LXXIII.
Jugurtha's irresolution. Metellus defeats him, LXXIV. The flight of
Jugurtha to Thala. The march of Metellus in pursuit of him, LXXV.
Jugurtha abandons Thala, and Metellus takes possession of it, LXXVI.
Metellus receives a deputation from Leptis, and sends a detachment
thither, LXXVII. The situation of Leptis, LXXVIII. The history of
the Philaeni, LXXIX. Jugurtha collects an army of Getulians, and gains
the support of Bocchus, King of Mauritania. The two kings proceed
toward Cirta, LXXX., LXXXI. Metellus marches against them, but hearing
that Marius is appointed to succeed him, contents himself with
endeavoring to alienate Bocchus from Jugurtha, and protracting the war
rather than prosecuting it, LXXXII., LXXXIII. The preparations of Marius
for his departure. His disposition toward the nobility. His popularity,
LXXXIV. His speech to the people, LXXXV. He completes his levies, and
arrives in Africa, LXXXVI. He opens the campaign, LXXXVII. The reception
of Metellus in Rome. The successes and plans of Marius. The applications
of Bocchus, LXXXVIII. Marius marches against Capsa, and takes it,
LXXXIX-XCI. He gains possession of a fortress which the Numidians thought
impregnable, XCII.-XCIV. The arrival of Sylla in the camp. His character,
XCV. His arts to obtain the favor of Marius and the soldiers, XCVI.
Jugurtha and Bocchus attack Marius, and are vigorously opposed, XCVII.,
XCVIII. Marius surprises them in the night, and routs them with great
slaughter, XCIX. Marius prepares to go into winter quarters. His
vigilance, and maintenance of discipline, C. He fights a second battle
with Jugurtha and Bocchus, and gains a second victory over them, CI. He
arrives at Cirta. He receives a deputation from Bocchus, and sends Sylla
and Manlius to confer with him, CII. Marina undertakes an expedition
Bocchus prepares to send ambassadors to Rome, who being stripped by
robbers, takes refuge in the Roman camp, and are entertained by Sylla
during the absence of Marius, CIII. Marius returns. The ambassadors
set out for Rome. The answer which they receive from the senate, CIV.
Bocchus desires a conference with Sylla; Sylla arrives at the camp of
Bocchus, CV.-CVII. Negotiations between Sylla and Bocchus, CVIII.,
CIX. The address of Bocchus to Sylla, CX. The reply of Sylla. The
subsequent transactions between them. The resolution of Bocchus to
betray Jugurtha, and the execution of it, CXI-CXIII. The triumph of
Marius, CXIV.

I. Mankind unreasonably complain of their nature, that, being weak and
short-lived, it is governed by chance rather than intellectual power;\footnote{I. Intellectual power—\emph{Virtute}. See the remarks on
\emph{virtus}, at the commencement of the Conspiracy of Catiline. A little
below, I have rendered \emph{via virtutis}, ``the path of true merit.''}
for, on the contrary, you will find, upon reflection, that there is
nothing more noble or excellent, and that to nature is wanting rather
human industry than ability or time.

The ruler and director of the life of man is the mind, which, when it
pursues glory in the path of true merit, is sufficiently powerful,
efficient, and worthy of honor,\footnote{Worthy of honor—\emph{Clarus}. ``A person may be called \emph{clarus}
either on account of his great actions and merits; or on account of
some honor which he has obtained, as the consuls were called
\emph{clarissimi viri}; or on account of great expectations which are
formed from him. But since the worth of him who is \emph{clarus} is known
by all, it appears that the mind is here called \emph{clarus} because its
nature is such that pre-eminence is generally attributed to it, and
the attention of all directed toward it.'' \emph{Dietsch}.} and needs no assistance from
fortune, who can neither bestow integrity, industry, or other good
qualities, nor can take them away. But if the mind, ensnared by
corrupt passions, abandons itself\footnote{Abandons itself—\emph{Pessum datus est}. Is altogether sunk and
overwhelmed.} to indolence and sensuality, when
it has indulged for a season in pernicious gratifications, and when
bodily strength, time, and mental vigor, have been wasted in sloth,
the infirmity of nature is accused, and those who are themselves in
fault impute their delinquency to circumstances.\footnote{Impute their delinquency to circumstances, etc.—\emph{Suam quisque
culpam ad negotia transferunt}. Men excuse their indolence and
inactivity, by saying that the weakness of their faculties, or the
circumstances in which they are placed, render them unable to
accomplish any thing of importance. But, says Seneca, \emph{Satis natura,
homini dedit roboris, si illo utamur;—nolle in causâ, non posse
praetenditur}. ``Nature has given men sufficient powers, if they will
but use them; but they pretend that they can not when the truth is
that they will not.'' ``\emph{Negotia} is a common word with Sallust, for
which other writers would use \emph{res, facta}.'' Gerlach. ``Cajus rei nos
ipsi sumus auctores, ejus culpam rebus externis attribuimus.''
\emph{Muller}. ``Auctores'' is the same as the [Greek: \emph{aitioi}].}

If man, however, had as much regard for worthy objects, as he has
spirit in the pursuit of what is useless,\footnote{Useless—\emph{Aliena}. Unsuitable, not to the purpose, not
contributing to the improvement of life.} unprofitable, and even
perilous, he would not be governed by circumstances more than he would
govern them, and would attain to a point of greatness, at which,
instead of being mortal,\footnote{Instead of being mortal—\emph{Pro mortalibus}. There are two senses
in which these words may be taken: \emph{as far as mortals can}, and
\emph{instead of being mortals}. Kritz and Dietsch say that the latter is
undoubtedly the true sense. Other commentators are either silent or
say little to the purpose. As for the translators, they have studied
only how to get over the passage delicately. The latter sense is
perhaps favored by what is said in c. 2, that ``the illustrious
achievements of the mind are, like the mind itself, immortal.''} he would be immortalized by glory.

II. As man is composed of mind and body, so, of all our concerns and
pursuits, some partake the nature of the body, and some that of the
mind. Thus beauty of person, eminent wealth, corporeal strength, and
all other things of this kind, speedily pass away; but the illustrious
achievements of the mind are, like the mind itself, immortal.

Of the advantages of person and fortune, as there is a beginning,
there is also an end; they all rise and fall,\footnote{II. They all rise and fall, etc.—\emph{Omnia orta occidunt, et
aucta senescunt}. This is true of things in general, but is here
spoken only of the qualities of the body, as De Brosses clearly
perceived.} increase and decay.
But the mind, incorruptible and eternal, the ruler of the human race,
actuates and has power over all things,\footnote{Has power over all things—\emph{Habet cuncta}. ``All things are in
its power.'' Dietsch. ``\emph{Sub ditione tenet}. So Jupiter, Ov. Met.
i. 197:
Quum mihi qui fulmen, qui vos habeoque regoque.''
\emph{Burnouf}. So Aristippus said, \emph{Habeo Laidem, non habeor a Laide},
[Greek: \emph{echo ouk echomai}]. Cic. Epist. ad Fam. ix. 26.} yet is itself free from
control.

The depravity of those, therefore, is the more surprising, who,
devoted to corporeal gratifications, spend their lives in luxury and
indolence, but suffer the mind, than which nothing is better or
greater in man, to languish in neglect and inactivity; especially when
there are so many and various mental employments by which the highest
renown may be attained.

III. Of these occupations, however, civil and military offices,\footnote{III. Civil and military offices—\emph{Magistratus et imperia}.
``Illo vocabule civilia, hoc militaria munera, significantur.''
\emph{Dietsch}.} and
all administration of public affairs, seem to me at the present time,
by no means to be desired; for neither is honor conferred on merit,
nor are those, who have gained power by unlawful means, the more
secure or respected for it. To rule our country or subjects\footnote{To rule our country or subjects, etc.—\emph{Nam vi quidem regere
patriam aut parentes}, etc. Cortius, Gerlach, Kritz, Dietsch, and
Muller are unanimous in understanding \emph{parentes} as the participle of
the verb \emph{pareo}. That this is the sense, says Gerlach, is
sufficiently proved by the conjunction \emph{aut}; for if Sallust had meant
\emph{parents}, he would have used \emph{ut}; and in this opinion Allen
coincides. Doubtless, also, this sense of the word suits extremely
well with the rest of the sentence, in which changes in government are
mentioned. But Burnouf, with Crispinus, prefers to follow Aldus
Manutius, who took the word in the other signification, supposing that
Sallust borrowed the sentiment from Plato, who says in his Epistle \emph{ad
Dionis Propinquos}: [Greek: \emph{Patera de hae maetera ouch osion
haegoumai prosbiazesthai, mae noso paraphrosunaes hechomenous. Bian de
patridi politeias metabolaes aeae prospherein, otan aneu phugon, kai
sphagaes andron, mae dunaton hae ginesthae taen ariostaen}.] And he
makes a similar observation in his Crito: [Greek: \emph{Pantachou
poiaetaen, o an keleuoi hae polis te, kai hae patris.—Biazesthai de
ouch osion oute maetera, oute patera poly de touton eti aetton taen
patrida}.] On which sentiments Cicero, ad Fam. i. 9, thus comments:
\emph{Id enim jubet idem ille Plato, quem ego auctorem vehementer sequor;
tantum contendere in republica quantum probare tuis civibus possis:
vim neque parenti, neque patriae afferre oportere}. There is also
another passage in Cicero, Cat. i. 3, which seems to favor this sense
of the word: \emph{Si te parentes timerent atque odissent tui, neque eos
ullâ ratione placare posses, ut opinor, ab eorum oculis aliquò
concederes; nunc te patria, quae communis est omnium nostrum parens
odit ac metuit}, etc. Of the first passage cited from Plato, indeed,
Sallust's words may seem to be almost a translation. Yet, as the
majority of commentators have followed Cortius, I have also followed
him. Sallust has the word in this sense in Jug., c. 102; \emph{Parentes
abunde habemus}. So Vell. Pat. ii. 108: \emph{Principatis constans ex
voluntate parentium}.} by
force, though we may have the ability, and may correct what is wrong,
is yet an ungrateful undertaking; especially as all changes in the
state lead to\footnote{Lead to—\emph{Portendant}. ``\emph{Portendere} in a \emph{pregnant sense},
meaning not merely to indicate, but \emph{quasi secum ferre}, to carry
along with them.'' \emph{Kritzius}.} bloodshed, exile, and other evils of discord; while
to struggle in ineffectual attempts, and to gain nothing, by wearisome
exertions, but public hatred, is the extreme of madness; unless when a
base and pernicious spirit, perchance, may prompt a man to sacrifice
his honor and liberty to the power of a party.

IV. Among other employments which are pursued by the intellect, the
recording of past events is of pre-eminent utility; but of its merits
I may, I think, be silent, since many have spoken of them, and since,
if I were to praise my own occupation, I might be considered as
presumptuously\footnote{IV, Presumptuously—\emph{Per insolentiam}. The same as
\emph{insolenter}, though some refer it, not to Sallust, but to \emph{quis
existumet,} in the sense of \emph{strangely,} i. e. \emph{foolishly or
ignorantly.} I follow Cortius's interpretation.} praising myself. I believe, too, that there will be
some, who, because I have resolved to live unconnected with political
affairs, will apply to my arduous and useful labors the name of
idleness; especially those who think it an important pursuit to court
the people, and gain popularity by entertainments. But if such persons
will consider at what periods I obtained office, what sort of men\footnote{At what periods I obtained office, what sort of men, etc.
—\emph{Quibus ego temporibus maqistratus adeptus sum, et quales viri,} etc.
—``Sallust obtained the quaestorship a few years after the conspiracy
of Catiline, about the time when the state was agitated by the
disorders of Clodius and his party. He was tribune of the people,
A.U.C. 701, the year in which Clodius was killed by Milo. He was
praetor in 708, when Caesar had made himself ruler. In the expression
\emph{quales viri,} etc., he alludes chiefly to Cato, who, when he stood
for the praetorship, was unsuccessful.'' \emph{Burnouf.} Kritzius defends
\emph{adeptus sum.}}
were then unable to obtain it, and what description of persons have
subsequently entered the senate,\footnote{What description of persons have subsequently entered the
senate—``Caesar chose the worthy and unworthy, as suited his own
purposes, to be members of the senate.'' \emph{Burnouf.}} they will think, assuredly, that
I have altered my sentiments rather from prudence than from indolence,
and that more good will arise to the state from my retirement, than
from the busy efforts of others.

I have often heard that Quintus Maximus,\footnote{Quintus Maximus—Quintus Fabius Maximus, of whom Ennius says,

Unus qui nobis cunctando restituit rem;

Non ponebat enim rumores ante salutem.} Publius Scipio,\footnote{Publius Scipio—Scipio Africanus the Elder, the conqueror of
Hannibal. See c. 5.} and
many other illustrious men of our country, were accustomed to observe,
that, when they looked on the images of their ancestors, they felt
their minds irresistibly excited to the pursuit of honor.\footnote{To the pursuit of honor—\emph{Ad vertutem. Virtus} in the same
sense as in \emph{virtutis viâ,} c. 1.} Not,
certainly, that the wax,\footnote{The wax—\emph{Ceram illam.} The images or busts of their ancestors,
which the nobility kept in the halls of their houses, were made of wax.
See Plin. H. N. xxxv., 2.} or the shape, had any such influence;
but, as they called to mind their forefathers' achievements, such a
flame was kindled in the breasts of those eminent persons, as could
not be extinguished till their own merit had equaled the fame and
glory of their ancestors.

But, in the present state of manners, who is there, on the contrary,
that does not rather emulate his forefathers in riches and extravagance,
than in virtue and labor? Even men of humble birth,\footnote{Men of humble birth—\emph{Homines novi}. See Cat., c. 23.} who formerly
used to surpass the nobility in merit, pursue power and honor rather
by intrigue and dishonesty, than by honorable qualifications; as if
the praetorship, consulate, and all other offices of the kind, were
noble and dignified in themselves, and not to be estimated according
to the worth of those who fill them.

But, in expressing my concern and regret at the manners of the state,
I have proceeded with too great freedom, and at too great length. I
now return to my subject.

V. I am about to relate the war which the Roman people carried on with
Jugurtha, King of the Numidians; first, because it was great, sanguinary,
and of varied fortune; and secondly, because then, for the first time,
opposition was offered to the power of the nobility; a contest which
threw every thing, religious and civil, into confusion,\footnote{V. Threw every thing, religious and civil, into confusion—\emph{Divina
et humana cuncta permiscuit}. ``All things, both divine and human, were
so changed, that their previous condition was entirely subverted.''
\emph{Dietsch}.} and was
carried to such a height of madness, that nothing but war, and the
devastation of Italy, could put an end to civil dissensions.\footnote{Civil dissensions—\emph{Studiis civilibus}. This is the sense in
which most commentators take \emph{studia}; and if this be right, the whole
phrase must be understood as I have rendered it. So Cortius; ``Ut non
prius finirentur [\emph{studia civilia}] nisi bello et vastitate Italiae.''
Sallust has \emph{studia partium} Jug. c. 42; and Gerlach quotes from Cic.
pro Marcell. c. 10: ``\emph{Non enim consiliis solis et studiis, sed armis
etiam et castris dissidebamus}''.} But
before I fairly commence my narrative, I will take a review of a few
preceding particulars, in order that the whole subject may be more
clearly and distinctly understood.

In the second Punic war, in which Hannibal, the leader of the
Carthaginians, had weakened the power of Italy more than any other
enemy\footnote{More than any other enemy—\emph{Maximè}.} since the Roman name became great,\footnote{Since the Roman name became great—\emph{Post magnitudinem nominis
Romani}. ``I know not why interpreters should find any difficulty in
this passage. I understand it to signify simply \emph{since} the Romans
became so great as they were in the time of Hannibal; for, \emph{before}
that period they had suffered even heavier calamities, especially
from the Gauls.'' \emph{Cortius}.} Masinissa, King of
the Numidians, being received into alliance by Publius Scipio, who,
from his merits was afterward surnamed Africanus, had performed for us
many eminent exploits in the field. In return for which services,
after the Carthaginians were subdued, and after Syphax,\footnote{Syphax—``He was King of the Masaesyli in Numidia; was at first
an enemy to the Carthaginians (Liv. xxiv. 48), and afterward their
friend (Liv. xxviii. 17). He then changed sides again, and made
a treaty with Scipio; but having at length been offered the hand of
Sophonisba, the daughter of Asdrubal, in marriage, he accepted it,
and returned into alliance with the Carthaginians. Being subsequently
taken prisoner by Masinissa and Laelius, the lieutenant of Scipio,
(Liv. xxx. 2) he was carried into Italy, and died at Tibur (Liv. xxx.
45).'' \emph{Burnouf}.} whose
power in Italy was great and extensive, was taken prisoner, the Roman
people presented to Masinissa, as a free gift, all the cities and
lands that they had captured. Masinissa's friendship for us,
accordingly, remained faithful and inviolate; his reign\footnote{His reign—\emph{Imperii}. Cortius thinks that the grant of the
Romans ceased with the life of Masinissa, and that his son, Micipsa,
reigned only over that part of Numidia which originally belonged to
his father. But in this opinion succeeding commentators have generally
supposed him to be mistaken.} and his
life ended together. His son, Micipsa, alone succeeded to his kingdom;
Mastanabal and Gulussa, his two brothers, having been carried off by
disease. Micipsa had two sons, Adherbal and Hiempsal, and had brought
up in his house, with the same care as his own children, a son of his
brother Mastanabal, named Jugurtha, whom Masinissa, as being the son
of a concubine, had left in a private station.

VI. Jugurtha, as he grew up, being strong in frame, graceful in
person, but, above all, vigorous in understanding, did not allow
himself to be enervated by pleasure and indolence, but, as is the
usage of his country, exercised himself in riding, throwing the
javelin, and contending in the race with his equals in age; and,
though he excelled them all in reputation, he was yet beloved by all.
He also passed much of his time in hunting; he was first, or among the
first, to wound the lion and other beasts; he performed very much, but
spoke very little of himself.

Micipsa, though he was at first gratified with these circumstances,
considering that the merit of Jugurtha would be an honor to his
kingdom, yet, when he reflected that the youth was daily increasing in
popularity, while he himself was advanced in age, and his children but
young, he was extremely disturbed at the state of things, and revolved
it frequently in his mind. The very nature of man, ambitious of power,
and eager to gratify its desires, gave him reason for apprehension, as
well as the opportunity afforded by his own age and that of his children,
which was sufficient, from the prospect of such a prize, to lead astray
even men of moderate desires. The affection of the Numidians, too, which
was strong toward Jugurtha, was another cause for alarm; among whom, if
he should cut off such a man, he feared that some insurrection or war
might arise.

VII. Surrounded by such difficulties, and seeing that a man, so
popular among his countrymen, was not to be destroyed either by force
or by fraud, he resolved, as Jugurtha was of an active disposition,
and eager for military reputation, to expose him to dangers in the
field, and thus make trial of fortune. During the Numantine war,\footnote{VII. During the Numantine war—\emph{Bella Numantino}. Numantia,
which stood near the source of the Durius or \emph{Douro} in Spain, was
so strong in its situation and fortifications, that it with stood the
Romans for fourteen years. See Florus, ii. 17,18; Vell. Pat. ii. 4.}
therefore, when he was sending supplies of horse and foot to the
Romans, he gave him the command of the Numidians, whom he dispatched
into Spain, hoping that he would certainly perish, either by an
ostentatious display of his bravery, or by the merciless hand of the
enemy. But this project had a very different result from that which he
had expected. For when Jugurtha, who was of an active and penetrating
intellect, had learned the disposition of Publius Scipio, the Roman
general, and the character of the enemy, he quickly rose, by great
exertion and vigilance, by modestly submitting to orders, and frequently
exposing himself to dangers, to such a degree of reputation, that he was
greatly beloved by our men, and extremely dreaded by the Numantines. He
was indeed, what is peculiarly difficult, both brave in action, and wise
in counsel; qualities, of which the one, from forethought, generally
produces fear, and the other, from confidence, rashness. The general,
accordingly, managed almost every difficult matter by the aid of
Jugurtha, numbered him among his friends, and grew daily more and more
attached to him, as a man whose advice and whose efforts were never
useless. With such merits were joined generosity of disposition, and
readiness of wit, by which he united to himself many of the Romans in
intimate friendship.

VIII. There were at that time, in our army, a number of officers, some
of low, and some of high birth, to whom wealth was more attractive
than virtue or honor; men who were attached to certain parties, and of
consequence in their own country; but, among the allies, rather
distinguished than respected. These persons inflamed the mind of
Jugurtha, of itself sufficiently aspiring, by assuring him, ``that if
Micipsa should die, he might have the kingdom of Numidia to himself;
for that he was possessed of eminent merit, and that anything might be
purchased at Rome.''

When Numantia, however, was destroyed, and Scipio had determined to
dismiss the auxiliary troops, and to return to Rome, he led Jugurtha,
after having honored him, in a public assembly, with the noblest
presents and applauses, into his own tent; where he privately
admonished him ``to court the friendship of the Romans rather by
attention to them as a body, than by practicing on individuals;\footnote{VIII. Rather by attention to them as a body, than by practicing
on individuals—\emph{Publicè quàm privatim.} ``Universae potius civitatis,
quàm privatorum gratiam quaerendo.'' \emph{Burnouf}. The words can only be
rendered periphrastically.}
to bribe no one, as what belonged to many could not without danger be
bought from a few; and adding that, if he would but trust to his own
merits, glory and regal power would spontaneously fall to his lot;
but, should he proceed too rashly, he would only, by the influence of
his money, hasten his own ruin.''

IX. Having thus spoken, he took leave of him, giving him a letter,
which he was to present to Micipsa, and of which the following was
the purport: ``The merit of your nephew Jugurtha, in the war against
Numantia, has been eminently distinguished; a fact which I am sure
will afford you pleasure. He is dear to us for his services, and we
shall strive, with our utmost efforts, to make him equally dear to the
senate and people of Rome. As a friend, I sincerely congratulate you;
you have a kinsman worthy of yourself, and of his grandfather
Masinissa.''

Micipsa, when he found, from the letter of the general, that what he
had already heard reported was true, being moved, both by the merit of
the youth and by the interest felt for him by Scipio, altered his
purpose, and endeavored to win Jugurtha by kindness. He accordingly,
in a short time,\footnote{IX. In a short time—\emph{Statim}. If what is said in c. 11 be
correct, that Jugurtha was adopted within three years of Micipsa's
death, his adoption did not take place till twelve years after the
taking of Numantia, which surrendered in 619, and Micipsa died in 634.
\emph{Statim} is therefore used with great latitude, unless we suppose
Sallust to mean that Micipsa signified to Jugurtha his intention to
adopt him immediately on his return from Numantia, and that the formal
ceremony of the adoption was delayed for some years.} adopted him as his son, and made him, by his
will, joint-heir with his own children.

A few years afterward, when, being debilitated by age and disease, he
perceived that the end of his life was at hand, he is said, in the
presence of his friends and relations, and of Adherbal and Hiempsal
his sons, to have spoken with Jugurtha in the following manner:

X. ``I received you, Jugurtha, at a very early age, into my kingdom,\footnote{X. I received you—into my kingdom—\emph{In meuum regnum accepi}.
By these words it is only signified that Micipsa received Jugurtha
into his palace so as to bring him up with his own children. The
critics who suppose that there is any allusion to the adoption, or
a pretended intention of it on the part of Micipsa, are evidently in
the wrong.}
at a time when you had lost your father, and were without prospects or
resources, expecting that, in return for my kindness, I should not be
less loved by you than by my own children, if I should have any. Nor
have my anticipations deceived me; for, to say nothing of your other
great and noble deeds, you have lately, on your return from Numantia,
brought honor and glory both to me and my kingdom; by your bravery,
you have rendered the Romans, from being previously our friends, more
friendly to us than ever; the name of our family is revived in Spain;
and, finally, what is most difficult among mankind, you have suppressed
envy by preeminent merit.\footnote{Pre-eminent merit—\emph{Gloriâ}. Our English word \emph{glory} is too
strong.}

And now, since nature is putting a period to my life, I exhort and
conjure you, by this right hand, and by the fidelity which you owe to
my kingdom,\footnote{By the fidelity which you owe to my kingdom—\emph{Per regni
fidem}. This seems to be the best of all the explanations that have
been offered of these words. ``Per fidem quam tu rex (futurus) mihi
regi praestare debes.'' \emph{Bournouf}. ``Per fidem quae decet in regno, \emph{i.
e.} regem.'' \emph{Dietsch}. ``Per eam fidem, qua esse decet eum qui regnum
obtinet. \emph{Kritzius}.} to regard these princes, who are your cousins by
birth, and your brothers by my generosity, with sincere affection; and
not to be more anxious to attach to yourself strangers, than to retain
the love of those connected with you by blood. It is not armies, or
treasures,\footnote{It is not armies, or treasures, etc.—[Greek: \emph{Ou tode to
chrusoun skaeptron to taen basileian diasozon estin, alla oi polloi
philoi skaeptron basileioin ulaethestaton kai asphalestaton}.] ``It is
not this golden scepter that can preserve a kingdom; but numerous
friends are to princes their trust and safest scepter.'' Xen. Cyrop.,
viii. 7,14.} that form the defenses of a kingdom, but friends, whom
you can neither command by force nor purchase with gold; for they are
acquired only by good offices and integrity. And who can be a greater
friend than one brother to another?\footnote{And who can be a greater friend than one brother to another?
—\emph{Quis autem amicior, quam frater fratri?} ``[Greek: \emph{Nomiz
adelphous tous alaethinous philous}] Menander.'' \emph{Wasse}.} Or what stranger will you find
faithful, if you are at enmity with your own family? I leave you a
kingdom, which will be strong if you act honorably, but weak, if you
are ill-affected to each other; for by concord even small states are
increased, but by discord, even the greatest fall to nothing.

But on you, Jugurtha, who are superior in age and wisdom, it is
incumbent, more than on your brothers, to be cautious that nothing of
a contrary tendency may arise; for, in all disputes, he that is the
stronger, even though he receive the injury, appears, because his
power is greater, to have inflicted it. And do you, Adherbal and
Hiempsal, respect and regard a kinsman of such a character; imitate
his virtues, and make it your endeavor to show that I have not adopted
a better son\footnote{That I have not adopted a better son, \&c—\emph{Ne ego meliores
liberos sumsissevidear quam genuisse}. As there is no allusion to
Micipsa's adoption of any other son than Jugurtha, Sallust's
expression \emph{liberos sumsisse} can hardly be defended. It is necessary
to give \emph{son} in the singular, in the translation.} than those whom I have begotten.''

XI. To this address, Jugurtha, though he knew that the king had spoken
insincerely,\footnote{XI. Had spoken insincerely—\emph{Ficta locutum}. Jugurtha saw
that Micipsa pretended more love for him than he really felt. Compare
c. 6,7.} and though he was himself revolving thoughts of a far
different nature, yet replied with good feeling, suitable to the
occasion. A few days afterward Micipsa died.

When the princes had performed his funeral with due magnificence, they
met together to hold a discussion on the general condition of their
affairs. Hiempsal, the youngest, who was naturally violent, and who
had previously shown contempt for the mean birth of Jugurtha, as being
inferior on his mother's side, sat down on the right hand of Adherbal,
in order to prevent Jugurtha from being the middle one of the three,
which is regarded by the Numidians as the seat of honor.\footnote{Which is regarded by the Numidians as the seat of honor—\emph{Quad
apud Numidas honori ducitur}. ``I incline,'' says Sir Henry Steuart,
``to consider those manuscripts as the most correct, in which the word
\emph{et} is placed immediately before \emph{apud, Quod et apud Numidas honori
ducitur}.'' Sir Henry might have learned, had he consulted the
commentators, that ``\emph{the word} et \emph{is placed immediately before}
apud'' in no manuscript; that Lipsius was the first who proposed its
insertion; and that Crispinus, the only editor who has received it
into his text, is ridiculed by Wasse for his folly. ``Lipsius,'' says
Cortius, ``cùm sciret apud Romanos etiam medium locum honoratiorem
fuisse, corrigit: \emph{quod et apud Numidas honori ducitur}. Sed quis
talia ab historico exegerit? Si de Numidis narrat, non facile aliquis
intulerit, aliter propterea fuisse apud Romanos.''} Being
urged by his brother, however, to yield to superior age, he at length
removed, but with reluctance, to the other seat.\footnote{To the other seat—\emph{In alteram partem}. We must suppose that
the three seats were placed ready for the three princes; that Adherbal
sat down first, in one of the outside seats; the one, namely, that
would be on the right hand of a spectator facing them; and that
Hiempsal immediately took the middle seat, on Adherbal's right hand,
so as to force Jugurtha to take the other outside one. Adherbal had
then to remove Hiempsal \emph{in alteram partem}, that is, to induce him to
take the seat corresponding to his own, on the other side of the
middle one.}

In the course of this conference, after a long debate about the
administration of the kingdom, Jugurtha suggested, among other
measures, ``that all the acts and decrees made in the last five years
should be annulled, as Micipsa, during that period, had been enfeebled
by age, and scarcely sound in intellect.''

Hiempsal replied, ``that he was exceedingly pleased with the proposal,
since Jugurtha himself, within the last three years, had been adopted
as joint-heir to the throne.'' This repartee sunk deeper into the mind
of Jugurtha than any one imagined. From that very time, accordingly,
being agitated with resentment and jealousy, he began to meditate and
concert schemes, and to think of nothing but projects for secretly
cutting off Hiempsal. But his plans proving slow in operation, and his
angry feelings remaining unabated, he resolved to execute his purpose
by any means whatsoever.

XII. At the first meeting of the princes, of which I have just spoken,
it had been resolved, in consequence of their disagreement, that the
treasures should be divided among them, and that limits should be set
to the jurisdiction of each. Days were accordingly appointed for both
these purposes, but the earlier of the two for the division of the
money. The princes, in the mean time, retired into separate places of
abode in the neighborhood of the treasury. Hiempsal, residing in the
town of Thirmida, happened to occupy the house of a man, who, being
Jugurtha's chief lictor,\footnote{Chief lictor—\emph{Proximus lictor}. ``The \emph{proximus lictor} was
he who, when the lictors walked before the prince or magistrate in a
regular line, one behind the other, was last, or next to the person on
whom they attended.'' \emph{Cortius}. He would thus be ready to receive the
great man's commands, and be in immediate communication with him. We
must suppose either that Sallust merely speaks in conformity with the
practice of the Romans, or, what is more probable, that the Roman
custom of being preceded by lictors had been adopted in Numidia.} had always been liked and favored by his
master. This man, thus opportunely presented as an instrument,
Jugurtha loaded with promises, and induced him to go to his house, as
if for the purpose of looking over it, and provide himself with false
keys to the gates; for the true ones used to be given to Hiempsal,
adding, that he himself, when circumstances should call for his
presence, would be at the place with a large body of men. This
commission the Numidian speedily executed, and, according to his
instructions, admitted Jugurtha's men in the night, who, as soon as
they had entered the house, went different ways in quest of the
prince; some of his attendants they killed while asleep, and others as
they met them; they searched into secret places, broke open those that
were shut, and filled the whole premises with uproar and tumult.
Hiempsal, after a time, was found concealed in the hut of a
maid-servant,\footnote{Hut of a maid-servant—\emph{Tugurio muliers ancillae} Rose renders
\emph{tugurio} ``a mean apartment,'' and other translators have given
something similar, as if they thought that the servant must have had a
room in the house. But she, and other Numidian servants, may have had
huts apart from the dwelling-house. \emph{Tugurium} undoubtedly signifies
\emph{a hut} in general.} where, in his alarm and ignorance of the locality,
he had at first taken refuge. The Numidians, as they had been ordered,
brought his head to Jugurtha.

XIII The report of so atrocious an outrage was soon spread through
Africa. Fear seized on Adherbal, and on all who had been subject to
Micipsa. The Numidians divided into two parties, the greater number
following Adherbal, but the more warlike, Jugurtha; who, accordingly,
armed as large a force as he could, brought several cities, partly by
force and partly by their own consent, under his power, and prepared
to make himself sovereign of the whole of Numidia. Adherbal, though he
had sent embassadors to Rome, to inform the senate of his brother's
murder and his own circumstances, yet, relying on the number of his
troops, prepared for an armed resistance. When the matter, however,
came to a contest, he was defeated, and fled from the field of battle
into our province,\footnote{XIII. Into our province—\emph{In Provinciam}. ``The word \emph{province},
in this place, signifies that part of Africa which, after the
destruction of Carthage, fell to the Romans by the right of conquest,
in opposition to the kingdom of Micipsa.'' \emph{Wasse}.} and from thence hastened to Rome.

Jugurtha, having thus accomplished his purposes,\footnote{Having thus accomplished his purposes—\emph{Patratis consiliis}.
After \emph{consiliis}, in all the manuscripts, occur the words \emph{postquam
omnis Numidia potiebatur}, which were struck out by Cortius, as being
\emph{turpissima glossa}. The recent editors, Gerlach, Kritz, Dietsch, and
Burnouf, have restored them.} and reflecting,
at leisure, on the crime which he had committed, began to feel a dread
of the Roman people, against whose resentment he had no hopes of
security but in the avarice of the nobility, and in his own wealth. A
few days afterward, therefore, he dispatched embassadors to Rome, with
a profusion of gold and silver, whom he directed, in the first place,
to make abundance of presents to his old friends, and then to procure
him new ones; and not to hesitate, in short; to effect whatever could
be done by bribery.

When these deputies had arrived at Rome, and had sent large presents,
according to the prince's direction, to his intimate friends,\footnote{His intimate friends—\emph{Hospitibus}. Persons probably with whom
he had been intimate at Numantia, or who had since visited him in
Numidia.} and
to others whose influence was at that time powerful, so remarkable a
change ensued, that Jugurtha, from being an object of the greatest
odium, grew into great regard and favor with the nobility; who, partly
allured with hope, and partly with actual largesses, endeavored, by
soliciting the members of the senate individually, to prevent any
severe measures from being adopted against him. When the embassadors
accordingly, felt sure of success, the senate, on a fixed day, gave
audience to both parties\footnote{The senate—gave audience to both parties—\emph{senatus utrisque
datur}. ``The embassadors of Jugurtha, and Adberbal in person, are
admitted into the senate-house to plead their cause.'' \emph{Burnouf}.}. On that occasion, Adherbal, as I have
understood, spoke to the following effect:

XIV. ``My father Micipsa, Conscript Fathers, enjoined me, on his
death-bed, to look upon the kingdom of Numidia as mine only by
deputation;\footnote{By deputation—\emph{Procuratione}. He was to consider himself only
the \emph{procurator}, manager, or deputed governor, of the kingdom.} to consider the right and authority as belonging to
you; to endeavor, at home and in the field, to be as serviceable to
the Roman people as possible; and to regard you as my kindred and
relatives:\footnote{Kindred—and relatives—\emph{Cognatorum—affinium}. \emph{Cognatus} is
a blood relation; \emph{affinis} is properly a relative by marriage.} saying that, if I observed these injunctions. I should
find, in your friendship, armies, riches, and all necessary defenses
of my realm. By these precepts I was proceeding to regulate my conduct,
when Jugurtha, the most abandoned of all men whom the earth contains,
setting at naught your authority, expelled me, the grandson of Masinissa,
and the hereditary\footnote{Hereditary—\emph{Ab stirpe}.} ally and friend of the Roman people, from my
kingdom and all my possessions.

Since I was thus to be reduced to such an extremity of wretchedness,
I could wish that I were able to implore your aid, Conscript Fathers,
rather for the sake of my own services than those of my ancestors; I
could wish, indeed, above all, that acts of kindness were due to me
from the Romans, of which I should not stand in need; and, next to
this,\footnote{Next to this—\emph{Secundum ea}. ``Priscianus, lib. xiii., de
praepositione agens, \emph{Secundum}, inquit, \_quando pro [Greek: \emph{kata ei
meta}] loco praepositionis est\_. Sallustius in Jugurthino: \emph{secundum
ea, uti deditis uterer}.—Videlicet hoc dicit, \emph{Secundum} in Sallustii
exemple, \emph{post} vel \emph{proximè} significare.'' \emph{Rivius}.} that, if I required your services, I might receive them as
my due. But as integrity is no defense in itself, and as I had no
power to form the character of Jugurtha,\footnote{As I had no power to form the character of Jugurtha—\emph{Neque mihi
in manu fuit, qualis Jugurtha foret}. ``\emph{In manu fuit} is simply
\emph{in potestate fuit}—Ter. Hec., iv. 4, 44: \emph{Uxor quid faciat in manu
non est meâ}.'' \emph{Cortius}.} I have fled to you,
Conscript Fathers, to whom, what is the most grievous of all things, I
am compelled to become a burden before I have been an assistance.

Other princes have been received into your friendship after having
been conquered in war, or have solicited an alliance with you in
circumstances of distress; but our family commenced its league with
the Romans in the war with Carthage, at a time when their faith was a
greater object of attraction than their fortune. Suffer not, then, O
Conscript Fathers, a descendent of that family to implore aid from you
in vain. If I had no other plea for obtaining your assistance but my
wretched fortune; nothing to urge, but that, having been recently a
king, powerful by birth, by character, and by resources, I am now
dishonored, afflicted,\footnote{Dishonored, afflicted—\_Deformatus aerumnis.} destitute, and dependent on the aid of
others, it would yet become the dignity of Rome to protect me from
injury, and to allow no man's dominions to be increased by crime. But
I am driven from those very territories which the Roman people gave to
my ancestors, and from which my father and grandfather, in conjunction
with yourselves, expelled Syphax and the Carthaginians. It is what you
bestowed that has been wrested from me; in my wrongs you are insulted.

Unhappy man that I am! Has your kindness, O my father Micipsa, come
to this, that he whom you made equal with your children, and a sharer
of your kingdom, should become, above all others,\footnote{Above all others—\emph{Potissimum}.} the destroyers
of your race? Shall our family, then, never be at peace? Shall we
always be harassed with war, bloodshed, and exile? While the
Carthaginians continued in power, we were necessarily exposed to all
manner of troubles; for the enemy were on our frontiers; you, our
friends, were at a distance; and all our dependence was on our arms.
But after that pest was extirpated, we were happy in the enjoyment of
tranquillity, as having no enemies but such as you should happen to
appoint us. But lo! on a sudden, Jugurtha, stalking forth with
intolerable audacity, wickedness, and arrogance, and having put to
death my brother, his own cousin, made his territory, in the first
place, the prize of his guilt; and next, being unable to ensnare me
with similar stratagems, he rendered me, when under your rule I
expected any thing rather than violence or war, an exile, as you see,
from my country and my home, the prey of poverty and misery, and safer
any where than in my own kingdom.

I was always of opinion, Conscript Fathers, as I had often heard my
father observe, that those who cultivated your friendship might indeed
have an arduous service to perform, but would be of all people the
most secure. What our family could do for you, it has done; it has
supported you in all your wars; and it is for you to provide for our
safety in time of peace. Our father left two of us, brothers; a third,
Jugurtha, he thought would be attached to us by the benefits conferred
upon him; but one of us has been murdered, and I, the other, have
scarcely escaped the hand of lawlessness.\footnote{One of us has been murdered, and I, the other, have scarcely
escaped the hand of lawlessness—\emph{Alter eorum necatus, alterius ipse
ego manus impias vix effugi}. This is the general reading, but it can
not be right. Adherbal speaks of himself and his brother as two
persons, and of Jugurtha as a third, and says that \emph{of those two} the
one (\emph{alter}) has been killed; he would then naturally proceed to
speak of himself as the other; \emph{i. e.} he would use the word \emph{alter}
concerning himself, not apply it to Jugurtha. Allen, therefore,
proposes to read \emph{alter necatus, alter manus impias vix effugi}. This
mode of correction strikes out too much; but there is no doubt that
the second \emph{alter} should be in the nominative case.} What course can I now
take? Unhappy that I am, to what place, rather than another, shall I
betake myself? All the props of our family are extinct; my father, of
necessity, has paid the debt of nature; a kinsman, whom least of all
men it became, has wickedly taken the life of my brother; and as for
my other relatives, and friends, and connections, various forms of
destruction have overtaken them. Seized by Jugurtha, some have been
crucified, and some thrown to wild beasts, while a few, whose lives
have been spared, are shut up in the darkness of the dungeon, and drag
on, amid suffering and sorrow, an existence more grievous than death
itself.

If all that I have lost, or all that, from being friendly, has become
hostile to me,\footnote{From being friendly, has become hostile to me—\emph{Ex necessariis
adversa facta sunt.} ``Si omnia mihi incolumia manerent, neque quidquam
rerum mearum (s. praesidiorum) amisissem, neque Jugurtha aliique mihi
ex necessariis inimici facti essent.''\emph{Kritzius}.} remained unchanged, yet, in case of any sudden
calamity, it is of you that I should still have to implore assistance,
to whom, from the greatness of your empire, justice and injustice in
general should be objects of regard. And at the present time, when I
am exiled from my country and my home, when I am left alone, and
destitute of all that is suitable to my dignity, to whom can I go, or
to whom shall: I appeal, but to you? Shall I go to nations and kings,
who, from our friendship with Rome, are all hostile to my family?
Could I go, indeed, to any place where there are not abundance of
hostile monuments of my ancestors? Will any one, who, has ever been at
enmity with you, take pity upon me?

Masinissa, moreover, instructed us, Conscript Fathers, to cultivate
no friendship but that of Rome, to adopt no new leagues or alliances,
as we should find, in your good-will, abundance of efficient support;
while, if the fortune of your empire should change, we must sink
together with it. But, by your own merits, and the favor of the gods,
you are great and powerful; the whole world regards you with favor and
yields to your power; and you are the better able, in consequence, to
attend to the grievances of your allies. My only fear is, that private
friendship for Jugurtha, too little understood, may lead any of you
astray; for his partisans, I hear, are doing their utmost in his
behalf, soliciting and importuning you individually, to pass no
decision against one who is absent, and whose cause is yet untried;
and saying that I state what is false, and only pretend to be an
exile, when I might, if I pleased, have remained still in my kingdom.
But would that I could see him,\footnote{But would that I could see him, etc.—\emph{Quod utinam illum—videam}.
The \emph{quod}, in \emph{quod utinam}, is the same as that in \emph{quod si}, which
we commonly translate, \emph{but if}. \emph{Quod}, in such expressions, serves
as a particle of connection, between what precedes and what follows it;
the Latins being fond of connection by means of relatives. See Zumpt's
Lat. Grammar on this point, Sect. 63, 82, Kenrick's translation.
Kritzius writes \emph{quodutinam}, \emph{quodsi}, \emph{quodnisi}, etc., as one word.
Cortius injudiciously interprets \emph{quod} in this passage as having
\emph{facientem} understood with it.} by whose unnatural crime I am thus
reduced to misery, pretending as I now pretend; and would that, either
with you or with the immortal gods, there may at length arise some
regard for human interests; for then assuredly will he, who is now
audacious and triumphant in guilt, be tortured by every kind of
suffering, and pay a heavy penalty for his ingratitude to my father,
for the murder of my brother, and for the distress which he has
brought upon myself.

And now, O my brother, dearest object of my affection, though thy
life has been prematurely taken from thee, and by a hand that should
have been the last to touch it, yet I think thy fate a subject for
rejoicing rather than lamentation, for, in losing life, thou hast not
been cut off from a throne, but from flight, expatriation, poverty,
and all those afflictions which now press upon me. But I, unfortunate
that I am, cast from the throne of my father into the depths of
calamity, afford an example of human vicissitudes, undecided what
course to adopt, whether to avenge thy wrongs, while I myself stand in
need of assistance, or to attempt the recovery of my kingdom, while my
life or death depends on the aid of others.\footnote{My life or death depends on the aid of others—\emph{Cujus vitae
necisque ex opibus alienis pendet}. On the aid of the Romans. Unless
they protected him, he expected to meet with the same fate as Hiempsal
at the hands of Jugurtha.}

Would that death could be thought an honorable termination to my
misfortunes, that I might not seem to live an object of contempt, if,
sinking under my afflictions, I tamely submit to injustice. But now I
can neither live with pleasure, nor can die without disgrace.\footnote{Without disgrace—\emph{Sine dedecore}. That is, if he did not succeed
in getting revenge on Jugurtha.} I
implore you, therefore, Conscript Fathers, by your regard for
yourselves,\footnote{By your regard for yourselves, etc.—I have here departed from
the text of Cortius, who reads \emph{per, vos, liberos atque parentes},
i. e. \emph{vos (obsecro) per liberos}, etc., as most critics would explain
it, though Cortius himself prefers taking \emph{vos} as the nominative case,
and joining it with \emph{subvenite}, which follows. Most other editions
have \emph{per vos, per liberos, atque parentes vestros}, to which I have
adhered. \emph{Per vos}, though an adjuration not used in modern times,
is found in other passages of the Roman writers. Thus Liv. xxix. 18:
\emph{Per vos, fidemque vestram}. Cic. pro Planc., c. 42; \emph{Per vos, per
fortunas vestras}.} for your children, and for your parents, and by the
majesty of the Roman people, to grant me succor in my distress, to
arrest the progress of injustice, and not to suffer the kingdom of
Numidia, which is your own property, to sink into ruin\footnote{To sink into ruin—\emph{Tabescere}. ``Paullatim interire.''
\emph{Cortius}. Lucret. ii. 1172: \emph{Omnia paullatim tabescere el ire Ad
capulum}. ``This speech,'' says Gerlach, ``though of less weighty
argument than the other speeches of Sallust, is composed with great
art. Neither the speaker nor his cause was adapted for the highest
flights of eloquence; but Sallust has shrouded Adherbal's weakness in
excellent language. That there is a constant recurrence to the same
topics, is no ground for blame; indeed, such recurrence could hardly
be avoided, for it is natural to all speeches in which the orator
earnestly labors to make his hearers adopt his own feelings and views.
The Romans were again and again to be supplicated, and again and again
to be reminded of the character and services of Masinissa, that they
might be induced, if not by the love of justice, yet by the dread of
censure, to relieve the distresses of his grandson…. He omits no
argument or representation that could move the pity of the Romans; and
if his abject prostration of mind appears more suitable to a woman
than a man, it is to be remembered that it is purposely introduced by
Sallust to exhibit the weakness of his character.''} through
villainy and the slaughter of our family.''

XV. When the prince had concluded his speech, the embassadors of
Jugurtha, depending more on their money than their cause, replied, in
a few words, ``that Hiempsal had been put to death by the Numidians for
his cruelty; that Adherbal, commencing war of his own accord, complained,
after he was defeated, of being unable to do injury; and that Jugurtha
entreated the senate not to consider him a different person from what
he had been known to be at Numantia, nor to set the assertions of his
enemy above his own conduct.''

Both parties then withdrew from the senate-house, and the senate
immediately proceeded to deliberate. The partisans of the embassadors,
with a great many others, corrupted by their influence, expressed
contempt for the statements of Adherbal, extolled with the highest
encomiums the merits of Jugurtha, and exerted themselves as
strenuously, with their interest and eloquence, in defense of the
guilt and infamy of another, as they would have striven for their own
honor. A few, however, on the other hand, to whom right and justice
were of more estimation than wealth, gave their opinion that Adherbal
should be assisted, and the murder of Hiempsal severely avenged. Of
all these the most forward was Aemilius Scaurus,\footnote{XV. Aemilius Scaurus—He was \emph{princeps senatus}(see c. 25),
and seems to be pretty faithfully characterized by Sallust as a man of
eminent abilities, but too avaricious to be strictly honest. Cicero,
who alludes to him in many passages with commendation (Off., i. 20,
30; Brut. 29; Pro Muraen., 7; Pro Fonteio, 7), mentions an anecdote
respecting him (De Orat. ii. 70), which shows that he had a general
character for covetousness. See Pliny, H. N, xxxvi. 14. Valerius
Maximus (iii. 7, 8) tells another anecdote of him, which shows that he
must have been held in much esteem, for whatever qualities, by the
public. Being accused before the people of having taken a bribe from
Mithridates, he made a few remarks on his own general conduct; and
added, ``Varius of Sucro says that Marcus Scaurus, being bribed with
the king's money, has betrayed the interests of the Roman people.
Marcus Scaurus denies that he is guilty of what is laid to his charge.
Which of the two do you believe?'' The people dismissed the accusation;
but the words of Scaurus may be regarded as those of a man rather
seeking to convey a notion of his innocence, than capable of proving
it. The circumstance which Cicero relates is this: Scaurus had
incurred some obloquy for having, as it was said, taken possession of
the property of a certain rich man, named Phyrgio Pompeius, without
being entitled to it by any will; and being engaged as an advocate in
some cause, Mommius, who was pleading on the opposite side, seeing a
funeral pass by at the time, said, ``Scaurus, yonder is a dead man, on
his way to the grave; if you can but get possession of his property!''
I mention these matters, because it has been thought that Sallust,
from some ill-feeling, represents Scaurus as more avaricious than he
really was.} a man of noble
birth and great energy, but factious, and ambitious of power, honor,
and wealth; yet an artful concealer of his own vices. He, seeing that
the bribery of Jugurtha was notorious and shameless, and fearing that,
as in such cases often happens, its scandalous profusion might excite
public odium, restrained himself from the indulgence of his ruling
passion.\footnote{His ruling passion—\emph{Consuetâ libidine}. Namely, avarice.}

XVI. Yet that party gained the superiority in the senate, which
preferred money and interest to justice. A decree was made, ``that ten
commissioners should divide the kingdom, which Micipsa had possessed,
between Jugurtha and Adherbal.'' Of this commission the leading person
was Lucius Opimius,\footnote{XVI. Lucius Opimius—His contention with the party of C. Gracchus
may be seen in any history of Rome. For receiving bribes from Jugurtha
he was publicly accused, and being condemned, ended his life, which
was protracted to old age, in exile and neglect. Cic. Brut. 33;
Planc. 28.} a man of distinction, and of great influence
at that time in the senate, from having in his consulship, on the
death of Caius Gracchus and Marcus Fulvius Flaccus, prosecuted the
victory of the nobility over the plebeians with great severity.

Jugurtha, though he had already counted Scaurus among his friends at
Rome, yet received him with the most studied ceremony, and, by
presents and promises, wrought on him so effectually, that he
preferred the prince's interest to his own character, honor, and all
other considerations. The rest of the commissioners he assailed in a
similar way, and gained over most of them; by a few only integrity was
more regarded than lucre. In the division of the kingdom, that part of
Numidia which borders on Mauretania, and which is superior in
fertility and population, was allotted to Jugurtha; of the other part,
which, though better furnished with harbors and buildings, was more
valuable in appearance than in reality, Adherbal became the possessor.

XVII. My subject seems to require of me, in this place, a brief
account of the situation of Africa, and of those nations in it with
whom we have had war or alliances. But of those tracts and countries,
which, from their heat, or difficulty of access, or extent of desert,
have been but little visited, I can not possibly give any exact
description. Of the rest I shall speak with all possible brevity.

In the division of the earth, most writers consider Africa as a third
part; a few admit only two divisions, Asia and Europe,\footnote{XVII. Only two divisions, Asia and Europe—Thus Varro, de L.
L. iv.13, ed. Bip. ``As all nature is divided into heaven and earth, so
the heaven is divided into regions, and the earth into Asia and
Europe.'' See Bronkh. ad Tibull., iv. 1, 176.} and include
Africa in Europe. It is bounded, on the west, by the strait connecting
our sea with the ocean;\footnote{The strait connecting our sea with the ocean—\emph{Fretum nostri
maris et oceani}. That is, the \emph{Fretum Gaditanum}, or Strait of
Gibraltar. By \emph{our see}, he means the Mediterranean. See Pomp. Mela,
i. 1.} on the east, by a vast sloping tract,
which the natives call the Catabathmos.\footnote{A vast sloping tract—Catabathmos—\emph{Declivem latitudinem,
quem locum Catabathmon incolae appellant. Catabathmus—vallis repente
convexa}, Plin. H. N. v 5. \emph{Catabathmus, vallis devexa in Aegyptum},
Pomp. Mela, i. 8. I have translated \emph{declivem latitudinem} in
conformity with these passages. \emph{Catabathmus}, a Greek word, means \emph{a
descent}. There were two, the \emph{major} and \emph{minor}; Sallust speaks of
the \emph{major}.} The sea is boisterous, and
deficient in harbors; the soil is fertile in corn, and good for
pasturage, but unproductive of trees. There is a scarcity of water
both from rain and from landsprings. The natives are healthy, swift of
foot, and able to endure fatigue. Most of them die by the gradual
decay of age,\footnote{Most of them die by the gradual decay of age—\emph{Plerosque
senectus dissolvit} ``A happy expression; since the effect of old age
on the bodily frame is not to break it in pieces suddenly, but to
dissolve it, as it were, gradually and imperceptibly.'' \emph{Burnouf}.} except such as perish by the sword or beasts of
prey; for disease finds but few victims. Animals of a venomous nature
they have in great numbers.

Concerning the original inhabitants of Africa, the settlers that
afterward joined them, and the manner in which they intermingled, I
shall offer the following brief account, which, though it differs from
the general opinion, is that which was interpreted to me from the Punic
volumes said to have belonged to King Hiempsal\footnote{King Hiempsal—``This is not the prince that was murdered by
Jugurtha, but the king who succeeded him; he was grandson of
Masinissa, son of Gulussa, and father of Juba. After Juba was killed
at Thapsus, Caesar reduced Numidia to the condition of a province, and
appointed Sallust over it, who had thus opportunities of gaining a
knowledge of the country, and of consulting the books written in the
language of it.'' \emph{Burnouf}.}, and which the
inhabitants of that country believe to be consistent with fact. For
the truth of the statement, however, the writers themselves must be
responsible.

XVIII. Africa, then, was originally occupied by the Getulians and
Libyans,\footnote{XVIII. Getulians and Lybians—\emph{Gaetuli et Libyes}, ``See
Pompon. Mel. i. 4; Plin. H. N. v. 4, 6, 8, v. 2, xxi. 13; Herod, iv.
159, 168.'' \emph{Gerlach}. The name \emph{Gaetuli}, is, however, unknown to
Herodotus. They lay to the south of Numidia and Mauretania. See
Strabo, xvii. 3. \emph{Libyes} is a term applied by the Greek writers
properly to the Africans of the North coast, but frequently to the
inhabitants of Africa in general.} rude and uncivilized tribes, who subsisted on the flesh
of wild animals, or, like cattle, on the herbage of the soil. They
were controlled neither by customs, laws, nor the authority of any
ruler; they wandered about, without fixed habitations, and slept in
the abodes to which night drove them. But after Hercules, as the
Africans think, perished in Spain, his army, which was composed of
various nations,\footnote{His army, which was composed of various nations—This seems
to have been an amplification of the adventure of Hercules with
Geryon, who was a king in Spain. But all stories that make Hercule
a leader of armies appear to be equally fabulous.} having lost its leader, and many candidates
severally claiming the command of it, was speedily dispersed. Of its
constituent troops, the Medes, Persians, and Armenians,\footnote{Medes, Persians and Armenians—De Brosses thinks that these
were not real Medes, etc., but that the names were derived from
certain companions of Hercules. The point is not worth discussion.} having
sailed over into Africa, occupied the parts nearest to our sea.\footnote{Our sea—The Mediterranean. See above, c. 17.}
The Persians, however, settled more toward the ocean,\footnote{More toward the Ocean—\emph{Intra oceanum magis}. ``\emph{Intra oceanum}
is differently explained by different commentators. Cortius, Muller
and Gerlach, understand the parts bounded by the ocean, lying close
upon it, and stretching toward the west; while Langius thinks that
the regions more remote from the Atlantic Ocean, and extending
toward the east, are meant. But Langius did not consider that those
who had inverted keels of vessels for cottages, could not have
strayed far from the ocean, but must have settled in parts
bordering upon it\_. And this is what is signified by \emph{intra oceanum}.
For \emph{intra aliquam rem} is not always used to denote what is actually
\emph{in a thing,} and circumscribed by its boundaries, but what approaches
toward it, and reaches close to it.'' Kritzius. He then instances
\emph{intra modum, intra legem; Hortensii scripta intra famam sunt},
Quintil. xi. 8, 8. But the best example which he produces is Liv. xxv.
11: \emph{Fossa ingens ducta, et vallum}, intra eam \emph{erigitur}. Cicero, in
Verr. iii. 89, has also, he notices, the same, expression, \emph{Locus}
intra oceanum \emph{jam nullus est—quò non nostrorum hominum libido
iniquitasque pervaserit}, i. e.. \emph{locus oceano conterminus}. Burnouf
absurdly follows Langius.} and used the
inverted keels of their vessels for huts, there being no wood in the
country, and no opportunity of obtaining it, either by purchase or
barter, from the Spaniards; for a wide sea, and an unknown tongue,
were barriers to all intercourse. These, by degrees, formed
intermarriages with the Getulians; and because, from constantly trying
different soils, they were perpetually shifting their abodes, they
called themselves NUMIDIANS.\footnote{Numidians—\emph{Numidas}. The same as \emph{Nomades}, or wanderers; a
term applied to pastoral nations, and which, as Kritzius observes,
the Africans must have had from the Greeks, perhaps those of Sicily.} And to this day the huts of the
Numidian boors, which they call \emph{mapalia}, are of an oblong shape,
with curved roofs; resembling the hulls of ships.

The Medes and Armenians connected themselves with the Libyans, who
dwelled near the African sea; while the Getulians lay more to the
sun,\footnote{More to the sun—\emph{sub sole magis}. I have borrowed this
expression from Rose. The Getulians were more southward.} not far from the torrid heats; and these soon built
themselves towns,\footnote{These soon built themselves towns—That is, the united Medes,
Armenians, and Libyans.} as, being separated from Spain only by a strait,
they proceeded to open an intercourse with its inhabitants. The name
of Medes the Libyans gradually corrupted, changing it, in their
barbarous tongue, into Moors.\footnote{Medes—into Moors—\emph{Mauris pro Medis}. A most improbable,
not to say impossible corruption.}

Of the Persians\footnote{Of the Persians—\emph{Persarum}. That is, of the Persians and
Getulians united.} the power rapidly increased; and at length, the
children, through excess of population, separating from the parents,
they took possession, under the name of Numidians, of those regions
bordering on Carthage which are now called Numidia. In process of
time, the two parties,\footnote{The two parties—\emph{Utrique}. The older Numidians, and the
younger, who had emigrated toward Carthage.} each assisting the other, reduced the
neighboring tribes, by force or fear, under their sway; but those who
had spread toward our sea, made the greater conquests: for the Lybians
are less warlike than the Getulians\footnote{Those who had spread toward our sea—for the Libyans are
less warlike than the Getulians—\emph{Magis hi, qui ad nostrum mare
processerant; quia Libyes quam Gaetuli minus bellicosi}. The Persians
and Getulians (under the name of Numidians), and their colonists, who
were more toward the Mediterranean, and were more warlike than the
Libyans (who were united with the Medes and Armenians) took from them
portions of their territories by conquest. This is clearly the sense,
as deducible from the preceding portion of the text.} At last nearly all lower
Africa\footnote{Lower Africa—\emph{Africa pars inferior}. The part nearest to the
sea. The ancients called the maritime parts of a country \emph{the lower
parts}, and the inland parts \emph{the higher}, taking the notion, probably,
from the course of the rivers. Lower Egypt was the part at the mouth of
the Nile.} was occupied by the Numidians; and all the conquered tribes
were merged in the nation and name of their conquerors.

XIX. At a later period, the Phoenicians, some of whom wished to lessen
their numbers at home, and others, ambitious of empire, engaged the
populace, and such as were eager for change, to follow them, founded
Hippo,\footnote{XIX. Hippo—``It is not Hippo Regius'' (now called \emph{Bona}) ``that is
meant, but another Hippo, otherwise called \emph{Diarrhytum} or \emph{Zarytum},
situate in Zengitana, not far from Utica. This is shown by the order
in which the places are named, as has already been observed by Cortius.''
\emph{Kritzius}.} Adrumetum, Leptis,\footnote{Leptis—There were two cities of this name. Leptis Major, now
\emph{Lebida}, lay between the two Syrtes; Leptis Minor, now \emph{Lempta},
between the smaller Sytis and Carthage. It is the latter that is meant
here, and in c. 77, 78.} and other cities, on the sea-coast;
which, soon growing powerful, became partly a support, and partly an
honor, to their parent state. Of Carthage I think it better to be
silent, than to say but little; especially as time bids me hasten to
other matters.

Next to the Catabathmos,\footnote{Next to the Catabathmos—\emph{Ad Catabathmon}. \emph{Ad} means, on the
side of the country toward the Catabathmos. ``Catabathmon \emph{initium}
ponens Sallustius ab eo \emph{discedit}.'' Kritzius.} then, which divides Egypt from Africa,
the first city along the sea-coast\footnote{Along the sea-coast—\emph{Secundo mari}. ``Si quis secundum mare
pergat'' \emph{Wasse.}} is Cyrene, a colony of
Theraeans;\footnote{Of Therseans—\emph{Theraeon}. From the island of Thera, one of the
Sporades, in the Aegean Sea, now called \emph{Santorin}. Battus was the
leader of the colony. See Herod., iv. 145; Strab., xvii. 8; Pind.
Pyth., iv.} after which are the two Syrtes,\footnote{Two Syrtes—See c. 78.} with Leptis\footnote{Leptis—That is, \emph{Leptis Major}. See above on this c.}
between them; then the Altars of the Philaeni,\footnote{Altars of the Philaeni—see c. 79.} which the
Carthaginians considered the boundary of their dominion on the side of
Egypt; beyond these are the other Punic towns. The other regions, as
far as Mauretania, the Numidians occupy; the Moors are nearest to
Spain. To the south of Numidia,\footnote{To the south of Numidia—\emph{Super Numidiam}. ``Ultra Numidiam,
meridiem versus.'' \emph{Burnouf}.} as we are informed, are the
Getulians, of whom some live in huts, and others lead a vagrant and
less civilized life; beyond these are the Ethiopians; and further on,
regions parched by the heat of the sun.

At the time of the Jugurthine war, most of the Punic towns, and the
territories which Carthage had lately possessed,\footnote{Had lately possessed—\emph{Novissimè habuerant}. In the interval
between the second and third Punic wars.} were under the
government of Roman praetors; a great part of the Getulians, and
Numidia as far as the river Mulucha, were subject to Jugurtha; while
the whole of the Moors were governed by Bocchus, a king who knew
nothing of the Romans but their name, and who, before this period,
was as little known to us, either in war or peace. Of Africa and its
inhabitants I have now said all that my narrative requires.

XX. When the commissioners, after dividing the kingdom, had left
Africa, and Jugurtha saw that, contrary to his apprehensions, he had
obtained the object of his crimes; he then being convinced of the
truth of what he had heard from his friends at Numantia, ``that all
things were purchasable at Rome,'' and being also encouraged by the
promises of those whom he had recently loaded with presents, directed
his views to the domain of Adherbal. He was himself bold and warlike,
while the other, at whose destruction he aimed, was quiet, unfit for
arms, of a mild temper, a fit subject for injustice, and a prey to
fear rather than an object of it. Jugurtha, accordingly, with a
powerful force, made a sudden irruption into his dominions, took
several prisoners, with cattle and other booty, set fire to the
buildings, and made hostile demonstrations against several places with
his cavalry. He then retreated, with all his followers, into his own
kingdom, expecting that Adherbal, roused by such provocation, would
avenge his wrongs by force, and thus furnish a pretext for war. But
Adherbal, thinking himself unable to meet Jugurtha in the field, and
relying on the friendship of the Romans more than on the Numidians,
merely sent embassadors to Jugurtha to complain of the outrage; and,
although they brought back but an insolent reply, yet he resolved to
endure any thing rather than have recourse to war, which, when he
attempted it before, had ended in his defeat. By such conduct the
eagerness of Jugurtha was not at all allayed; for he had now, indeed,
in imagination, possessed himself of all Adherbal's dominions. He
therefore renewed hostilities, not, as before, with a predatory band,
but at the head of a large army which he had collected, and openly
aspired to the sovereignty of all Numidia. Wherever he marched, he
ravaged the towns and the fields, drove off booty, and raised
confidence in his own men and dismay among the enemy.

XXI. Adherbal, when he found that matters had arrived at such a point,
that he must either abandon his dominions, or defend them by force of
arms, collected an army from necessity, and advanced to meet Jugurtha.
Both armies took up\footnote{XXI. Both armies took up, etc.—I have omitted the word
\emph{interim} at the beginning of this sentence, as it would be worse than
useless in the translation. It signifies, \emph{during the interval before
the armies came to an engagement}; but this is sufficiently expressed
at the termination of the sentence.} their position near the town of Cirta\footnote{Cirta—Afterward named \emph{Sittianorum Colonia}, from P. Sittius
Nucerinus (mentioned in Cat., c. 21), who assisted Caesar in the
African war, and was rewarded by him with the possession of this
city and its lands. It is now called \emph{Constantina}, from Constantine
the Great, who enlarged and restored it when it had fallen into decay.
Strabo describes it, xvii. 3.}, at
no great distance from the sea; but, as evening was approaching,
encamped without coming to an engagement. But when the night was far
advanced, and twilight was beginning to appear\footnote{Twilight was beginning to appear—\emph{Obscuro etiam tum lumine}.
Before day had fairly dawned.}, the troops of
Jugurtha, at a given signal, rushed into the camp of the enemy, whom
they routed and put to flight, some half asleep and others resuming
their arms. Adherbal, with a few of his cavalry, fled to Cirta; and,
had there not been a number of Romans\footnote{Romans—\emph{Togatorum}. Romans, with, perhaps, some of the
allies, engaged in merchandise, or other peaceful occupations, and
therefore wearing the \emph{toga}. They are called \emph{Italici} in c. 26.} in the town, who repulsed
his Numidian pursuers from the walls, the war between the two princes
would have been begun and ended on the same day.

Jugurtha proceeded to invest the town, and attempted to storm it with
the aid of mantelets, towers, and every kind of machines; being
anxious above all things, to take it before the ambassadors could
arrive at Rome, who, he was informed, had been dispatched thither by
Adherbal before the battle was fought. But as soon as the senate heard
of their contention, three young men\footnote{Three young men—\emph{Tres adolescentes}. Cortius includes these
words in brackets, regarding them as the insertion of some sciolist. But
a sciolist, as Burnouf observes, would hardly have thought of inserting
\emph{tres adolescentes}. The words occur in all the MSS., and are pretty
well confirmed by what is said below, c. 25, that when the senate next
sent a deputation, they took care to make it consist of \emph{majores natu,
nobiles}. See on \emph{adolescens}, Cat., c. 38.} were sent as deputies into
Africa, with directions to go to both of the princes, and to announce
to them, in the words of the senate and people of Rome, ``that it was
their will and resolution that they should lay down their arms, and
settle their disputes rather by arbitration than by the sword; since
to act thus would be to the honor both of the Romans and themselves.''

XXII. These deputies soon arrived in Africa, using the greater
dispatch, because, while they were preparing for their journey, a
report was spread at Rome of the battle which had been fought, and of
the siege of Cirta; but this report told much less than the truth\footnote{XXII. Told much less than the truth—\emph{Sed is rumor clemens erat}.
``It fell below the truth, not telling the whole of the atrocity that
had been committed.'' \emph{Gruter.} ``Priscian (xviii. 26) interprets
\emph{clemens} 'non nimius,' alluding to this passage of Sallust.''
\emph{Kritzius.} All the later commentators have adopted this interpretation,
except Burnouf, who adopts the supposition of Ciacconius, that a \emph{vague
and uncertain} rumor is meant.}
Jugurtha, having given them an audience, replied, ``that nothing was of
greater weight with him, nothing more respected, than the authority of
the senate; that it had been his endeavor, from his youth, to deserve
the esteem of all men of worth; that he had gained the favor of
Publius Scipio, a man of the highest eminence, not by dishonorable
practices, but by merit; that, for the same good qualities, and not
from want of heirs to the throne, he had been adopted by Micipsa; but
that, the more honorable and spirited his conduct had been, the less
could his feelings endure injustice; that Adherbal had formed designs
against his life on discovering which, he had counteracted his malice;
that the Romans would act neither justly nor reasonably, if they
withheld from him the common right of nations;\footnote{XXIV. Pays no regard—\emph{Neque—in animo habeat.} This letter of
Adherbal's, both in matter and tone, is very similar to his speech
in c. 14.} and, in conclusion,
that he would soon send embassadors to Rome to explain the whole of
his proceedings.'' On this understanding, both parties separated. Of
addressing Adherbal the deputies had no opportunity.

XXIII. Jugurtha, as soon as he thought that they had quitted Africa,
surrounded the walls of Cirta, which, from the nature of its
situation, he was unable to take by assault, with a rampart and a
trench; he also erected towers, and manned them with soldiers; he made
attempts on the place, by force or by stratagem, day and night; he
held out bribes, and some times menaces, to the besieged; he roused
his men, by exhortations, to efforts of valor, and resorted, with the
utmost perseverance, to every possible expedient.

Adherbal, on the other hand, seeing that his affairs were in a
desperate condition, that his enemy was determined on his ruin, that
there was no hope of succor, and that the siege, from want of
provisions, could not long be protracted, selected from among those
who had fled with him to Cirta, two of his most resolute supporters,
whom he induced, by numerous promises, and an affecting representation
of his distress, to make their way in the night, through the enemy's
lines, to the nearest point of the coast, and from thence to Rome.

XXIV. The Numidians, in a few days executed their commission; and a
letter from Adherbal was read in the senate, of which the following
was the purport:

``It is not through my own fault, Conscript Fathers, that I so often
send requests to you; but the violence of Jugurtha compels me; whom so
strong a desire for my destruction has seized, that he pays no
regard\footnote{XXIV. Pays no regard—\emph{Neque—in animo habeat.} This letter of
Adherbal's, both in matter and tone, is very similar to his speech
in c. 14.} either to you or to the immortal gods; my blood he covets
beyond every thing. Five months, in consequence, have I, the ally and
friend of the Roman people, been besieged with an armed force; neither
the remembrance of my father Micipsa's benefits, nor your decrees, are
of any avail for my relief; and whether I am more closely pressed by
the sword, or by famine, I am unable to say.

From writing further concerning Jugurtha, my present condition deters
me; for I have experienced, even before,\footnote{I have experienced, even before—\emph{Jam antea expertus sum.}
He means, in the result of his speech to the senate.} that little credit is
given to the unfortunate. Yet I can perceive that his views extend
further than to myself, and that he does not expect to possess, at the
same time, your friendship and my kingdom; which of the two he thinks
the more desirable, must be manifest to every one. For, in the first
place, he murdered my brother Hiempsal; and, in the next, expelled me
from my dominions; which, however, may be regarded as our own wrongs,
and as having no reference to you. But now he occupies your kingdom
with an army; he keeps me, whom you appointed a king over the
Numidians, in a state of blockade; and in what estimation he holds the
words of your embassadors, my perils may serve to show. What then is
left, except your arms, that can make an impression upon him?

I could wish, indeed, that what I now write, as well as the complaints
which I lately made before the senate, were false, rather than that my
present distresses should confirm the truth of my statements. But
since I am born to be an example of Jugurtha's villainy, I do not now
beg a release from death or distress, but only from the tyranny of an
enemy, and from bodily torture. Respecting the kingdom of Numidia,
which is your own property, determine as you please, but if the memory
of my grandfather Masinissa is still cherished by you, deliver me, I
entreat you, by the majesty of your empire, and by the sacred ties of
friendship, from the inhuman hands of Jugurtha.''

XXV. When this letter was read, there were some who thought that an
army should be dispatched into Africa, and relief afforded to
Adherbal, as soon as possible; and that the senate, in the mean time,
should give judgment on the conduct of Jugurtha, in not having obeyed
the embassadors. But by the partisans of Jugurtha, the same that had
before supported his cause, effectual exertions were made to prevent
any decree from being passed; and thus the public interest, as is too
frequently the case, was defeated by private influence.

An embassy was, however, dispatched into Africa, consisting of men of
advanced years, and of noble birth, and who had filled the highest
offices of the state; among whom was Marcus Scaurus, already mentioned,
a man who had held the consulship, and who was at that time chief of
the senate\footnote{XXV. Chief of the senate—\emph{Princeps senatûs.} ``He whose name
was first entered in the censors' books was called \emph{Princeps Senatûs},
which title used to be given to the person who of those alive had been
censor first (\emph{qui primus censor, ex iis qui viverent, fuisset}), but
after the year 544, to him whom the censors thought most worthy, Liv.,
xxvii. 13. This dignity, although it conferred no command or emolument,
was esteemed the very highest, and was usually retained for life, Liv.,
xxxiv. 44; xxxix. 52. It is called \emph{Principatus}; and hence afterward
the Emperor was named \emph{Princeps}, which word properly denotes rank, and
not power.'' Adam's Rom. Antiq., p. 3.}. These embassadors, as their business was an affair of
public odium, and as they were urged by the entreaties of the Numidians,
embarked in three days; and having soon arrived at Utica, sent a letter
from thence to Jugurtha, desiring him ``to come to the province as
quickly as possible, as they were deputed by the senate to meet him.''

Jugurtha, when he found that men of eminence, whose influence at Rome
he knew to be powerful, were come to put a stop to his proceedings,
was at first perplexed, and distracted between fear and cupidity. He
dreaded the displeasure of the senate, if he should disobey the
embassadors; while his eager spirit, blinded by the lust of power,
hurried him on to complete the injustice which he had begun. At length
the evil incitements of ambition prevailed\footnote{At length the evil incitements of ambition prevailed—\emph{Vicit
tamen in avido ingenio pravum consilium.} ``Evil propensities gained
the ascendency in his ambitious disposition.''}. He accordingly drew
his army round the city of Cirta, and endeavored, with his utmost
efforts, to force an entrance; having the strongest hopes, that, by
dividing the attention of the enemy's troops, he should be able, by
force or artifice, to secure an opportunity of success. When his
attempts, however, were unavailing, and he found himself unable, as
he had designed, to get Adherbal into his power before he met the
embassadors, fearing that, by further delay, he might irritate
Scaurus, of whom he stood in great dread, he proceeded with a small
body of cavalry into the Province. Yet, though serious menaces were
repeated to him in the name of the senate, because he had not desisted
from the siege, nevertheless, after spending a long time in conference,
the embassadors departed without making any impression upon him.

XXVI. When news of this result was brought to Cirta, the Italians\footnote{XXVI. The Italians—\emph{Italici.} See c. 21.},
by whose exertions the city had been defended, and who trusted that, if
a surrender were made, they would be able, from respect to the greatness
of the Roman power, to escape without personal injury, advised Adherbal
to deliver himself and the city to Jugurtha, stipulating only that his
life should be spared, and leaving all other matters to the care of the
senate. Adherbal, though he thought nothing less trustworthy than the
honor of Jugurtha, yet, knowing that those who advised could also compel
him if he resisted, surrendered the place according to their desire.
Jugurtha immediately proceeded to put Adherbal to death with torture,
and massacred all the inhabitants that were of age, whether Numidians
or Italians, as each fell in the way of his troops.

XXVII. When this outrage was reported at Rome, and became a matter of
discussion in the senate, the former partisans of Jugurtha applied
themselves, by interrupting the debates and protracting the time,
sometimes exerting their interest, and sometimes quarreling with
particular members, to palliate the atrocity of the deed. And had not
Caius Memmius, one of the tribunes of the people elect, a man of
energy, and hostile to the power of the nobility, convinced the people
of Rome that an attempt was being made, by the agency of a small
faction, to have the crimes of Jugurtha pardoned, it is certain that
the public indignation against him would have passed off under the
protraction of the debates; so powerful was party interest, and the
influence of Jugurtha's money. When the senate, however, from
consciousness of misconduct, became afraid of the people, Numidia and
Italy, by the Sempronian law,\footnote{XXVII. By the Sempronian law—\emph{Lege Semproniâ.} This was
the \emph{Lex Sempronia de Provinciis.} In the early ages of the republic,
the provinces were decreed by the senate to the consuls after they
were elected; but by this law, passed A.U.C. 631, the senate fixed on
two provinces for the future consuls before their election (Cic. Pro
Dom., 9; De Prov. Cons., 2), which they, after entering on their
office, divided between themselves by lot or agreement. The law was
passed by Caius Gracchus. See Adam's Rom. Antiq., p. 105.} were appointed as provinces to the
succeeding consuls, who were declared to be Publius Scipio Nasica\footnote{Publius Scipio Nasica—``The great-grandson of him who was
pronounced by the senate to be \emph{vir optimus}; and son of him who,
though holding no office at the time, took part in putting to death
Tiberius Gracchus. He was consul with Bestia, A.U.C. 643, and died in
his consulship. Cic. Brut., 34.'' \emph{Burnouf.}},
and Lucius Bestia Calpurnius\footnote{Lucius Bestia Calpurnius—``He had been on the side of the
nobility against the Gracchi, and was therefore in favor with the
senate. After his consulship he was accused and condemned by the
Mamilian law (c. 40), for having received money from Jugurtha, Cic.
Brut. c. 34. De Brosses thinks that he was the grandfather of that
Bestia who was engaged in the conspiracy of Catilina.'' Burnouf.\_}. Numidia fell to Calpurnius, and Italy
to Scipio. An army was then raised to be sent into Africa; and pay, and
all other necessaries of war, were decreed for its use.

XXVIII. When Jugurtha received this news, which was utterly at
variance with his expectations, as he had felt convinced that all
things were purchasable at Rome, he sent his son, with two of his
friends, as deputies to the senate, and directed them, like those whom
he had sent on the murder of Hiempsal, to attack every body with
bribes. Upon the approach of these deputies to Rome, the senate was
consulted by Bestia, whether they would allow them to be admitted
within the gates; and the senate decreed, ``that, unless they came to
surrender Jugurtha's kingdom and himself, they must quit Italy within
the ten following days.'' The consul directed this decree to be
communicated to the Numidians, who consequently returned home without
effecting their object.

Calpurnius, in the mean time, having raised an army, chose for his
officers men of family and intrigue, hoping that whatever faults he
might commit, would be screened by their influence; and among these
was Scaurus, of whose disposition and character we have already
spoken. There were, indeed, in our consul Calpurnius, many excellent
qualities, both mental and personal, though avarice interfered with
the exercise of them; he was patient of labor, of a penetrating
intellect, of great foresight, not inexperienced in war, and extremely
vigilant against danger and surprise.

The troops were conducted through Italy to Rhegium, from thence to
Sicily, and from Sicily into Africa; and Calpurnius's first step,
after collecting provisions, was to invade Numidia with spirit, where
he took many prisoners, and several towns, by force of arms.

XXIX. But when Jugurtha began, through his emissaries, to tempt him
with bribes, and to show the difficulties of the war which he had
undertaken to conduct, his mind, corrupted with avarice, was easily
altered. His accomplice, however, and manager in all his schemes, was
Scaurus; who, though he had at first, when most of his party were
corrupted, displayed violent hostility to Jugurtha, yet was afterward
seduced, by a vast sum of money, from integrity and honor to injustice
and perfidy. Jugurtha, however, at first sought only to purchase a
suspension of hostilities, expecting to be able, during the interval,
to make some favorable impression, either by bribery or by interest,
at Rome; but when he heard that Scaurus was co-operating with
Calpurnius, he was elated with great hopes of regaining peace, and
resolved upon a conference with them in person respecting the terms of
it. In the mean time, for the sake of giving confidence \footnote{XXIX. For the sake of giving confidence—\emph{Fidei causâ.} ``In
order that Jugurtha might have confidence in Bestia, Sextius the
quaestor was sent as a sort of hostage into one of Jugurtha's towns.''
\emph{Cortius.}} to
Jugurtha, Sextus the quaestor was dispatched by the consul to Vaga,
one of the prince's towns; the pretext for his journey being the
receiving of corn, which Calpurnius had openly demanded from Jugurtha's
emissaries, on the ground that a truce was observed through their delay
to make a surrender. Jugurtha then, as he had determined, paid a visit
to the consul's camp, where, having made a short address to the council,
respecting the odium cast upon his conduct, and his desire for a
capitulation, he arranged other matters with Bestia and Scaurus in
secret; and the next day, as if by an evident majority of voices\footnote{As if by an evident majority of voices—\emph{Quasi per saturam
exquisitis sententiis.} ``The opinions being taken in a confused
manner,'' or, as we say, \emph{in the lump}. The sense manifestly is, that
there was (or was said to be) such a preponderating majority in
Jugurtha's favor, that it was not necessary to ask the opinion of each
individual in order. \emph{Satura}, which some think to be always an
adjective, with \emph{lanx} understood, though \emph{lanx}, according to
Scheller, is never found joined with it in ancient authors, was \emph{a
plate filled with various kinds of fruit, such as was annually offered
to the gods.} ``Lanx plena diversis frugibus in templum Cereris
infertur, quae satura nomine appellatur,'' Acron. ad Hor. Sat., i. 1,
\emph{init}. ``Lanx. referta variis multisque primitiis, sacris Cereris
inferebatur,'' Diomed., iii. p. 483.''Satura, cibi genus ex variis rebus
conditum,'' Festus \emph{sub voce}. See Casaubon. de Rom. Satirâ, ii. 4;
Kritzius ad h. l., and Scheller's Lex. v., \emph{Satur.} In the Pref. to
Justinian's Pandects, that work is called \emph{opus sparsim et quasi per
saturam collectum, utíle cum inutilibus mixtim.}},
he was formally allowed to surrender. But, as was demanded in the
hearing of the council, thirty elephants, a considerable number of
cattle and horses, and a small sum of money, were delivered into the
hands of the quaestor. Calpurnius then returned to Rome to preside at
the election of magistrates\footnote{To preside at the election of magistrates—\emph{Ad magistratus
rogandos.} The presiding magistrate had \emph{to ask} the consent of the
people, saying \emph{Velitis, jubeatis—rogo Quirites.}}, and peace was observed throughout
Numidia and the Roman army.

XXX. When rumor had made known the affairs transacted in Africa, and
the mode in which they had been brought to pass, the conduct of the
consul became a subject of discussion in every place and company at
Rome. Among the people there was violent indignation; as to the
senators, whether they would ratify so flagitious a proceeding, or
annul the act of the consul, was a matter of doubt. The influence of
Scaurus, as he was said to be the supporter and accomplice of Bestia,
was what chiefly restrained the senate from acting with justice and
honor. But Caius Memmius, of whose boldness of spirit, and hatred to
the power of the nobility, I have already spoken, excited the people
by his harangues, during the perplexity and delay of the senators, to
take vengeance on the authors of the treaty; he exhorted them not to
abandon the public interest or their own liberty; he set before them
the many tyrannical and violent proceedings of the nobles, and omitted
no art to inflame the popular passions. But as the eloquence of
Memmius, at that period, had great reputation and influence I have
thought proper to give in full\footnote{XXX. To give in full—\emph{Perscribere.} ``To write at length.''
The reader might suppose, at first, that Sallust transcribed this
speech from some publication; but in that case, as Burnouf observes,
he would rather have said \emph{ascribere.} Besides, the following
\emph{hujuscemodi} shows that Sallust did not profess to give the exact
words of Memmius. And the speech is throughout marked with Sallustian
phraseology. ``The commencement of it, there is little doubt, is
imitated from Cato, of whose speech \emph{De Lusitanis} the following
fragment is extant in Aul. Gell. xiii. 24: \emph{Multa me dehortata sunt
huc prodere, anni, aetas, vox vires, senectus.}'' Kritzius.} one out of many of his speeches;
and I take, in preference to others, that which he delivered in the
assembly of the people, after the return of Bestia, in words to the
following effect:

XXXI. ``Were not my zeal for the good of the state, my fellow-citizens,
superior to every other feeling, there are many considerations which
would deter me from appearing in your cause; I allude to the power of
the opposite party, your own tameness of spirit, the absence of all
justice, and, above all, the fact that integrity is attended with more
danger than honor. Indeed, it grieves me to relate, how, during the
last fifteen years\footnote{XXXI. During the last fifteen years—\emph{His annis quindecim.}
``It was at this time, A.U.C. 641, twenty-two years since the death of
Tiberius Gracchus, and ten since that of Caius; Sallust, or Memmius,
not to appear to make too nice a computation, takes a mean.''
\emph{Burnouf.} The manuscripts, however, vary; some read \emph{fifteen}, and
others \emph{twelve}. Cortius conjectured \emph{twenty}, as a rounder number,
which Kritzius and Dietsch have inserted in their texts. \emph{Twenty} is
also found in the Editio Victoriana, Florence, 1576.}, you have been a sport to the arrogance of an
oligarchy; how dishonorably, and how utterly unavenged, your defenders
have perished\footnote{Your defenders have perished—\emph{Perierint vestri defensores.}
Tiberius and Caius Gracchus, and their adherents.}; and how your spirit has become degenerate by sloth
and indolence; for not even now, when your enemies are in your power,
will you rouse yourselves to action, but continue still to stand in
awe of those to whom you should be a terror.

Yet, notwithstanding this state of things, I feel prompted to make an
attack on the power of that faction. That liberty of speech\footnote{Liberty of speech—\emph{Libertatem.} Liberty of speech is evidently
intended.},
therefore, which has been left me by my father, I shall assuredly
exert against them; but whether I shall use it in vain, or for your
advantage, must, my fellow-citizens, depend upon yourselves. I do not,
however, exhort you, as your ancestors have often done, to rise in
arms against injustice.

There is at present no need of violence, no need of secession; for
your tyrants must work their fall by their own misconduct.

After the murder of Tiberius Gracchus, whom they accused of aspiring
to be king, persecutions were instituted against the common people of
Rome; and after the slaughter of Caius Gracchus and Marcus Fulvius,
many of your order were put to death in prison. But let us leave these
proceedings out of the question; let us admit that to restore their
rights to the people, was to aspire to sovereignty; let us allow that
what can not be avenged without shedding the blood of citizens, was
done with justice. You have seen with silent indignation, however, in
past years, the treasury pillaged; you have seen kings, and free
people, paying tribute to a small party of Patricians, in whose hands
were both the highest honors and the greatest wealth; but to have
carried on such proceedings with impunity, they now deem but a small
matter; and, at last, your laws and your honor, with every civil and
religious obligation\footnote{Every civil and religions obligation—\emph{Divina et humana omnia.}
``They offended against the laws, when they took bribes from an
enemy; against the honor of Rome when they did what was unworthy of
it, and greatly to its injury; and against gods and men, against all
divine and human obligations, when they granted to a wicked prince not
only impunity, but even rewards, for his crimes.'' \emph{Dietsch.}}, have been sacrificed for the benefit of
your enemies. Nor do they, who have done these things, show either
shame or contrition, but parade proudly before your faces, displaying
their sacerdotal dignities, their consulships, and some of them their
triumphs, as if they regarded them as marks of honor, and not as
fruits of their dishonesty. Slaves, purchased with money\footnote{Slaves purchased with money, etc.—\emph{Servi, aere parati,} etc.
This is taken from another speech of Cato, of which a portion is
preserved in Aul. Gell. x. 3: \emph{Servi injurias nimis aegre ferunt; quid
illos bono genere natos, magnâ virtute praeditos, animi habuisse atque
habituros, dum vivent?} ``Slaves are apt to be too impatient of
injuries; and what feelings do you think that men of good family, and
of great merit, must have had, and will have as long as they live?''}, will
not submit to unjust commands from their masters; yet you, my
fellow-citizens, who are born to empire, tamely endure oppression.

But who are these that have thus taken the government into their
hands? Men of the most abandoned character, of blood-stained hands, of
insatiable avarice, of enormous guilt, and of matchless pride; men by
whom integrity, reputation, public spirit\footnote{Public spirit—\emph{Pietas.} Under this word are included all
duties that we ought to perform to those with whom we are intimately
connected, or on whom we are dependent, as our parents, our country,
and the gods. I have borrowed my translation of the word from Rose.}, and indeed every
thing, whether honorable or dishonorable, is converted to a means of
gain. Some of them make it their defense that they have killed
tribunes of the people; others, that they have instituted unjust
prosecutions; others, that they have shed your blood; and thus, the
more atrocities each has committed, the greater is his security; while
your oppressors, whom the same desires, the same aversions, and the
same fears, combine in strict union (a union which among good men is
friendship, but among the bad confederacy in guilt), have excited in
you, through your want of spirit, that terror which they ought to feel
for their own crimes.

But if your concern to preserve your liberty were as great as their
ardor to increase their power of oppression, the state would not be
distracted as it is at present; and the marks of favor which proceed
from you\footnote{The marks of favor which proceed from you—\emph{Beneficia vestra.}
Offices of state, civil and military.}, would be conferred, not on the most shameless, but on
the most deserving. Your forefathers, in order to assert their rights
and establish their authority, twice seceded in arms to Mount
Aventine; and will not you exert yourselves, to the utmost of your
power, in defense of that liberty which you received from them? Will
you not display so much the more spirit in the cause, from the
reflection that it is a greater disgrace to lose\footnote{A greater disgrace to lose, etc.—\emph{Quod majus dedecus est
parta amitere quam omnino non paravisse.} [Greek: Aischion de echontas
aphairethaenai ae ktomenous atychaesai] Thucyd. ii. 62.} what has been
gained, than not to have gained it at all?

But some will ask me, 'What course of conduct, then, would you advise
us to pursue?' I would advise you to inflict punishment on those who
have sacrificed the interests of their country to the enemy; not,
indeed, by arms, or any violence (which would be more unbecoming,
however, for you to inflict than for them to suffer), but by
prosecutions, and by the evidence of Jugurtha himself, who, if he has
really surrendered, will doubtless obey your summons; whereas, if he
shows contempt for it, you will at once judge what sort of a peace or
surrender it is, from which springs impunity to Jugurtha for his
crimes, immense wealth to a few men in power, and loss and infamy to
the republic.

But perhaps you are not yet weary of the tyranny of these men;
perhaps these times please you less than those\footnote{These times please you less than those, etc.—\emph{Illa quam
haec tempora magis placent,} etc. ``\emph{Those times}, which immediately
succeeded the deaths of the Gracchi, and which were distinguished for
the tyranny of the nobles, and the humiliation of the people; \emph{these
times}, in which the people have begun to rouse their spirit and exert
their liberty.'' \emph{Burnouf.}} when kingdoms,
provinces, laws, rights, the administration of justice, war and peace,
and indeed every thing civil and religious, was in the hands of an
oligarchy; while you, that is, the people of Rome, though unconquered
by foreign enemies, and rulers of all nations around, were content
with being allowed to live; for which of you had spirit to throw off
your slavery? For myself, indeed, though I think it most disgraceful
to receive an injury without resenting it, yet I could easily allow
you to pardon these basest of traitors, because they are your
fellow-citizens, were it not certain that your indulgence would end in
your destruction. For such is their presumption, that to escape
punishment for their misdeeds will have but little effect upon them,
unless they be deprived, at the same time, of the power of doing
mischief; and endless anxiety will remain for you, if you shall have
to reflect that you must either be slaves or preserve your liberty by
force of arms.

Of mutual trust, or concord, what hope is there? They wish to be
lords, you desire to be free; they seek to inflict injury, you to
repel it; they treat your allies as enemies, your enemies as allies.
With feelings so opposite, can peace or friendship subsist between
you? I warn, therefore, and exhort you, not to allow such enormous
dishonesty to go unpunished. It is not an embezzlement of the public
money\footnote{Embezzlement of the public money—\emph{Peculatus aerarii.} ``Peculator,
qui furtum facit pecuniae publicae.'' Ascon. Pedian. in Cic. Verr i.} that has been committed; nor is it a forcible extortion of
money from your allies; offenses which, though great, are now, from
their frequency, considered as nothing; but the authority of the
senate, and your own power, have been sacrificed to the bitterest of
enemies, and the public interest has been betrayed for money, both at
home and abroad; and unless these misdeeds be investigated, and
punishment be inflicted on the guilty, what remains for us but to live
the slaves of those who committed them? For those who do what they
will with impunity are undoubtedly kings.\footnote{Kings—I have substituted the plural for the singular. ``No
name was more hated at Rome than that of a king; and no sentiment,
accordingly, could have been better adapted to inflame the minds of
Memmius's hearers, than that which he here utters.'' \emph{Dietsch.}}

I do not, however, wish to encourage you, O Romans, to be better
satisfied at finding your fellow-citizens guilty than innocent, but
merely to warn you not to bring ruin on the good, by suffering the bad
to escape. It is far better, in any government, to be unmindful of a
service than of an injury; for a good man, if neglected, only becomes
less active; but a bad man, more daring. Besides, if the crimes of the
wicked are suppressed,\footnote{If the crimes of the wicked are suppressed, etc.—\emph{Si injuriae
non sint, haud saepe auxilii egeas.} ``Some foolishly interpret
\emph{auxilium} as signifying \emph{auxilium tribunicium}, the aid of the
tribunes; but it is evident to me that Sallust means \emph{aid against
the injuries of bad men}, i.e. revenge or punishment.'' \emph{Kritzius.} ``If
injuries are repressed, or prevented, there will be less need for the
help of good men and it will be of less consequence if they become
inactive.'' \emph{Dietsch.}} the state will seldom need extraordinary
support from the virtuous.''

XXXII. By repeating these and similar sentiments, Memmius prevailed on
the people to send Lucius Cassius,\footnote{XXXII. Lucius Cassius—This is the man from whom came the
common saying \emph{cui bono?} ``Lucius Cassius, whom the Roman people
thought the most accurate and wisest of judges, was accustomed
constantly to inquire, in the progress of a cause, \emph{cui bono fuisset},
of what advantage any thing had been.'' Cic. pro Rosc. Am. 80. ``His
tribunal,'' says Valerius Maximus (iii. 7), ``was called, from his
excessive severity, the rock of the accused.'' It was probably on
account of this quality in his character that he was now sent into
Numidia.} who was then praetor, to
Jugurtha, and to bring him, under guarantee of the public faith\footnote{Under guarantee of the public faith—\emph{Interpositâ fide publicâ.}
See Cat. 47, 48. So a little below, \emph{fidem suam interponit. Interpono}
is ``to pledge.''},
to Rome, in order that, by the prince's evidence, the misconduct of
Scaurus and the rest, whom they charged with having taken bribes,
might more easily be made manifest.

During the course of these proceedings at Rome, those whom Bestia had
left in Numidia in command of the army, following the example of their
general, had been guilty of many scandalous transactions. Some, seduced
by gold, had restored Jugurtha his elephants; others had sold him his
deserters; others had ravaged the lands of those at peace with us; so
strong a spirit of rapacity, like the contagion of a pestilence, had
pervaded the breasts of all.

Cassius, when the measure proposed by Memmius had been carried, and
while all the nobility were in consternation, set out on his mission
to Jugurtha, whom, alarmed as he was, and despairing of his fortune,
from a sense of guilt, he admonished ``that since he had surrendered
himself to the Romans, he had better make trial of their mercy than
their power.'' He also pledged his own word, which Jugurtha valued not
less than that of the public, for his safety. Such, at that period,
was the reputation of Cassius.

XXXIII. Jugurtha, accordingly, accompanied Cassius to Rome, but
without any mark of royalty, and in the garb, as much as possible, of
a suppliant\footnote{XXXIII. In the garb, as much as possible, of a suppliant—\emph{Cultu
quam maximè miserabili}. ``In such a garb as accused persons, or
suppliants, were accustomed to adopt, when they wished to excite
compassion, putting on a mean dress, and allowing their hair and beard
to grow.'' \emph{Burnouf.}}; and, though he felt great confidence on his own
part, and was supported by all those through whose power or villainy
he had accomplished his projects, he purchased, by a vast bribe, the
aid of Caius Baebius, a tribune of the people, by whose audacity he
hoped to be protected against the law, and against all harm.

An assembly of the people being convoked, Memmius, although they were
violently exasperated against Jugurtha, (some demanding that he should
be cast into prison, others that, unless he should name his
accomplices in guilt, he should be put to death, according to the
usage of their ancestors, as a public enemy), yet, regarding rather
their character than their resentment, endeavored to calm their
turbulence and mitigate their rage; and assured them that, as far as
depended on him, the public faith should not be broken. At length,
when silence was obtained, he brought forward Jugurtha, and addressed
them. He detailed the misdeeds of Jugurtha at Rome and in Numidia, and
set forth his crimes toward his father and brothers; and admonished
the prince, ``that the Roman people, though they were well aware by
whose support and agency he had acted, yet desired further testimony
from himself; that, if he disclosed the truth, there was great hope
for him in the honor and clemency of the Romans; but if he concealed
it, he would certainly not save his accomplices, but ruin himself and
his hopes forever.''

XXXIV. But when Memmius had concluded his speech, and Jugurtha was
expected to give his answer, Caius Baebius, the tribune of the people,
whom I have just noticed as having been bribed, enjoined the prince to
hold his peace\footnote{XXXIV. Enjoined the prince to hold his peace—A single tribune
might, by such intervention, offer an effectual opposition to almost
any proceeding. On the great power of the tribunes, see Adam's Rom.
Ant., under the head ``Tribunes of the People.''}; and though the multitude, who formed the
assembly, were desperately enraged, and endeavored to terrify the
tribune by outcries, by angry looks, by violent gestures, and by every
other act to which anger prompts\footnote{Every other act to which anger prompts—\emph{Aliis omnibus, qua
ira fieri amat.} ``These words have given rise to wonderful
hallucinations; for Quintilian, ix. 3. 17, having observed that many
expressions of Sallust are borrowed from the Greek, as \emph{Vulgus amat
fieri,} all interpreters, from Cortius downward, have thought that the
structure of Sallust's words must be Greek, and have taken \emph{ira,} in
this passage, for an ablative, and \emph{quae} for a nominative plural.
Gerlach has even gone so far as to take liberties with the words cited
By Quintilian, and to correct them, please the gods, into \emph{quae in
Vulgus amat fieri.} But how could there have been such want of
penetration in learned critics, such deficiency in the knowledge of
the two languages, that, when the imitation of the Greek, noticed by
Quintilian, has reference merely to the word [Greek: philei], \emph{amat},
they should think of extending it to the dependence of a singular verb
on a neuter plural? With truth, indeed, though with much simplicity,
does Gerlach observe, that you will in vain seek for instances of this
mode of expression in other writers.'' \emph{Kritzius.} Dietsch agrees with
Kritzius; and there will, I hope, be no further doubt that \emph{quae} is
the accusative and \emph{ira} the nominative; the sense being, ``which anger
loves or desires to be done.'' Another mode of explanation has been
suggested, namely, to understand \emph{multitudo} as the nominative case to
\emph{amat}, making \emph{ira} the ablative; but this method is far more
cumbersome, and less in accordance with the style of Sallust. The
words quoted by Quintilian do not refer, as Cortius erroneously
supposes, to this passage, but to some part of Sallust's works that is
now lost.}, his audacity was at last
triumphant. The people, mocked and set at naught, withdrew from the
place of assembly; and the confidence of Jugurtha, Bestia, and the
others, whom this investigation had alarmed, was greatly augmented.
XXXV. There was at this period in Rome a certain Numidian named
Massiva, a son of Gulussa and grandson of Masinissa, who, from having
been, in the dissensions among the princes, opposed to Jugurtha, had
been obliged, after the surrender of Cirta and the murder of Adherbal,
to make his escape out of Africa. Spurius Albinus, who was consul with
Quintus Minucius Rufus the year after Bestia, prevailed upon this man,
as he was of the family of Masinissa, and as odium and terror hung
over Jugurtha for his crimes, to petition the senate for the kingdom
of Numidia. Albinus, being eager for the conduct of a war, was
desirous that affairs should be disturbed\footnote{XXXV. Should be disturbed—\emph{Movere} is the reading of Cortius;
\emph{moveri} that of most other editors, in conformity with most of the MSS.
and early editions.}, rather than sink into
tranquillity; especially as, in the division of the provinces, Numidia
had fallen to himself, and Macedonia to Minucius.

When Massiva proceeded to carry these suggestions into execution,
Jugurtha, finding that he had no sufficient support in his friends, as
a sense of guilt deterred some, and evil report or timidity others,
from coming forward in his behalf, directed Bomilcar, his most
attached and faithful adherent, to procure by the aid of money, by
which he had already effected so much, assassins to kill Massiva; and
to do it secretly if he could; but, if secrecy should be impossible,
to cut him off in any way whatsoever. This commission Bomilcar soon
found means to execute; and, by the agency of men versed in such
service, ascertained the direction of his journeys, his hours of
leaving home, and the times at which he resorted to particular places
\footnote{The times at which he resorted to particular places—\emph{Loca
atque tempora cuncta.} ``All his places and times.'' There can be no
doubt that the sense is what I have given in the text.}, and, when all was ready, placed his assassins in ambush. One of
their number sprung upon Massiva, though with too little caution, and
killed him; but being himself caught, he made, at the instigation of
many, and especially of Albinus the consul, a full confession.
Bomilcar was accordingly committed for trial, though rather on the
principles of reason and justice than in accordance with the law of
nations\footnote{In accordance with the law of nations, etc.—As the public faith
had been pledged to Jugurtha for his security, his retinue was on the
same footing as that of embassadors, the persons of whose attendants
are considered as inviolable as their own, as long as they commit no
offense against the laws of the country in which they are resident. If
any such offense is committed by an attendant of an embassador, an
application is usually made by the government to the embassador to
deliver him up for trial. Bomilcar seems to have been apprehended
without any application having been made to Jugurtha; as, in our own
country, the Portuguese embassador's brother, who was one of his
retinue, was apprehended and executed for a murder, by Oliver
Cromwell. See, on this point, Grotius De Jure Bell, et Pac., xviii, 8;
Vattel, iv. 9; Burlamaqui on Politic Law, part iv. ch. 15. Jugurtha,
says Vattel, should have given up Bomilcar; but such was not
Jugurtha's object.}, as he was in the retinue of one who had come to Rome on
a pledge of the public faith for his safety. But Jugurtha, though
clearly guilty of the crime, did not cease to struggle against the
truth, until he perceived that the infamy of the deed was too strong
for his interest or his money. For which reason, although, at the
commencement of the proceedings\footnote{At the commencement of the proceedings—\emph{In priori actione.}
That is, when Bomilcar was apprehended and charged with the murder.}, he had given fifty of his
friends as bail for Bomilcar, yet, thinking more of his kingdom than
of the sureties, he sent him off privately into Numidia; for he feared
that if such a man should be executed, his other subjects would be
deterred from obeying him\footnote{His other subjects would be deterred from obeying him—\emph{Reliquos
popularis metus invaderet parendi sibi.} ``Fear of obeying him should
take possession of his other subjects.''}. A few days after, he himself departed,
having been ordered by the senate to quit Italy. But, as he was going
from Rome, he is said, after frequently looking back on it in silence,
to have at last exclaimed, ``That it was a venal city, and would soon
perish, if it could but find a purchaser!''\footnote{That it was a venal city, etc.—\emph{Urbem venalem,} etc. I
consider, with Cortias, that this is the proper way of taking these
words. Some would render them \emph{O venal city,} etc., because Livy,
Epit. lxiv., has \emph{O urbem venalem,} but this seems to require that the
verb should be in the second person; and it is probable that in Livy
we should either eject the O or read \emph{inveneris.} Florus, iii. 1,
gives the words in the same way as Sallust.}

XXXVI. The war being now renewed, Albinus hastened to transport
provisions, money, and other things necessary for the army, into
Africa, whither he himself soon followed, with the hope that, before
the time of the comitia, which was not far distant, he might be able,
by an engagement, by capitulation, or by some other method, to bring
the contest to a conclusion.

Jugurtha, on the other hand, tried every means of protracting the war,
continually inventing new causes for delay; at one time he promised to
surrender, at another he feigned distrust; he retreated when Albinus
attacked him, and then, lest his men should lose courage, attacked in
return, and thus amused the consul with alternate procrastinations of
war and of peace.

There were some, at that time, who thought that Albinus understood
Jugurtha's object, and who believed that so ready a protraction of the
war, after so much haste at the commencement, was to be attributed
less to tardiness than to treachery. However this might be, Albinus,
when time passed on, and the day of the comitia approached, left his
brother Aulus in the camp as propraetor\footnote{XXXVI. As propraetor—\emph{Pro praetore.} With the power of
lieutenant-general.}, and returned to Rome.

XXXVII. The republic, at this time, was grievously distracted by the
contentions of the tribunes. Two of them, Publius Lucullus and Lucius
Annius, were struggling against the will of their colleagues, to
prolong their term of office; and this dispute put off the comitia
throughout the year\footnote{XXXVII. Throughout the year—\emph{Totius anni.} That is, all that
remained of the year.}. In consequence of this delay, Aulus, who, as
I have just said, was left as propraetor in the camp, conceiving hopes
either of finishing the war, or of extorting money from Jugurtha by
the terror of his army, drew out his troops in the month of January,
from their winter-quarters into the field, and by forced marches,
during severe weather, made his way to the town of Suthul, where
Jugurtha's treasures were deposited. And though this place, both from
the inclemency of the season, and from its advantageous situation,
could neither be taken nor besieged; for around its walls, which were
built on the edge of a steep hill\footnote{On the edge of a steep hill—\emph{In praerupti montis extremo.
``In extremo} a scholiast rightly interprets \emph{in margine},'' Gerlach.
Cortius, whom Langius follows, considers that \emph{in extremo} means \emph{at
the bottom}; a notion which Kritzius justly condemns; for, as Gerlach
asks, what would that have to do withthe strength of the place? Muller
would have us believe that \emph{in extremo} means \emph{at the top}; but if
Sallust had meant to say that the city was at the top, he would hardly
have chosen the word \emph{extremus} for the purpose. Doubtless, as Gerlach
observes, the city was on the top of the hill, which was broad enough
to hold it; but the words \emph{in extremo} signify that the walls were
even with the side of the hill. Of the site of the town of Suthul no
traces are now to be found.}, a marshy plain, flooded by the
rains of winter, had been converted into a lake; yet Aulus, either as
a feint to strike terror into Jugurtha, or blinded by avarice, began
to move forward his vineae\footnote{Vineae—Defenses made of hurdles or other wood, and often
covered with raw hides, to defend the soldiers who worked the
battering-ram. The word that comes nearest to \emph{vineae} in our language
is \emph{mantelets}. Before this word, in many editions, occurs the phrase
\emph{ob thesauros oppidi potiundi}, which Cortius, whom I follow, omits.}, to cast up a rampart, and to hasten
all necessary preparations for a siege.

XXXVIII. Jugurtha, seeing the propraetor's vanity and ignorance,
artfully strengthened his infatuation; he sent him, from time to time,
deputies with submissive messages, while he himself, as if desirous to
escape, led his army away through woody defiles and cross-roads. At
length he succeeded in alluring Aulus, by the prospect of a surrender
on conditions, to leave Suthul, and pursue him, as if in full retreat,
into the remoter parts of the country. Meanwhile, by means of skillful
emissaries, he tampered night and day with our men, and prevailed on
some of the officers, both of infantry and cavalry, to desert to him
at once, and upon others to quit their posts at a given signal, that
their defection might thus be less observed\footnote{XXXIII. That their defection might thus be less observed—\emph{Ita
delicta occultiora fore.} Cortius transferred these words to this place
from the end of the preceding sentence; Kritzius and Dietsch have
restored them to their former place. Gerlach thinks them an intruded
gloss.}. Having prepared
matters according to his wishes, he suddenly surrounded the camp of
Aulus, in the dead of night, with a vast body of Numidians. The Roman
soldiers were alarmed with an unusual disturbance; some of them seized
their arms, others hid themselves, others encouraged those that were
afraid; but consternation prevailed every where; for the number of the
enemy was great, the sky was thick with clouds and darkness, the
danger was indiscernible, and it was uncertain whether it were safer
to flee or to remain. Of those whom I have just mentioned as being
bribed, one cohort of Ligurians, with two troops of Thracian horse,
and a few common soldiers, went over to Jugurtha; and the chief
centurion\footnote{The chief centurion—\_Centurio primi pili.\_There were sixty
centurions in a Roman legion; the one here meant was the first, or
oldest, centurion of the Triarii, or Pilani.} of the third legion allowed the enemy an entrance at
the very post which he had been appointed to defend, and at which all
the Numidians poured into the camp. Our men fled disgracefully, the
greater part having thrown away their arms, and took possession of a
neighboring hill. Night, and the spoils of the camp, prevented the
enemy from making full use of this victory. On the following day,
Jugurtha, coming to a conference with Aulus, told him, ``that though he
held him hemmed in by famine and the sword, yet that, being mindful of
human vicissitudes, he would, if they would make a treaty with him,
allow them to depart uninjured; only that they must pass under the
yoke, and quit Numidia within ten days.'' These terms were severe and
ignominious; but, as death was the alternative\footnote{As death was the alternative—\_Quia mortis metu mutabant.
Neither manuscripts nor critics are agreed about this passage. Cortius,
from a suggestion of Palmerius, adopted \emph{mutabant}; most other editors
have \emph{mutabantur}; but both are to be taken in the same sense; for
\emph{mutabant} is equivalent to \emph{mutabant se}. Cortius's interpretation
appears the most eligible: ``Permutabantur cum metuendâ morte,'' i.e.
there were those conditions on one side, and death on the other, and
if they did not accept the conditions, they must die. Kritzius
fancifully and strangely interprets, \emph{propter mortis metum se mutabant,
i.e., alia videbantur atque erant}, or the acceptance of the terms
appeared excusable to the soldiers, because they were threatened with
death if they did not accept them. It is worth while to notice the
variety of readings exhibited in the manuscripts collated by Cortius:
ten exhibit \emph{mutabantur}; three, \emph{minitabantur}; three, \emph{multabantur};
three, \emph{tenebantur}; one, \emph{tenebatur}; one, \emph{cogebantur}; one,
\emph{cogebatur}; one, \emph{angustiabantur}; one, \emph{urgebantur}; and one \emph{mortis
metuebant pericula}. There is also, he adds, in some copies, \emph{mutabant},
which the Bipont editors and Müller absurdly adopted.}, peace was
concluded as Jugurtha desired.

XXXIX. When this affair was made known at Rome, consternation and
dismay pervaded the city; some were concerned for the glory of the
republic; others, ignorant of war, trembled for their liberty. But
all were indignant at Aulus, and especially those who had been
distinguished in the field, because, with arms in his hands, he had
sought safety in disgrace rather than in resistance. The consul
Albinus, apprehending, from the delinquency of his brother, odium and
danger to himself, consulted the senate on the treaty which had been
made, but, at the same time, raised recruits for the army, sent for
auxiliaries to the allies and Latins, and made general preparations
for war. The senate, as was just, decreed, ``that no treaty could be
made without their own consent and that of the people.''

The consul, though he was hindered by the influence of the tribunes
from taking with him the force which he had raised, set out in a few
days for the province of Africa, where the whole army, being
withdrawn, according to the agreement, from Numidia, had gone into
winter-quarters. When he arrived there, although he longed to pursue
Jugurtha, and diminish the odium that had fallen on his brother, yet,
when he saw the state of the troops, whom, besides the flight and
relaxation of discipline, licentiousness, and debauchery had
corrupted, he determined, under all the circumstances of the
case\footnote{XXXIX. Under all the circumstances of the case—\emph{Ex copiâ rerum.}
From the number of things which he had to consider.}, to attempt nothing.

XL. At Rome, in the mean time, Caius Mamilius Limetanus, one of the
tribunes, proposed that the people should pass a bill for instituting
an inquiry into the conduct of those by whose influence Jugurtha had
set at naught the decrees of the senate, as well as of those who,
whether as embassadors or commanders, had received money from him, or
who had restored to him his elephants and deserters, or had made any
compacts with the enemy relative to peace or war. To this bill some,
who were conscious of guilt, and others, who apprehended danger from
the jealousy of parties, secretly raised obstructions through the
agency of friends, and especially of men among the Latins and Italian
allies\footnote{XL. The Latins and Italian allies—\emph{Per homines nominis Latini,
et socios Italicos.} ``The right of voting was not extended to all
the Latin people till A.U.C. 664, and the Italian allies did not
obtain it till some years afterward.'' \emph{Kritzius.} So that at this
period, which was twenty years earlier, their influence could only be
employed in an underhand way. Compare c.42.}, since they could not openly resist it, without admitting
that these and similar practices met their approbation. But as to the
people, it is incredible what eagerness they displayed, and with what
spirit they approved, voted, and passed the bill, though rather from
hatred to the nobility, against whom these severe measures were
directed, than from concern for the republic; so violent was the fury
of party.

While the rest of the delinquents were in trepidation, Marcus Scaurus
\footnote{Marcus Scaurus—See c. 15. That he was appointed on this
occasion, is an evident proof of his commanding influence.}, whom I have previously noticed as Bestia's lieutenant,
contrived, amid the exultation of the populace, the dismay of his own
party, and the continued agitation in the city, to have himself
elected one of the three commissioners who were appointed by the bill
of Mamilius to carry it into execution. But the investigation,
notwithstanding, was conducted \footnote{The parents and children of the soldiers, etc.—} with great rigor and violence,
under the influence of common rumor and popular caprice; for the
insolence of success, which had often distinguished the nobility, on
this occasion characterized the people.

XLI. The prevalence of parties among the people, and of factions in
the senate, and of all evil practices attendant on them, had its
origin at Rome, a few years before, during a period of tranquillity,
and amid the abundance of all that mankind regarded as desirable. For,
before the destruction of Carthage, the senate and people managed the
affairs of the republic with mutual moderation and forbearance; there
were no contests among the citizens for honor or ascendency; but the
dread of an enemy kept the state in order. When that fear, however,
was removed from their minds, licentiousness and pride, evils which
prosperity loves to foster, immediately began to prevail; and thus
peace, which they had so eagerly desired in adversity, proved, when
they had obtained it, more grievous and fatal than adversity itself.
The patricians carried their authority, and the people their liberty,
to excess; every man took, snatched, and seized\footnote{XLI. Took, snatched, and seized—Ducere, \emph{trahere, rapere}.
``\emph{Ducere} conveys the notion of cunning and fraud; \emph{trahere} of some
degree of force; \emph{rapere} of open violence.'' \emph{Muller}. The words chiefly
refer to offices in the state, as is apparent from what follows.} what he could.
There was a complete division into two factions, and the republic was
torn in pieces between them. Yet the nobility still maintained an
ascendency by conspiring together; for the strength of the people,
being disunited and dispersed among a multitude, was less able to
exert itself. Things were accordingly directed, both at home and in
the field, by the will of a small number of men, at whose disposal
were the treasury, the provinces, offices, honors, and triumphs; while
the people were oppressed with military service and with poverty, and
the generals divided the spoils of war with a few of their friends.
The parents and children of the soldiers, meantime, if they
chanced to dwell near a powerful neighbor, were driven from their
homes. Thus avarice, leagued with power, disturbed, violated, and
wasted every thing, without moderation or restraint; disregarding
alike reason and religion, and rushing headlong, as it were, to its
own destruction. For whenever any arose among the nobility\footnote{Among the nobility—\emph{Ex nobilitate.} Cortius injudiciously
omits those words. The reference is to the Gracchi.}, who
preferred true glory to unjust power, the state was immediately in a
tumult, and civil discord spread with as much disturbance as attends a
convulsion of the earth.

XLII. Thus when Tiberius and Caius Gracchus, whose forefathers had
done much to increase the power of the state in the Punic and other
wars, began to vindicate the liberty of the people, and to expose the
misconduct of the few, the nobility, conscious of guilt, and seized
with alarm, endeavored, sometimes by means of the allies and
Latins\footnote{By means of the allies and Latins—See on, c. 40.}, and sometimes by means of the equestrian order, whom the
hope of coalition with the patricians had detached from the people, to
put a stop to the proceedings of the Gracchi; and first they killed
Tiberius, and a few years after Caius, who pursued the same measures
as his brother, the one when he was tribune, and the other when he was
one of a triumvirate for settling colonies; and with them they cut off
Marcus Fulvius Flaccus. In the Gracchi, indeed, it must be allowed
that, from their ardor for victory, there was not sufficient prudence.
But to a reasonable man it is more agreeable to submit\footnote{But to a reasonable man it is more agreeable to submit,
etc.—\emph{Sed bono vinci satius est, quam malo more injuriam vincere.
Bono,} sc. \emph{viro}. ``That is, if the nobility had been truly worthy
characters, they would rather have yielded to the Gracchi, than have
revenged any wrong that they had received from them in an unprincipled
manner.'' \emph{Dietsch.} Thus this is a reflection on the nobles; in which
notion of the passage Allen concurs with Dietsch. Others, as Cortius,
think it a reflection on the too great violence of the Gracchi. The
brevity with which Sallust had expressed himself makes it difficult to
decide. Kritzius, who thinks that the remark is in praise of the
Gracchi, supplies the ellipse thus: ``Sane concedi debet Gracchis non
satis moderatum animum fuiase; \emph{quae res ipsis adeo interitum
attulit}; sed \emph{sic quoque egregii viri putandi sunt; nam} bono vinci,''
etc. Langius and Burnouf join \emph{bono} with \emph{more}, but do not differ
much in their interpretations of the passage from that given by
Dietsch.} to
injustice than to triumph over it by improper means. The nobility,
however, using their victory with wanton extravagance, exterminated
numbers of men by the sword or by exile, yet rather increased, for the
time to come, the dread with which they were regarded, than their real
power. Such proceedings have often ruined powerful states; for of two
parties, each strives to suppress the other by any means whatever, and
take vengeance with undue severity on the vanquished.

But were I to attempt to treat of the animosities of parties, and of
the morals of the state, with minuteness of detail, and suitably to
the vastness of the subject, time would fail me sooner than matter. I
therefore return to my subject.

XLIII. After the treaty of Aulus, and the disgraceful flight of our
army, Quintus Metellus and Marcus Silanus, the consuls elect, divided
the provinces between them; and Numidia fell to Metellus, a man of
energy, and, though an opponent of the popular party, yet of a
character uniformly irreproachable\footnote{XLIII. Of a character uniformly irreproachable—\emph{Famâ tamen
aequabili et inviolatâ. Aequabilis} is uniform, always the same,
keeping an even tenor.}. He, as soon as he entered on
his office, regarded all other things as common to himself and his
colleague\footnote{Regarded all things as common to himself and his colleague—\emph{Ali
omnia sibi cum collegâ ratus.} ``Other matters, unconnected with the war
against Jugurtha, he thought that he would have to manage in
conjunction with his colleague, and that, consequently, he might give
but partial attention to them; but that the war in Numidia was
committed to his sole care.'' \emph{Cortius.} Other interpretations of these
words have been suggested; but they are fanciful and unworthy of notice.}, but directed his chief attention to the war which he
was to conduct. Distrusting, therefore, the old army, he began to
raise new troops, to procure auxiliaries from all parts, and to
provide arms, horses, and other military requisites, besides
provisions in abundance, and every thing else which was likely to be
of use in a war varied in its character, and demanding great
resources. To assist in accomplishing these objects, the allies and
Latins, by the appointment of the senate, and different princes\footnote{Princes—\emph{Reges.} Who these were, the commentators have not
attempted to conjecture.}
of their own accord, sent supplies; and the whole state exerted itself
in the cause with the greatest zeal. Having at length prepared and
arranged every thing according to his wishes, Metellus set out for
Numidia, attended with sanguine expectations on the part of his
fellow-citizens, not only because of his other excellent qualities,
but especially because his mind was proof against gold; for it was
through the avarice of our commanders, that, down to this period, our
affairs in Numidia had been ruined, and those of the enemy rendered
prosperous.

XLIV. When he arrived in Africa, the command of the army was resigned
to him by Albinus, the proconsul\footnote{XLIV. By Spurius Albinus, the proconsul.—\emph{A Spurio Albino
proconsule}. This is the general reading. Cortius has, \emph{Spurii Albini
pro consule}, with which we may understand \emph{agentis} or \emph{imperantis},
but can hardly believe it to be what Sallust wrote. Kritzius reads,
\emph{Spurii Albini proconsulis}.}; but it was an army spiritless
and unwarlike; incapable of encountering either danger or fatigue;
more ready with the tongue than with the sword; accustomed to plunder
our allies, while itself was the prey of the enemy; unchecked by
discipline, and void of all regard to its character. The new general,
accordingly, felt more anxiety from the corrupt morals of the men,
than confidence or hope from their numbers. He determined, however,
though the delay of the comitia had shortened his summer campaign, and
though he knew his countrymen to be anxious for the result of his
proceedings, not to commence operations, until, by a revival of the
old discipline, he had brought the soldiers to bear fatigue. For
Albinus, dispirited by the disaster of his brother Aulus and his army,
and having resolved not to leave the province during the portion of
the summer that he was to command, had kept the soldiers, for the most
part, in a stationary camp\footnote{In a stationary camp—\emph{Stativis castris}. In contradistinction
to that which the soldiers formed at the end of a day's march.}, except when stench, or want of
forage, obliged them to remove. But neither had the camp been
fortified\footnote{But neither had the camp been fortified, etc.—\emph{Sed neque
muniebantur ea} (se. castra), \emph{neque more militari vigiliae
deducebantur}. ``The words \emph{sed neque muniebantur ea} are wanting in
almost all the manuscripts, as well as in all the editions, except
that of Cyprianus Popma.'' \emph{Kritzius}. Gerlach, however, had,
previously to Kritz, inserted them in his text though in brackets;
for he supposed them to be a mere conjecture of some scribe, who was
not satisfied with a single \emph{neque}. But they have been found in a
codex of Fronto, by Angelo Mai, and have accordingly been received
as genuine by Kritz and Dietsch. Potter and Burnouf have omitted the
\emph{ea}, thinking, I suppose, that in such a position it could hardly be
Sallust's; but the verb requires a nominative case to prevent it from
being referred to the following \emph{vigiliae}.}, nor the watches kept, according to military usage;
every one had been allowed to leave his post when he pleased. The
camp-followers, mingled with the soldiers, wandered about day and
night, ravaging the country, robbing the houses, and vying with each
other in carrying off cattle and slaves, which they exchanged with
traders for foreign wine\footnote{Foreign wine—\emph{Vino advectitio} Imported. Africa does not
abound in wine.} and other luxuries; they even sold the
corn, which was given them from the public store, and bought bread
from day to day; and, in a word, whatever abominations, arising from
idleness and licentiousness, can be expressed or imagined, and even
more, were to be seen in that army.

XLV. But I am assured that Metellus, in these difficult circumstances,
no less than in his operations against the enemy, proved himself a
great and wise man; so just a medium did he observe between an
affectation of popularity and an excessive enforcement of discipline.
His first measure was to remove incentives to idleness, by a general
order that no one should sell bread, or any other dressed provisions,
in the camp; that no sutlers should follow the army; and that no
common soldier should have a servant, or beast of burden, either in a
camp or on a march. He made the strictest regulations, too, with
regard to other things.\footnote{XLV. With regard to other things—Caeteris. Cortius, whom
Gerlach follows, considers this word as referring to the men or
officers; but Kritzius and Dietsch, with better judgment, understand
\emph{rebus}.} He moved his camp daily, exercising the
soldiers by marches across the country; he fortified it with a rampart
and a trench, exactly as if the enemy had been at hand; he placed
numerous sentinels\footnote{Numerous sentinels—\emph{Vigilias crebras}. At short intervals,
says Kritzius, from each other.} by night, and went the rounds with his
officers; and, when the army was on the march; he would be at one time
in the front, at another in the rear, and at another in the center, to
see that none quitted their ranks, that the men kept close to their
standards, and that every soldier carried his provisions and his arms.
Thus by preventing rather than punishing irregularities, he in a short
time rendered his army effective.

XLVI. Jugurtha, meantime, having learned from his emissaries how
Metellus was proceeding, and having heard, when he was in Rome, of the
integrity of the consul's character, began to despair of his plans,
and at length actually endeavored to effect a capitulation. He
therefore sent deputies to the consul with proposals of submission,
stipulating only for his own life and that of his children, and
offering to surrender every thing else to the Romans. But Metellus had
already learned by experience, that the Numidians were a faithless
race, of unsettled disposition, and fond of change; and he accordingly
applied himself to each of the deputies separately, and after
gradually sounding them, and finding them proper instruments for his
purpose, prevailed on them, by large promises, to deliver Jugurtha
into his hands; bringing him alive, if they could, or dead, if to take
him alive was impracticable. In public, however, he directed that such
an answer should be given to the king as would be agreeable to his
wishes.

A few days afterward, he led the army, which was now vigorous and
resolute, into Numidia, where, instead of any appearance of war, he
found the cottages full of people, and the cattle and laborers in the
fields, while the officers of Jugurtha came from the towns and
villages\footnote{XLVI. Villages—\emph{Mapalibus}. See c. xviii. The word is here
used for a collection of huts, a village.} to meet him, offering to supply him with corn, to convey
provisions for him, and to do whatever might be required of them.
Metellus, notwithstanding, made no diminution in the caution with
which he marched, but kept as much upon the defensive as if an enemy
had been at hand; and he dispatched scouts to explore the country,
thinking that these signs of submission were but pretense, and that
the Numidians were watching an opportunity for treachery. He himself,
with some light-armed cohorts, and a select body of slingers and
archers, advanced always in the front; while Caius Marius, his
lieutenant-general, at the head of the cavalry, had charge of the
rear. The auxiliary horse, distributed among the tribunes of the
legions and prefects of the cohorts, he placed on the flanks, so that,
with the aid of the light troops mixed with them, they might repel the
enemy whenever an approach should be made. For such was the subtlety
of Jugurtha, and such his knowledge of the country and the art of war,
that it was doubtful whether he was more formidable absent or present,
offering peace or threatening hostilities.

XLVII. There lay, not far from the route which Metellus was pursuing,
a city of the Numidians named Vaga, the most celebrated place for
trade in the whole kingdom, in which many Italian merchants were
accustomed to reside and traffic. Here the consul, to try the
disposition of the inhabitants, and, should they allow him, to take
advantage of the situation of the place\footnote{Here the consul, to try the disposition of the inhabitants,
and, should they allow him, to take advantage of the situation of the
place, etc.—\emph{Huc consul, simul tentandi gratiâ, et si paterentur,
opportuniatis loci, praesidium imponit.} This is a \emph{locus
veratissimus}, about which no editor has satisfied himself. I have
deserted Cortius and followed Dietsch, who seems to have settled the
passage, on the basis of Havercamp's text, with more judgment than any
other commentator. Cortius read, \emph{Huc consul simul tentandi gratiâ, si
paterent opportunitates loci}, etc., taking \emph{opportuniatates} in the
sense of \emph{munitiones}, ``defenses;'' but would Sallust have said \emph{that
Metellus put a garrison in the place, to try if its defenses would be
open to him?} Havercamp's reading is, \emph{simul tentan si gratiâ, et si
paterentur opportunitates loci}, etc. Palmerius conjectured \emph{simul
tentandi gratiâ, si paterentur; et opportunitate loci,} which Gerlach
and Kritsius adopt, except that they change the place of the \emph{et}, and
put it before \emph{si}. Allen thinks that he has amended the passage by
reading \emph{Huc consul, simul si paterentur tentandi, et opportunitatis
loci, gratiâ;} but this conjecture is liable to similar objection with
that of Cortius. Other varieties of reading it is needless to notice.
But it is observable that four manuscripts, as Kritzius remarks, have
\emph{propter opportunitates,} which led me long ago to suppose that the
true reading must be \emph{simul tentandi gratiâ, simul propter
opportunitates loci. Simul propter} might easily have been corrupted
into \emph{si paterentur}.}, established a garrison,
and ordered the people to furnish him with corn, and other necessaries
for war; thinking, as circumstances indeed suggested, that the
concourse of merchants, and frequent arrival of supplies\footnote{Frequent arrival of supplies—\emph{Commeatum.} ``Frumenti et omnium
rerum quarum in bello usus est, largam copiam.'' \emph{Kritzius}. I follow
the text of Cortius (retaining the words \emph{juvaturum ecercitum})
which Kritzius sufficiently justifies. There is a variety of readings,
but all much the same in sense.}, would
add strength to his army, and further the plans which he had already
formed.

In the midst of these proceedings, Jugurtha, with extraordinary
earnestness\footnote{Extraordinary earnestness—\emph{Impensius modo.} Cortius and
Kritzius interpret this \emph{modo} as the ablative case of \emph{modus; i. e.
quam modus erat,} or \emph{supra modum;} but Dietsch and Burnouf question
the propriety of this interpretation, and consider the \emph{modo} to be
the same as that in \emph{tantummodo, dummodo,} etc. The same expression
occurs again in c. 75.}, sent deputies to sue for peace, offering to resign
every thing to Metellus, except his own life and that of his children.
These, like the former, the consul first reduced to treachery, and
then sent back; the peace which Jugurtha asked, he neither granted nor
refused, but waited, during these delays, the performance of the
deputies' promises. XLVIII. Jugurtha, on comparing the words of
Metellus with his actions, perceived that he was assailed with his own
artifices; for though peace was offered him in words, a most vigorous
war was in reality pursued against him; one of his strongest cities
was wrested from him; his country was explored by the enemy, and the
affections of his subjects alienated. Being compelled, therefore, by
the necessity of circumstances, he resolved to try the fortune of a
battle. Having, with this view, informed himself of the exact route of
the enemy, and hoping for success from the advantage of the ground, he
collected as large a force of every kind as he could, and, marching by
cross-roads, got in advance of Metellus' army.

There was, in that part of Numidia, of which, on the division of the
kingdom, Adherbal had become possessor, a river named Muthul, flowing
from the south; and, about twenty miles from it, was a range of
mountains running parallel with the stream\footnote{XLVIII. Running parallel with the stream—\emph{Tractu pari.} It
may be well to illustrate this and the following chapter by a copy of
the lines which Cortius has drawn, ``to excite,'' as he says, ``the
imagination of his readers:''}, wild and
uncultivated; but from the center of it stretched a kind of hill,
reaching to a vast distance, covered with wild olives, myrtles, and
other trees, such as grow in a dry and sandy soil. The plain, which
lay between the mountains and the Muthul, was uninhabited from want of
water, except the parts bordering on the river, which were planted
with trees, and full of cattle and inhabitants.

XLIX. On this hill, which I have just mentioned, stretching in a
transverse direction\footnote{XLIX. In a transverse direction—\emph{Transverso itinere}. It lay on
the flank of the Romans as they marched toward the river, \emph{in dextero
latere}, c.49, fin.}, Jugurtha took post with his line drawn out
to a great length. The command of the elephants, and of part of the
infantry, he committed to Bomilcar, and gave him instructions how to
act. He himself, with the whole of the cavalry and the choicest of the
foot, took his station nearer to the range of mountains. Then, riding
round among the several squadrons and battalions, he exhorted and
conjured them to call to mind their former prowess and triumphs, and
to defend themselves and their country from Roman rapacity; saying
that they would have to engage with those whom they had already
conquered and sent under the yoke, and that, though their commander
was changed, there was no alteration in their spirit. He added, that
he had provided for his men every thing becoming a general; that he
had chosen the higher ground, where they, being well acquainted with
the country\footnote{Well acquainted with the country—\emph{Prudentes.} ``Periti loci
et regionis.'' \_Cortius.\_Or it may mean knowing what they were to do,
while the enemy would be \emph{imperiti,} surprised and perplexed.}, would contend with adversaries ignorant of it; nor
would they engage, inferior in numbers and skill, with a larger or
more experienced force; and that they should, therefore, be ready,
when the signal should be given, to fall vigorously on the Romans, as
that day would either crown\footnote{Would crown—\emph{Confirmaturum}. Would establish, settle, put the
last hand to them.} all their labors and victories, or be
a prelude to the most grievous calamities. He also addressed himself,
individually, to any one whom he had rewarded with money or honors for
military desert, reminding him of his favors, and pointing him out as
an example to the rest; and finally he excited all his men, some in
one way and some in another, by threats or entreaties, according to
the different dispositions of each.

Metellus, who was still ignorant of the enemy's position, was now
seen\footnote{Was seen—\emph{Conspicitur.} This is the reading adopted by Cortius,
Müller, and Allen, as being that of all the manuscripts. Havercamp,
Kritzius, and Dietsch admitted into their texts, on the sole authority
of Donatus ad Ter. Eun. ii. 3, \emph{conspicatur}, i.e. (Metellus) \emph{catches
sight} of the enemy. The latter reading, perhaps, makes a better
connection.} descending the mountain with his army. He was at first
doubtful what the strange appearance before him indicated; for the
Numidians, both cavalry and infantry, had taken post among the wood,
not entirely concealing themselves, by reason of the lowness of the
trees, yet rendering it uncertain\footnote{Rendering it uncertain—\emph{Incerti.} Presenting such an
appearance that a spectator could not be certain what they were.} what they were, as both
themselves and their standards were screened as well by the nature of
the ground as by artifice; but soon perceiving that there were men in
ambush, he halted awhile, and, having altered the arrangement of his
troops, he drew up those in the right wing, which was nearest to the
enemy, in three lines\footnote{He drew up these in the right wing—in three lines—\emph{In
dextero latere triplicibus subsidiis aciem instruxit.} In the other
passages in which Sallust has the word \emph{subsidia} (Cat. c. 59), he
uses it for \emph{the lines behind the front.} Thus he says of Catiline,
\emph{Octo cohortes in fronte constituit; reliqua signa in subsidiis
arctiùs collocat;} and of Petreius, \emph{Cohortes veteranas—in fronte;
post eas reliquum exercitum in subsidiis locat.} But whether he uses
the word in the same sense here; whether we might, as Cortius thinks
(whom Gerlach and Dietsch follow), call the division of Metellus's
troops \emph{quadruple} instead of \emph{triple,} or whether he arranged them as
De Brosses and others suppose, in the usual disposition of Hastati,
Principes, and Triarii, who shall place beyond dispute? The probability,
however, if Sallust is consistent with himself in his use of the word,
lies with Cortius. Gerlach refers to Caesar, De Bell, Civ., iii. 89:
``\emph{Celeriter ex tertiâ acie singulas cohortes detraxit, atque ex his
quartam instituit;} but this does not illustrate Sallust's use of the
word \emph{subsidia}: Caesar forms a fourth \emph{acies}; Metellus draws up one
\emph{acies} triplicibus subsidia''.}; he distributed the slingers and archers
among the infantry, posted all the cavalry on the flanks, and having
made a brief address, such as time permitted, to his men, he led them
down, with the front changed into a flank\footnote{With the front changed into a flank—\emph{Transversis principiis.}
He made the whole army wheel to the left, so that what was their front
line, or \emph{principia,} as they faced the enemy on the hill, became their
flank as they marched from the mountain toward the river.}, toward the plain.

L. But when he observed that the Numidians remained quiet, and did not
offer to descend from the hill, he became apprehensive that his army,
from the season of the year and the scarcity of water, might be
overcome with thirst, and therefore sent Rutilius, one of his
lieutenant-generals, with the light-armed cohorts and a detachment of
cavalry, toward the river, to secure ground for an encampment,
expecting that the enemy, by frequent charges and attacks on his
flank, would endeavor to impede his march, and, as they despaired of
success in arms, would try the effect of fatigue and thirst on his
troops.

He then continued to advance by degrees, as his circumstances and the
ground permitted, in the same order in which he had descended from the
range of mountains. He assigned Marius his post behind the front
line\footnote{L. Behind the front line—\emph{Post principia.} The \emph{principia}
are the same as those mentioned in the preceding note, that is, the
front line when the army faced that of Jugurtha on the hill, but which
presented its flank to the enemy when the army was on its march. So
that Marius commanded in the center (``in medio agmine,'' says Dietsch),
while Metellus took the lead with the cavalry of the left wing. See
the following note.}, and took on himself the command of the cavalry on the left
wing, which, on the march, had become the van\footnote{Cavalry on the left wing—which, on the march, had become
the van—\emph{Sinistrae ulae equitibus—qui in agmine principes facti
erant.} When Metellus halted (c. 49, fin.), and drew up his troops
fronting the hill on which Jugurtha was posted, he placed all his
cavalry in the wings; consequently, when the army wheeled to the left,
and marched forward, the cavalry of the left wing became the van.}.

When Jugurtha perceived that the rear of the Roman army had passed his
first line, he took possession of that part of the mountain from which
Metellus had descended, with a body of about two thousand infantry,
that it might not serve the enemy, if they were driven back, as a
place of retreat, and afterward as a post of defense; and then,
ordering the signal to be given, suddenly commenced his attack. Some
of his Numidians made havoc in the rear of the Romans, while others
assailed them on the right and left wings; they all advanced and
charged furiously, and every where threw the consul's troops into
confusion. Even those of our men who made the stoutest resistance,
were baffled by the enemy's versatile method of fighting, and wounded
from a distance, without having the power of wounding in return, or of
coming to close combat; for the Numidian cavalry, as they had been
previously instructed by Jugurtha, retreated whenever a troop of
Romans attempted to pursue them, but did not keep in a body, or
collect themselves into one place, but dispersed as widely as
possible. Thus, being superior in numbers, if they could not deter the
Romans from pursuing, they surrounded them, when disordered, on the
rear or flank, or, if the hill seemed more convenient for retreat than
the plain, the Numidian horses, being accustomed to the brushwood,
easily made their way among it, while the difficulty of the ascent,
and want of acquaintance with the ground, impeded those of the Romans.

LI. The aspect of the whole struggle\footnote{LI. Of the whole struggle—\emph{Totius negotii.} That is, on the side
of the Romans.} was indeed various,
perplexing, direful, and lamentable; the men, separated from their
comrades, were partly fleeing, partly pursuing; neither standards nor
ranks were regarded, but wherever danger pressed, there they made a
stand and defended themselves; arms and weapons, horses and men,
enemies, and fellow-countrymen, were all mingled in confusion; nothing
was done by direction or command, but chance ordered every thing.
Though the day, therefore, was now far advanced, the event of the
contest was still uncertain. At last, however, when all were faint
with exertion and the heat of the day, Metellus, observing that the
Numidians were less vigorous in their charges, drew his troops
together by degrees, restored order among them, and led four cohorts
of the legions against the enemy's infantry, of whom a great number,
overcome with fatigue, had seated themselves on the high ground. He at
the same time entreated and exhorted his men not to lose courage, nor
to suffer a flying enemy to be victorious; adding that they had
neither camp nor citadel to which they could flee, but that their only
dependence was on their arms. Nor was Jugurtha, in the mean time,
inactive; he rode round among his troops, cheered them, renewed the
contest, and, at the head of a select body, made every possible effort
for victory; supporting his own men, charging such of the enemy as
wavered, and repressing with missiles such as he saw remaining
unshaken.

LII. Thus did these two commanders, both eminent men, maintain the
contest against each other. In personal ability they were equal, but
in circumstances unequal. Metellus had resolute troops, but a
disadvantageous position; Jugurtha had every thing in his favor except
men. At last the Romans, seeing that they had no place of refuge, that
the enemy allowed no opportunity for a regular engagement, and that
the evening was fast approaching, forced their way, according to the
orders which were given, up the hill. The Numidians were thus driven
from their position, routed, and put to flight; a few of them were
slain, but their speed, and the enemy's ignorance of the country\footnote{LII. The enemy's ignorance of the country—\emph{Regio hostibus ignara.
Ignara} for \emph{ignota;} a country unknown to the enemy.},
saved the greater number of them.

Meanwhile Bomilcar, who, as I have said before, was appointed by
Jugurtha over the elephants and a part of the infantry, having seen
Rutilius pass by him, led down his men gradually into the plain, and
while Rutilius hastened to the river, to which he had been dispatched,
quietly drew them up in such order as circumstances required; not
omitting, at the same time, to watch every movement of the enemy. When
he learned that Rutilius had taken his position, and seemed free from
apprehension of danger, and heard, at the same time, an increasing
noise where Jugurtha was engaged, fearing lest the lieutenant-general,
taking the alarm, should go to the support of his countrymen in
difficulties, he, in order to intercept his march, increased the
extent of his lines, which, from distrust of the bravery of his men,
he had previously condensed, and advanced in this order toward
Rutilius' camp.

LIII. The Romans, on a sudden, observed a vast cloud of dust, which,
as the ground, thickly covered with brushes, obstructed their view,
they at first supposed to be only sand raised by the wind; but at
length, when they saw that it continued uniform, and approached nearer
and nearer as the line advanced, they understood the real cause of it,
and, hastily seizing their arms, drew up, as their commander directed,
before the camp. When the enemy came up, both sides rushed to the
encounter with loud shouts. But the Numidians maintained the contest
only as long as they trusted for support to their elephants; for, when
they saw the animals entangled in the boughs of the trees, and
dispersed or surrounded by the enemy, they betook themselves to
flight, and most of them, having thrown away their arms, escaped, by
favor of the hill, or of the night, which was now coming on, without
injury. Of the elephants, four were taken, and the rest, to the number
of forty, were killed.

The Romans, though fatigued and exhausted\footnote{LIII. Fatigued and exhausted—\emph{Fessi lassique.} I am once more
obliged to desert Cortius, who reads \emph{laetique}. The sense, as Kritzius
and Dietsch observe, shows that \emph{laeti} can not be the reading, for
there must evidently be a complete antithesis between the two parts of
the sentence; an antithesis which would be destroyed by the introduction
of \emph{laeti}. Gerlach, though he retains \emph{laeti} in his text, condemns it
in his notes.} with their march, the
construction of their camp, and the engagement, yet, as Metellus was
longer in coming than they expected, advanced to meet him in regular
and steady order. The subtlety of the Numidians, indeed, allowed them
neither rest nor relaxation. But as the two parties drew together, in
the obscurity of the night, each occasioned, by a noise like that of
enemies approaching, alarm and trepidation in the other; and, had not
parties of horse, sent forward from both sides, ascertained the truth,
a fatal disaster was on the point of happening from the mistake.
However, in place of fear, joy quickly succeeded; the soldiers met
with mutual congratulations, relating their adventures, or listening
to those of others, and each extolling his own achievements to the
skies. For thus it is with human affairs; in success, even cowards may
boast; while defeat lowers the character even of heroes.

LIV. Metellus remained four days in the same camp. He carefully
provided for the recovery of the wounded, rewarded, in military
fashion, such as had distinguished themselves in the engagements, and
praised and thanked them all in a public address; exhorting them to
maintain equal resolution in their future labors, which would be less
arduous, as they had fought sufficiently for victory, and would now
have to contend only for spoil. In the mean time he dispatched
deserters, and other eligible persons, to ascertain where Jugurtha
was, or what he was doing; whether he had but few followers, or a
large army; and how he conducted himself under his defeat. The prince,
he found, had retreated to places full of wood, well defended by
nature, and was there collecting an army, which would be more numerous
indeed than the former, but inactive and inefficient, as being
composed of men better acquainted with husbandry and cattle than with
war. This had happened from the circumstance, that, in case of flight,
none of the Numidian troops, except the royal cavalry, follow their
king; the rest disperse, wherever inclination leads them; nor is this
thought any disgrace to them as soldiers, such being the custom of the
people.

Metellus, therefore, seeing that Jugurtha's spirit was still
unsubdued; that a war was being renewed, which could only be
conducted\footnote{LIV. Which could only be conducted, etc.—\emph{Quod, nisi ex illius
lubidine, geri non posset.} Cortius omits the \emph{non} before \emph{posset},
but almost every other editor, except Allen, has retained it, from a
conviction of necessity.} according to the prince's pleasure; and that he was
struggling with the enemy on unequal terms, as the Numidians suffered
a defeat with less loss than his own men gained a victory, he resolved
to manage the contest, not by pitched battles or regular warfare, but
in another method. He accordingly marched into the richest parts of
Numidia, captured and burned many fortresses and towns, which were
insufficiently or wholly undefended, put the youth to the sword, and
gave up every thing else as plunder to his soldiers. From the terror
caused by these proceedings, many persons were given up as hostages to
the Romans; corn, and other necessaries, were supplied in abundance;
and garrisons were admitted wherever Metellus thought fit.

These measures alarmed Jugurtha much more than the loss of the late
battle; for he, whose whole security lay in flight, was compelled to
pursue; and he who could not defend his own part of the kingdom, was
obliged to make war in that which was occupied by others. Under these
circumstances, however\footnote{Under these circumstances, however—\emph{Ex copiâ tamen.} With
\emph{copiâ} we must understand \emph{consiliorum} or \emph{rerum,} as at the end of
c. 39. All the manuscripts, except two, have \emph{inopiâ}, which editors
have justly rejected as inconsistent with the sense.}, he adopted what seemed the most eligible
plan. He ordered the main body of his army to continue stationary;
while he himself, with a select troop of cavalry, went in pursuit of
Metellus, and coming upon him unperceived, by means of night marches
and by-roads, he fell upon such of the Roman as were straggling about,
of whom the greater number, being unarmed, were slain, and several
others made prisoners; not one of them, indeed, escaped unharmed; and
the Numidians, before assistance could arrive from the camp, fled, as
they had been ordered, to the nearest hills.

LV. In the mean time great joy appeared at Rome when the proceedings
of Metellus were reported, and when it was known how he was conducting
himself and his army conformably to the ancient discipline; how, on
adverse ground, he had gained a victory by his valor; how he was
securing possession of the enemy's territory; and how he had driven
Jugurtha, when elated by the weakness of Aulus, to depend for safety
on the desert or on flight. For these successes, accordingly, the
senate decreed a thanksgiving\footnote{LV. A thanksgiving—\emph{Supplicia.} The same as \emph{supplicatio,} on
which the reader may consult Adam's Rom. Ant., or Dr. Smith's
Dictionary.} to the immortal gods; the city,
which had been full of anxiety, and apprehensive as to the event of
the war, was now filled with joy; and the fame of Metellus was raised
to the utmost height.

The consul's eagerness to gain a complete victory was thus increased;
he exerted himself in every possible way, taking care, at the same
time, to give the enemy no opportunity of attacking him to advantage.
He remembered that envy is the concomitant of glory, and thus, the
more renowned he became, the greater was his caution and
circumspection. He never went out to plunder, after the sudden attack
of Jugurtha, with his troops in scattered parties; when corn or forage
was sought, a body of cohorts, with the whole of the cavalry, were
stationed as a guard. He himself conducted part of the army, and
Marius the rest. The country was wasted, however, more by fire than by
spoliation. They had separate camps, not far from each other; whenever
there was occasion for force, they formed a union; but, that
desolation and terror might spread the further, they acted separately.
Jugurtha, meanwhile, continued to follow them along the hills,
watching for a favorable opportunity or situation for an attack. He
destroyed the forage, and spoiled the water, which was scarce,
wherever he found that the enemy were coming. He presented himself
sometimes to Metellus, and sometimes to Marius; he would attack their
rear upon a march, and instantly retreat to the hills; he would
threaten sometimes one point, and sometimes another, neither giving
battle nor allowing rest, but making it his great object to retard the
progress of the enemy.

LVI. The Roman commander, finding himself thus harassed by artifices,
and allowed no opportunity of coming to a general engagement, resolved
on laying siege to a large city, named Zama, which was the bulwark of
that part of the kingdom in which it was situate; expecting that
Jugurtha, as a necessary consequence, would come to the relief of his
subjects in distress, and that a battle would then follow. But the
king, being apprised by some deserters of the consul's design, reached
the place, by rapid marches, before him, and exhorted the inhabitants
to defend their walls, giving them, as a reinforcement, a body of
deserters; a class of men, who, of all the royal forces, were the most
to be trusted, inasmuch as they dared not be guilty of treachery\footnote{LVI. Dared not be guilty of treachery—\emph{Fallere nequibant.}
``Through dread of the severest punishments if they should fall into
the hands of the Romans. Valerius Maximus, ii. 7, speaks of deserters
having been deprived of their hands by Quintus Fabius Maximus; of
others who were crucified or beheaded by the elder Africanus; of
others who were thrown to wild beasts by Africanus the younger; and of
others who were sentenced by Paulus Aemilius to be trampled to death
by elephants. Hence it appears that the punishment of deserters was
left to the pleasure of the general.'' \emph{Burnouf}.}.
He also promised to support them, whenever it should be necessary,
with his whole army.

Having taken these precautions, he retired into the deserts of the
interior; where he soon after learned that Marius, with a few cohorts,
had been dispatched from the line of march to bring provisions from
Sicca\footnote{Sicca—It stood on the banks of the Bagradas, at some distance
from the coast, and contained a celebrated Temple of Venus. Val. Max.,
ii. 6. D'Anville thinks it the same as the modern \emph{Kef.}}, a town which had been the first to revolt from him after
his defeat. To this place he hastened by night, accompanied by a
select body of cavalry, and attacked the Romans at the gate, just as
they were leaving the city; calling to the inhabitants, at the same
time, with a loud voice, to surround the cohorts in the rear; adding,
that Fortune had given them an opportunity for a glorious exploit; and
that, if they took advantage of it, he would henceforth enjoy his
kingdom, and they their liberty, without fear. And had not Marius
hastened to advance the standards, and to escape from the town, it is
certain that all, or the greater part of the inhabitants, would have
changed their allegiance; so great is the fickleness which the
Numidians exhibit in their conduct. The soldiers of Jugurtha, animated
for a time by their king, but finding the enemy pressing them with
superior force, betook themselves, after losing a few of their number,
to flight.

LVII. Marius arrived at Zama. This town, built on a plain, was better
fortified by art than by nature. It was well supplied with
necessaries, and contained plenty of arms and men. Metellus, having
made arrangements suitable for the time and the place, encompassed the
whole city with his army, assigning to each of his officers his post
of command. At a given signal, a loud shout was raised on every side,
but without exciting the least alarm in the Numidians, who awaited the
attack full of spirit and resolution. The assault was consequently
commenced; the Romans were allowed to act each according to his
inclination; some annoyed the enemy with slings and stones from a
distance; others came close up to the walls, and attempted to
undermine or scale them, desiring to engage in close combat with the
besieged. The Zamians, on the other hand, rolled down stones, and
hurled burning stakes, javelins\footnote{LVII. Javelins—\emph{Pila.} This \emph{pilum} may have been, as Müller
suggests, similar to the \emph{falarica} which Livy (xxi. 8) says that the
Saguntines used against their besiegers. \emph{Falarica erat Saguntinis,
missile telum hastili abiegno—id, sicut in pilo, quadratum stuppa
circumligabant, linebantque pice:—quod cum medium accensum mitteretur},
etc. Of Sallust's other words, in the latter part of this sentence,
the sense is clear, but the readings of different editors are extremely
various. Cortius and Gerlach have \emph{sudes, pila praeterea picem sulphure
et taedâ mixtam ardentia mittere:} but it can scarcely be believed that
Sallust wrote \emph{picem—taedâ mixtam.} Havercamp gives \emph{pice et sulphure
taedam mixtam ardentia mittere,} which has been adopted by Kritzius and
Dietsch, except that they have changed \emph{ardentia,} on the authority of
some of the manuscripts, into \emph{ardenti}.}, and wood smeared with pitch and
sulphur, on the nearest assailants. Nor was caution a sufficient
protection to those who kept aloof; for darts, discharged from engines
or by the hand, inflicted wounds on most of them; and thus the brave and
the timid, though of unequal merit, were exposed to equal danger.

LVIII. While the struggle was thus continued at Zama, Jugurtha, at the
head of a large force, suddenly attacked the camp of the Romans, and,
through the remissness of those left to guard it, who expected any
thing rather than an attack, effected an entrance at one of the gates.
Our men, struck with sudden consternation, acted each on his own
impulse; some fled, others seized their arms; and many of them were
wounded or slain. About forty, however, out of the whole number
mindful of the honor of Rome, formed themselves into a body, and took
possession of a slight eminence, from which they could not be
dislodged by the utmost efforts of the enemy, but hurled back the
darts discharged at them, and, as they were few against many, not
without execution. If the Numidians came near them, they displayed
their courage, and slaughtered, repulsed, and dispersed them, with the
greatest fury. Metellus, meanwhile, who was vigorously pursuing the
siege, heard a noise, as of enemies, in his rear, and, turning round
his horse, perceived a party of soldiers in flight toward him; a
certain proof that they were his own men. He instantly, therefore,
dispatched the whole of the cavalry to the camp, and immediately
afterward Caius Marius, with the cohorts of the allies, entreating him
with tears, by their mutual friendship, and by his regard for the
public welfare, to allow no stain to rest on a victorious army, and
not to let the enemy escape with impunity. Marius soon executed his
orders. Jugurtha, in consequence, after being embarrassed in the
intrenchments of the camp, while some of his men threw themselves over
the ramparts, and others, in their haste, obstructed each other at the
gates, fled, with considerable loss, to his strongholds, Metellus, not
succeeding in his attempt on the town, retired with his forces, at the
approach of night, into his camp.

LIX. On the following day, before he marched out to resume the siege,
he ordered the whole of his cavalry to take their station before the
camp, on the side where the approach of Jugurtha was to be apprehended;
assigning the gates, and adjoining posts, to the charge of the tribunes.
He then marched toward the town, and commenced an assault upon the walls
as on the day before. Jugurtha, meanwhile, issuing from his concealment,
suddenly attacked our men in the camp, of whom those stationed in advance
were for the moment alarmed and thrown into confusion; but the rest soon
came to their support; nor would the Numidians have longer maintained
their ground, had not their foot, which were mingled with the cavalry,
done great execution in the struggle; for the horse, relying on the
infantry, did not, as is common in actions of cavalry, charge and then
retreat, but pressed impetuously forward, disordering and breaking the
ranks, and thus, with the aid of the light-armed foot, almost succeeded
in giving the army a defeat\footnote{LIX. And thus, with the aid of the light-armed foot, almost
succeeded in giving the enemy a defeat—\emph{Ita expeditis peditibus suis
hostes paene victos dare.} Cortius, Kritzius, and Allen, concur in
regarding \emph{expeditis peditibus} as an ablative of the instrument, i.e.
as equivalent to \emph{per expeditos pedites} and \emph{victos dare} as nothing
more than \emph{vincere.} This appears to be the right mode of explanation;
but most of the translators, French as well as English, have taken
\emph{expeditis peditibus} as a dative, and given to the passage the sense
that ``the cavalry delivered up the enemy, when nearly conquered, to be
dispatched by the light-armed foot.''}.

LX. The conflict at Zama, at the same time, was continued with great
fury. Wherever any lieutenant or tribune commanded, there the men
exerted themselves with the utmost vigor. No one seemed to depend for
support on others, but every one on his own exertions. The townsmen,
on the other side, showed equal spirit. Attacks, or preparations for
defense, were made in all quarters\footnote{LX. Attacks, or preparations for defense, were made in all
quarters—\emph{Oppugnare aut parare omnibus locis.} There is much
discussion among the critics whether these verbs are to be referred to
the besiegers or the besieged. Cortius and Gerlach attribute
\emph{oppugnare} to the Romans, and \emph{parare} to the men of Zama; a
distinction which Kritzius justly condemns. There can be little doubt
that they are spoken of both parties equally.}. All appeared more eager to
wound their enemies than to protect themselves. Shouts, mingled with
exhortations, cries of joy, and the clashing of arms, resounded
through the heavens. Darts flew thick on every side. If the besiegers,
however, in the least relaxed their efforts, the defenders of the
walls immediately turned their attention to the distant engagement of
the cavalry; they were to be seen sometimes exhibiting joy, and
sometimes apprehension, according to the varying fortune of Jugurtha,
and, as if they could be heard or seen by their friends, uttering
warnings or exhortations, making signs with their hands, and moving
their bodies to and fro, like men avoiding or hurling darts. This
being noticed by Marius, who commanded on that side of the town, he
artfully relaxed his efforts, as if despairing of success, and allowed
the besieged to view the battle at the camp unmolested. Then, while
their attention was closely fixed on their countrymen, he made a
vigorous assault on the wall, and the soldiers mounting their scaling
ladders, had almost gained the top, when the townsmen rushed to the
spot in a body, and hurled down upon them stones, firebrands, and
every description of missiles. Our men made head against these
annoyances for a while, but at length, when some of the ladders were
broken, and those who had mounted them dashed to the ground, the rest
of the assailants retreated as they could, a few indeed unhurt, but
the greater number miserably wounded. Night put an end to the efforts
of both parties.

LXI. When Metellus saw that all his attempts were vain; that the town
was not to be taken; that Jugurtha was resolved to abstain from
fighting, except from an ambush, or on his own ground, and that the
summer was now far advanced, he withdrew his army from Zama, and
placed garrisons in such of the cities that had revolted to him as
were sufficiently strong in situation or fortifications. The rest of
his forces he settled in winter quarters, in that part of our province
nearest to Numidia\footnote{LXI. The rest of his forces—in that part of our province
nearest to Numidia—\emph{Caeterum exercitum in provinciam, quae proxima
est Numidiae, hiemandi gratiâ collocat.} ``The words \emph{quae proxima est
Numidiae} Cortius would eject as superfluous and spurious. But it is
to be understood that Metellus did not distribute his troops through
the whole of the province, but in that part which is nearest to
Numidia, in order that they might be easily assembled in case of an
attack of the enemy or any other emergency. There is, therefore, no
need to read with the Bipont edition and Müller, \emph{qua proxima,} etc.
though this is in itself not a bad conjecture.'' \emph{Kritzius}.}.

This season of repose, however, he did not, like other commanders,
abandon to idleness and luxury; but as the war had been but slowly
advanced by fighting, he resolved to try the effect of treachery on
the king through his friends, and to employ their perfidy instead of
arms. He accordingly addressed himself with large promises, to
Bomilcar, the same nobleman who had been with Jugurtha at Rome, and
who had fled from thence, notwithstanding he had given bail, to escape
being tried for the murder of Massiva; selecting this person for his
instrument, because, from his great intimacy with Jugurtha, he had the
best opportunities of betraying him. He prevailed on him, in the first
place, to come to a conference with him privately, when, having given
him his word, ``that, if he should deliver up Jugurtha, alive or dead,
the senate would grant him a pardon, and the full possession of his
property,'' he easily brought him over to his purpose, especially as he
was naturally faithless, and also apprehensive that, if peace were
made with the Romans, he himself would be surrendered to justice by
the terms of it.

LXII. Bomilcar took the earliest opportunity of addressing Jugurtha,
at a time when he was full of anxiety, and lamenting his ill success.
He exhorted and implored him, with tears in his eyes, to take at
length some thought for himself and his children, as well as for the
people of Numidia, who had so much claim upon him. He reminded him
that they had been, defeated in every battle; that the country was
laid waste; that numbers of his subjects had been captured or slain;
that the resources of the kingdom were greatly reduced; that the valor
of his soldiers, and his own fortune, had been already sufficiently
tried; and that he should beware, lest, if he delayed to consult for
his people, his people should consult for themselves. By these and
similar appeals, he prevailed with Jugurtha to think of a surrender.
Embassadors were accordingly sent to the Roman general, announcing
that Jugurtha was ready to submit to whatever he should desire, and to
trust himself and his kingdom unconditionally to his honor. Metellus,
on receiving this statement, summoned such of his officers as were of
senatorial rank, from their winter quarters; of whom, with, others
whom he thought eligible, he formed a council. By a resolution of this
assembly, in conformity with ancient usage, he demanded of Jugurtha,
through his embassadors, two hundred thousand pounds' weight of
silver, all his elephants, and a portion of his horses and arms. These
requisitions being immediately complied with, he next desired that all
the deserters should be brought to him in chains. A large number of
them were accordingly brought; but a few, when the surrender first
began to be mentioned, had fled into Mauretania to king Bocchus.

When Jugurtha, however, after being thus despoiled of arms, men and
money, was summoned to appear in person at Tisidium\footnote{LXII. Was summoned to appear in person at Tisidium, etc.
—\emph{Cum ipse ad imperandum Tisidium vocaretur.} The gerund is used, as
grammarians say, in a passive sense. ``The town of Tisidium is nowhere
else mentioned. Strabo (xvii. 3, p. 488, Ed. Tauch.) speaks of a place
named [Greek: \emph{Tisiaioi}], which was utterly destroyed, and not a
vestige of it left.'' \emph{Gerlach}.}, to await the
consul's commands, he began again to change his mind, dreading, from a
consciousness of guilt, the punishment due to his crimes. Having spent
several days in hesitation, sometimes, from disgust at his ill
success, believing any thing better than war, and sometimes
considering with himself how grievous would be the fall from
sovereignty to slavery, he at last determined, notwithstanding that he
had lost so many and so valuable means of resistance, to commence
hostilities anew.

At Rome, meanwhile, the senate, having been consulted about the
provinces, had decreed Numidia to Metellus.

LXIII. About the same time, as Caius Marius, who happened to be at
Utica, was sacrificing to the gods\footnote{LXIII. Sacrificing to the gods—\emph{Per hostias dis supplicante.}
Supplicating or worshiping the gods with sacrifices, and trying to
learn their intentions as to the future by inspection of the entrails.
``Marius was either a sincere believer in the absurd superstitions and
dreams of the soothsayers, or pretended to be so, from a knowledge of
the nature of mankind, who are eager to listen to wonders, and are
ore willing to be deceived than to be taught.'' \emph{Burnouf.} See Plutarch,
Life of Marius. He could interpret omens for himself, according to
Valerius Maximus, i. 5.}, an augur told him that great
and wonderful things were presaged to him; that he might therefore
pursue whatever designs he had formed, trusting to the gods for
success; and that he might try fortune as often as he pleased, for
that all his undertakings would prosper. Previously to this period an
ardent longing for the consulship had possessed him; and he had,
indeed, every qualification for obtaining it, except antiquity of
family; he had industry, integrity, great knowledge of war, and a
spirit undaunted in the field; he was temperate in private life,
superior to pleasure and riches, and ambitious only of glory.
Having been born at Arpinum, and brought up there during his boyhood,
he employed himself, as soon as he was of age to bear arms, not in the
study of Greek eloquence, nor in learning the refinements of the city,
but in military service; and thus, amid the strictest discipline, his
excellent genius soon attained full vigor. When he solicited the
people, therefore, for the military tribuneship, he was well known by
name, though most were strangers to his face, and unanimously elected
by the tribes. After this office he attained others in succession, and
conducted himself so well in his public duties, that he was always
deemed worthy of a higher station than he had reached. Yet, though
such had been his character hitherto (for he was afterward carried
away by ambition), he had not ventured to stand for the consulship.
The people, at that time, still disposed of\footnote{The people—disposed of, etc.—\emph{Etiam tum alios magistratus
plebes, consulatum nobilitas, inter se per manus tradebat.} The
commentators have seen the necessity of understanding a verb with
\emph{plebes.} Kritzius suggests \emph{habebat;} Gerlach \emph{grebat} or \emph{accipiebat}.} other civil offices,
but the nobility transmitted the consulship from hand to hand among
themselves. Nor had any commoner appeared, however famous or
distinguished by his achievements, who would not have been thought
unworthy of that honor, and, as it were, a disgrace to it\footnote{A disgrace to it—\emph{Pollutus.} He was considered, as it were,
unclean. See Cat., c. 23, \emph{fin}.}.

LXIV. But when Marius found that the words of the augur pointed in the
same direction as his own inclinations prompted him, he requested of
Metellus leave of absence, that he might offer himself a candidate for
the consulship. Metellus, though eminently distinguished by virtue,
honor, and other qualities valued by the good, had yet a haughty and
disdainful spirit, the common failing of the nobility. He was at
first, therefore, astonished at so extraordinary an application,
expressed surprise at Marius's views, and advised him, as if in
friendship, ``not to indulge such unreasonable expectations, or elevate
his thoughts above his station; that all things were not to be coveted
by all men; that his present condition ought to satisfy him; and,
finally, that he should be cautious of asking from the Roman people
what they might justly refuse him.'' Having made these and similar
remarks, and finding that the resolution of Marius was not at all
affected by them, he told him ``that he would grant what he desired as
soon as the public business would allow him''.\footnote{LXIV. As soon as the public business would allow him—\emph{Ubi
primum potuisset per negotia publica.} As soon as he could through
(regard to) the public business.} On Marius repeating
his request several times afterward, he is reported to have said,
``that he need not be in a hurry to go, as he would be soon enough if
he became a candidate with his own son.''\footnote{With his own son—\emph{Cum filio suo.} With the son of Metellus.
He tells Marius that it would be soon enough for him to stand for
the consulship in twenty-three years' time, the legitimate age for
the consulship being forty-three.} Metellus's son was then
on service in the camp with his father\footnote{In the camp with his father—\emph{Contubernio patris.} He was
among the young noblemen in the consul's retinue, who were sent out
to see military service under him. This was customary. See Cic. Pro
Cael. Pro Planc. 11.}, and was about twenty
years old.

This taunt served only to rouse the feelings of Marius, as well for
the honor at which he aimed, as against Metellus. He suffered himself
to be actuated, therefore, by ambition and resentment, the worst of
counselors. He omitted nothing henceforward, either in deeds or words,
that could increase his own popularity. He allowed the soldiers, of
whom he had the command in the winter quarters, more relaxation of
discipline than he had ever granted them before. He talked of the war
among the merchants, of whom there was a great number at Utica,
censoriously with respect to Metellus, and vauntingly with regard to
himself; saying ``that if but half of the army were granted him, he
would in a few days have Jugurtha in chains; but that the war was
purposely protracted by the consul, because, being a man of vanity and
regal pride, he was too fond of the delights of power.'' All these
assertions appeared the more credible to the merchants, as, by the
long continuance of the war, they had suffered in their fortunes; and
to impatient minds no haste is sufficient.

LXV. There was then in our army a Numidian named Gauda, the son of
Mastanabal, and grandson of Masinissa, whom Micipsa, in his will, had
appointed next heir to his immediate successors. This man had been
debilitated by ill-health, and, from the effect of it, was somewhat
impaired in his understanding. He had petitioned Metellus to allow him
a seat, like a prince, next to himself, and a troop of horse for a
bodyguard; but Metellus had refused him both; the seat, because it was
granted only to those whom the Roman people had addressed as kings,
and the guard, because it would be an indignity to Roman cavalry to
act as guards to a Numidian. While Gauda was discontented at these
refusals, Marius paid him a visit, and prompted him, with his
assistance, to seek revenge for the affronts put upon him by the
general; inflating his mind, which was as weak as his body,\footnote{LXV. Which was as weak as his body—\emph{Ob morbos—parum valido.}
Sallust had already expressed this a few lines above.} with
flattering speeches, telling him that he was a prince, a great man,
and the grandson of Masinissa; that if Jugurtha were taken or killed,
he would immediately become king of Numidia; and that this event might
soon happen, if he himself were sent as consul to the war.

Thus partly the influence of Marius himself, and partly the hope of
obtaining peace, induced Gauda, as well as most of the Roman knights,
both soldiers and merchants,\footnote{Merchants—\emph{Negotiatores.} ``Every one knows that Romans of
equestrian dignity were accustomed to trade in the provinces.''
\emph{Burnouf}.} to write to their friends at Rome,
in a style of censure, respecting Metellus's management of the war,
and to intimate that Marius should be appointed general. The consulship,
accordingly, was solicited for him by numbers of people, with the most
honorable demonstrations in his favor.\footnote{With the most honorable demonstrations in his favor
—\emph{Honestissimâ suffragatione.} ``\emph{Suffragatio} was the zealous
recommendation of those who solicited the votes of their
fellow-citizens in favor of some candidate. See Festus, s.v.
\emph{Suffragatores,} p. 266, Lindem.'' \emph{Dietsch.} It was honorable, in
the case of Marius, as it was without bribery, and seemed to have
the good of the republic in view.} It happened that the
people too, at this juncture, having just triumphed over the nobility
by the Mamilian law,\footnote{The Mamilian law—See c. 40.} were eager to raise commoners to office.
Hence every thing was favorable to Marius's views.

LXVI. Jugurtha, meantime, who, after relinquishing his intention to
surrender, had renewed the war, was now hastening the preparations for
it with the utmost diligence. He assembled an army; he endeavored, by
threats or promises, to recover the towns that had revolted from him;
he fortified advantageous positions;\footnote{LXVI. Advantageous positions—\emph{Suos locos.} Places favorable
for his views. See Kritzius on c. 54.} he repaired or purchased
arms, weapons, and other necessaries, which he had given up on the
prospect of peace; he tried to seduce the slaves of the Romans, and
even tempted with bribes the Romans themselves who occupied the
garrisons; he, indeed, left nothing untried or neglected, but put
every engine in motion.

Induced by the entreaties of their king, from whom, indeed, they had
never been alienated in affection, the leading inhabitants of Vacca, a
city in which Metellus, when Jugurtha began to treat for peace, had
placed a garrison, entered into a conspiracy against the Romans. As
for the common people of the town, they were, as is generally the
case, and especially among the Numidians, of a fickle disposition,
factious and turbulent, and therefore already desirous of a change,
and adverse to peace and quiet. Having arranged their plans, they
fixed upon the third day following for the execution of them, because
that day, being a festival, celebrated throughout Africa, would
promise merriment and dissipation rather than alarm. When the time
came, they invited the centurions and military tribunes, with Titus
Turpilius Silanus, the governor of the town, to their several houses,
and butchered them all, except Turpilius, at their banquets; and then
fell upon the common soldiers, who, as was to be expected on such a
day, when discipline was relaxed, were wandering about without their
arms. The populace followed the example of their chiefs, some of them
having been previously instructed to do so, and others induced by a
liking for such disorders, and, though ignorant of what had been done
or intended, finding sufficient gratification in tumult and variety.

LXVII. The Roman soldiers, perplexed with sudden alarm, and not
knowing what was best for them to do, were in trepidation. At the
citadel,\footnote{LXVII. Were in trepidation. At the citadel, etc.—I have
translated this passage in conformity with the texts of Gerlach,
Kritzius, Dietsch, Müller, and Allen, who put a point between
\emph{trepidare} and \emph{ad arcem}. Cortina, Havercamp, and Burnouf have
\emph{trepidare ad arcem}, without any point. Which method gives the better
sense, any reader can judge.} where their standards and shields were, was posted a
guard of the enemy; and the city-gates, previously closed, prevented
escape. Women and children, too, on the roofs of the houses,\footnote{On the roofs of the houses—\emph{Pro tectis aedificiorum}. In
front of the roofs of the houses; that is, at the parapets. ``In prima
tectorum parte.'' \emph{Kritzius}. The roofs were flat.}
hurled down upon them, with great eagerness, stones and whatever else
their position furnished. Thus neither could such twofold danger be
guarded against, nor could the bravest resist the feeblest; the worthy
and the worthless, the valiant and the cowardly, were alike put to
death unavenged. In the midst of this slaughter, while the Numidians
were exercising every cruelty, and the town was closed on all sides,
Turpilius was the only one, of all the Italians, that escaped unhurt.
Whether his flight was the consequence of compassion in his entertainer,
of compact, or of chance, I have never discovered; but since, in such a
general massacre, he preferred inglorious safety to an honorable name,
he seems to have been a worthless and infamous character.\footnote{Worthless and infamous character—\emph{Improbus intestabilisque}.
These words are taken from the twelve tables of the Roman law: See
Aul. Gell. vi. 7, xv. 3. Horace, in allusion to them, has \emph{intestabilis
et sacer}, Sat. ii. 3.181, \emph{Intestabilis} signified a person to be of
so infamous a character that he was not allowed to give evidence in a
court of justice.}

LXVIII. When Metellus heard of what had happened at Vacca, he retired
for a time, overpowered with sorrow, from the public gaze; but at
length, as indignation mingled with his grief, he hastened, with the
utmost spirit, to take vengeance for the outrage. He led forth, at
sunset, the legion that was in winter quarters with him, and as many
Numidian horse as he could, and arrived, about the third hour on the
following day, at a certain plain surrounded by rising grounds. Here
he acquainted the soldiers, who were now exhausted with the length of
their march, and averse to further exertion,\footnote{LXVIII. Averse to further exertion—\emph{Tum abnuentes omnia}.
Most of the translators have understood by these words that the troops
refused to obey orders; but Sallust's meaning is only that they
expressed, by looks and gestures, their unwillingness to proceed.} that the town of
Vacca was not above a mile distant, and that it became them to bear
patiently the toil that remained, with the hope of exacting revenge for
their countrymen, the bravest and most unfortunate of men. He likewise
generously promised them the whole of the plunder. Their courage being
thus revived, he ordered them to resume their march, the cavalry
maintaining an extended line in front, and the infantry, with their
standards concealed, keeping the closest order behind.

LXIX. The people of Vacca, perceiving an army coming toward them,
judged rightly at first that it was Metellus, and shut their gates;
but, after a while, when they saw that their fields were not laid
waste, and that the front consisted of Numidian cavalry, they imagined
that it was Jugurtha, and went out with great joy to meet him. A
signal being immediately given, both cavalry and infantry commenced an
attack; some cut down the multitude pouring from the town, others
hurried to the gates, others secured the towers, revenge and the hope
of plunder prevailing over their weariness. Thus Vacca triumphed only
two days in its treachery; the whole city, which was great and
opulent, was given up to vengeance and spoliation. Turpilius, the
governor, whom we mentioned as the only person that escaped, was
summoned by Metellus to answer for his conduct, and not being able to
clear himself, was condemned, as a native of Latium,\footnote{LXIX As a native of Latium—\emph{Nam is civis ex Latio erat}.
``As he was a Latin, he was not protected by the Porcian law (see Cat.,
c. 51), though how far this law had power in the camp, is not agreed.''
\emph{Allen}. Gerlach thinks that it had the same power in the camp as
elsewhere, with reference to Roman citizens. But Roman citizenship was
not extended to the Latins till the end of the Social War, A.U.C. 662.
Plutarch, however, in his Life of Caius Gracchus (c. 9), speaks of
Livius Drusus having been abetted by the patricians in proposing a law
for exempting the Latin soldiers from being flogged, about thirty
years earlier; and it seems to have been passed, but, from this
passage of Sallust, appears not to have remained in force. Lipsius
touches on this obscure point in his \emph{Militia Romana}, v. 18, but
settles nothing. Plutarch, in his Life of Marius, c. 8, says that
Turpilius was an old retainer of the family of Metellus, whom he
attended, in this war, as \emph{prafectus fabrûm}, or master of the
artificers; that, being afterward appointed governor of Vacea, he
exercised his office with great justice and humanity, that his life
was spared by Jugurtha at the solicitation of the inhabitants; that,
when he was brought to trial, Metellus thought him innocent, and that
he would not have been condemned but for the malice of Marius, who
exasperated the other members of the council against him. He adds,
that after his death, his innocence became apparent, and that Marius
boasted of having planted in the breast of Metellus an avenging fury,
that would not fail to torment him for having put to death the
innocent friend of his family. Hence Sir Henry Steuart has accused
Sallust of wilfully misrepresenting the character of Turpilius, as
well as the whole transaction. But as much credit is surely due to
Sallust as to Plutarch.} to be
scourged and put to death.

LXX. About this time, Bomilcar, at whose persuasion Jugurtha had
entered upon the capitulation which he had discontinued through fear,
being distrusted by the king, and distrusting him in return, grew
desirous of a change of government. He accordingly meditated schemes
for Jugurtha's destruction, racking his invention night and day. At
last, to leave nothing untried, he sought an accomplice in Nabdalsa, a
man of noble birth and great wealth, who was in high regard and favor
with his countrymen, and who, on most occasions, used to command a
body of troops distinct from those of the king, and to transact all
business to which Jugurtha, from fatigue, or from being occupied with
more important matters, was unable to attend;\footnote{LXX. To which Jugurtha—was unable to attend—\emph{Quae Jugartha,
fesso, aut majoribus astricto, superaverant}. ``Which had remained to
(or been too much for) Jugurtha, when weary, or engaged in more
important affairs.''} employments by
which he had gained both honors and wealth. By these two men in
concert, a day was fixed for the execution of their treachery;
succeeding matters they agreed to settle as the exigences of the
moment might require. Nabdalsa then proceeded to join his troops,
which he kept in readiness, according to orders, among the winter
quarters of the Romans,\footnote{Among the winter-quarters of the Romans—\emph{Inter hiberna
Romanorum}.It is stated in c. 61, as Kritzius observes, that Metellus,
when he put his army into winter-quarters, had, at the same time,
placed garrisons in such of Jugurtha's towns as had revolted to him.
The forces of the Romans being thus dispersed, Nabdalsa might justly
be said to have his army \emph{inter hiberna}, ``\emph{among} their
winter-quarters.''} to prevent the country from being ravaged
by the enemy with impunity.

But as Nabdalsa, growing alarmed at the magnitude of the undertaking,
failed to appear at the appointed time, and allowed his fears to
hinder their plans, Bomilcar, eager for their execution, and
disquieted at the timidity of his associate, lest he should relinquish
his original intentions and adopt some new course, sent him a letter
by some confidential person, in which he ``reproached him with
pusillanimity and irresolution, and conjured him by the gods, by whom
he had sworn, not to turn the offers of Metellus to his own
destruction;'' assuring him ``that the fall of Jugurtha was approaching;
that the only thing to be considered was whether he should perish by
their hand or by that of Metellus; and that, in consequence, he might
consider whether to choose rewards, or death by torture.''

LXXI. It happened that when this letter was brought, Nabdalsa,
overcome with fatigue, was reposing on his couch, where, after reading
Bomilcar's letter, anxiety at first, and afterward, as is usual with a
troubled mind, sleep overpowered him. In his service there was a
certain Numidian, the manager of his affairs, a person who possessed
his confidence and esteem, and who was acquainted with all his designs
except the last. He, hearing that a letter had arrived, and supposing
that there would be occasion, as usual, for his assistance or
suggestions, went into the tent, and, while his master was asleep,
took up the letter thrown carelessly upon the cushion behind his
head,\footnote{LXXI. Behind his head—\emph{Super caput}. On the back of the
bolster that supported his head; part of which might be higher than
the head itself.} and read it; and, having thus discovered the plot, set off
in haste to Jugurtha. Nabdalsa, who awoke soon after, missing the
letter, and hearing of the whole affair, and how it had happened, at
first attempted to pursue the informer, but finding that pursuit was
vain, he went himself to Jugurtha to try to appease him; saying that
the disclosure which he intended to make, had been anticipated by the
perfidy of his servant; and beseeching him with tears, by his
friendship, and by his own former proofs of fidelity, not to think
that he could be guilty of such treachery.

LXXII. To these entreaties the king replied with a mildness far
different from his real feelings. After putting to death Bomilcar, and
many others whom he knew to be privy to the plot, he refrained from
any further manifestation of resentment, lest an insurrection should
be the consequence of it. But after this occurrence he had no peace
either by day or by night; he thought himself safe neither in any
place, nor with any person, nor at any time; he feared his subjects
and his enemies alike; he was always on the watch, and was startled at
every sound; he passed the night sometimes in one place, and sometimes
in another, and often in places little suited to royal dignity; and
sometimes, starting from his sleep, he would seize his arms and raise
an alarm. He was indeed so agitated by extreme terror, that he
appeared under the influence of madness.

LXXIII. Metellus, hearing from some deserters of the fate of Bomilcar,
and the discovery of the conspiracy, made fresh preparations for
action, and with the utmost dispatch, as if entering upon an entirely
new war. Marius, who was still importuning him for leave of absence,
he allowed to go home; thinking that as he served with reluctance, and
bore him personal enmity, he was not likely to prove a very useful
officer.

The common people at Rome, having learned the contents of the letters
written from Africa concerning Metellus and Marius, had listened to
the accounts given of both with eagerness. But the noble birth of
Metellus, which had previously been a motive for paying him honor, had
now become a cause of unpopularity; while the obscurity of Marius's
origin had procured him favor. In regard to both, however, party
feeling had more influence than the good or bad qualities of either.
The factious tribunes,\footnote{LXXIIL The factious tribunes—\emph{Seditiosi magistratus}.} too, inflamed the populace, charging
Metellus, in their harangues, with offenses worthy of death, and
exaggerating the excellent qualities of Marius. At length the people
were so excited that all the artisans and rustics, whose whole
subsistence and credit depended on their labor, quitting their several
employments, attended Marius in crowds, and thought less of their own
wants than of his exaltation. Thus the nobility being borne down, the
consulship, after the lapse of many years,\footnote{After the lapse of many years—\emph{Post multas tempestates}.
Apparently the period since A.U.C. 611, when Quintus Pompeius, who, as
Cicero says (in Verr. ii. 5), was \emph{humile atque obscuro loco natus},
obtained the consulship; that is, a term of forty-three or forty-four
years.} was once more given to
a man of humble birth. And afterward, when the people were asked by
Manilius Mancinus, one of their tribunes, whom they would appoint to
carry on the war against Jugurtha, they, in a full assembly, voted it
to Marius. The senate had previously decreed it to Metellus; but that
decree was thus rendered abortive.\footnote{That decree was thus rendered abortive—\emph{Ea res frustra fuit}.
By a \emph{lex Sempronia}, a law of Caius Gracchus, it was enacted that the
senate should fix the provinces for the future consuls before the
\emph{comitia} for electing them were held. But from Jug. c. 26, it appears
that the consuls might settle by lot, or by agreement between
themselves, which of those two provinces each of them should take. How
far the senate were allowed or accustomed in general, to interfere in
the arrangement, it is not easy to discover: but on this occasion they
had taken on themselves to pass a resolution in favor of the patrician.
Lest similar scenes, however, to those of the Sempronian times should
be enacted, they yielded the point to the people.}

LXXIV. During this period, Jugurtha, as he was bereft of his friends
(of whom he had put to death the greater number, while the rest, under
the influence of terror, had fled partly to the Romans, and partly to
Bocchus), as the war, too, could not be carried on without officers,
and as he thought it dangerous to try the faith of new ones after such
perfidy among the old, was involved in doubt and perplexity; no
scheme, no counsel, no person could satisfy him; he changed his route
and his captains daily; he hurried sometimes against the enemy, and
sometimes toward the deserts; depended at one time on flight, and at
another on resistance; and was unable to decide whether he could less
trust the courage or the fidelity of his subjects. Thus, to whatever
direction he turned his thoughts, the prospect was equally
disheartening.

In the midst of his irresolution, Metellus suddenly made his
appearance with his army. The Numidians were assembled and drawn up by
Jugurtha, as well as time permitted; and a battle was at once
commenced. Where the king commanded in person, the struggle was
maintained for some time; but the rest of his force was routed and put
to flight at the first onset. The Romans took a considerable number of
standards and arms, but not many prisoners; for, in almost every
battle, their feet afforded more security to the Numidians than their
swords.

LXXV. In consequence of this defeat, Jugurtha, feeling less confidence
in the state of his affairs than ever, retreated with the deserters,
and part of his cavalry, first into the deserts, and afterward to
Thala,\footnote{LXXV. Thala—The river on which this town stood is not named by
Sallust, but it appears to have been the Bagrada. It seems to have been
nearly destroyed by the Romans, after the defeat of Juba, in the time
of Julius Caesar; though Tacitus, Ann. iii. 21, mentions it as having
afforded a refuge to the Romans in the insurrection of the Numidian
chief, Tacfarinas. D'Anville and Dr. Shaw, \emph{Travels in Bombay}, vol.
i. pt. 2, ch. 5, think it the same with Telepte, now \emph{Ferre-anah}; but
this is very doubtful. See Cellar. iv. 5. It was in ruins in the time
of Strabo.} a large and opulent city, where lay the greater portion of
his treasures, and where there was magnificent provision for the
education of his children. When Metellus was informed of this,
although he knew that there was, between Thala and the nearest river,
a dry and desert region fifty miles broad, yet, in the hope of
finishing the war if he should gain possession of the town, he
resolved to surmount all difficulties, and to conquer even Nature
herself. He gave orders that the beasts of burden, therefore, should
be lightened of all the baggage excepting ten days' provision; and
that they should be laden with skins and other utensils for holding
water. He also collected from the fields as many laboring cattle as he
could find, and loaded them with vessels of all sorts, but chiefly
wooden, taken from the cottages of the Numidians. He directed such of
the neighboring people, too, as had submitted to him after the retreat
of Jugurtha, to bring him as much water as they could carry,
appointing a time and a place for them to be in attendance. He then
loaded his beasts from the river, which, as I have intimated, was the
nearest water to the town, and, thus provided, set out for Thala.

When he came to the place at which he had desired the Numidians to
meet him, and had pitched and fortified his camp, so copious a fall of
rain is said to have happened, as would have furnished more than
sufficient water for his whole army. Provisions, too, were brought him
far beyond his expectations; for the Numidians, like most people after
a recent surrender, had done more than was required of them.\footnote{Had done more than was required of them—\emph{Officia intenderant.}
``Auxit \emph{intenditque} saevitiam exacerbatus indicio filii sui Drusi''
Suet. Tib. 62.} The
men, however, from a religious feeling, preferred using the
rain-water; the fall of which greatly increased their courage, for
they thought of themselves the peculiar care of the gods. On the next
day, to the surprise of Jugurtha, they arrived at Thala. The
inhabitants, who thought themselves secured by difficulties of the
approach to them, were astonished at so strange and unexpected a
sight, but, nevertheless, prepared for their defense. Our men showed
equal alacrity on their side.

LXXVI. But Jugurtha himself, believing that Metellus, who, by his
exertions, had triumphed over every obstacle, over arms, deserts,
seasons, and finally over Nature herself that controls all, nothing
was impossible, fled with his children, and a great portion of his
treasure, from the city during the night. Nor did he ever, after this
time, continue\footnote{LXXVI. Nor did he ever—continue, etc.—\emph{Neque postea—moratus,
simulabat}, etc.—Most editors take \emph{moratus} for \emph{morans}; Allen
places a colon after it, as if it were for \emph{moratus est}.} more than one day or night in any place;
pretending to be hurried away by business, but in reality dreading
treachery, which he thought he might escape by change of residence, as
schemes of such a kind are the result of leisure and opportunity.

Metellus, seeing that the people of Thala were determined on
resistance, and that the town was defended both by art and situation,
surrounded the walls with a rampart and a trench. He then directed his
machines against the most eligible points, threw up a mound, and
erected towers upon it to protect\footnote{And erected towns upon it to protect, etc.—\emph{Et super aggerem
impositis turribus epus et administros tutari}. ``And protected
the work and the workmen with towers placed on the mound.'' \emph{Impositis
turribus} is not the ablative absolute, but the ablative of the
instrument.} the works and the workmen. The
townsmen, on the other hand, were exceedingly active and diligent; and
nothing was neglected on either side. At last the Romans, though
exhausted with much previous fatigue and fighting, got possession,
forty days after their arrival, of the town, and the town only; for
all the spoil had been destroyed by the deserters; who, when they saw
the walls shaken by the battering-ram, and their own situation
desperate, had conveyed the gold and silver, and whatever else is
esteemed valuable, to the royal palace, where, after being sated with
wine and luxuries, they destroyed the treasures, the building, and
themselves, by fire, and thus voluntarily submitted to the sufferings
which, in case of being conquered, they dreaded at the hands of the
enemy.

LXXVII. At the very time that Thala was taken, there came to Metellus
embassadors from the city of Leptis,\footnote{LXXVII. Leptis—Leptis Major, now \emph{Lebida}. In c. 19, Leptis Minor
is meant.} requesting him to send them
a garrison and a governor; saying ``that a certain Hamilcar, a man of
rank, and of a factious disposition, against whom the magistrates and
the laws were alike powerless, was trying to induce them to change
sides; and that unless he attended to the matter promptly, their own
safety,\footnote{Their own safety—\emph{Suam salutem}: i.e. the safety of the people of
Leptis.} and the allies of Rome, would be in the utmost danger.''
For the people at Leptis, at the very commencement of the war with
Jugurtha, had sent to the consul Bestia, and afterward to Rome,
desiring to be admitted into friendship and alliance with us. Having
been granted their request, they continued true and faithful adherents
to us, and promptly executed all orders from Bestia, Albinus, and
Metellus. They therefore readily obtained from the general the aid
which they solicited; and four cohorts of Ligurians were dispatched to
Leptis, with Caius Annius to be governor of the place.

LXXVIII. This city was built by a party of Sidonians, who, as I have
understood, being driven from their country through civil dissensions,
came by sea into those parts of Africa. It is situated between the two
Syrtes, which take their name from their nature\footnote{LXXVIII. Which take their name from their nature—\emph{Quibus
nomen ex re inditum.} From [Greek: \emph{surein}], \emph{to draw,} because the
stones and sand were drawn to and fro by the force of the wind and
tide. But it has been suggested that this etymology is probably false;
it is less likely that their name should be from the Greek than from
the Arabic, in which \emph{sert} signifies a desert tract or region, a term
still applied to the desert country bordering on the Syrtea. See
Ritter, Allgem. vergleich, Geog. vol. i. p. 929. The words which, in
Havercamp, close this description of the Syrtes, ``Syrtes ab tractu
nominatae'', and which Gruter and Putschius suspected not to be
Sallust's, Cortius omitted; and his example has been followed by
Muller and Burnouf; Gerlach, Kritzius, and Dietsch, have retained
them. Gerlach, however, thinks them a gloss, though they are found in
every manuscript but one.} These are two
gulfs almost at the extremity of Africa\footnote{Almost at the extremity of Africa—\emph{Prope in extremâ Africâ.}
``By \emph{extremâ Africâ} Gerlach rightly understands the eastern part of
Africa, bordering on Egypt, and at a great distance from Numidia.''
\emph{Kritzius}.} of unequal size, but of
similar character. Those parts of them next to the land are very deep;
the other parts sometimes deep and sometimes shallow, as chance may
direct; for when the sea swells, and is agitated by the winds, the
waves roll along with them mud, sand, and huge stones; and thus the
appearance of the gulfs changes with the direction of the wind.

Of this people, the language alone\footnote{The language alone—\emph{Lingua modò}.} has been altered by their
intermarriages with the Numidians; their laws and customs continue for
the most part Sidonian; which they have preserved with the greater
ease, through living at so great a distance from the king's
dominions.\footnote{From the king's dominions—\emph{Ab imperio regis.} ``Understand
Masinissa's, Micipsa's, or Jugurtha's.'' \emph{Burnouf}.} Between them and the populous parts of Numidia lie
vast and uncultivated deserts.

LXXIX. Since the affairs of Leptis have led me into these regions, it
will not be foreign to my subject to relate the noble and singular act
of two Carthaginians, which the place has brought to my recollection.

At the time when the Carthaginians were masters of the greater part of
Africa, the Cyrenians were also a great and powerful people. The
territory that lay between them was sandy, and of a uniform
appearance, without a stream or a hill to determine their respective
boundaries; a circumstance which involved them in a severe and
protracted war. After armies and fleets had been routed and put to
flight on both sides, and each people had greatly weakened their
opponents, fearing lest some third party should attack both victors
and vanquished in a state of exhaustion, they came to an agreement,
during a short cessation of arms, ``that on a certain day deputies
should leave home on either side, and that the spot where they should
meet should be the common boundary between the two states.'' From
Carthage, accordingly, were dispatched two brothers, who were named
Philaeni,\footnote{LXXIX. Philaeni—The account of these Carthaginian brothers
with a Greek name, \emph{philainoi, praise-loving}, is probably a fable.
Cortius thinks that the inhabitants, observing two mounds rising above
the surrounding level, fancied they must have been raised, not by
nature, but by human labor, and invented a story to account for their
existence. ``The altars,'' according to Mr. Rennell (Geog. of Herod., p.
640), ``were situated about seven ninths of the way from Carthage to
Cyrene; and the deception,'' he adds, ``would have been too gross, had
it been pretended that the Carthaginian party had traveled seven parts
in nine, while the Cyrenians had traveled no more than two such parts
of the way.'' Pliny (II. N. v. 4) says that the altars were of sand;
Strabo (lib. iii.) says that in his time they had vanished. Pomponius
Mela and Valerius Maximus repeat the story, but without adding any
thing to render it more probable.} and who traveled with great expedition. The deputies of
the Cyrenians proceeded more slowly; but whether from indolence or
accident I have not been informed. However, a storm of wind in these
deserts will cause obstruction to passengers not less than at sea; for
when a violent blast, sweeping over a level surface devoid of
vegetation,\footnote{Devoid of vegetation—\emph{Nuda gignentium}. So c. 93, \emph{cunota
gignentium natura}. Kritzius justly observes that \emph{gignentia} is not
to be taken in the sense of \emph{genita}, as Cortius and others interpret,
but in its own active sense; the ground was bare \emph{of all that was
productive}, or \emph{of whatever generates any thing}. This interpretation
is suggested by Perizonius ad Sanctu Minerv. i. 15.} raises the sand from the ground, it is driven onward
with great force, and fills the mouth and eyes of the traveler, and
thus, by hindering his view, retards his progress. The Cyrenian
deputies, finding that they had lost ground, and dreading punishment
at home for their mismanagement, accused the Carthaginians of having
left home before the time; quarreling about the matter, and preferring
to do any thing rather than submit. The Philaeni, upon this, asked
them to name any other mode of settling the controversy, provided it
were equitable; and the Cyrenians gave them their choice, ``either that
they should be buried alive in the spot which they claimed as the
boundary for their people, or that they themselves, on the same
conditions, should be allowed to go forward to whatever point they
should think proper.'' The Philaeni, having accepted the conditions,
sacrificed themselves\footnote{Sacrificed themselves—\emph{Seque vitamque—condonavere}.
``Nihil aliud est quàm \emph{vitam suam}, sc.[Greek: \emph{eu dia dyoin}].''
\emph{Allen}.} to the interest of their country, and were
interred alive. The people of Carthage consecrated altars to the
brothers on the spot; and other honors were instituted to them at
home. I now return to my subject.

LXXX. After the loss of Thala, Jugurtha, thinking no place sufficiently
secure against Metellus, fled with a few followers into the country of
the Getulians, a people savage and uncivilized, and, at that period,
unacquainted with even the name of Rome. Of these barbarians he collected
a great multitude, and trained them by degrees to march in ranks, to
follow standards, to obey the word of command, and to perform other
military exercises. He also gained over to his interest, by large
presents and larger promises, the intimate friends of king Bocchus, and
working upon the king by their means, induced him to commence war
against the Romans. This was the more practicable and easy, because
Bocchus, at the commencement of hostilities with Jugurtha, had sent an
embassy to Rome to solicit friendship and alliance; but a faction,
blinded by avarice, and accustomed to sell their votes on every question
honorable or dishonorable,\footnote{LXXX. Sell—honorable or dishonorable—\emph{Omnia honesta atque
inhonesta vendere}. See Cat. c. 30. They had been bribed by Jugurtha
to use their influence against Bocchus.} had caused his advances to be rejected,
though they were of the highest consequence to the war recently begun.
A daughter of Bocchus, too, was married to Jugurtha,\footnote{A daughter of Bocchus, too, was married to Jugurtha—\emph{Jugurthae
filia Bocchi nupserat}. Several manuscripts and old editions have
\emph{Boccho}, making Bocchus the son-in-law of Jugurtha. But Plutarch
(Vit. Mar. c. 10, Sull. c. 8) and Florus (iii. 1) agree in speaking
of him as Jugurtha's father-in-law. Bocchus was doubtless an older man
than Jugurtha, having a grown up son, Volux, c. 105. Castilioneus and
Cortius, therefore, saw the necessity of reading \emph{Bocchi}, and, other
editors have followed them, except Gerlach, ``who,'' says Kritzius, ``has
given \emph{Bocchi} in his larger, and \emph{Boccho} in his smaller and more
recent edition, in order that readers using both may have an
opportunity of making a choice.''} but such a
connection, among the Numidians and Moors, is but lightly regarded;
for every man has as many wives as he pleases, in proportion to his
ability to maintain them; some ten, others more, but the kings most of
all. Thus the affection of the husband is divided among a multitude;
no one of them becomes a companion to him,\footnote{No one of them becomes a companion to him—\_Nulla pro sociâ
obtinet The use of \emph{obtinet} absolutely, or with the word dependent on
it understood, prevails chiefly among the later Latin writers. Livy,
however, has \emph{fama obtinuit}, xxi. 46. ``The \emph{tyro} is to be reminded,''
says Dietsch, ``that \emph{obtinet} is not the same as \emph{habetar}, but is
always for \emph{locum obtinet}.''} but all are equally
neglected.

LXXXI. The two kings, with their armies,\footnote{LXXXI. The two kings, with their armies—The text has only
\emph{exercitus}.} met in a place settled
by mutual agreement, where, after pledges of amity were given and
received, Jugurtha inflamed the mind of Bocchus by observing ``that the
Romans were a lawless people, of insatiable covetousness, and the
common enemies of mankind; that they had the same motive for making
war on Bocchus as on himself and other nations, the lust of dominion;
that all independent states were objects of hatred to them; at present,
for instance, himself; a little before, the Carthaginians had been so,
as well as king Perses; and that, in future, as any sovereign became
conspicuous for his power, so would he assuredly be treated as an enemy
by the Romans.''

Induced by these and similar considerations, they determined to march
against Cirta, where Metellus had deposited his plunder, prisoners,
and baggage. Jugurtha supposed that, if he took the city, there would
be ample recompense for his exertions; or that, if the Roman general
came to succor his adherents, he would have the opportunity of
engaging him in the field. He also hastened this movement from policy,
to lessen Bocchus's chance of peace;\footnote{To lessen Bocchus's chance of peace—\emph{Bocchi pacem
imminuere}. He wished to engage Bocchus in some act of hostility
against the Romans, so as to render any coalition between them
impossible.} lest, if delay should be
allowed, he should decide upon something different from war.

LXXXII. Metellus, when he heard of the confederacy of the kings, did
not rashly, or in every place, give opportunities of fighting, as he
had been used to do since Jugurtha had been so often defeated, but,
fortifying his camp, awaited the approach of the kings at no great
distance from Cirta; thinking it better, when he should have learned
something of the Moors,\footnote{LXXXII. Should have learned something of the Moors
—\emph{Cognitis Mauris, i.e.} after knowing something of the Moors,
\emph{and not before}. \emph{Cognitis militibus} is used in the same way in c.
39; and Dietsch says that \emph{amicitia Jugurthae parum cognita} is for
\emph{nondum cognita}, c. 14.} as they were new enemies in the field,
to give battle on an advantage. In the mean time he was informed, by
letters from Rome, that the province of Numidia was assigned to Marius,
of whose election to the consulship he had already heard.

Being affected at these occurrences beyond what was proper and
decorous, he could neither restrain his tears nor govern his tongue;
for though he was a man eminent in other respects, he had too little
firmness in bearing trouble of mind. His irritation was by some
imputed to pride; others said that a noble spirit was wounded by
insult; many thought him chagrined because victory, just attained, was
snatched from his grasp. But to me it is well known that he was more
troubled at the honor bestowed on Marius than at the injustice done to
himself; and that he would have shown much less uneasiness if the
province of which he was deprived had been given to any other than
Marius.

LXXXIII. Discouraged, therefore, by such a mortification, and thinking
it folly to promote another man's success at his own hazard, he sent
deputies to Bocchus, entreating him ``not to become an enemy to the
Romans without cause;'' and observing ``that he had a fine opportunity
of entering into friendship and alliance with them, which were far
preferable to war; that though he might have confidence in his
resources, he ought not to change certainties for uncertainties; that
a war was easily begun, but discontinued with difficulty; that its
commencement and conclusion were not dependent on the same party; that
any one, even a coward, might commence hostilities, but that they
could be broken off only when the conqueror thought proper; and that
he should therefore consult for his interest and that of his kingdom,
and not connect his own prosperous circumstances with the ruined
fortunes of Jugurtha.''

To these representations the king mildly answered, ``that he desired
peace, but felt compassion for the condition of Jugurtha, to whom if
similar proposals were made, all would easily be arranged.'' Metellus,
in reply to this request of Bocchus, sent deputies with overtures, of
which the King approved some, and rejected others. Thus, in sending
messengers to and fro, the time passed away, and the war, according to
the consul's desire, was protracted without being advanced.

LXXXIV. Marius, who, as I said before, had been made consul with great
eagerness on the part of the populace, began, though he had always
been hostile to the patricians, to inveigh against them, after the
people gave him the province of Numidia, with great frequency and
violence; he attacked them sometimes individually and sometimes in a
body; he said that he had snatched from them the consulship as spoils
from vanquished enemies; and uttered other remarks laudatory to
himself and offensive to them. Meanwhile he made the provision for the
war his chief object; he asked for reinforcements for the legions; he
sent for auxiliaries from foreign states, kings, and allies; he also
enlisted all the bravest men from Latium, most of whom were known to
him by actual service, some few only by report, and induced, by
earnest solicitation, even discharged veterans\footnote{LXXXIV. Discharged veterans—\emph{Homines emeritis stipendiis.}
Soldiers who had completed their term of service.} to accompany him.
Nor did the senate, though adverse to him, dare to refuse him any
thing; the additions to the legions they had voted even with
eagerness, because military service was thought to be unpopular with
the multitude, and Marius seemed likely to lose either the means of
warfare\footnote{Means of warfare—\emph{Usum belli.} That is \emph{ea quae belli usus
posceret}, troops and supplies.}, or the favor of the people. But such expectations were
entertained in vain, so ardent was the desire of going with Marius
that had seized on almost all. Every one cherished the fancy\footnote{Cherished the fancy—\emph{Animis trahebant. ``Trahere animo} is
always to revolve in the mind, not to let the thought of a thing
escape from the mind.'' \emph{Kritzius}.} that
he should return home laden with spoil, crowned with victory, or
attended with some similar good fortune. Marius himself, too, had
excited them in no small degree by a speech; for, when all that he
required was granted, and he was anxious to commence a levy, he called
an assembly of the people, as well to encourage them to enlist, as to
inveigh, according to his practice, against the nobility. He spoke, on
the occasion, as follows:

LXXXV. ``I am aware, my fellow-citizens, that most men do not appear as
candidates before you for an office, and conduct themselves in it when
they have obtained it, under the same character; that they are at
first industrious, humble, and modest, but afterward lead a life of
indolence and arrogance. But to me it appears that the contrary should
be the case; for as the whole state is of greater consequence than the
single office of consulate or praetorship, so its interests ought to
be managed\footnote{LXXXV. Its interests ought to be managed, etc.—\emph{Majore curâ
illam administrari quàm haec peti debere.} Cortius injudiciously
omits the word \emph{illam}. No one has followed him but Allen.} with greater solicitude than these magistracies are
sought. Nor am I insensible how great a weight of business I am,
through your kindness, called upon to sustain. To make preparations
for war, and yet to be sparing of the treasury; to press those into
the service whom I am unwilling to offend; to direct every thing at
home and abroad; and to discharge these duties when surrounded by the
envious, the hostile,\footnote{Hostile—\emph{Occursantis.} Thwarting, opposing.} and the factious, is more difficult, my
fellow-citizens, than is generally imagined. In addition to this, if
others fail in their undertakings, their ancient rank, the heroic
actions of their ancestors, the power of their relatives and
connections, their numerous dependents, are all at hand to support
them; but as for me, my whole hopes rest upon myself, which I must
sustain by good conduct and integrity; for all other means are
unavailing.

I am sensible, too, my fellow-citizens, that the eyes of all men are
turned upon me; that the just and good favor me, as my services are
beneficial to the state, but that the nobility seek occasion to attack
me. I must therefore use the greater exertion, that you may not be
deceived in me,\footnote{That you may not be deceived in me—\emph{Ut neque vos capiamini.}
``This verb is undoubtedly used in this passage for \emph{decipere}.
Compare Tibull. Eleg. iii. 6, 45: \emph{Nec vos aut capiant pendentia
brachia collo, Aut fallat blandâ sordida tingua prece.} Cic. Acad.
iv. 20: \emph{Sapientis vim maximam esse cavere, ne capiatur.}'' Gerlach.} and that their views may be rendered abortive. I
have led such a life, indeed, from my boyhood to the present hour,
that I am familiar with every kind of toil and danger; and that
exertion, which, before your kindness to me, I practiced gratuitously,
it is not my intention to relax after having received my reward. For
those who have pretended to be men of worth only to secure their
election,\footnote{To secure their election—\emph{Per ambitionem. Ambire} is to
canvass for votes; to court the favor of the people.} it may be difficult to conduct themselves properly in
office: but to me, who have passed my whole life in the most honorable
occupations, to act well has from habit become nature.

You have commanded me to carry on the war against Jugurtha; a
commission at which the nobility are highly offended. Consider with
yourselves, I pray you, whether it would be a change for the better,
if you were to send to this, or to any other such appointment, one of
yonder crowd of nobles\footnote{Of yonder crowd of nobles—\emph{ex illo globo nobilitatis. Illo,}
[Greek: \emph{deiktikos}].}, a man of ancient family, of innumerable
statues, and of no military experience; in order, forsooth, that in so
important an office, and being ignorant of every thing connected with
it, he may exhibit hurry and trepidation, and select one of the people
to instruct him in his duty. For so it generally happens, that he whom
you have chosen to direct, seeks another to direct him. I know some,
my fellow-citizens, who, after they have been elected\footnote{I know some—who after they have been elected, etc.—``At
whom Marina directs this observation, it is impossible to tell.
Gerlach referring to Cic. Quest. Acad. ii. 1, 2, thinks that Lucullus
is meant. But if he supposes that Lucullus was present \emph{to the mind of
Marius} when he spoke, he is egregiously deceived, for Marius was
forty years antecedent to Lucullus. It is possible, however, that
\emph{Sallust}, thinking of Lucullus when he wrote Marius's speech, may
have fallen into an anachronism, and have attributed to Marius, whose
character he had assumed, an observation which might justly have been
made in his own day.'' \emph{Kritzius}.} consuls,
have begun to read the acts of their ancestors, and the military
precepts of the Greeks; persons who invert the order of things;\footnote{Persons who invert the order of things—\emph{Homines Praeposteri.}
Men who do that last which should be done first.}
for though to discharge the duties of the office\footnote{For though to discharge the duties of the office, etc.—\emph{Nam
gerere, quam fieri, tempore posterius, re atque usu prius est.} With
\emph{gerere} is to be understood \emph{consulatum}; with \emph{fieri, consulem.}
This is imitated from Demosthenes, Olynth. iii.: [Greek: \emph{To gar
prattein ton legein kai cheirotonein, usteron on tae taxei, proteron
tae dynamei kai kreitton esti}.] ``Acting is posterior in order to
speaking and voting, but prior and superior in effect.''} is posterior, in
point of time, to election, it is, in reality and practical
importance, prior to it.

Compare now, my fellow-citizens, me, who am \emph{a new man,} with those
haughty nobles.\footnote{With those haughty nobles—\emph{Cum illorum superbui. Virtus
Scipiades et mitis sapientia Laeli.}} What they have but heard or read, I have
witnessed or performed. What they have learned from books, I have
acquired in the field; and whether deeds or words are of greater
estimation, it is for you to consider.

They despise my humbleness of birth; I contemn their imbecility. My
condition\footnote{My condition \emph{Mihi fortuna}. ``That is, my lot, or condition,
in which I was born, in which I had no hand in producing.'' \emph{Dietsch}.} is made an objection to me; their misconduct is a
reproach to them. The circumstance of birth,\footnote{The circumstance of birth, etc. \emph{Naturam unam et communem
omnium existumo}. ``Nascendi sortem'' is the explanation which Dietsch
gives to \emph{naturam}. One man is \emph{born} as well as another, but the
difference between men is made by their different modes of action; a
difference which the nobles falsely suppose to proceed from fortune.
``Voltaire, Mohammed, Act.I., sce. iv., has expressed the sentiment of
Sallust exactly:} indeed, I consider
as one and the same to all; but think that he who best exerts himself
is the noblest. And could it be inquired of the fathers,\footnote{And could it be inquired of the fathers, etc.—\emph{Ac, si jam
ex patribus Alibini aut Bestiae quaeri posset}, etc. \emph{Patres}, in this
passage, is not, as Anthon imagines, the same as \emph{majores}; as is
apparent from the word \emph{gigni}. The fathers of Albinus and Bestia were
probably dead at the time that Marius spoke. The passage which Anthon
quotes from Plutarch to illustrate \emph{patres}, is not applicable, for
the word there is [Greek: \emph{pragonoi: Epunthaneto ton paronton, ei mae
kai tous ekeinon oiontai progonous auto mallon an emxasthai
paraplaesious ekgonous apolitein, ate dae maed autous di eugeneian,
all ap aretaes kai kalon ergon endoxous genomenous}.] Vit. Mar. c. 9.
``He would then ask the people whether they did not think that the
ancestors of those men would have wished rather to leave a posterity
like him, since they themselves had not risen to glory by their high
birth, but by their virtue and heroic achievements?'' \emph{Langhorne}.} of
Albinus and Bestia, whether they would rather be the parents of them
or of me, what do you suppose that they would answer, but that they
would wish the most deserving to be their offspring! If the patricians
justly despise me, let them also despise their own ancestors, whose
nobility, like mine, had its origin in merit. They envy me the honor
that I have received; let them also envy me the toils, the
abstinence,\footnote{Abstinence—\emph{Innocentiae}. Abstinence from all vicious indulgence.} and the perils, by which I obtained that honor.

But they, men eaten up with pride, live as if they disdained all the
distinctions that you can bestow, and yet sue for those distinctions
as if they had lived so as to merit them. Yet those are assuredly
deceived, who expect to enjoy, at the same time, things so
incompatible as the pleasures of indolence and the rewards of
honorable exertion.\footnote{Honorable exertion—\emph{Virtutis}. See notes on Cat. c. 1, and
Jug. c. 1.}

When they speak before you, or in the senate, they occupy the
greatest part of their orations in extolling their ancestors;\footnote{They occupy the greatest part of their orations in extolling their
ancestors—\emph{Plerâque oratione majores suos extollunt.} ``They extol their
ancestors in the greatest part of their speech.''}
for, they suppose that, by recounting the heroic deeds of their
forefathers, they render themselves more illustrious. But the reverse
of this is the case; for the more glorious were the lives of their
ancestors, the more scandalous is their own inaction. The truth,
indeed, is plainly this, that the glory of ancestors sheds a light on
their posterity,\footnote{The glory of ancestors sheds a light on their posterity, Juvenal,
viii.138:} which suffers neither their virtues nor their
vices to be concealed. Of this light, my fellow-citizens, I have no
share; but I have, what confers much more distinction, the power of
relating my own actions. Consider, then, how unreasonable they are;
what they claim to themselves for the merit of others, they will not
grant to me for my own; alleging, forsooth, that I have no statues,
and that my distinction is newly-acquired; but it is surely better to
have acquired such distinction myself than to bring disgrace on that
received from others.

I am not ignorant, that, if they were inclined to reply to me, they
would make an abundant display of eloquent and artful language. Yet,
since they attack both you and myself on occasion of the great favor
which you have conferred upon me, I did not think proper to be silent
before them, lest any one should construe my forbearance into a
consciousness of demerit. As for myself, indeed, nothing that is said
of me, I feel assured,\footnote{I feel assured—\emph{Ex animi sententiâ}. ``It was a common form of strong
asseveration.'' \emph{Gerlach.}} can do me injury; for what is true, must
of necessity speak in my favor; what is false, my life and character
will refute. But since your judgment, in bestowing on me so
distinguished an honor and so important a trust, is called in
question, consider, I beseech you, again and again, whether you are
likely to repent of what you have done. I can not, to raise your
confidence in me, boast of the statues, or triumphs, or consulships of
my ancestors; but, if it be thought necessary, I can show you spears,\footnote{Spears—\emph{Hastas}. ``A \emph{hasta pura}, that is a spear without iron, was
anciently the reward of a soldier the first time that he conquered in
battle, Serv. ad Virg. Aen. vi. 760; it was afterward given to one who had
struck down an enemy in a sally or skirmish, Lips. ad Polyb. de Milit. Rom.
v.17.'' \emph{Burnouf}.}
a banner,\footnote{A banner—\emph{Vexillum}. ``Standards were also military rewards.
Vopiscus relates that ten \emph{hastae purae}, and four standards of two
colors, were presented to Aurelian. Suetonius (Aug. 25) says that Agrippa
was presented by Augustus, after his naval victory, with a standard of the
color of the sea. These standards therefore, were not, as Badius Ascensius
thinks, always taken from the enemy; though this was sometimes the case,
as appears from Sil. Ital. x.v. 261:} caparisons\footnote{Caparisons—\emph{Phaleras}. ``Sil. Ital. xv. 255:} for horses, and other military rewards;
besides the scars of wounds on my breast. These are my statues; this
is my nobility; honors, not left, like theirs, by inheritance, but
acquired amid innumerable toils and dangers.

My speech, they say, is inelegant; but that I have ever thought of
little importance. Worth sufficiently displays itself; it is for my
detractors to use studied language, that they may palliate base
conduct by plausible words. Nor have I learned Greek; for I had no
wish to acquire a tongue that adds nothing to the valor\footnote{Valor—\emph{Virtutem.} ``The Greeks, those illustrious instructors
of the world, had not been able to preserve their liberty; their
learning therefore had not added to their valor. \emph{Virtus}, in this
passage, is evidently \emph{fortitudo bellica}, which, in the opinion of
Marius, was \emph{the only virtue}.'' Burnouf. See Plutarch, Vit. Mar. c. 2.} of those
who teach it. But I have gained other accomplishments, such as are of
the utmost benefit to a state; I have learned to strike down an enemy;
to be vigilant at my post;\footnote{To be vigilant at my post—\emph{Praesidia agitare}. Or ``to keep
guard at my post.'' ``\emph{Praesidia agitare} signifies nothing more than to
protect a party of foragers or the baggage, or to keep guard round a
besieged city.'' \emph{Vortius}.} to fear nothing but dishonor; to bear
cold and heat with equal endurance; to sleep on the ground; and to
sustain at the same time hunger and fatigue. And with such rules of
conduct I shall stimulate my soldiers, not treating them with rigor
and myself with indulgence, nor making their toils my glory. Such a
mode of commanding is at once useful to the state, and becoming to a
citizen. For to coerce your troops with severity, while you yourself
live at ease, is to be a tyrant, not a general.

It was by conduct such as this, my fellow-citizens, that your
ancestors made themselves and the republic renowned. Our nobility,
relying on their forefathers' merits, though totally different from
them in conduct, disparage us who emulate their virtues; and demand of
you every public honor, as due, not to their personal merit, but to
their high rank. Arrogant pretenders, and utterly unreasonable! For
though their ancestors left them all that was at their disposal, their
riches, their statues, and their glorious names, they left them not,
nor could leave them, their virtue; which alone, of all their
possessions, could neither be communicated nor received.

They reproach me as being mean, and of unpolished manners, because,
forsooth, I have but little skill in arranging an entertainment, and
keep no actor,\footnote{Keep no actor—\emph{Histrionem nullum—habeo}. ``Luxuriae peregrinae
origo ab exercitu Asiatico (Manlii sc. Vulsonis, A.U.C. 563) invecta
in urbem est.——Tum psaltriae sambucistriaeque et convivalia
\emph{ludionum} obiectamenta, addita epulis.'' Liv. xxxix. 6. ``By this army
returning from Asia was the origin of foreign luxury imported into the
city.——At entertainments—were introduced players on the harp and
timbrel, with \emph{buffoons} for the diversion of the guests.'' \emph{Baker}.
Professor Anthon, who quotes this passage, says that \emph{histrio} ``here
denotes a buffoon kept for the amusement of the company.'' But such is
not the meaning of the word \emph{histrio}. It signifies one who in some way
\emph{acted}, either by dancing and gesticulation, or by reciting, perhaps
to the music of the \_sambucistriae or other minstrels. See Smith's
Dict. of Gr. and Rom. Ant. Art. \emph{Histrio}, sect. 2. Scheller's Lex.
sub. vv. \emph{Histrio, Ludio}, and \emph{Salto}. The emperors had whole companies
of actors, \emph{histriones aulici}, for their private amusement. Suetonius
says of Augustus (c. 74) that at feasts he introduced \emph{acroamata et
histriones}. See also Spartian. \emph{Had}. c. 19; Jul.Capitol. \emph{Verus}, c.8.} nor give my cook\footnote{My cook—\emph{Coquum}. Livy, in the passage just cited from him, adds
\emph{tum coquus villisimum antiquis mancipium, et estimatione et usu in
pretio esse; ut quod ministerium fuerat, ars haberi coepta}. ``The cook,
whom the ancients considered as the meanest of their slaves both in
estimation and use, became highly valuable.'' \emph{Baker}.} higher wages than my
steward; all which charges I must, indeed, acknowledge to be just; for
I learned from my father, and other venerable characters, that vain
indulgences belong to women, and labor to men; that glory, rather than
wealth, should be the object of the virtuous; and that arms and armor,
not household furniture, are marks of honor. But let the nobility, if
they please, pursue what is delightful and dear to them; let them
devote themselves to licentiousness and luxury; let them pass their
age as they have passed their youth, in revelry and feasting, the
slaves of gluttony and debauchery; but let them leave the toil and
dust of the field, and other such matters, to us, to whom they are
more grateful than banquets. This, however, they will not do; for when
these most infamous of men have disgraced themselves by every species
of turpitude, they proceed to claim the distinctions due to the most
honorable. Thus it most unjustly happens that luxury and indolence,
the most disgraceful of vices, are harmless to those who indulge in
them, and fatal only to the innocent commonwealth.

As I have now replied to my calumniators, as far as my own character
required, though not so fully as their flagitiousness deserved, I
shall add a few words on the state of public affairs. In the first
place, my fellow-citizens, be of good courage with regard to Numidia;
for all that hitherto protected Jugurtha, avarice, inexperience, and
arrogance\footnote{Avarice, inexperience, and arrogance—\emph{Avaritiam, imperitiam,
superbiam}. ``The President De Brosses and Dotteville have observed,
that Marius, in these words, makes an allusion to the characters of
all the generals that had preceded him, noticing at once the avarice
of Calpurnius, the inexperience of Albinus, and the pride of Metellus.''
\emph{Le Brun}.}, you have entirely removed. There is an army in it,
too, which is well acquainted with the country, though, assuredly,
more brave than fortunate; for a great part of it has been destroyed
by the avarice or rashness of its commanders. Such of you, then, as
are of military age, co-operate with me, and support the cause of your
country; and let no discouragement, from the ill-fortune of others, or
the arrogance of the late commanders, affect any one of you. I myself
shall be with you, both on the march and in the battle, both to direct
your movements and to share your dangers. I shall treat you and myself
on every occasion alike; and, doubtless, with the aid of the gods, all
good things, victory, spoil, and glory, are ready to our hands; though,
even if they were doubtful or distant, it would still become every able
citizen to act in defense of his country. For no man, by slothful
timidity, has escaped the lot of mortals\footnote{For no man, by slothful timidity, has escaped the lot of
mortals—\_Etenim ignavia nemo immortalis factus. The English
translators have rendered this phrase as if they supposed the sense to
be, ``No man has gained immortal renown by inaction.'' But this is not
the signification. What Marius means, is, that \emph{no man, however
cautiously and timidly he may avoid danger, has prolonged his life to
immortality}. Taken in this sense, the words have their proper
connection with what immediately follows: \emph{neque quisquam parens
liberis, uti aeterni forent, optavit}. The sentiment is the same as in
the verse of Horace: \emph{Mors et fugacem persequitur virum}; or in these
lines of Tyrtaeus:}; nor has any parent wished
for his children\footnote{Nor has any parent wished for his children, etc.—[Greek:
\emph{Ou gar athanatous sphisi paidas euchontai genesthai, all' agathous kai
eukleeis}.] ``Men do not pray that they may have children that will
never die, but such as will be good and honorable.'' Plato, Menex. 20.
``This speech, differing from the other speeches in Sallust both in
words and thoughts, conveys a clear notion of that fierce and
objurgatory eloquence which was natural to the rude manners and bold
character of Marius. It is a speech which can not be called polished
and modulated, but must rather be termed rough and ungraceful. The
phraseology is of an antique cast, and some of the wordscoarse.——But
it is animated and fervid, rushing on like a torrent; and by language
of such a character and structure, the nature and manners of Marius are
excellently represented.'' \emph{Gerlach}.} that they might live forever, but rather that they
might act in life with virtue and honor. I would add more, my
fellow-citizens, if words could give courage to the faint-hearted; to
the brave I think that I have said enough.''

LXXXVI. After having spoken to this effect, Marius, when he found that
the minds of the populace were excited, immediately freighted vessels
with provisions, pay, arms, and other necessaries, and ordered Aulus
Manlius, his lieutenant-general, to set sail with them. He himself, in
the mean time, proceeded to enlist soldiers, not after the ancient
method, or from the classes\footnote{LXXXVI. Not after the ancient method, or from the classes—\emph{Non
more majorum, neque ex classibus}. By the regulation of Servius Tullius,
who divided the Roman people into six classes, the highest class
consisting of the wealthiest, and the others decreasing downward in
regular gradation, none of the sixth class, who were not considered as
having any fortune, but were \emph{capite censi}, ``rated by the head,'' were
allowed to enlist in the army. The enlistment of the lower order,
commenced, it is said, by Marius, tended to debase the army, and to
render it a fitter tool for the purposes of unprincipled commanders.
See Aul. Gell., xvi. 10.}, but taking all that were willing to
join him, and the greater part from the lowest ranks. Some said that
this was done from a scarcity of better men, and others from the
consul's desire to pay court\footnote{Desire to pay court—\emph{Per ambitionem}.} to the poorer class, because it was
by that order of men that he had been honored and promoted; and,
indeed, to a man grasping at power, the most needy are the most
serviceable, persons to whom their property (as they have none) is not
an object of care, and to whom every thing lucrative appears honorable.
Setting out, accordingly, for Africa, with a somewhat larger force than
had been decreed, he arrived in a few days at Utica. The command of the
army was resigned to him by Publius Rutilius, Metullus's
lieutenant-general; for Metullus himself avoided the sight of Marius,
that he might not see what he could not even endure to hear mentioned.

LXXXVII. Marius, having filled up his legions\footnote{LXXXVII. Having filled up his legions, etc. Their numbers had been
thinned in actions with the enemy, and Metellus perhaps took home some
part of the army which did not return to it.} and auxiliary
cohorts, marched into a part of the country which was fertile and
abundant in spoil, where, whatever he captured, he gave up to his
soldiers. He then attacked such fortresses or towns as were ill
defended by nature or with troops, and ventured on several
engagements, though only of a light character, in different places.
The new recruits, in process of time, began to join in an encounter
without fear; they saw that such as fled were taken prisoners or
slain; that the bravest were the safest; that liberty, their country,
and parents,\footnote{Their country and parents, etc—\emph{Patriam parentesque}, etc.
Sallust means to say that the soldiers would see such to be the general
effect and result of vigorous warfare; not that they had any country or
parents to protect in Numidia. But the observation has very much of the
rhetorician in it.} are defended, and glory and riches acquired, by
arms. Thus the new and old troops soon became as one body, and the
courage of all was rendered equal.

The two kings, when they heard of the approach of Marius, retreated,
by separate routes, into parts that were difficult of access; a plan
which had been proposed by Jugurtha, who hoped that, in a short time,
the enemy might be attacked when dispersed over the country, supposing
that the Roman soldiers, like the generality of troops, would be less
careful and observant of discipline when the fear of danger was removed.

LXXXVIII. Metellus, meanwhile, having taken his departure for Rome,
was received there, contrary to his expectation, with the greatest
feelings of joy, being equally welcomed, since public prejudice had
subsided, by both the people and the patricians.

Marius continued to attend, with equal activity and prudence, to his
own affairs and those of the enemy. He observed what would be
advantageous, or the contrary, to either party; he watched the
movements of the kings, counteracted their intentions and stratagems,
and allowed no remissness in his own army, and no security in that of
the enemy. He accordingly attacked and dispersed, on several
occasions, the Getulians and Jugurtha on their march, as they were
carrying off spoil from our allies;\footnote{LXXXVIII. From our allies—\emph{Ex sociis nostris}. The people of the
province.} and he obliged the king
himself, near the town of Cirta, to take flight without his arms\footnote{Obliged the king himself—to take flight without his arms
\emph{Ipsumque regem—armis exuerat}. He attacked Jugurtha so suddenly and
vigorously that he was compelled to flee, leaving his arms behind him.}
But finding that such enterprises merely gained him honor, without
tending to terminate the war, he resolved on investing, one after
another, all the cities, which, by the strength of their garrisons or
situation, were best suited either to support the enemy, or to resist
himself; so that Jugurtha would either be deprived of his fortresses,
if he suffered them to be taken, or be forced to come to an engagement
in their defense. As to Bocchus, he had frequently sent messengers to
Marius, saying that he desired the friendship of the Roman people, and
that the consul need fear no act of hostility from him. But whether he
merely dissembled, with a view to attack us unexpectedly with greater
effect, or whether, from fickleness of disposition he habitually
wavered between war and peace, was never fairly ascertained.

LXXXIX. Marius, as he had determined, proceeded to attack the
fortified towns and places of strength, and to detach them, partly by
force, and partly by threats or offers of reward, from the enemy. His
operations in this way, however, were at first but moderate; for he
expected that Jugurtha, to protect his subjects, would soon come to an
engagement. But finding that he kept at a distance, and was intent on
other affairs, he thought it was time to enter upon something of
greater importance and difficulty. Amid the vast deserts there lay a
great and strong city, named Capsa, the founder of which is said to have
been the Libyan Hercules.\footnote{: LXXXIX. The Libyan Hercules—\emph{Hercules Libys}. ``He is one of
the forty and more whom Varro mentions, and who, it is probable, were
leaders of trading expeditions or colonies. See \emph{supra}, c. 18. A
Libyan Hercules is mentioned by Solinus xxvii.'' \emph{Bernouf}.} Its inhabitants were exempted from taxes
by Jugurtha, and under mild government, and were consequently regarded
as the most faithful of his subjects. They were defended against enemies,
not only by walls, magazines of arms, and bodies of troops, but still
more by the difficulty of approaching them; for, except the parts
adjoining the walls, all the surrounding country is waste and
uncultivated, destitute of water, and infested with serpents, whose
fierceness, like that of other wild animals, is aggravated by want of
food; while the venom of such reptiles, deadly in itself, is exacerbated
by nothing so much as by thirst. Of this place Marius conceived a strong
desire\footnote{Marius conceived a strong desire—\emph{Marium maxima cupido
invaserat}. ``A strong desire had seized Marius.''} to make himself master, not only from its importance for the
war, but because its capture seemed an enterprise of difficulty; for
Metellus had gained great glory by taking Thala, a town similarly
situated and fortified; except that at Thala there were several springs
near the walls, while the people of Capsa had only one running stream,
and that within the town, all the water which they used beside being
rain-water. But this scarcity, both here and in other parts of Africa,
where the people live rudely and remote from the sea, was endured with
the greater ease, as the inhabitants subsist mostly on milk and wild
beasts' flesh,\footnote{Wild beasts' flesh—\emph{Ferinâ carne}. Almost all our
translators have rendered this ``venison.'' But the Africans lived on
the flesh of whatever beasts they took in the chase.} and use no salt, or other provocatives of appetite,
their food being merely to satisfy hunger or thirst, and not to encourage
luxury or excess.

XC. The consul,\footnote{XC. The consul, etc.—Here is a long and awkward parenthesis.
I have adhered to the construction of the original. The ``yet,'' \emph{tamen},
that follows the parenthesis, refers to the matter included in it.} having made all necessary investigations, and
relying, I suppose, on the gods (for against such difficulties he
could not well provide by his own forethought, as he was also
straitened for want of corn, because the Numidians apply more to
pasturage than agriculture, and had conveyed, by the king's order,
whatever corn had been raised into fortified places, while the ground
at the time, it being the end of summer, was parched and destitute of
vegetation), yet, under the circumstances, conducted his arrangements
with great prudence. All the cattle, which had been taken for some
days previous, he consigned to the care\footnote{He consigned to the care, etc.—\emph{Equitibus auxiliariis agendum
attribuit}. ``He gave to be driven by the auxiliary cavalry.''} of the auxiliary cavalry;
and directed Aulus Manlius, his lieutenant-general, to proceed with
the light-armed cohorts to the town of Lares,\footnote{The town of Lares—\emph{Oppidum Laris}. Cortius seems to have
been right in pronouncing \emph{Laris} to be an accusative plural. Gerlach
observes that Lares occurs in the Itinerary of Antonius and in St.
Augustine, Adv. Donatist., vi. 28.} where he had
deposited provisions and pay for the army, telling him that, after
plundering the country, he would join him there in a few days. Having
by this means concealed his real design, he proceeded toward the river
Tana.

XCI. On his march he distributed daily, to each division of the
infantry and cavalry, an equal portion of the cattle, and gave orders
that water-bottles should be made of their hides; thus compensating,
at once, for the scarcity of corn, and providing, while all remained
ignorant of his intention, utensils which would soon be of service. At
the end of six days, accordingly, when he arrived at the river, a
large number of bottles had been prepared. Having pitched his camp,
with a slight fortification, he ordered his men to take refreshment,
and to be ready to resume their march at sunset; and, having laid aside
all their baggage, to load themselves and their beasts only with water.
As soon as it seemed time, he quitted the camp, and, after marching the
whole night,\footnote{XCI. After marching the whole night.—He seems to have
marched in the night for the sake of coolness.} encamped again.

The same course he pursued on the following night, and on the third,
long before dawn, he reached a hilly spot of ground, not more than two
miles distant from Capsa, where he waited, as secretly as possible,
with his whole force. But when daylight appeared, and many of the
Numidians, having no apprehensions of an enemy, went forth out of the
town, he suddenly ordered all the cavalry, and with them the lightest
of the infantry, to hasten forward to Capsa, and secure the gates. He
himself immediately followed, with the utmost ardor, restraining his
men from plunder.

When the inhabitants perceived that the place was surprised, their
state of consternation and extreme dread, the suddenness of the
calamity, and the consideration that many of their fellow-citizens
were without the walls in the power of the enemy, compelled them to
surrender. The town, however, was burned; of the Numidians, such as
were of adult age, were put to the sword; the rest were sold, and the
spoil divided among the soldiers. This severity, in violation of the
usages of war, was not adopted from avarice or cruelty in the consul,
but was exercised because the place was of great advantage to Jugurtha,
and difficult of access to us, while the inhabitants were a fickle and
faithless race, to be influenced neither by kindness nor by terror.

XCII. When Marius had achieved so important an enterprise, without any
loss to his troops, he who was great and honored before became still
greater and still more honored. All his undertakings,\footnote{XCII. All his undertakings, etc.—\emph{Omnia non bene consulta
in virtutem trahebantur}. ``All that he did rashly was attributed to
his \emph{consciousness of} extraordinary power.'' If they could not praise
his prudence, they praised his resolution and energy.} however
ill-concerted, were regarded as proofs of superior ability; his
soldiers, kept under mild discipline, and enriched with spoil,
extolled him to the skies; the Numidians dreaded him as some thing
more than human; and all, indeed, allies as well as enemies, believed
that he was either possessed of supernatural power, or had all things
directed for him by the will of the gods.

After his success in this attempt, he proceeded against other towns; a
few, where they offered resistance, he took by force; a greater number,
deserted in consequence of the wretched fate of Capsa, he destroyed by
fire; and the whole country was filled with mourning and slaughter.

Having at length gained possession of many places, and most of them
without loss to his army, he turned his thoughts to another enterprise,
which, though not of the same desperate character as that at Capsa, was
yet not less difficult of execution.\footnote{Difficult of execution—\emph{Difficilem}. There seemed to be as
many impediments to success as in the affair at Capsa, though the
undertaking was not of so perilous a nature.} Not far from the river Mulucha,
which divided the kingdoms of Jugurtha and Bocchus, there stood, in the
midst of a plain,\footnote{In the midst of a plain—\emph{Inter caeteram planitiem}. By
\emph{caeteram} he signifies that \emph{the rest} of the ground, except the part
on which the fort stood, was plain and level.} a rocky hill, sufficiently broad at the top for
a small fort; it rose to a vast height, and had but one narrow ascent
left open, the whole of it being as steep by nature as it could have
been rendered by labor and art. This place, as there were treasures of
the king in it, Marius directed his utmost efforts to take.\footnote{Directed his utmost efforts to take—\emph{Summâ vi capere
intendit}. It is to be observed that \emph{summâ vi} refers to \emph{intendit},
not to \emph{capere}. \emph{Summâ ope} animum \emph{intendit ut caperet}.} But
his views were furthered more by fortune than by his own contrivance.
In the fortress there were plenty of men and arms for its defense,
as well as an abundant store of provisions, and a spring of water;
while its situation was unfavorable for raising mounds, towers, and
other works; and the road to it, used by its inhabitants, was extremely
steep, with a precipice on either side. The vineae were brought up with
great danger, and without effect; for, before they were advanced any
considerable distance, they were destroyed with fire or stones. And from
the difficulties of the ground, the soldiers could neither stand in front
of the works, nor act among the vineae,\footnote{Among the vineae—\emph{Inter vineas}. ``\emph{Inter}, for which Müller,
from a conjecture of Glareanus, substituted \emph{intra} is supported by
all the manuscripts, and ought not to be altered, although \emph{intra}
would have been more exact, as the signification of \emph{inter} is of
greater extent, and includes that of \emph{intra}. \emph{Inter} is used when
a thing is inclosed on each side; \emph{intra}, when it is inclosed on
all sides. If the soldiers, therefore, are considered as surrounded
with the \emph{vineae}, they should be described as \emph{intra vineas}; but
as there is no reason why they may not also be contemplated as being
inclosed only laterally by the \emph{vineae}, the phrase \emph{inter vineas}
may surely in that case be applied to them. Gronovius and Drakenborch
ad Liv., i. 10, have observed how often these propositions are
interchanged when referred to \emph{time}.'' Kritzius. On \emph{vineae}, see
c. 76.} without danger; the boldest
of them were killed or wounded, and the fear of the rest increased.

XCIII. Marius having thus wasted much time and labor, began seriously
to consider whether he should abandon the attempt as impracticable, or
wait for the aid of Fortune, whom he had so often found favorable.
While he was revolving the matter in his mind, during several days and
nights, in a state of much doubt and perplexity, it happened that a
certain Ligurian, a private soldier in the auxiliary cohorts,\footnote{XCIII. A certain Ligurian—in the auxiliary cohorts—The
Ligurians were not numbered among the Italians or \emph{socii} in the Roman
army, but attached to it only as auxiliaries.}
having gone out of the camp to fetch water, observed, near that part
of the fort which was furthest from the besiegers, some snails
crawling among the rocks, of which, when he had picked up one or two,
and afterward more, he gradually proceeded, in his eagerness for
collecting them, almost to the top of the hill. When he found this
part deserted, a desire, incident to the human mind, of seeing what he
had never seen,\footnote{A desire—of seeing what he had never seen—\emph{More humani
ingenii, cupido ignara visundi invadit}. This is the reading of
Cortius, to which Muller and Allen adhere. Gerlach inserted in his
text, \emph{More humani ingeni, cupidio difficilia faciundi animum vortit};
which Kritzius, Orelli, and Dietsch, have adopted, and which Cortius
acknowledged to be the reading of the generality of the manuscripts,
except that they vary as to the last two words, some having
\emph{animad vortit}. The sense of this reading will be, ``the desire of
doing something difficult, which is natural to the human mind, drew
off his thoughts from gathering snails, and led him to contemplate
something of a more arduous character.'' But the reading of Cortius
gives so much better a sense to the passage, that I have thought
proper to follow it. Burnouf, with Havercamp and the editions
antecedent to Cortius, reads \emph{more humanae cupidinis ignara visundi
animum vortit}, of which the first five words are taken from a
quotation of Aulus Gellius, ix. 12, who, however, may have transcribed
them from some other part of Sallust's works, now lost.} took violent possession of him. A large oak
chanced to grow out among the rocks, at first, for a short distance,
horizontally,\footnote{Horizontally—\emph{Prona}. This word here signifies \emph{forward}, not
\emph{downward}, as Anthon and others interpret, for trees growing out
of a rock or bank will not take a \emph{descending} direction.} and then, as nature directs all vegetables,\footnote{As nature directs all vegetables—\emph{Quò cuncta gignentium
natura fert}. It is to be observed that the construction is \emph{natura
fert cuncta gignentium}, for \emph{cuncta gignentia}. On \emph{gignentia}, i.e.
vegetable, or \emph{whatever produces any thing}, see c. 79, and Cat., c.
53.}
turning and shooting upward. Raising himself sometimes on the boughs
of this tree, and sometimes on the projecting rocks, the Ligurian, as
all the Numidians were intently watching the besiegers, took a full
survey of the platform of the fortress. Having observed whatever he
thought it would afterward prove useful to know, he descended the same
way, not unobservantly, as he had gone up, but exploring and noticing
all the peculiarities of the path. He then hastened to Marius,
acquainted him with what he had done, and urged him to attack the fort
on that side where he had ascended, offering himself to lead the way
and the attempt. Marius sent some of those about him, along with the
Ligurian, to examine the practicability of his proposal, who,
according to their several dispositions, reported the affair as
difficult or easy. The consul's hopes, however, were somewhat
encouraged; and he accordingly selected, from his band of trumpeters
and bugle-men, five of the most nimble, and with them four centurions
for a guard;\footnote{Four centurions for a guard—\emph{Praesidio qui forent, quatuor
centuriones}. It is a question among the commentators whether the
centurions were attended by their centuries or not; Cortius thinks
that they were not, as ten men were sufficient to cause an alarm in
the fortress, which was all that Marius desired. But that Cortius is
in the wrong, and that there were common soldiers with the centurions,
appears from the following considerations: 1. Marius would hardly have
sent, or Sallust have spoken of, \emph{four} men as a guard to \emph{six}. 2.
Why should centurions only have been selected, and not common soldiers
as well as their officers? 3. An expression in the following chapter,
\emph{laqueis—quibus allevati milites facilius escenderent}, seems to
prove that there were others present besides the centurions and the
trumpeters. The word \emph{milites} is indeed wanting in the text of
Cortius, but appears to have been omitted by him merely to favor his
own notion as to the absence of soldiers, for he left it out, as
Kritzius says, \emph{summâ libidine, ne uno quidem codice assentiente},
``purely of his own will, and without the authority of a single
manuscript.'' Taking a fair view of the passage, we seem necessarily
led to believe that the centurions were attended by a portion, if not
the whole, of their companies. See the following note.} all of whom he directed to obey the Ligurian,
appointing the next day for commencing the experiment.

XCIV. When, according to their instructions, it seemed time to set
out, the Ligurian, after preparing and arranging every thing,
proceeded to the place of ascent. Those who commanded the
centuries,\footnote{XCIV. Those who commanded the centuries—\emph{Illi qui centuriis
praeerant}. This is the reading of several manuscripts, and of almost
all the editions before that of Kritzius, and may be tolerated if we
suppose that the centurions were attended by their men, and that
Sallust, in speaking of the change of dress, meant \emph{to include the
men}, although he specifies only the officers. Yet it is difficult
to conceive why Sallust should have used such a periphrase for
\emph{centuriones}. Seven of the manuscripts, however, have \emph{qui adsensuri
erant}, which Kritzius and Dietsch have adopted. Two have \emph{qui ex
centuriis praeerant}. Allen, not unhappily, conjectures, \emph{qui
praesidio erant}. Cortius suspected the phrase, \emph{qui centuriis
praeerant}, and thought it a transformation of the words \emph{qui
adscensuris praeerat}, which somebody had written in the margin as an
explanation of the following word \emph{duce}, and which were afterward
altered and thrust into the text.} being previously instructed by the guide, had changed
their arms and dress, having their heads and feet bare, that their
view upward, and their progress among the rocks, might be less
impeded;\footnote{Progress—might be less impeded—\emph{Nisus—faciliùs foret}. The
adverb for the adjective. So in the speech of Adherbal, c. 14, \emph{ut
tutius essem}.} their swords were slung behind them, as well as their
shields, which were Numidian, and made of leather, both for the sake
of lightness, and in order that, if struck against any object, they
might make less noise. The Ligurian went first, and tied to the rocks,
and whatever roots of trees projected through age, a number of ropes,
by which the soldiers supporting themselves might climb with the
greatest ease. Such as were timorous, from the extraordinary nature of
the path, he sometimes pulled up by the hand; when the ascent was
extremely rugged, he sent them on singly before him without their
arms, which he then carried up after them; whatever parts appeared
unsafe,\footnote{Unsafe—\emph{Dubia nisu}. ``Not to be depended upon for support.''
\emph{Nisu} is the old dative for \emph{nisui}.} he first tried them himself, and, by going up and down
repeatedly in the same place, and then standing aside, he inspired the
rest with courage to proceed. At length, after uninterrupted and
harassing exertion they reached the fortress, which, on that side, was
undefended, for all the occupants, as on other days, were intent on
the enemy in the opposite quarter.

Though Marius had kept the attention of the Numidians, during the
whole day, fixed on his attacks, yet, when he heard from his scouts
how the Ligurian had succeeded, he animated his soldiers to fresh
exertions, and he himself, advancing beyond the vineae, and causing a
testudo to be formed,\footnote{Causing a testudo to be formed—\emph{Testudine actâ}. The soldiers
placed their shields over their heads, and joined them close together,
forming a defense like the shell of a tortoise.} came up close under the walls, annoying the
enemy, at the same time, with his engines, archers, and slingers, from
a distance.

But the Numidians, having often before overturned and burned the
vineae of the Romans, no longer confined themselves within the
fortress, but spent day and night before the walls, railing at the
Romans, upbraiding Marius with madness, threatening our soldiers with
being made slaves to Jugurtha, and exhibiting the utmost audacity on
account of their successful defense. In the mean time, while both the
Romans and Numidians were engaged in the struggle, the one side
contending for glory and dominion, the other for their very existence,
the trumpets suddenly sounded a blast in the rear of the enemy, at
which the women and children, who had gone out to view the contest,
were the first to flee; next those who were nearest to the wall, and
at length the whole of the Numidians, armed and unarmed, retreated
within the fort. When this had happened, the Romans pressed upon the
enemy with increased boldness, dispersing them, and at first only
wounding the greater part, but afterward making their way over the
bodies of those who fell, thirsting for glory, and striving who should
be first to reach the wall; not a single individual being detained by
the plunder. Thus the rashness of Marius, rendered successful by
fortune, procured him renown from his very error.

XCV. During the progress of this affair, Lucius Sylla, Marius's
quaestor, arrived in the camp with a numerous body of cavalry, which
he had been left at Rome to raise among the Latins and allies.

Of so eminent a man, since my subject brings him to my notice, I think
it proper to give a brief account of the character and manners; for I
shall in no other place allude to his affairs;\footnote{XCV. For I shall in no other place allude to his affairs—\emph{Neque
enim allo loco de Sullae rebus dicturi sumus.} ``These words show that
Sallust, at this time, had not thought of writing \emph{Histories}, but
that he turned his attention to that pursuit after he had finished
the Jugurthine war. For that he spoke of Sylla in his large history
is apparent from several extant fragments of it, and from Plutarch,
who quotes Sallust, Vit. Syll., c. 3.'' \emph{Kritzius.}} and Lucius
Sisenna,\footnote{Lucius Sisenna—He wrote a history of the civil wars between
Sylla and Marius, Vell. Paterc. ii. 9. Cicero alludes to his style
as being jejune and puerile, Brut., c. 64, De Legg. i. 2. About a
undred and fifty fragments of his history remain.} who has treated that subject the most ably and accurately
of all writers, seems to me to have spoken with too little freedom.
Sylla, then, was of patrician descent, but of a family almost sunk in
obscurity by the degeneracy of his forefathers. He was skilled, equally
and profoundly, in Greek and Roman literature. He was a man of large
mind, fond of pleasure, but fonder of glory. His leisure was spent in
luxurious gratifications, but pleasure never kept him from his duties,
except that he might have acted more for his honor with regard to his
wife\footnote{Except that he might have acted more for his honor with
regard to his wife—\emph{Nisi quod de uxore potuit honestius consuli.} As
these words are vague and indeterminate, it is not agreed among the
critics and translators to what part of Sylla's life Sallust refers.
I suppose, with Rupertus, Aldus, Manutius, Crispinus, and De Brosses,
that the allusion is to his connection with Valeria, of which the
history is given by Plutarch in his life of Sylla, which the English
reader may take in Langhorne's translation: ``A few months after
Metella's death, he presented the people with a show of gladiators;
and as, at that time, men and women had no separate places, but sat
promiscuously in the theater, a woman of great beauty, and of one of
the best families, happened to sit near Sylla. She was the daughter of
Messala, and sister to the orator Hortensius; her name was Valeria;
and she had lately been divorced from her husband. This woman, coming
behind Sylla, touched him, and took off a little of the nap of his
robe, and then returned to her place. Sylla looked at her, quite
amazed at her familiarity, when she said, 'Wonder not, my lord, at
what I have done; I had only a mind to share a little in your good
fortune.' Sylla was far from being displeased; on the contrary, it
appeared that he was flattered very agreeably, for he sent to ask her
name, and to inquire into her family and character. Then followed an
interchange of amorous regards and smiles, which ended in a contract
and marriage. The lady, perhaps, was not to blame. But Sylla, though
he got a woman of reputation, and great accomplishments, yet came into
the match upon wrong principles. Like a youth, he was caught with soft
looks and languishing airs, things that are wont to excite the lowest
of the passions.'' Others have thought that Sallust refers to Sylla's
conduct on the death of his wife Metella, above mentioned, to whom, as
she happened to fall sick when he was giving an entertainment to the
people, and as the priest forbade him to have his house defiled with
death on the occasion, he unfeelingly sent a bill of divorce, ordering
her to be carried out of the house while the breath was in her.
Cortius, Kritz, and Langius. think that the allusion is to Sylla'a
general faithlessness to his wives, for he had several; as if Sallust
had used the singular for the plural, \emph{uxore} for \emph{uxoribus}, or
\emph{reuxoriâ}; but if Sallust meant to allude to more than one wife, why
should he have restricted himself to the singular?}. He was eloquent and subtle, and lived on the easiest terms
with his friends.\footnote{Lived on the easiest terms with his friends—\emph{Facilis
amicitiâ} The critics are in doubt about the sense of this phrase. I
have given that which Dietsch prefers, who says that a man \emph{facilis
amicitiâ} is ``one who easily grants his friends all that they desire,
exacts little from them, and is no severe censor of their morals.''
Cortius explains it \emph{facilis ad amicitiam}, and Facciolati, in his
Lexicon, \emph{facilè sibi amicos parans}, but these interpretations, as
Kritzius observes, are hardly suitable to the ablative case.} His depth of thought in disguising his
intentions, was incredible; he was liberal of most things, but
especially of money. And though he was the most fortunate \footnote{Most fortunate—\emph{Felicissumo}. Alluding, perhaps, to the
title of Felix, which he assumed after his great victory over Marius.} of all
men before his victory in the civil war, yet his fortune was never
beyond his desert;\footnote{His desert—\emph{Industriam}. That is, the efforts which he made to
attain distinction.} and many have expressed a doubt whether his
success or his merit were the greater. As to his subsequent acts, I
know not whether more of shame, or of regret must be felt at the
recital of them.

XCVI. When Sylla came with his cavalry into Africa, as has just been
stated, and arrived at the camp of Marius, though he had hitherto been
unskilled and undisciplined in the art of war, he became, in a short
time, the most expert of the whole army. He was besides affable to the
soldiers; he conferred favors on many at their request, and on others
of his own accord, and was reluctant to receive any in return. But he
repaid other obligations more readily than those of a pecuniary
nature; he himself demanded repayment from no one; but rather made it
his object that as many as possible should be indebted to him. He
conversed, jocosely as well as seriously, with the humblest of the
soldiers; he was their frequent companion at their works, on the
march, and on guard. Nor did he ever, as is usual with depraved
ambition, attempt to injure the character of the consul, or of any
deserving person.

His sole aim, whether in the council or the field, was to suffer none
to excel him; to most he was superior. By such conduct he soon became
a favorite both with Marius and with the army.

XCVII. Jugurtha, after he had lost the city of Capsa, and other strong
and important places, as well as a vast sum of money, dispatched
messengers to Bocchus, requesting him to bring his forces into Numidia
as soon as possible, and stating that the time for giving battle was
at hand. But finding that he hesitated, and was balancing the
inducements to peace and war, he again corrupted his confidants, as on
a previous occasion, with presents, and promised the Moor himself a
third part of Numidia, should either the Romans be driven from Africa,
or the war brought to an end without any diminution of his own
territories. Being allured by this offer, Bocchus joined Jugurtha with
a large force.

The armies of the kings being thus united, they attacked Marius, on
his march to his winter quarters, when scarcely a tenth part of the
day remained\footnote{XCVII. When scarcely a tenth part of the day remained—\emph{Vix
decimâ parte die reliquâ.} A remarkably exact specification of the time.}, expecting that the night, which was now coming on,
would be a shelter to them if they were beaten, and no impediment if
they should conquer, as they were well acquainted with the country,
while either result would be worse for the Romans in the dark. At the
very moment, accordingly, that Marius heard from various quarters\footnote{From various quarters—\emph{Ex multis.} From his scouts, who came in
from all sides.}
of the enemy's approach, the enemy themselves were upon him, and
before the troops could either form themselves or collect the baggage,
before they could receive even a signal or an order, the Moorish and
Getulian horse, not in line, or any regular array of battle, but in
separate bodies, as chance had united them, rushed furiously on our
men; who, though all struck with a panic, yet, calling to mind what
they had done on former occasions, either seized their arms, or
protected those who were looking for theirs, while some, springing on
their horses, advanced against the enemy. But the whole conflict was
more like a rencounter with robbers than a battle; the horse and foot
of the enemy, mingled together without standards or order, wounded
some of our men, and cut down others, and surprised many in the rear
while fighting stoutly with those in front; neither valor nor arms
were a sufficient defense, the enemy being superior in numbers, and
covering the field on all sides. At last the Roman veterans, who were
necessarily well experienced in war,\footnote{The Roman veterans, who were necessarily well experienced
in war,—The reading of Cortius is, \emph{Romani veteres, novique, et ob
ea scientes belli;} which he explains by supposing that the new
recruits \emph{were joined with} the veterans, and that both united were
consequently well skilled in war, citing, in support of his
supposition, a passage in c. 87: \emph{Sic brevi spatio} novi veteresqua
\emph{coaluere, et virtus omnium aequalis facta.} And Ascensius had
previously given a similar explanation, \emph{quod etiam veterani
adessent.} But many later critics have not been induced to believe
that Cortius's reading will bear any such interpretation; and
accordingly Kritzius, Dietsch, and Orelli, have ejected \emph{novique}; as
indeed Ciaeconius and Ursinus had long before recommended. Muller,
Burnouf, and Allen, retain it, adopting Cortius's interpretation.
Gerlach also retains it, but not without hesitation. But it is very
remarkable that it occurs in all the manuscripts but one, which has
\emph{Romani veteres boni scientes erant ut quos locus,} etc.} formed themselves, wherever
the nature of the ground or chance allowed them to unite, in circular
bodies, and thus secured on every side, and regularly drawn up,
withstood the attacks of the enemy.

XCVIII. Marius, in this desperate emergency, was not more alarmed or
disheartened than on any previous occasion, but rode about with his
troop of cavalry, which he had formed of his bravest soldiers rather
than his nearest friends, in every quarter of the field, sometimes
supporting his own men when giving way, sometimes charging the enemy
where they were thickest, and doing service to his troops with his
sword, since, in the general confusion, he was unable to command with
his voice.

The day had now closed, yet the barbarians abated nothing of their
impetuosity, but, expecting that the night would be in their favor,
pressed forward, as their kings had directed them, with increased
violence. Marius, in consequence, resolved upon a measure suited to
his circumstances, and, that his men might have a place of retreat,
took possession of two hills contiguous to each other, on one of
which, too small for a camp, there was an abundant spring of water,
while the other, being mostly elevated and steep, and requiring little
fortification, was suited for his purpose as a place of encampment. He
then ordered Sylla, with a body of cavalry, to take his station for
the night on the eminence containing the spring, while he himself
collected his scattered troops by degrees, the enemy being not less
disordered\footnote{\emph{Neque minus hostibus conturbatis}. If the enemy had not been
in as much disorder as himself, Marius would hardly have been able to
effect his retreat.}, and led them all at a quick march\footnote{\emph{Pleno gradu}.—``By the \emph{militaris gradus} twenty miles were
completed in five hours of a summer day; by the \emph{plénusus}, which is
quicker, twenty-four miles were traversed in the same time.'' Veget. i.9.} up the other
hill. Thus the kings, obliged by the strength of the Roman position,
were deterred from continuing the combat; yet they did not allow their
men to withdraw to a distance, but, surrounding both hills with a
large force, encamped without any regular order. Having then lighted
numerous fires, the barbarians, after their custom, spent most of the
night in merriment, exultation, and tumultuous clamor, the kings,
elated at having kept their ground, conducting themselves as
conquerors. This scene, plainly visible to the Romans, under cover of
the night and on the higher ground, afforded great encouragement to
them.

XCIX. Marius, accordingly, deriving much confidence from the
imprudence of the enemy, ordered the strictest possible silence to be
kept, not allowing even the trumpets, as was usual, to be sounded when
the watches were changed;\footnote{XCIX. When the watches were changed—\_Per vigilias: i. e.
at the end of each watch, when the guards were relieved. ``The nights,
by the aid of a clepsydra, were divided into four watches, the
termination of each being marked by the blast of a trumpet or horn.
See Viget. in. 8: \emph{A tubicine omnes vigiliae committuntur; et finitis
horis a cornicine revocantur}.'' Kritzius He also refers to Liv. vii.
35; Lucan. viii. 24; Tacit. Hist. v. 22.} cavalry, and legions, to sound all and
then, when day approached, and the enemy were fatigued and just
sinking to sleep, he ordered the sentinels, with the trumpeters of the
auxiliary cohorts,\footnote{Auxiliary cohorts—\emph{Cohortium}. I have added the word \emph{auxiliary}.
That they were the cohorts of the auxiliaries or allies is apparent,
as the word \emph{legionum} follows. Kritzius indeed thinks otherwise,
supposing that the cohorts had particular trumpeters, distinct from
those of the whole legion. But for this notion there seems to be no
sufficient ground. Sallust speaks of the \emph{cohortes sociorum}, c. 58,
and \emph{cohortes Ligurum}, c. 100.} their instruments at once, and the soldiers,
at the same time, to raise a shout, and sally forth from the camp\footnote{Sally forth from the camp—\emph{Portis erumpere}. Sallust uses
the common phrase for issuing from the camp. It can hardly be
supposed, that the Romans had formed a regular camp with gates during
the short time that they had been upon the hill, especially as they
had fled to it in great disorder.}
upon the enemy. The Moors and Getulians, suddenly roused by the
strange and terrible noise, could neither flee, nor take up arms,
could neither act, nor provide for their security, so completely had
fear, like a stupor,\footnote{Stupor—\emph{Vecordia}. A feeling that deprived them of all sense.} from the uproar and shouting, the absence of
support, the charge of our troops, and the tumult and alarm, seized
upon them all. The whole of them were consequently routed and put to
flight; most of their arms, and military standards, were taken; and
more were killed in this than in all former battles, their escape
being impeded by sleep and the sudden alarm.

C. Marius now continued the route, which he had commenced, toward his
winter quarters, which, for the convenience of getting provisions, he
had determined to fix in the towns on the coast. He was not, however,
rendered careless or presumptuous by his victory, but marched with his
army in form of a square\footnote{C. in form of a square—\emph{Quadrato agmine}. ``A hollow square,
with the baggage in the center; see Serv. ad Verg. Aen. xii.121. …
Such an \emph{agmen} Sallust, in c. 46, calls \emph{munitum}, as it was
prepared to defend itself against the enemy, from whatever quarter
they might approach.'' \emph{Kritzius}.} just as if he were in sight of the
enemy. Sylla, with his cavalry, was on the right; Aulus Manlius, with
the slingers and archers, and Ligurian cohorts, had the command on the
left; the tribunes, with the light-armed infantry, the consul had
placed in the front and rear. The deserters, whose lives were of
little value, and who were well acquainted with the country, observed
the route of the enemy. Marius himself, too, as if no other were
placed in charge, attended to every thing, went through the whole of
the troops, and praised or blamed them according to their desert. He
was always armed and on the alert, and obliged his men to imitate his
example. He fortified his camp with the same caution with which he
marched; stationing cohorts of the legions to watch the gates, and the
auxiliary cavalry in front, and others upon the rampart and lines. He
went round the posts in person, not from suspicion that his orders
would not be observed, but that the labor of the soldiers, shared
equally by their general, might be endured by them with cheerfulness.
\footnote{Might be endured by them with cheerfulness \emph{Volentibus
esset.} A Greek phrase, \emph{Boulomenois eiae.}} Indeed, Marius, as well at this as at other periods of the war,
kept his men to their duty rather by the dread of shame\footnote{Dread of shame—\emph{Pudore.} Inducing each to have a regard to
his character.} than of
severity; a course which many said was adopted from desire of popularity,
but some thought it was because he took pleasure in toils to which he had
been accustomed from his youth, and in exertions which other men call
perfect miseries. The public interest, however, was served with as much
efficiency and honor as it could have been under the most rigorous
command.

CI. At length, on the fourth day of his march, when he was not far
from the town of Cirta, his scouts suddenly made their appearance from
all quarters at once; a circumstance by which the enemy was known to
be at hand. But as they came in from different points, and all gave
the same account, the consul, doubting in what form to draw up his
army, made no alteration in it, but halted where he was, being already
prepared for every contingency. Jugurtha's expectations, in consequence,
disappointed him; for he had divided his force into four bodies, trusting
that one of them, assuredly,\footnote{CI. Trusting that one of them, assuredly, etc.—\emph{Ratua es
omnibus aque aliquos ab tergo hostibus ventures. By aequo Sallust}
signifies that each of the four bodies would have an equal chance of
coming on the rear of the Romans.} would surprise the Romans in the rear.
Sylla, meanwhile, with whom they first came in contact, having cheered
on his men, charged the Moors, in person and with his officers,\footnote{In person and with his officers—\emph{Ipse aliique.} ``The
\emph{alii}, are the \emph{praefecti equitum,} officers of the cavalry.''
\emph{Kritzius.}}
with troop after troop of cavalry, in the closest order possible; while
the rest of his force, retaining their position, protected themselves
against the darts thrown from a distance, and killed such of the enemy
as fell into their hands.

While the cavalry was thus engaged, Bocchus, with his infantry, which
his son Volux had brought up, and which, from delay on their march,
had not been present in the former battle, assailed the Romans in the
rear. Marius was at that moment occupied in front, as Jugurtha was
there with his largest force. The Numidian king, hearing of the
arrival of Bocchus, wheeled secretly about, with a few of his
followers, to the infantry,\footnote{Wheeled secretly about, with a few of his followers, to the
infantry—\emph{Clam—ad pedites convertit}. What infantry are meant, the
commentators can not agree, nor is there any thing in the narrative on
which a satisfactory decision can be founded. As the arrival of
Bocchus is mentioned immediately before, Cortius supposes that the
infantry of Bocchus are signified; and it may be so; but to whatever
party the words wore addressed, they were intended to be heard by the
Romans, or for what purpose were they spoken in Latin? Jugurtha may
have spoken the words in both languages, and this, from what follows,
would appear to have been the case, for both sides understood him.
\emph{Quod ubi milites} (evidently the Roman soldiers) \emph{accepere—simul
barbari animos tollere}, etc. The \emph{clam} signifies that Jugurtha
turned about, or wheeled off, so as to escape the notice of Marius,
with whom he had been contending.} and exclaimed in Latin, which he had
learned to speak at Numantia, ``that our men wore struggling in vain;
for that he had just slain Marius with his own hand;'' showing, at the
same time, his sword besmeared with blood, which he had, indeed,
sufficiently stained by vigorously cutting down our infantry\footnote{By vigorously cutting down our infantry—\emph{Satis impigre
occiso pedite nostro}. ``A ces mots il leur montra son épée teinte du
sang des notres, dont il venait, en effet, de faire une assez cruelle
boucherie.'' \emph{De Brosses}. Of the other French translators, Beauzée and
Le Brun render the passage in a similar way; Dotteville and Durean
Delamalle, as well as all our English translators, take \emph{pedite} as
signifying \emph{only one soldier}. Sir Henry Steuart even specifies that
it was ``a legionary soldier.'' The commentators, I should suppose, have
all regarded the word as having a plural signification; none of them,
except Burnouf, who expresses a needless doubt, say any thing on the
point.}.

When the soldiers heard this, they felt a shock, though rather at the
horror of such an event, than from belief in him who asserted it; the
barbarians, on the other hand, assumed fresh courage, and advanced
with greater fury on the disheartened Romans, who were just on the
point of taking to flight, when Sylla, having routed those to whom he
had been opposed, fell upon the Moors in the flank. Bocchus instantly
fled. Jugurtha, anxious to support his men, and to secure a victory so
nearly won, was surrounded by our cavalry, and all his attendants,
right and left, being slain, had to force a way alone, with great
difficulty, through the weapons of the enemy. Marius, at the same
time, having put to flight the cavalry, came up to support such of his
men as he had understood to be giving ground. At last the enemy were
defeated in every quarter. The spectacle on the open plains was then
frightful;\footnote{The spectacle on the open plains was then frightful—\emph{Tum
spectaculum horribile campis patentibus}, etc. The idea of this
passage was probably taken, as Ciacconius intimates, from a
description in Xenophon, Agesil. ii. 12, 14, part of which is quoted
by Longinus, Sect. 19, as an example of the effect produced by the
omission of conjunctions: [Greek: \emph{Kai symbalontes tas aspidas
eothounto, emachonto, apekteinon, apethnaeskon Epei ge maen elaexen
hae machae, paraen dae theasasthai entha synepeson allaelois, taen men
gaen aimati pe, ormenaen, nekrous de peimenous philious kai polemious
met allaelon, aspidas de diatethrummenas, dorata syntethrausmena,
egchoipidia gumna kouleon ta men chamai, ta d'en somasi, ta d'eti meta
cheiras}.] ``Closing their shields together, they pushed, they fought,
… But when the battle was over, you might have seen, where they had
fought, the ground clotted with blood, the corpses of friends and
enemies mingled together, and pierced shields, broken lances, and
swords without their sheaths, strewed on the ground, sticking in the
dead bodies, or still remaining in the hands that had wielded them
when alive.'' Tacitus, Agric. c. 37. has copied this description of
Sallust, as all the commentators have remarked: \emph{Tum vero patentibus
locis grande et atrox spectaculum. Sequi, vulnerare, capere, atque
eosdem, oblatis aliis, trucidare…. Passim, arma et corpora, et
laceri artus, et cruenta humus}. ``The sight on the open field was then
striking and horrible; they pursued, they inflicted wounds, they took
… Every where were seen arms and corpses, mangled limbs, and the
ground stained with blood.''} some were pursuing, others fleeing; some were being
slain, others captured; men and horses were dashed to the earth; many,
who were wounded, could neither flee nor remain at rest, attempting to
rise, and instantly falling back; and the whole field, as far as the
eye could reach, was strewed with arms and dead bodies, and the
intermediate spaces saturated with blood.

CII. At length the consul, now indisputably victor, arrived at the
town of Cirta, whither he had at first intended to go. To this place,
on the fifth day after the second defeat of the barbarians, came
messengers from Bocchus, who, in the king's name, requested of Marius
to send him two persons in whom he had full confidence, as he wished
to confer with them on matters concerning both the interest of the
Roman people and his own. Marius immediately dispatched Sylla and
Aulus Manlius; who, though they went at the king's invitation, thought
proper, notwithstanding, to address him first, in the hope of altering
his sentiments, if he were unfavorable to peace, or of strengthening
his inclination, if he were disposed to it. Sylla, therefore, to whose
superiority, not in years but in eloquence, Manlius yielded
precedence, spoke to Bocchus briefly as follows:

``It gives us great pleasure, King Bocchus, that the gods have at
length induced a man, so eminent as yourself, to prefer peace to war,
and no longer to stain your own excellent character by an alliance
with Jugurtha, the most infamous of mankind; and to relieve us, at the
same time, from the disagreeable necessity of visiting with the same
punishment your errors and his crimes. Besides, the Roman people, even
from the very infancy\footnote{Besides, the Roman people, even from the very infancy—The
reading of this passage, before the edition of Cortius, was this: \emph{Ad
hoc, populo Romano jam a principio inopi melius visum amicos, quam
servos, quaerere}. Gruter proposed to read \emph{Ad hoc populo Romano inopi
melius est visum}, etc., whence Cortius made \emph{Ad hoc, populo Romano jam
inopi visum}, etc. But the Bipont editors, observing that \emph{inopi} was
not quite consistent with \emph{quaerere servos}, altered the passage to
\emph{Ad hoc, populo Romano jam à principio reipublicae melius visum},
etc., which seems to be the best emendation that has been proposed,
and which I have accordingly followed. Kritzius and Dietsch adopt it,
except that they omit \emph{reipublicae}, and put nothing in the place of
\emph{inopi}. Gerlach retains \emph{inopi}, on the principle of ``quo
insolentius, eo verius,'' and it may, after all, be genuine. Cortius
omitted \emph{melius} on no authority but his own.} of their state, have thought it better to
seek friends than slaves, thinking it safer to rule over willing than
forced subjects. But to you no friendship can be more suitable than
ours; for, in the first place, we are at a distance from you, on which
account there will be the less chance of misunderstanding between us,
while our good feeling for you will be as strong as if we were near;
and, secondly, because, though we have subjects in abundance, yet
neither we, nor any other nation, can ever have a sufficiency of
friends. Would that such had been your inclination from the first; for
then you would assuredly, before this time, have received from the
Roman people more benefits than you have now suffered evils. But since
Fortune has the chief control in human affairs, and it has pleased her
that you should experience our force as well as our favor, now, when,
she gives you this fair opportunity, embrace it without delay, and
complete the course which you have begun. You have many and excellent
means of atoning, with great ease, for past errors by future services.
Impress this, however, deeply on your mind, that the Roman people are
never outdone in acts of kindness; of their power in war you have
already sufficient knowledge.''

To this address Bocchus made a temperate and courteous reply, offering
a few observations, at the same time, in extenuation of his error; and
saying ``that he had taken arms, not with any hostile feeling, but to
defend his own dominions, as part of Numidia, out of which he had
forcibly driven Jugurtha,\footnote{Out of which he had forcibly driven Jugurtha—\emph{Unde ut
Jugurtham expulerit [expulerat]} There is here some obscurity. The
manuscripts vary between \emph{expulerit} and \emph{expulerat}. Cortius, and
Gerlaen in his second edition, adopt \emph{expulerat}, which they of
necessity refer to Marius; but to make Bocchus speak thus, is, as
Kritzius says, to make him speak very foolishly and arrogantly.
Kritzius himself, accordingly, adopts \emph{expulerit}, and supposes that
Bocchus invents a falsehood, in the belief that the Romans wouldhave
no means of detecting it. But Bocchus may have spoken truth, referring,
as Müller suggests, to some previous transactions between him and
Jugurtha, to which Sallust does not elsewhere allude.} was his by right of conquest, and he
could not allow it to be laid waste by Marius; that when he formerly
sent embassadors to the Romans, he was refused their friendship; but
that he would say nothing more of the past, and would, if Marius gave
him permission, send another embassy to the senate.'' But no sooner was
this permission granted, than the purpose of the barbarian was altered
by some of his friends, whom Jugurtha, hearing of the mission of Sylla
and Manlius, and fearful of what was intended by it, had corrupted
with bribes.

CIII. Marius, in the mean time, having settled his army in winter
quarters, set out, with the light-armed cohorts and part of the
cavalry, into a desert part of the country, to besiege a fortress of
Jugurtha's, in which he had placed a garrison consisting wholly of
Roman deserters. And now again Bocchus, either from reflecting on what
he had suffered in the two engagements, or from being admonished by
such of his friends as Jugurtha had not corrupted, selected, out of
the whole number of his adherents, five persons of approved integrity
and eminent abilities, whom he directed to go, in the first place, to
Marius, and afterward to proceed, if Marius gave his consent, as
embassadors to Rome, granting them full powers to treat concerning his
affairs, and to conclude the war upon any terms whatsoever. These five
immediately set out for the Roman winter-quarters, but being beset and
spoiled by Getulian robbers on the way, fled, in alarm and ill
plight,\footnote{In ill plight—\emph{Sine decore}.} to Sylla, whom the consul, when he went on his expedition,
had left as pro-praetor with the army. Sylla received them, not, as they
had deserved, like faithless enemies, but with the greatest ceremony and
munificence; from which the barbarians concluded that what was said of
Roman avarice was false, and that Sylla, from his generosity, must be
their friend. For interested bounty,\footnote{For interested bounty—\emph{Largitio}. ``The word signifies liberal
treatment of others with a view to our own interest; without any real
goodwill.'' \emph{Müller}. ``He intends a severe stricture on his own age,
and the manners of the Romans.'' \emph{Dietsch}.} in those days, was still
unknown to many; by whom every man who was liberal was also thought
benevolent, and all presents were considered to proceed from kindness.
They therefore disclosed to the quaestor their commission from Bocchus,
and asked him to be their patron and adviser; extolling, at the same
time, the power, integrity, and grandeur of their monarch, and adding
whatever they thought likely to promote their objects, or to procure
the favor of Sylla. Sylla promised them all that they requested; and,
being instructed how to address Marius and the senate, they tarried in
the camp about forty days.\footnote{About forty days. Waiting, apparently, for the return of Marius.}

CIV. When Marius, having failed in the object\footnote{CIV. Having failed in the object, etc.—\emph{Infecto, quo
intenderat, negotio.} Though this is the reading of most of the
manuscripts, Kritzius, Müller, and Dietach, read \emph{confecto}, as if
Marius could not have failed in his attempt.} of his expedition,
returned to Cirta, and was informed of the arrival of the embassadors,
he desired both them and Sylla to come to him, together with Lucius
Bellienus, the praetor from Utica, and all that were of senatorial rank
in any part of the country, with whom he discussed the instructions of
Bocchus to his embassadors; to whom permission to proceed to Rome was
granted by the consul. In the mean time a truce was asked, a request
to which assent was readily expressed by Sylla and the majority; the
few, who advocated harsher measures, were men inexperienced in human
affairs, which, unstable and fluctuating, are always verging to
opposite extremes.\footnote{Are always verging to opposite extremes.—\emph{Semper in adversa
mutant}. Rose renders this ``are always changing, and constantly for
the worse;'' and most other translators have given something similar.
But this is absurd; for every one sees that all changes in human
affairs are not for the worse. \emph{Adversa} is evidently to be taken in
the sense which I have given.}

The Moors having obtained all that they desired, three of them started
for Rome with Oneius Octavius Rufus, who, as quaestor, had brought pay
for the army to Africa; the other two returned to Bocchus, who heard
from them, with great pleasure, their account both of other
particulars, and especially of the courtesy and attention of Sylla.

To his three embassadors that went to Rome, when, after a deprecatory
acknowledgment that their king had been in error, and had been led
astray by the treachery of Jugurtha, they solicited for him friendship
and alliance, the following answer was given: ``The senate and people
of Rome are wont to be mindful of both services and injuries; they
pardon Bocchus, since he repents of his fault, and will grant him
their alliance and friendship when he shall have deserved them.''

CV. When this reply was communicated to Bocchus, he requested Marius,
by letter, to send Sylla to him, that, at his discretion,\footnote{CV. At his discretion—\emph{Arbitratu}. Kritzius observes that
this word comprehends the notion of plenary powers to treat and
decide: \emph{der mit unbeschränkter Vollmacht unterhandeln könnte}.}
measures might be adopted for their common interest. Sylla was
accordingly dispatched, attended with a guard of cavalry, infantry,
and Balearic slingers, besides some archers and a Pelignian cohort,
who, for the sake of expedition, were furnished with light arms,
which, however, protected them, as efficiently as any others, against
the light darts of the enemy. As he was on his march, on the fifth day
after he set out, Volux, the son of Bocchus, suddenly appeared on the
open plain with a body of cavalry, which amounted in reality to not
more than a thousand, but which, as they approached in confusion and
disorder, presented to Sylla and the rest the appearance of a greater
number, and excited apprehensions of hostility. Every one, therefore,
prepared himself for action, trying and presenting\footnote{Presenting—\emph{Intendere}. The critics are in doubt to what
to refer this word; some have thought of understanding \emph{animum};
Cortius, Wasse, and Müller, think it is meant only of the bows of the
archers; Kritzius, Burnouf, and Allen, refer it, apparently with
better judgment, to the \emph{arma} and \emph{tela} in general.} his arms and
weapons; some fear was felt among them, but greater hope, as they were
now conquerors, and were only meeting those whom they had often
overcome. After a while, however, a party of horse sent forward to
reconnoiter, reported, as was the case, that nothing but peace was
intended.

CVI. Volux, coming forward, addressed himself to Sylla, saying that he
was sent by Bocchus his father to meet and escort him. The two parties
accordingly formed a junction, and prosecuted their journey, on that
day and the following, without any alarm. But when they had pitched
their camp, and evening had set in, Volux came running, with looks of
perplexity, to Sylla, and said that he had learned from his scouts
that Jugurtha was at hand, entreating and urging him, at the same
time, to escape with him privately in the night. Sylla boldly replied,
``that he had no fear of Jugurtha, an enemy so often defeated; that he
had the utmost confidence in the valor of his troops; and that, even
if certain destruction were at hand, he would rather keep his ground,
than save, by deserting his followers, a life at best uncertain, and
perhaps soon to be lost by disease.'' Being pressed, however, by Volux,
to set forward in the night, he approved of the suggestion, and
immediately ordered his men to dispatch their supper,\footnote{CVI. To dispatch their supper—\emph{Coenatos esse}. ``The perfect is
not without its force; it signifies that Sylla wished his orders to
be performed with the greatest expedition.'' \emph{Kritzius}. He orders them
\emph{to have done} supper.} to light as
many fires as possible in the camp, and to set out in silence at the
first watch.

When they were all fatigued with their march during the night, and
Sylla was preparing, at sunrise, to pitch his camp, the Moorish
cavalry announced that Jugurtha was encamped about two miles in
advance. At this report, great dismay fell upon our men; for they
believed themselves betrayed by Volux, and led into an ambuscade. Some
exclaimed that they ought to take vengeance on him at once, and not
suffer such perfidy to remain unpunished.

CVII. But Sylla, though he had similar thoughts, protected the Moor
from violence; exhorting his soldiers to keep up their spirits; and
saying, ``that a handful of brave men had often fought successfully
against a multitude; that the less anxious they were to save their
lives in battle, the greater would be their security; and that no man,
who had arms in his hands, ought to trust for safety to his unarmed
heels, or to turn to the enemy, in however great danger, the
defenseless and blind parts of his body''.\footnote{CVII. And blind parts of his body—\emph{Caecum corpus}. Imitated
from Xenephon, Cyrop. iii. 3, 45: [Greek: \emph{Moron gar to kratein
boulomenous, ta tuphla, tou somatos, kai aopla, tauta enantia
tattein tois polemiois pheugontas}.] ``It is folly for those that
desire to conquer to turn the blind, unarmed, and handless parts of
the body, to the enemy in flight.''} Having then called
almighty Jupiter to witness the guilt and perfidy of Bocchus, he
ordered Volux, as being an instrument of his father's hostility,\footnote{At being an instrument of his father's hostility—\emph{Quoniam
hostilia faceret}. ``Since he wished to deceive the Romans by pretended
friendship.'' \emph{Müller}.}
to quit the camp.

Volux, with tears in his eyes, entreated him to entertain no such
suspicions; declaring ``that nothing in the affair had been caused by
treachery on his part, but all by the subtlety of Jugurtha, to whom
his line of march had become known through his scouts. But as Jugurtha
had no great force with him, and as his hopes and resource were
dependent on his father Bocchus, he assuredly would not attempt any
open violence, when the son of Bocchus would himself be a witness of
it. He thought it best for Sylla, therefore, to march boldly through
the middle of his camp, and that as for himself, he would either send
forward his Moors, or leave them where they were, and accompany Sylla
alone.'' This course, under such circumstances, was adopted; they set
forward without delay, and, as they came upon Jugurtha unexpectedly,
while he was in doubt and hesitation how to act, they passed without
molestation. In a few days afterward, they arrived at the place to
which their march was directed.

CVIII. There was, at this time, in constant and familiar intercourse
with Bocchus, a Numidian named Aspar, who had been sent to him by
Jugurtha, when he heard of Sylla's intended interview, in the
character of embassador, but secretly to be a spy on the Mauretanian
king's proceedings. There was also with him a certain Dabar, son of
Massugrada, one of the family of Masinissa,\footnote{CVIII. Of the family of Masinissa—\emph{Ex gente Masinissae.}
Massugrada was the son of Masinissa by a concubine.} but of inferior birth
on the maternal side, as his father was the son of a concubine. Dabar,
for his many intellectual endowments, was liked and esteemed by
Bocchus, who, having found him faithful\footnote{Faithful—\emph{Fidum}. After this word, in the editions of Cortius,
Kritzius, Gerlach, Allen, and Dietsch, follows \emph{Romanis} or \emph{esse
Romanis}. These critics defend \emph{Romanis} on the plea that a dative
is necessary after \emph{fidum}, and that it was of importance, as
Castilioneus observes that Dabar should be well disposed toward the
Romans, and not have been corrupted, like many other courtiers of
Bocchus, by the bribes of Jugurtha. Glarcanus, Badius Ascensius, the
Bipont editors, and Burnouf, with, most of the translators, omit
\emph{Romanis}, and I have thought proper to imitate their example.} on many former occasions,
sent him forthwith to Sylla, to say that ``he was ready to do whatever
the Romans desired; that Sylla himself should appoint the place, day,
and hour,\footnote{Place, day, and hour—\emph{Diem, locum, tempus.} Not only the
day, but the time of the day.} for a conference; that he kept all points, which he had
settled with him before, inviolate;\footnote{That he kept all points, which he had settled with him
before, inviolate—\emph{Consulta sese omnia cum illo integra habere}.
Kritzius justly observes that most editors, in interpreting this
passage, have erroneously given to \emph{consulta} the sense of
\emph{consulenda}; and that the sense is, ``that all that he had arranged
with Sylla before, remained unaltered, and that he was not drawn from
his resolutions by the influence of Jugurtha.''} and that he was not to fear
the presence of Jugurtha's embassador as any restraint\footnote{And that he was not to fear the presence of Jugurtha's
embassador, as any restraint, etc.—\emph{Neu Jugurthae legatum
pertimesceret, quo res communis licentius gereretur}. There is some
difficulty in this passage. Burnouf makes the nearest approach to a
satisfactory explanation of it. ``Sylla,'' says he, ``was not to fear the
envoy of Jugurtha, \emph{quo}, on which account (equivalent to \emph{eoque}, and
on that account, \emph{i. e.} on account of his freedom from apprehension)
their common interests would be more freely arranged.'' Yet it appears
from what follows that fear of Jugurtha's envoy \emph{could not be
dismissed}, and that there could be no freedom of discussion in his
presence, as Sylla was to say but little before him, and to speak more
at large at a private meeting. These considerations have induced
Kritzius to suppose that the word \emph{remoto}, or something similar, has
been lost after \emph{quo}. The Bipont editors inserted \emph{cautum esse}
before \emph{quo}, which is without authority, and does not at all assist
the sense.} on the
discussion of their common interests, since, without admitting him, he
could have no security against Jugurtha's treachery''. I find, however,
that it was rather from African duplicity\footnote{African duplicity—\emph{Punica fide}. ``\emph{Punica fides} was a
well-known proverbial expression for treachery and deceit. The origin
of it is perhaps attributable not so much to fact, as to the implacable
hatred of the Romans toward the Carthaginians.'' \emph{Burnouf}.} than from the motives
which he professed, that Bocchus thus allured both the Romans and
Jugurtha with the hopes of peace; that he frequently debated with
himself whether he should deliver Jugurtha to the Romans, or Sylla to
Jugurtha; and that his inclination swayed him against us, but his
fears in our favor.

CIX. Sylla replied, ``that he should speak on but few particulars
before Aspar, and discuss others at a private meeting, or in the
presence of only a few;'' dictating, at the same time, what answer
should be returned by Bocchus.\footnote{CIX. What answer should be returned by Bocchus—That is, in
the presence of Aspar.} Afterward, when they met, as
Bocchus had desired, Sylla stated, ``that he had come, by order of the
consul, to inquire whether he would resolve on peace or on war.''
Bocchus, as he had been previously instructed by Sylla, requested him
to come again at the end of ten days, since he had as yet formed no
determination, but would at that time give a decisive answer. Both
then retired to their respective camps.\footnote{Both then retired to their respective camps—\emph{Deinde ambo in
sua castra digressi}. Both, \emph{i. e.} Bocchus and Sylla, not Aspar and
Sylla, as Cortius imagines.} But when the night was
far advanced, Sylla was secretly sent for by Bocchus. At their
interview, none but confidential interpreters were admitted on either
side, together with Dabar, the messenger between them, a man of honor,
and held in esteem by both parties. The king at once commenced thus:

CX. ``I never expected that I, the greatest monarch in this part of the
world, and the richest of all whom I know, should ever owe a favor to
a private man. Indeed, Sylla, before I knew you, I gave assistance to
many who solicited me, and to others without solicitation, and stood
in need of no man's assistance.

But at this loss of independence, at which others are wont to repine,
I am rather inclined to rejoice. It will be a pleasure to me\footnote{CX. It will be a pleasure to me—\emph{Fuerit mihi}. Some editions,
as that of Langius, the Bipont, and Burnouf's, have \emph{fuerit mihi
pretium}. Something of the kind seems to be wanting. ``Res in bonis
numeranda fuerit mihi.'' \emph{Burnouf}. Allen, who omits \emph{pretium},
interprets, ``Grata mihi egestas sit, quae ad tuam, amicitiam
coufugiat;'' but who can deduce this sense from the passage, unless he
have \emph{pretium}, or something similar, in his mind?} to
have once needed your friendship, than which I hold nothing dearer to
my heart. Of the sincerity of this assertion you may at once make
trial, take my arms, my soldiers, my money, or whatever you please,
and use it as your own. But do not suppose, as long as you live, that
your kindness to me has been fully requited; my sense of it will
always remain undiminished, and you shall, with my knowledge, wish for
nothing in vain. For, as I am of opinion, it is less dishonorable to a
prince to be conquered in battle than to be surpassed in generosity.

With respect to your republic, whose interests you are sent to guard,
hear briefly what I have to say. I have neither made war upon the
Roman people, nor desired that it should be made; I have merely
defended my territories with arms against an armed force. But from
hostilities, since such is your pleasure, I now desist. Prosecute the
war with Jugurtha as you think proper. The river Malucha, which was
the boundary between Miscipsa and me, I shall neither pass myself, nor
suffer Jugurtha to come within it. And if you shall ask any thing
besides, worthy of me and of yourself, you shall not depart with a
refusal.''

CXI. To this speech Sylla replied, as far as concerned himself,
briefly and modestly; but spoke, with regard to the peace and their
common concerns, much more at length. He signified to the king ``that
the senate and people of Rome, as they had the superiority in the
field, would think themselves little obliged by what he promised; that
he must do something which would seem more for their interest than his
own; and that for this there was now a fair opportunity, since he had
Jugurtha in his power, for, if he delivered him to the Romans, they
would feel greatly indebted to him, and their friendship and alliance,
as well as that part of Numidia which he claimed,\footnote{CXI. That part of Numidia which he claimed—\emph{Numidiae partem
quam nunc peteret}. See the second note on c. 102. Bocchus continues,
in his speech in the preceding chapter, to signify that a part of
Numidia belonged to him.} would readily
be granted him.'' Bocchus at first refused to listen to the proposal,
saying that affinity, the ties of blood,\footnote{The ties of blood—\emph{Cognationem}. To this blood-relationship
between him and Jugurtha no allusion is elsewhere made.} and a solemn league,
connected him with Jugurtha; and that he feared, if he acted
insincerely, he might alienate the affections of his subjects, by whom
Jugurtha was beloved, and the Romans disliked. But at last, after
being frequently importuned, his resolution gave way,\footnote{His resolution gave way—\emph{Lenitur}. Cortius whom Gerlach and
Müller follow, reads \emph{leniter}, but, with Kritzius and Gerlach, I
prefer the verb to the adverb; which, however, is found in the greater
number of the manuscripts.} and he
engaged to do every thing in accordance with Sylla's wishes. They then
concerted measures for conducting a pretended treaty of peace, of
which Jugurtha, weary of war, was extremely desirous. Having settled
their plans, they separated.

CXII. On the next day Bocchus sent for Aspar, Jugurtha's envoy, and
acquainted him that he had ascertained from Sylla, through Dabar, that
the war might be concluded on certain conditions; and that he should
therefore make inquiry as to the sentiments of his king. Aspar
proceeded with joy to Jugurtha's camp, and having received full
instructions from him, returned in haste to Bocchus at the end of
eight days, with intelligence ``that Jugurtha was eager to do whatever
might be required, but that he put little confidence in Marius, as
treaties of peace, concluded with Roman generals, had often before
proved of no effect; that if Bocchus, however, wished to consult the
interests of both,\footnote{CXII. Interests of both—\emph{Ambobus}. Both himself and Jugurtha.} and to have an established peace, he should
endeavor to bring all parties together to a conference, as if to
settle the conditions, and then deliver Sylla into his hands, for when
he had such a man in his power, a treaty would at once be concluded by
order of the senate and people of Rome; as a man of high rank, who had
fallen into the hands of the enemy, not from want of spirit, but from
zeal for the public interest, would not be left in captivity''.

CXIII. The Moor, after long meditation on these suggestions, at length
expressed his assent to them, but whether in pretense or sincerity I
have not been able to discover. But the inclinations of kings, as they
are violent, are often fickle, and at variance with themselves. At
last, after a time and place were fixed for coming to a conference
about the treaty, Bocchus addressed himself at one time to Sylla and
at another to the envoy of Jugurtha, treating them with equal
affability, and making the same professions to both. Both were in
consequence equally delighted, and animated with the fairest
expectations. But on the night preceding the day appointed for the
conference, the Moor, after first assembling his friends, and then,
on a change of mind, dismissing them, is reported to have had many
anxious struggles with himself, disturbed alike in his thoughts and
his gestures, which, even when he was silent, betrayed the secret
agitation of his mind. At last, however, he ordered that Sylla should
be sent for, and, according to his desire, laid an ambush for
Jugurtha.

As soon as it was day, and intelligence was brought that Jugurtha was
at hand, Bocchus, as if to meet him and do him honor, went forth,
attended by a few friends, and our quaestor, as far as a little hill,
which was full in the view of the men who were placed in ambush. To
the same spot came Jugurtha with most of his adherents, unarmed,
according to agreement; when immediately, on a signal being given, he
was assailed on all sides by those who were lying in wait. The others
were cut to pieces, and Jugurtha himself was delivered bound to Sylla,
and by him conducted to Marius.

CXIV. At this period war was carried on unsuccessfully by our generals
Quintus Caepio and Marcus Manlius, against the Gauls; with the terror
of which all Italy was thrown into consternation. Both the Romans of
that day, indeed, and their descendants, down to our own times,
maintained the opinion that all other nations must yield to their
valor, but that they contended with the Gauls, not for glory, but
merely in self-defense. But after the war in Numidia was ended, and
it was announced that Jugurtha was coming in chains to Rome, Marius,
though absent from the city, was created consul, and Gaul decreed to
him as his province. On the first of January he triumphed as consul,
with great glory. At that time\footnote{CXIV. At that time—\emph{Eâ tempestate}. ``In many manuscripts
is found \emph{ex eâ tempestate}, by which the sense is wholly perverted.
Sallust signifies that Marius did not continue always deserving of
such honor; for, as is said in c. 63, 'he was afterward carried
headlong by ambition.''' \emph{Kritzius}.} the hopes and dependence of the
state were placed on him.
